% Recovered from the previously completed PDF; no new mathematical content.
% Prose is editable; formula boxes retain original symbol positions.
% Build from repository root: xelatex -no-shell-escape -halt-on-error -output-directory=... print/fudan/NAME.tex
\documentclass[11pt,a4paper]{article}
\usepackage[margin=23mm,headheight=14pt,headsep=8.5mm,footskip=10mm]{geometry}
\usepackage{amsmath,amssymb}
\DeclareMathAlphabet{\mathsf}{OT1}{cmss}{m}{n}
\usepackage[no-math]{fontspec}
\usepackage{xeCJK}
\setmainfont{Latin Modern Roman}
\setCJKmainfont{FandolSong-Regular.otf}[Path=assets/fonts/fandol/,BoldFont=FandolSong-Bold.otf,ItalicFont=FandolSong-Regular.otf,Script=CustomDefault]
\xeCJKsetup{PunctStyle=quanjiao,CheckSingle=true}
\usepackage{setspace}\setstretch{1.12}
\setlength{\parindent}{2em}\setlength{\parskip}{3pt}
\emergencystretch=6em
\hfuzz=9pt
\usepackage{titlesec}
\titleformat{\section}{\centering\normalfont\Large\bfseries}{}{0pt}{}
\titlespacing*{\section}{0pt}{1.8\baselineskip}{1.0\baselineskip}
\usepackage{fancyhdr}\pagestyle{fancy}\fancyhf{}
\fancyhead[L]{\small 数学手记·高等代数}\fancyhead[R]{\small 16S03–16S09}
\fancyfoot[C]{\thepage}\renewcommand{\headrulewidth}{0.4pt}
\usepackage[hidelinks,unicode]{hyperref}
\hypersetup{pdftitle={16S03–16S09：题目、解答与小结},pdfauthor={Vincent Arqain}}
\newcommand{\RB}[4]{\raisebox{0pt}[#2bp][#3bp]{\hbox to #1bp{#4\hss}}}
\newcommand{\RG}[5]{\hbox to 0pt{\kern#1bp\raisebox{#2bp}[0pt][0pt]{\fontsize{#3bp}{#3bp}\csname font#4\endcsname\char#5}\hss}}
\newcommand{\RT}[4]{\hbox to 0pt{\kern#1bp\raisebox{#2bp}[0pt][0pt]{\fontsize{#3bp}{#3bp}\selectfont#4}\hss}}
\newcommand{\RS}[4]{\hbox to 0pt{\kern#1bp\raisebox{#2bp}[0pt][0pt]{\fontsize{#3bp}{#3bp}\selectfont\ensuremath{#4}}\hss}}
\newcommand{\RR}[4]{\hbox to 0pt{\kern#1bp\raisebox{#2bp}[0pt][0pt]{\rule{#3bp}{#4bp}}\hss}}
\newcommand{\fontmi}{\usefont{OML}{cmm}{m}{it}}
\newcommand{\fontsy}{\usefont{OMS}{cmsy}{m}{n}}
\newcommand{\fontex}{\usefont{OMX}{cmex}{m}{n}}
\newcommand{\fontrm}{\usefont{OT1}{cmr}{m}{n}}
\newcommand{\fontbf}{\usefont{OT1}{cmr}{bx}{n}}
\newcommand{\fontam}{\usefont{U}{msa}{m}{n}}
\newcommand{\fontbb}{\usefont{U}{msb}{m}{n}}
\widowpenalty=10000\clubpenalty=10000\raggedbottom
\begin{document}
\begin{center}{\LARGE\bfseries 16S03–16S09\par}\vspace{8pt}复旦大学高等代数习题集·题目与解答\end{center}
% Original page 1.
复旦大学高等代数习题集题目与证明

题目选自谢启鸿《复旦大学高等代数习题集》（2023版），§1.5，第24–27页。

记号：\allowbreak{}\hskip.1em \RB{6.184}{4.903}{0.607}{%
\RG{0.000}{0.003}{10.9091}{mi}{110}%
}\hskip.1em\allowbreak{} 为正整数；对数域\allowbreak{}\hskip.1em \RB{6.901}{7.302}{0.500}{%
\RG{0.000}{0.003}{10.9091}{mi}{76}%
}\hskip.1em\allowbreak{} ，\allowbreak{}\hskip.1em \RB{28.823}{7.969}{2.986}{%
\RG{0.000}{0.003}{10.9091}{mi}{77}\RG{9.657}{-1.491}{7.9701}{mi}{110}\RG{14.801}{0.003}{10.9091}{rm}{40}\RG{18.672}{0.003}{10.9091}{mi}{76}\RG{25.446}{0.003}{10.9091}{rm}{41}%
}\hskip.1em\allowbreak{} 表示\allowbreak{}\hskip.1em \RB{6.901}{7.302}{0.500}{%
\RG{0.000}{0.003}{10.9091}{mi}{76}%
}\hskip.1em\allowbreak{} 上全体\allowbreak{}\hskip.1em \RB{6.184}{4.903}{0.607}{%
\RG{0.000}{0.003}{10.9091}{mi}{110}%
}\hskip.1em\allowbreak{} 阶方阵。\allowbreak{}\hskip.1em \RB{9.327}{7.302}{2.064}{%
\RG{0.000}{0.003}{10.9091}{mi}{73}\RG{4.376}{-1.491}{7.9701}{mi}{110}%
}\hskip.1em\allowbreak{} 为\allowbreak{}\hskip.1em \RB{6.184}{4.903}{0.607}{%
\RG{0.000}{0.003}{10.9091}{mi}{110}%
}\hskip.1em\allowbreak{} 阶单位矩阵，\allowbreak{}\hskip.1em \RB{7.866}{7.521}{0.716}{%
\RG{0.000}{0.003}{10.9091}{mi}{79}%
}\hskip.1em\allowbreak{} 为相应阶的零矩阵；\allowbreak{}\hskip.1em \RB{13.547}{9.112}{0.500}{%
\RG{0.000}{0.003}{10.9091}{mi}{66}\RG{8.050}{3.615}{7.9701}{rm}{84}%
}\hskip.1em\allowbreak{} 、\allowbreak{}\hskip.1em \RB{23.510}{7.411}{0.607}{%
\RG{0.000}{0.003}{10.9091}{rm}{100}\RG{5.535}{0.003}{10.9091}{rm}{101}\RG{9.955}{0.003}{10.9091}{rm}{116}\RG{15.484}{0.003}{10.9091}{mi}{66}%
}\hskip.1em\allowbreak{} 、\allowbreak{}\hskip.1em \RB{17.457}{7.302}{0.607}{%
\RG{0.000}{0.003}{10.9091}{rm}{116}\RG{3.872}{0.003}{10.9091}{rm}{114}\RG{9.431}{0.003}{10.9091}{mi}{66}%
}\hskip.1em\allowbreak{} 分别表示方阵\allowbreak{}\hskip.1em \RB{8.026}{7.302}{0.500}{%
\RG{0.000}{0.003}{10.9091}{mi}{66}%
}\hskip.1em\allowbreak{} 的转置、行列式与迹。

\section*{2016S03}
（1）\allowbreak{}\hskip.1em \RB{6.555}{5.309}{0.611}{%
\RG{0.000}{-0.002}{10.9091}{mi}{110}%
}\hskip.1em\allowbreak{} 元非零复系数多项式\allowbreak{}\hskip.1em \RB{59.689}{8.680}{3.240}{%
\RG{0.000}{-0.002}{10.9091}{mi}{102}\RG{6.430}{-0.002}{10.9091}{rm}{40}\RG{10.584}{-0.002}{10.9091}{mi}{120}\RG{16.654}{-1.639}{7.9701}{rm}{49}\RG{20.972}{-0.002}{10.9091}{mi}{59}\RG{25.751}{-0.002}{10.9091}{mi}{58}\RG{30.540}{-0.002}{10.9091}{mi}{58}\RG{35.318}{-0.002}{10.9091}{mi}{58}\RG{40.107}{-0.002}{10.9091}{mi}{59}\RG{44.896}{-0.002}{10.9091}{mi}{120}\RG{50.969}{-1.639}{7.9701}{mi}{110}\RG{56.112}{-0.002}{10.9091}{rm}{41}%
}\hskip.1em\allowbreak{} 的零点集

\[
\RB{215.160}{8.683}{3.237}{%
\RG{116.780}{4.505}{7.9701}{mi}{110}\RG{0.000}{0.001}{10.9091}{mi}{90}\RG{8.095}{0.001}{10.9091}{rm}{40}\RG{12.249}{0.001}{10.9091}{mi}{102}\RG{18.679}{0.001}{10.9091}{rm}{41}\RG{25.857}{0.001}{10.9091}{rm}{61}\RG{37.202}{0.001}{10.9091}{sy}{102}\RG{42.656}{0.001}{10.9091}{rm}{40}\RG{46.811}{0.001}{10.9091}{mi}{97}\RG{56.736}{0.001}{10.9091}{mi}{59}\RG{61.515}{0.001}{10.9091}{mi}{58}\RG{66.304}{0.001}{10.9091}{mi}{58}\RG{71.082}{0.001}{10.9091}{mi}{58}\RG{75.871}{0.001}{10.9091}{mi}{59}\RG{80.660}{0.001}{10.9091}{mi}{97}\RG{91.414}{0.001}{10.9091}{rm}{41}\RG{98.598}{0.001}{10.9091}{sy}{50}\RG{108.901}{0.001}{10.9091}{bb}{67}\RG{124.953}{0.001}{10.9091}{rm}{58}\RG{130.951}{0.001}{10.9091}{mi}{102}\RG{137.381}{0.001}{10.9091}{rm}{40}\RG{141.535}{0.001}{10.9091}{mi}{97}\RG{151.461}{0.001}{10.9091}{mi}{59}\RG{156.239}{0.001}{10.9091}{mi}{58}\RG{161.028}{0.001}{10.9091}{mi}{58}\RG{165.806}{0.001}{10.9091}{mi}{58}\RG{170.595}{0.001}{10.9091}{mi}{59}\RG{175.385}{0.001}{10.9091}{mi}{97}\RG{186.138}{0.001}{10.9091}{rm}{41}\RG{193.316}{0.001}{10.9091}{rm}{61}\RG{204.662}{0.001}{10.9091}{rm}{48}\RG{210.002}{0.001}{10.9091}{sy}{103}\RG{52.417}{-1.636}{7.9701}{rm}{49}\RG{86.270}{-1.636}{7.9701}{mi}{110}\RG{147.143}{-1.636}{7.9701}{rm}{49}\RG{180.995}{-1.636}{7.9701}{mi}{110}%
}
\]
称为\allowbreak{}\hskip.1em \RB{12.831}{8.183}{0.705}{%
\RG{0.000}{0.003}{10.9091}{bb}{67}\RG{7.880}{3.962}{7.9701}{mi}{110}%
}\hskip.1em\allowbreak{} 中的一个超曲面。将线性同构\allowbreak{}\hskip.1em \RB{62.300}{10.656}{3.235}{%
\RG{0.000}{0.003}{10.9091}{mi}{77}\RG{10.334}{-1.633}{7.9701}{mi}{110}\RG{15.477}{0.003}{10.9091}{rm}{40}\RG{19.631}{0.003}{10.9091}{bb}{67}\RG{27.510}{0.003}{10.9091}{rm}{41}\RG{34.738}{3.201}{10.9091}{sy}{24}\RG{34.827}{-0.452}{10.9091}{rm}{61}\RG{46.298}{0.003}{10.9091}{bb}{67}\RG{54.177}{3.962}{7.9701}{mi}{110}\RG{58.865}{6.529}{5.9776}{rm}{50}%
}\hskip.1em\allowbreak{} 看作等同，证明：所有不可对角化的\allowbreak{}\hskip.1em \RB{6.555}{5.315}{0.605}{%
\RG{0.000}{0.004}{10.9091}{mi}{110}%
}\hskip.1em\allowbreak{} 阶复矩阵包含在\allowbreak{}\hskip.1em \RB{16.003}{10.656}{0.704}{%
\RG{0.000}{0.004}{10.9091}{bb}{67}\RG{7.879}{3.963}{7.9701}{mi}{110}\RG{12.568}{6.529}{5.9776}{rm}{50}%
}\hskip.1em\allowbreak{} 的一个超曲面中。（2）设\allowbreak{}\hskip.1em \RB{43.047}{8.679}{3.623}{%
\RG{0.000}{-0.003}{10.9091}{mi}{65}\RG{11.390}{-0.003}{10.9091}{rm}{61}\RG{23.074}{-0.003}{10.9091}{rm}{40}\RG{27.234}{-0.003}{10.9091}{mi}{97}\RG{32.841}{-1.639}{7.9701}{mi}{105}\RG{35.474}{-1.639}{7.9701}{mi}{106}\RG{39.471}{-0.003}{10.9091}{rm}{41}%
}\hskip.1em\allowbreak{} 为\allowbreak{}\hskip.1em \RB{6.555}{5.308}{0.612}{%
\RG{0.000}{-0.003}{10.9091}{mi}{110}%
}\hskip.1em\allowbreak{} 阶复矩阵。证明：存在矩阵\allowbreak{}\hskip.1em \RB{67.380}{8.679}{3.623}{%
\RG{0.000}{-0.003}{10.9091}{mi}{65}\RG{8.007}{-0.003}{10.9091}{rm}{40}\RG{12.161}{-0.003}{10.9091}{mi}{116}\RG{16.018}{-0.003}{10.9091}{rm}{41}\RG{23.557}{-0.003}{10.9091}{rm}{61}\RG{35.251}{-0.003}{10.9091}{rm}{40}\RG{39.399}{-0.003}{10.9091}{mi}{97}\RG{45.006}{-1.639}{7.9701}{mi}{105}\RG{47.639}{-1.639}{7.9701}{mi}{106}\RG{51.635}{-0.003}{10.9091}{rm}{40}\RG{55.790}{-0.003}{10.9091}{mi}{116}\RG{59.647}{-0.003}{10.9091}{rm}{41}\RG{63.804}{-0.003}{10.9091}{rm}{41}%
}\hskip.1em\allowbreak{} ，其每个元素都是\allowbreak{}\hskip.1em \RB{4.013}{7.316}{0.612}{%
\RG{0.000}{-0.003}{10.9091}{mi}{116}%
}\hskip.1em\allowbreak{}的多项式，满足\allowbreak{}\hskip.1em \RB{44.228}{8.680}{3.240}{%
\RG{0.000}{-0.002}{10.9091}{mi}{65}\RG{8.008}{-0.002}{10.9091}{rm}{40}\RG{12.165}{-0.002}{10.9091}{rm}{48}\RG{17.510}{-0.002}{10.9091}{rm}{41}\RG{24.688}{-0.002}{10.9091}{rm}{61}\RG{36.026}{-0.002}{10.9091}{mi}{65}%
}\hskip.1em\allowbreak{} ，且当\allowbreak{}\hskip.1em \RB{23.613}{7.753}{0.807}{%
\RG{0.000}{-0.002}{10.9091}{mi}{116}\RG{6.884}{-0.002}{10.9091}{mi}{62}\RG{18.226}{-0.002}{10.9091}{rm}{48}%
}\hskip.1em\allowbreak{} 充分小时，\allowbreak{}\hskip.1em \RB{19.596}{8.680}{3.240}{%
\RG{0.000}{-0.002}{10.9091}{mi}{65}\RG{8.008}{-0.002}{10.9091}{rm}{40}\RG{12.162}{-0.002}{10.9091}{mi}{116}\RG{16.019}{-0.002}{10.9091}{rm}{41}%
}\hskip.1em\allowbreak{} 均可对角化。

证明：（1）记\allowbreak{}\hskip.1em \RB{9.662}{7.954}{2.206}{%
\RG{0.000}{0.003}{10.9091}{mi}{73}\RG{4.711}{-1.633}{7.9701}{mi}{110}%
}\hskip.1em\allowbreak{} 为\allowbreak{}\hskip.1em \RB{6.555}{5.314}{0.606}{%
\RG{0.000}{0.003}{10.9091}{mi}{110}%
}\hskip.1em\allowbreak{} 阶单位矩阵。对任意\allowbreak{}\hskip.1em \RB{99.081}{8.685}{3.617}{%
\RG{0.000}{0.003}{10.9091}{mi}{66}\RG{14.733}{0.003}{10.9091}{rm}{61}\RG{29.111}{0.003}{10.9091}{rm}{40}\RG{33.262}{0.003}{10.9091}{mi}{98}\RG{37.804}{-1.633}{7.9701}{mi}{105}\RG{40.436}{-1.633}{7.9701}{mi}{106}\RG{44.433}{0.003}{10.9091}{rm}{41}\RG{54.655}{0.003}{10.9091}{sy}{50}\RG{67.995}{0.003}{10.9091}{mi}{77}\RG{78.329}{-1.633}{7.9701}{mi}{110}\RG{83.472}{0.003}{10.9091}{rm}{40}\RG{87.626}{0.003}{10.9091}{bb}{67}\RG{95.504}{0.003}{10.9091}{rm}{41}%
}\hskip.1em\allowbreak{} ，设其特征多项式\allowbreak{}\hskip.1em \RB{99.996}{8.686}{3.235}{%
\RG{0.000}{0.004}{10.9091}{mi}{112}\RG{5.361}{-1.633}{7.9701}{mi}{66}\RG{12.003}{0.004}{10.9091}{rm}{40}\RG{16.157}{0.004}{10.9091}{mi}{122}\RG{21.606}{0.004}{10.9091}{rm}{41}\RG{28.784}{0.004}{10.9091}{rm}{61}\RG{40.130}{0.004}{10.9091}{rm}{100}\RG{46.064}{0.004}{10.9091}{rm}{101}\RG{50.810}{0.004}{10.9091}{rm}{116}\RG{54.966}{0.004}{10.9091}{rm}{40}\RG{59.118}{0.004}{10.9091}{mi}{122}\RG{64.562}{0.004}{10.9091}{mi}{73}\RG{69.277}{-1.633}{7.9701}{mi}{110}\RG{76.845}{0.004}{10.9091}{sy}{0}\RG{87.755}{0.004}{10.9091}{mi}{66}\RG{96.420}{0.004}{10.9091}{rm}{41}%
}\hskip.1em\allowbreak{} 的全部根为\allowbreak{}\hskip.1em \RB{45.658}{8.075}{2.602}{%
\RG{0.000}{0.004}{10.9091}{mi}{21}\RG{6.232}{-1.633}{7.9701}{rm}{49}\RG{10.550}{0.004}{10.9091}{mi}{59}\RG{15.328}{0.004}{10.9091}{mi}{58}\RG{20.117}{0.004}{10.9091}{mi}{58}\RG{24.895}{0.004}{10.9091}{mi}{58}\RG{29.684}{0.004}{10.9091}{mi}{59}\RG{34.474}{0.004}{10.9091}{mi}{21}\RG{40.707}{-1.633}{7.9701}{mi}{110}%
}\hskip.1em\allowbreak{} ，按代数重数计。定义

\[
\RB{121.135}{9.838}{12.910}{%
\RG{50.815}{9.089}{9.0910}{ex}{89}\RG{114.190}{4.501}{7.9701}{rm}{50}\RG{0.000}{-0.003}{10.9091}{rm}{1}\RG{8.902}{-0.003}{10.9091}{rm}{40}\RG{13.055}{-0.003}{10.9091}{mi}{66}\RG{21.721}{-0.003}{10.9091}{rm}{41}\RG{28.900}{-0.003}{10.9091}{rm}{61}\RG{73.002}{-0.003}{10.9091}{rm}{40}\RG{77.157}{-0.003}{10.9091}{mi}{21}\RG{88.897}{-0.003}{10.9091}{sy}{0}\RG{99.806}{-0.003}{10.9091}{mi}{21}\RG{110.036}{-0.003}{10.9091}{rm}{41}\RG{118.508}{-0.003}{10.9091}{mi}{59}\RG{83.387}{-1.639}{7.9701}{mi}{105}\RG{106.038}{-1.639}{7.9701}{mi}{106}\RG{40.244}{-10.926}{7.9701}{rm}{49}\RG{44.107}{-10.926}{7.9701}{sy}{20}\RG{50.118}{-10.926}{7.9701}{mi}{105}\RG{52.751}{-10.926}{7.9701}{mi}{60}\RG{58.758}{-10.926}{7.9701}{mi}{106}\RG{62.304}{-10.926}{7.9701}{sy}{20}\RG{68.313}{-10.926}{7.9701}{mi}{110}%
}
\]
其中\allowbreak{}\hskip.1em \RB{25.714}{7.752}{0.612}{%
\RG{0.000}{-0.003}{10.9091}{mi}{110}\RG{9.406}{-0.003}{10.9091}{rm}{61}\RG{20.741}{-0.003}{10.9091}{rm}{49}%
}\hskip.1em\allowbreak{} 时的空乘积规定为\allowbreak{}\hskip.1em \RB{4.973}{7.752}{0.503}{%
\RG{0.000}{-0.003}{10.9091}{rm}{49}%
}\hskip.1em\allowbreak{} 。右端是这些根的对称多项式，由对称多项式基本定理，它是\allowbreak{}\hskip.1em \RB{11.644}{5.309}{2.618}{%
\RG{0.000}{-0.002}{10.9091}{mi}{112}\RG{5.362}{-1.638}{7.9701}{mi}{66}%
}\hskip.1em\allowbreak{} 各系数的多项式。行列式展开表明，\allowbreak{}\hskip.1em \RB{11.643}{5.309}{2.618}{%
\RG{0.000}{-0.002}{10.9091}{mi}{112}\RG{5.361}{-1.638}{7.9701}{mi}{66}%
}\hskip.1em\allowbreak{} 的每个系数都是各\allowbreak{}\hskip.1em \RB{10.678}{8.069}{3.622}{%
\RG{0.000}{-0.002}{10.9091}{mi}{98}\RG{4.542}{-1.638}{7.9701}{mi}{105}\RG{7.174}{-1.638}{7.9701}{mi}{106}%
}\hskip.1em\allowbreak{} 的多项式。因此存在

\[
\RB{110.831}{8.683}{3.619}{%
\RG{0.000}{0.001}{10.9091}{mi}{70}\RG{11.429}{0.001}{10.9091}{sy}{50}\RG{21.732}{0.001}{10.9091}{bb}{67}\RG{29.611}{0.001}{10.9091}{rm}{91}\RG{32.641}{0.001}{10.9091}{mi}{120}\RG{48.371}{0.001}{10.9091}{rm}{58}\RG{54.361}{0.001}{10.9091}{rm}{49}\RG{62.740}{0.001}{10.9091}{sy}{20}\RG{74.256}{0.001}{10.9091}{mi}{105}\RG{77.899}{0.001}{10.9091}{mi}{59}\RG{82.678}{0.001}{10.9091}{mi}{106}\RG{90.748}{0.001}{10.9091}{sy}{20}\RG{102.263}{0.001}{10.9091}{mi}{110}\RG{108.640}{0.001}{10.9091}{rm}{93}\RG{38.712}{-1.635}{7.9701}{mi}{105}\RG{41.344}{-1.635}{7.9701}{mi}{106}%
}
\]
使\allowbreak{}\hskip.1em \RB{105.202}{8.686}{3.617}{%
\RG{0.000}{0.004}{10.9091}{rm}{1}\RG{8.902}{0.004}{10.9091}{rm}{40}\RG{13.056}{0.004}{10.9091}{mi}{66}\RG{21.721}{0.004}{10.9091}{rm}{41}\RG{28.900}{0.004}{10.9091}{rm}{61}\RG{40.245}{0.004}{10.9091}{mi}{70}\RG{48.643}{0.004}{10.9091}{rm}{40}\RG{52.800}{0.004}{10.9091}{rm}{40}\RG{56.951}{0.004}{10.9091}{mi}{98}\RG{61.492}{-1.633}{7.9701}{mi}{105}\RG{64.125}{-1.633}{7.9701}{mi}{106}\RG{68.123}{0.004}{10.9091}{rm}{41}\RG{72.276}{-1.633}{7.9701}{rm}{49}\RG{76.140}{-1.633}{7.9701}{sy}{20}\RG{82.151}{-1.633}{7.9701}{mi}{105}\RG{84.784}{-1.633}{7.9701}{mi}{59}\RG{86.929}{-1.633}{7.9701}{mi}{106}\RG{90.472}{-1.633}{7.9701}{sy}{20}\RG{96.483}{-1.633}{7.9701}{mi}{110}\RG{101.626}{0.004}{10.9091}{rm}{41}%
}\hskip.1em\allowbreak{} 。

取\allowbreak{}\hskip.1em \RB{96.763}{8.679}{3.241}{%
\RG{0.000}{-0.003}{10.9091}{mi}{68}\RG{12.196}{-0.003}{10.9091}{rm}{61}\RG{23.530}{-0.003}{10.9091}{rm}{100}\RG{29.465}{-0.003}{10.9091}{rm}{105}\RG{32.432}{-0.003}{10.9091}{rm}{97}\RG{37.778}{-0.003}{10.9091}{rm}{103}\RG{43.265}{-0.003}{10.9091}{rm}{40}\RG{47.422}{-0.003}{10.9091}{rm}{49}\RG{52.757}{-0.003}{10.9091}{mi}{59}\RG{57.542}{-0.003}{10.9091}{rm}{50}\RG{62.884}{-0.003}{10.9091}{mi}{59}\RG{67.662}{-0.003}{10.9091}{mi}{58}\RG{72.451}{-0.003}{10.9091}{mi}{58}\RG{77.230}{-0.003}{10.9091}{mi}{58}\RG{82.019}{-0.003}{10.9091}{mi}{59}\RG{86.808}{-0.003}{10.9091}{mi}{110}\RG{93.187}{-0.003}{10.9091}{rm}{41}%
}\hskip.1em\allowbreak{} ，则

\[
\RB{130.357}{9.841}{12.907}{%
\RG{51.316}{9.091}{9.0910}{ex}{89}\RG{103.822}{4.504}{7.9701}{rm}{50}\RG{0.000}{0.000}{10.9091}{rm}{1}\RG{8.902}{0.000}{10.9091}{rm}{40}\RG{13.056}{0.000}{10.9091}{mi}{68}\RG{22.222}{0.000}{10.9091}{rm}{41}\RG{29.401}{0.000}{10.9091}{rm}{61}\RG{73.503}{0.000}{10.9091}{rm}{40}\RG{77.658}{0.000}{10.9091}{mi}{105}\RG{83.726}{0.000}{10.9091}{sy}{0}\RG{94.635}{0.000}{10.9091}{mi}{106}\RG{99.667}{0.000}{10.9091}{rm}{41}\RG{111.171}{0.000}{10.9091}{sy}{54}\RG{111.171}{0.000}{10.9091}{rm}{61}\RG{122.505}{0.000}{10.9091}{rm}{48}\RG{127.849}{0.000}{10.9091}{mi}{58}\RG{40.745}{-10.924}{7.9701}{rm}{49}\RG{44.608}{-10.924}{7.9701}{sy}{20}\RG{50.619}{-10.924}{7.9701}{mi}{105}\RG{53.252}{-10.924}{7.9701}{mi}{60}\RG{59.259}{-10.924}{7.9701}{mi}{106}\RG{62.805}{-10.924}{7.9701}{sy}{20}\RG{68.814}{-10.924}{7.9701}{mi}{110}%
}
\]
故\allowbreak{}\hskip.1em \RB{8.562}{7.926}{0.504}{%
\RG{0.000}{-0.004}{10.9091}{mi}{70}%
}\hskip.1em\allowbreak{} 不是零多项式。若\allowbreak{}\hskip.1em \RB{45.628}{8.678}{3.242}{%
\RG{0.000}{-0.004}{10.9091}{rm}{1}\RG{8.902}{-0.004}{10.9091}{rm}{40}\RG{13.055}{-0.004}{10.9091}{mi}{66}\RG{21.721}{-0.004}{10.9091}{rm}{41}\RG{28.906}{-0.004}{10.9091}{sy}{54}\RG{28.906}{-0.004}{10.9091}{rm}{61}\RG{40.240}{-0.004}{10.9091}{rm}{48}%
}\hskip.1em\allowbreak{} ，则\allowbreak{}\hskip.1em \RB{8.606}{7.947}{0.504}{%
\RG{0.000}{-0.004}{10.9091}{mi}{66}%
}\hskip.1em\allowbreak{} 有\allowbreak{}\hskip.1em \RB{6.555}{5.307}{0.613}{%
\RG{0.000}{-0.004}{10.9091}{mi}{110}%
}\hskip.1em\allowbreak{} 个互异的特征值；不同特征值的特征向量线性无关，因而可组成\allowbreak{}\hskip.1em \RB{12.831}{8.177}{0.710}{%
\RG{0.000}{-0.003}{10.9091}{bb}{67}\RG{7.880}{3.956}{7.9701}{mi}{110}%
}\hskip.1em\allowbreak{} 的一组基，\allowbreak{}\hskip.1em \RB{8.606}{7.948}{0.503}{%
\RG{0.000}{-0.003}{10.9091}{mi}{66}%
}\hskip.1em\allowbreak{} 可对角化。由此，所有不可对角化矩阵均属于超曲面\allowbreak{}\hskip.1em \RB{24.224}{8.680}{3.240}{%
\RG{0.000}{-0.002}{10.9091}{mi}{90}\RG{8.095}{-0.002}{10.9091}{rm}{40}\RG{12.250}{-0.002}{10.9091}{mi}{70}\RG{20.647}{-0.002}{10.9091}{rm}{41}%
}\hskip.1em\allowbreak{} 。当\allowbreak{}\hskip.1em \RB{25.714}{7.753}{0.611}{%
\RG{0.000}{-0.002}{10.9091}{mi}{110}\RG{9.406}{-0.002}{10.9091}{rm}{61}\RG{20.741}{-0.002}{10.9091}{rm}{49}%
}\hskip.1em\allowbreak{} 时，\allowbreak{}\hskip.1em \RB{27.737}{7.927}{0.502}{%
\RG{0.000}{-0.002}{10.9091}{mi}{70}\RG{11.429}{-0.002}{10.9091}{rm}{61}\RG{22.764}{-0.002}{10.9091}{rm}{49}%
}\hskip.1em\allowbreak{} ，其零点集为空集，与所有一阶矩阵均可对角化一致。（2）取上面定义的\allowbreak{}\hskip.1em \RB{9.107}{7.952}{0.500}{%
\RG{0.000}{0.001}{10.9091}{mi}{68}%
}\hskip.1em\allowbreak{} ，构造

\[
\RB{102.079}{8.686}{3.235}{%
\RG{0.000}{0.004}{10.9091}{mi}{65}\RG{8.008}{0.004}{10.9091}{rm}{40}\RG{12.162}{0.004}{10.9091}{mi}{116}\RG{16.019}{0.004}{10.9091}{rm}{41}\RG{23.197}{0.004}{10.9091}{rm}{61}\RG{34.543}{0.004}{10.9091}{rm}{40}\RG{38.700}{0.004}{10.9091}{rm}{49}\RG{46.462}{0.004}{10.9091}{sy}{0}\RG{57.371}{0.004}{10.9091}{mi}{116}\RG{61.228}{0.004}{10.9091}{rm}{41}\RG{65.382}{0.004}{10.9091}{mi}{65}\RG{75.813}{0.004}{10.9091}{rm}{43}\RG{86.546}{0.004}{10.9091}{mi}{116}\RG{90.408}{0.004}{10.9091}{mi}{68}\RG{99.572}{0.004}{10.9091}{mi}{58}%
}
\]
各矩阵元素均为\allowbreak{}\hskip.1em \RB{4.013}{7.316}{0.612}{%
\RG{0.000}{-0.003}{10.9091}{mi}{116}%
}\hskip.1em\allowbreak{} 的多项式，且\allowbreak{}\hskip.1em \RB{44.228}{8.679}{3.241}{%
\RG{0.000}{-0.003}{10.9091}{mi}{65}\RG{8.008}{-0.003}{10.9091}{rm}{40}\RG{12.165}{-0.003}{10.9091}{rm}{48}\RG{17.510}{-0.003}{10.9091}{rm}{41}\RG{24.688}{-0.003}{10.9091}{rm}{61}\RG{36.026}{-0.003}{10.9091}{mi}{65}%
}\hskip.1em\allowbreak{} 、\allowbreak{}\hskip.1em \RB{45.133}{8.679}{3.241}{%
\RG{0.000}{-0.003}{10.9091}{mi}{65}\RG{8.007}{-0.003}{10.9091}{rm}{40}\RG{12.164}{-0.003}{10.9091}{rm}{49}\RG{17.509}{-0.003}{10.9091}{rm}{41}\RG{24.688}{-0.003}{10.9091}{rm}{61}\RG{36.025}{-0.003}{10.9091}{mi}{68}%
}\hskip.1em\allowbreak{} 。令\allowbreak{}\hskip.1em \RB{68.469}{8.679}{3.241}{%
\RG{0.000}{-0.003}{10.9091}{mi}{113}\RG{5.127}{-0.003}{10.9091}{rm}{40}\RG{9.281}{-0.003}{10.9091}{mi}{116}\RG{13.138}{-0.003}{10.9091}{rm}{41}\RG{20.317}{-0.003}{10.9091}{rm}{61}\RG{31.662}{-0.003}{10.9091}{rm}{1}\RG{40.564}{-0.003}{10.9091}{rm}{40}\RG{44.718}{-0.003}{10.9091}{mi}{65}\RG{52.725}{-0.003}{10.9091}{rm}{40}\RG{56.879}{-0.003}{10.9091}{mi}{116}\RG{60.736}{-0.003}{10.9091}{rm}{41}\RG{64.892}{-0.003}{10.9091}{rm}{41}%
}\hskip.1em\allowbreak{} 。由（1），\allowbreak{}\hskip.1em \RB{35.418}{8.679}{3.230}{%
\RG{0.000}{-0.003}{10.9091}{mi}{113}\RG{8.158}{-0.003}{10.9091}{sy}{50}\RG{18.461}{-0.003}{10.9091}{bb}{67}\RG{26.340}{-0.003}{10.9091}{rm}{91}\RG{29.370}{-0.003}{10.9091}{mi}{116}\RG{33.227}{-0.003}{10.9091}{rm}{93}%
}\hskip.1em\allowbreak{} ，而\allowbreak{}\hskip.1em \RB{79.274}{8.680}{3.240}{%
\RG{0.000}{-0.002}{10.9091}{mi}{113}\RG{5.128}{-0.002}{10.9091}{rm}{40}\RG{9.285}{-0.002}{10.9091}{rm}{49}\RG{14.630}{-0.002}{10.9091}{rm}{41}\RG{21.809}{-0.002}{10.9091}{rm}{61}\RG{33.154}{-0.002}{10.9091}{rm}{1}\RG{42.056}{-0.002}{10.9091}{rm}{40}\RG{46.201}{-0.002}{10.9091}{mi}{68}\RG{55.368}{-0.002}{10.9091}{rm}{41}\RG{62.552}{-0.002}{10.9091}{sy}{54}\RG{62.552}{-0.002}{10.9091}{rm}{61}\RG{73.887}{-0.002}{10.9091}{rm}{48}%
}\hskip.1em\allowbreak{} ，故\allowbreak{}\hskip.1em \RB{5.311}{5.309}{2.618}{%
\RG{0.000}{-0.002}{10.9091}{mi}{113}%
}\hskip.1em\allowbreak{} 不是零多项式。

设\allowbreak{}\hskip.1em \RB{150.601}{8.683}{3.237}{%
\RG{0.000}{0.001}{10.9091}{mi}{69}\RG{12.542}{0.001}{10.9091}{rm}{61}\RG{24.845}{0.001}{10.9091}{sy}{102}\RG{30.300}{0.001}{10.9091}{mi}{114}\RG{39.405}{0.001}{10.9091}{sy}{50}\RG{50.674}{0.001}{10.9091}{bb}{82}\RG{62.547}{0.001}{10.9091}{rm}{58}\RG{69.511}{0.001}{10.9091}{mi}{114}\RG{78.609}{0.001}{10.9091}{mi}{62}\RG{90.919}{0.001}{10.9091}{rm}{48}\RG{96.260}{0.001}{10.9091}{mi}{59}\RG{105.249}{0.001}{10.9091}{mi}{113}\RG{110.384}{0.001}{10.9091}{rm}{40}\RG{114.539}{0.001}{10.9091}{mi}{114}\RG{119.649}{0.001}{10.9091}{rm}{41}\RG{127.798}{0.001}{10.9091}{rm}{61}\RG{140.104}{0.001}{10.9091}{rm}{48}\RG{145.443}{0.001}{10.9091}{sy}{103}%
}\hskip.1em\allowbreak{} 。非零一元多项式只有有限个根，所以\allowbreak{}\hskip.1em \RB{8.715}{7.930}{0.500}{%
\RG{0.000}{0.001}{10.9091}{mi}{69}%
}\hskip.1em\allowbreak{} 为有限集。若\allowbreak{}\hskip.1em \RB{32.447}{7.931}{0.520}{%
\RG{0.000}{0.002}{10.9091}{mi}{69}\RG{12.160}{0.002}{10.9091}{rm}{61}\RG{24.081}{0.002}{10.9091}{bb}{63}%
}\hskip.1em\allowbreak{} ，取\allowbreak{}\hskip.1em \RB{25.650}{8.258}{0.607}{%
\RG{0.000}{0.002}{10.9091}{mi}{14}\RG{8.754}{0.002}{10.9091}{rm}{61}\RG{20.678}{0.002}{10.9091}{rm}{49}%
}\hskip.1em\allowbreak{} ；若\allowbreak{}\hskip.1em \RB{32.447}{8.313}{2.855}{%
\RG{0.000}{0.002}{10.9091}{mi}{69}\RG{12.160}{0.002}{10.9091}{sy}{54}\RG{12.160}{0.002}{10.9091}{rm}{61}\RG{24.081}{0.002}{10.9091}{bb}{63}%
}\hskip.1em\allowbreak{} ，取\allowbreak{}\hskip.1em \RB{49.064}{8.258}{0.607}{%
\RG{0.000}{0.002}{10.9091}{mi}{14}\RG{8.754}{0.002}{10.9091}{rm}{61}\RG{20.678}{0.002}{10.9091}{rm}{109}\RG{29.580}{0.002}{10.9091}{rm}{105}\RG{32.547}{0.002}{10.9091}{rm}{110}\RG{40.349}{0.002}{10.9091}{mi}{69}%
}\hskip.1em\allowbreak{} 。两种情形下均有\allowbreak{}\hskip.1em \RB{26.062}{8.258}{0.804}{%
\RG{0.000}{0.002}{10.9091}{mi}{14}\RG{8.749}{0.002}{10.9091}{mi}{62}\RG{20.675}{0.002}{10.9091}{rm}{48}%
}\hskip.1em\allowbreak{} ，且对每个\allowbreak{}\hskip.1em \RB{40.188}{8.684}{3.236}{%
\RG{0.000}{0.002}{10.9091}{mi}{116}\RG{6.888}{0.002}{10.9091}{sy}{50}\RG{17.191}{0.002}{10.9091}{rm}{40}\RG{21.347}{0.002}{10.9091}{rm}{48}\RG{26.686}{0.002}{10.9091}{mi}{59}\RG{31.464}{0.002}{10.9091}{mi}{14}\RG{36.612}{0.002}{10.9091}{rm}{41}%
}\hskip.1em\allowbreak{} ，\allowbreak{}\hskip.1em \RB{37.046}{8.684}{3.236}{%
\RG{0.000}{0.002}{10.9091}{mi}{113}\RG{5.127}{0.002}{10.9091}{rm}{40}\RG{9.281}{0.002}{10.9091}{mi}{116}\RG{13.139}{0.002}{10.9091}{rm}{41}\RG{20.323}{0.002}{10.9091}{sy}{54}\RG{20.323}{0.002}{10.9091}{rm}{61}\RG{31.658}{0.002}{10.9091}{rm}{48}%
}\hskip.1em\allowbreak{} 。于是\allowbreak{}\hskip.1em \RB{19.596}{8.684}{3.236}{%
\RG{0.000}{0.002}{10.9091}{mi}{65}\RG{8.007}{0.002}{10.9091}{rm}{40}\RG{12.162}{0.002}{10.9091}{mi}{116}\RG{16.019}{0.002}{10.9091}{rm}{41}%
}\hskip.1em\allowbreak{} 有\allowbreak{}\hskip.1em \RB{6.555}{5.313}{0.607}{%
\RG{0.000}{0.002}{10.9091}{mi}{110}%
}\hskip.1em\allowbreak{} 个互异特征值，从而可对角化。\allowbreak{}\hskip.1em \RB{8.366}{7.996}{0.500}{%
\RG{0.000}{0.002}{10.9091}{am}{3}%
}\hskip.1em\allowbreak{}

% Original page 2.
\section*{2016S04}
设\allowbreak{}\hskip.1em \RB{6.555}{5.313}{0.607}{%
\RG{0.000}{0.002}{10.9091}{mi}{110}%
}\hskip.1em\allowbreak{} 阶方阵\allowbreak{}\hskip.1em \RB{8.202}{8.291}{0.500}{%
\RG{0.000}{0.002}{10.9091}{mi}{65}%
}\hskip.1em\allowbreak{} 满足\allowbreak{}\hskip.1em \RB{45.534}{8.684}{3.236}{%
\RG{0.000}{0.002}{10.9091}{mi}{102}\RG{6.430}{0.002}{10.9091}{rm}{40}\RG{10.583}{0.002}{10.9091}{mi}{65}\RG{18.592}{0.002}{10.9091}{rm}{41}\RG{25.770}{0.002}{10.9091}{rm}{61}\RG{37.114}{0.002}{10.9091}{mi}{79}%
}\hskip.1em\allowbreak{} ，其中

\[
\RB{294.433}{17.153}{11.523}{%
\RG{258.645}{11.816}{7.9701}{mi}{109}\RG{265.480}{11.816}{7.9701}{sy}{0}\RG{271.491}{11.816}{7.9701}{rm}{49}\RG{260.434}{9.089}{9.0910}{ex}{88}\RG{54.144}{4.502}{7.9701}{mi}{109}\RG{103.431}{4.502}{7.9701}{mi}{109}\RG{110.267}{4.502}{7.9701}{sy}{0}\RG{116.277}{4.502}{7.9701}{rm}{49}\RG{0.000}{-0.002}{10.9091}{mi}{102}\RG{6.430}{-0.002}{10.9091}{rm}{40}\RG{10.584}{-0.002}{10.9091}{mi}{120}\RG{16.654}{-0.002}{10.9091}{rm}{41}\RG{23.833}{-0.002}{10.9091}{rm}{61}\RG{35.176}{-0.002}{10.9091}{mi}{97}\RG{48.074}{-0.002}{10.9091}{mi}{120}\RG{63.858}{-0.002}{10.9091}{rm}{43}\RG{74.591}{-0.002}{10.9091}{mi}{97}\RG{97.361}{-0.002}{10.9091}{mi}{120}\RG{123.019}{-0.002}{10.9091}{rm}{43}\RG{133.752}{-0.002}{10.9091}{sy}{1}\RG{138.596}{-0.002}{10.9091}{sy}{1}\RG{143.451}{-0.002}{10.9091}{sy}{1}\RG{148.904}{-0.002}{10.9091}{rm}{43}\RG{159.637}{-0.002}{10.9091}{mi}{97}\RG{169.562}{-0.002}{10.9091}{mi}{120}\RG{178.057}{-0.002}{10.9091}{rm}{43}\RG{188.789}{-0.002}{10.9091}{mi}{97}\RG{198.714}{-0.002}{10.9091}{mi}{59}\RG{225.319}{-0.002}{10.9091}{sy}{106}\RG{228.349}{-0.002}{10.9091}{mi}{97}\RG{241.246}{-0.002}{10.9091}{sy}{106}\RG{247.306}{-0.002}{10.9091}{mi}{62}\RG{277.172}{-0.002}{10.9091}{sy}{106}\RG{280.203}{-0.002}{10.9091}{mi}{97}\RG{288.895}{-0.002}{10.9091}{sy}{106}\RG{291.926}{-0.002}{10.9091}{mi}{58}\RG{40.784}{-1.638}{7.9701}{mi}{109}\RG{80.198}{-1.638}{7.9701}{mi}{109}\RG{87.033}{-1.638}{7.9701}{sy}{0}\RG{93.043}{-1.638}{7.9701}{rm}{49}\RG{165.244}{-1.638}{7.9701}{rm}{49}\RG{194.396}{-1.638}{7.9701}{rm}{48}\RG{233.956}{-1.638}{7.9701}{mi}{109}\RG{285.810}{-1.639}{7.9701}{mi}{105}\RG{260.747}{-10.871}{7.9701}{mi}{105}\RG{263.379}{-10.871}{7.9701}{rm}{61}\RG{269.386}{-10.871}{7.9701}{rm}{48}%
}
\]
证明：矩阵方程\allowbreak{}\hskip.1em \RB{82.019}{9.290}{1.345}{%
\RG{0.000}{-0.005}{10.9091}{rm}{50}\RG{5.341}{-0.005}{10.9091}{mi}{88}\RG{17.487}{-0.005}{10.9091}{rm}{43}\RG{28.220}{-0.005}{10.9091}{mi}{65}\RG{36.228}{-0.005}{10.9091}{mi}{88}\RG{48.981}{-0.005}{10.9091}{rm}{61}\RG{60.320}{-0.005}{10.9091}{mi}{88}\RG{70.040}{-0.005}{10.9091}{mi}{65}\RG{78.050}{3.954}{7.9701}{rm}{50}%
}\hskip.1em\allowbreak{} 只有零解：

证明：以下在复数域上证明，实数情形亦包含在内。系数条件给出\allowbreak{}\hskip.1em \RB{34.238}{8.311}{2.856}{%
\RG{0.000}{0.000}{10.9091}{mi}{97}\RG{5.607}{-1.636}{7.9701}{mi}{109}\RG{16.719}{0.000}{10.9091}{sy}{54}\RG{16.719}{0.000}{10.9091}{rm}{61}\RG{28.850}{0.000}{10.9091}{rm}{48}%
}\hskip.1em\allowbreak{} ；非零常数多项式不能零化正阶矩阵，所以\allowbreak{}\hskip.1em \RB{30.065}{7.756}{1.994}{%
\RG{0.000}{0.001}{10.9091}{mi}{109}\RG{12.975}{0.001}{10.9091}{sy}{21}\RG{25.092}{0.001}{10.9091}{rm}{49}%
}\hskip.1em\allowbreak{} 。任取\allowbreak{}\hskip.1em \RB{8.202}{8.290}{0.500}{%
\RG{0.000}{0.001}{10.9091}{mi}{65}%
}\hskip.1em\allowbreak{} 的特征值\allowbreak{}\hskip.1em \RB{6.336}{8.072}{0.619}{%
\RG{0.000}{0.001}{10.9091}{mi}{21}%
}\hskip.1em\allowbreak{} 及相应的非零特征向量\allowbreak{}\hskip.1em \RB{5.486}{5.312}{0.608}{%
\RG{0.000}{0.001}{10.9091}{mi}{118}%
}\hskip.1em\allowbreak{} 。由\allowbreak{}\hskip.1em \RB{39.647}{8.291}{0.618}{%
\RG{0.000}{0.002}{10.9091}{mi}{65}\RG{8.007}{0.002}{10.9091}{mi}{118}\RG{16.594}{0.002}{10.9091}{rm}{61}\RG{27.932}{0.002}{10.9091}{mi}{21}\RG{34.161}{0.002}{10.9091}{mi}{118}%
}\hskip.1em\allowbreak{} 得\allowbreak{}\hskip.1em \RB{69.125}{8.684}{3.236}{%
\RG{0.000}{0.002}{10.9091}{mi}{102}\RG{6.430}{0.002}{10.9091}{rm}{40}\RG{10.584}{0.002}{10.9091}{mi}{65}\RG{18.592}{0.002}{10.9091}{rm}{41}\RG{22.746}{0.002}{10.9091}{mi}{118}\RG{31.332}{0.002}{10.9091}{rm}{61}\RG{42.671}{0.002}{10.9091}{mi}{102}\RG{49.100}{0.002}{10.9091}{rm}{40}\RG{53.254}{0.002}{10.9091}{mi}{21}\RG{59.486}{0.002}{10.9091}{rm}{41}\RG{63.640}{0.002}{10.9091}{mi}{118}%
}\hskip.1em\allowbreak{} ，再由\allowbreak{}\hskip.1em \RB{45.535}{8.684}{3.236}{%
\RG{0.000}{0.002}{10.9091}{mi}{102}\RG{6.430}{0.002}{10.9091}{rm}{40}\RG{10.583}{0.002}{10.9091}{mi}{65}\RG{18.592}{0.002}{10.9091}{rm}{41}\RG{25.770}{0.002}{10.9091}{rm}{61}\RG{37.115}{0.002}{10.9091}{mi}{79}%
}\hskip.1em\allowbreak{} 得\allowbreak{}\hskip.1em \RB{40.726}{8.684}{3.236}{%
\RG{0.000}{0.002}{10.9091}{mi}{102}\RG{6.430}{0.002}{10.9091}{rm}{40}\RG{10.584}{0.002}{10.9091}{mi}{21}\RG{16.815}{0.002}{10.9091}{rm}{41}\RG{23.993}{0.002}{10.9091}{rm}{61}\RG{35.339}{0.002}{10.9091}{rm}{48}%
}\hskip.1em\allowbreak{} 。若\allowbreak{}\hskip.1em \RB{31.810}{8.684}{3.225}{%
\RG{0.000}{0.002}{10.9091}{sy}{106}\RG{3.030}{0.002}{10.9091}{mi}{21}\RG{9.261}{0.002}{10.9091}{sy}{106}\RG{15.316}{0.002}{10.9091}{sy}{21}\RG{26.837}{0.002}{10.9091}{rm}{49}%
}\hskip.1em\allowbreak{} ，则

\[
\RB{108.788}{17.153}{94.636}{%
\RG{54.726}{16.362}{9.0910}{ex}{12}\RG{94.294}{16.362}{9.0910}{ex}{12}\RG{57.757}{11.816}{7.9701}{mi}{109}\RG{64.592}{11.816}{7.9701}{sy}{0}\RG{70.602}{11.816}{7.9701}{rm}{49}\RG{54.726}{10.907}{9.0910}{ex}{12}\RG{94.294}{10.907}{9.0910}{ex}{12}\RG{59.545}{9.089}{9.0910}{ex}{88}\RG{54.726}{5.452}{9.0910}{ex}{12}\RG{94.294}{5.452}{9.0910}{ex}{12}\RG{33.068}{4.502}{7.9701}{mi}{109}\RG{91.207}{4.502}{7.9701}{mi}{105}\RG{0.000}{-0.002}{10.9091}{sy}{106}\RG{3.030}{-0.002}{10.9091}{mi}{97}\RG{15.927}{-0.002}{10.9091}{sy}{106}\RG{20.771}{-0.002}{10.9091}{sy}{106}\RG{23.807}{-0.002}{10.9091}{mi}{21}\RG{30.037}{-0.002}{10.9091}{sy}{106}\RG{43.388}{-0.002}{10.9091}{rm}{61}\RG{54.726}{-0.002}{9.0910}{ex}{12}\RG{76.284}{-0.002}{10.9091}{mi}{97}\RG{84.977}{-0.002}{10.9091}{mi}{21}\RG{94.294}{-0.002}{9.0910}{ex}{12}\RG{8.638}{-1.639}{7.9701}{mi}{109}\RG{81.892}{-1.639}{7.9701}{mi}{105}\RG{54.726}{-5.457}{9.0910}{ex}{12}\RG{94.294}{-5.457}{9.0910}{ex}{12}\RG{59.859}{-10.871}{7.9701}{mi}{105}\RG{62.490}{-10.871}{7.9701}{rm}{61}\RG{68.497}{-10.871}{7.9701}{rm}{48}\RG{54.903}{-21.012}{7.9701}{mi}{109}\RG{61.738}{-21.012}{7.9701}{sy}{0}\RG{67.749}{-21.012}{7.9701}{rm}{49}\RG{56.692}{-23.740}{9.0910}{ex}{88}\RG{102.294}{-28.326}{7.9701}{mi}{105}\RG{43.388}{-32.830}{10.9091}{sy}{20}\RG{73.431}{-32.830}{10.9091}{sy}{106}\RG{76.460}{-32.830}{10.9091}{mi}{97}\RG{85.154}{-32.830}{10.9091}{sy}{106}\RG{89.998}{-32.830}{10.9091}{sy}{106}\RG{93.033}{-32.830}{10.9091}{mi}{21}\RG{99.263}{-32.830}{10.9091}{sy}{106}\RG{82.068}{-34.467}{7.9701}{mi}{105}\RG{57.006}{-43.699}{7.9701}{mi}{105}\RG{59.636}{-43.699}{7.9701}{rm}{61}\RG{65.644}{-43.699}{7.9701}{rm}{48}\RG{76.303}{-53.841}{7.9701}{mi}{109}\RG{83.138}{-53.841}{7.9701}{sy}{0}\RG{89.148}{-53.841}{7.9701}{rm}{49}\RG{78.092}{-56.568}{9.0910}{ex}{88}\RG{67.195}{-61.155}{7.9701}{mi}{109}\RG{43.388}{-65.660}{10.9091}{sy}{20}\RG{54.897}{-65.660}{10.9091}{sy}{106}\RG{57.934}{-65.660}{10.9091}{mi}{21}\RG{64.164}{-65.660}{10.9091}{sy}{106}\RG{94.830}{-65.660}{10.9091}{sy}{106}\RG{97.861}{-65.660}{10.9091}{mi}{97}\RG{106.553}{-65.660}{10.9091}{sy}{106}\RG{103.468}{-67.296}{7.9701}{mi}{105}\RG{78.405}{-76.528}{7.9701}{mi}{105}\RG{81.036}{-76.528}{7.9701}{rm}{61}\RG{87.043}{-76.528}{7.9701}{rm}{48}\RG{87.794}{-86.904}{7.9701}{mi}{109}\RG{43.388}{-91.408}{10.9091}{mi}{60}\RG{54.726}{-91.408}{10.9091}{sy}{106}\RG{57.757}{-91.408}{10.9091}{mi}{97}\RG{70.654}{-91.408}{10.9091}{sy}{106}\RG{75.498}{-91.408}{10.9091}{sy}{106}\RG{78.533}{-91.408}{10.9091}{mi}{21}\RG{84.763}{-91.408}{10.9091}{sy}{106}\RG{95.084}{-91.408}{10.9091}{mi}{59}\RG{63.364}{-93.044}{7.9701}{mi}{109}%
}
\]
产生矛盾。因此，\allowbreak{}\hskip.1em \RB{8.202}{8.293}{0.500}{%
\RG{0.000}{0.004}{10.9091}{mi}{65}%
}\hskip.1em\allowbreak{} 的每个特征值均满足\allowbreak{}\hskip.1em \RB{31.633}{8.686}{3.224}{%
\RG{0.000}{0.004}{10.9091}{sy}{106}\RG{3.030}{0.004}{10.9091}{mi}{21}\RG{9.261}{0.004}{10.9091}{sy}{106}\RG{15.322}{0.004}{10.9091}{mi}{60}\RG{26.660}{0.004}{10.9091}{rm}{49}%
}\hskip.1em\allowbreak{} 。

由复矩阵的上三角化定理，存在可逆矩阵\allowbreak{}\hskip.1em \RB{8.573}{7.948}{0.503}{%
\RG{0.000}{-0.003}{10.9091}{mi}{80}%
}\hskip.1em\allowbreak{} ，使\allowbreak{}\hskip.1em \RB{92.674}{9.292}{3.623}{%
\RG{0.000}{-0.003}{10.9091}{mi}{84}\RG{11.499}{-0.003}{10.9091}{rm}{61}\RG{23.564}{-0.003}{10.9091}{mi}{80}\RG{31.932}{3.956}{7.9701}{sy}{0}\RG{37.943}{3.956}{7.9701}{rm}{49}\RG{42.261}{-0.003}{10.9091}{mi}{65}\RG{50.269}{-0.003}{10.9091}{mi}{80}\RG{62.393}{-0.003}{10.9091}{rm}{61}\RG{74.459}{-0.003}{10.9091}{rm}{40}\RG{78.610}{-0.003}{10.9091}{mi}{116}\RG{82.468}{-1.639}{7.9701}{mi}{105}\RG{85.101}{-1.639}{7.9701}{mi}{106}\RG{89.098}{-0.003}{10.9091}{rm}{41}%
}\hskip.1em\allowbreak{} 为上三角矩阵。记\allowbreak{}\hskip.1em \RB{89.072}{8.680}{3.240}{%
\RG{0.000}{-0.002}{10.9091}{mi}{21}\RG{6.232}{-1.639}{7.9701}{mi}{105}\RG{12.347}{-0.002}{10.9091}{rm}{61}\RG{23.685}{-0.002}{10.9091}{mi}{116}\RG{27.543}{-1.639}{7.9701}{mi}{105}\RG{30.176}{-1.639}{7.9701}{mi}{105}\RG{36.889}{-0.002}{10.9091}{rm}{40}\RG{41.046}{-0.002}{10.9091}{rm}{49}\RG{49.414}{-0.002}{10.9091}{sy}{20}\RG{60.929}{-0.002}{10.9091}{mi}{105}\RG{67.604}{-0.002}{10.9091}{sy}{20}\RG{79.119}{-0.002}{10.9091}{mi}{110}\RG{85.495}{-0.002}{10.9091}{rm}{41}%
}\hskip.1em\allowbreak{} ，则这些对角元是\allowbreak{}\hskip.1em \RB{8.202}{8.287}{0.502}{%
\RG{0.000}{-0.002}{10.9091}{mi}{65}%
}\hskip.1em\allowbreak{} 的全部特征值，故\allowbreak{}\hskip.1em \RB{34.718}{8.680}{3.229}{%
\RG{0.000}{-0.002}{10.9091}{sy}{106}\RG{3.030}{-0.002}{10.9091}{mi}{21}\RG{9.261}{-1.639}{7.9701}{mi}{105}\RG{12.346}{-0.002}{10.9091}{sy}{106}\RG{18.407}{-0.002}{10.9091}{mi}{60}\RG{29.745}{-0.002}{10.9091}{rm}{49}%
}\hskip.1em\allowbreak{} 。设\allowbreak{}\hskip.1em \RB{9.642}{7.949}{0.502}{%
\RG{0.000}{-0.002}{10.9091}{mi}{88}%
}\hskip.1em\allowbreak{} 满足原方程，令\allowbreak{}\hskip.1em \RB{93.666}{9.294}{3.622}{%
\RG{0.000}{-0.002}{10.9091}{mi}{89}\RG{11.602}{-0.002}{10.9091}{rm}{61}\RG{22.940}{-0.002}{10.9091}{mi}{80}\RG{31.308}{3.958}{7.9701}{sy}{0}\RG{37.319}{3.958}{7.9701}{rm}{49}\RG{41.637}{-0.002}{10.9091}{mi}{88}\RG{51.357}{-0.002}{10.9091}{mi}{80}\RG{62.759}{-0.002}{10.9091}{rm}{61}\RG{74.094}{-0.002}{10.9091}{rm}{40}\RG{78.252}{-0.002}{10.9091}{mi}{121}\RG{83.460}{-1.638}{7.9701}{mi}{105}\RG{86.093}{-1.638}{7.9701}{mi}{106}\RG{90.090}{-0.002}{10.9091}{rm}{41}%
}\hskip.1em\allowbreak{} ，并记\allowbreak{}\hskip.1em \RB{46.871}{9.294}{3.622}{%
\RG{0.000}{-0.002}{10.9091}{mi}{84}\RG{7.744}{3.958}{7.9701}{rm}{50}\RG{15.093}{-0.002}{10.9091}{rm}{61}\RG{26.428}{-0.002}{10.9091}{rm}{40}\RG{30.586}{-0.002}{10.9091}{mi}{117}\RG{36.665}{-1.638}{7.9701}{mi}{105}\RG{39.298}{-1.638}{7.9701}{mi}{106}\RG{43.295}{-0.002}{10.9091}{rm}{41}%
}\hskip.1em\allowbreak{} 、\allowbreak{}\hskip.1em \RB{9.662}{7.949}{2.211}{%
\RG{0.000}{-0.002}{10.9091}{mi}{73}\RG{4.711}{-1.638}{7.9701}{mi}{110}%
}\hskip.1em\allowbreak{} 为\allowbreak{}\hskip.1em \RB{6.555}{5.309}{0.611}{%
\RG{0.000}{-0.002}{10.9091}{mi}{110}%
}\hskip.1em\allowbreak{} 阶单位矩阵。原方程化为

\[
\RB{90.606}{9.844}{3.235}{%
\RG{83.661}{4.508}{7.9701}{rm}{50}\RG{0.000}{0.004}{10.9091}{rm}{40}\RG{4.156}{0.004}{10.9091}{rm}{50}\RG{9.496}{0.004}{10.9091}{mi}{73}\RG{21.774}{0.004}{10.9091}{rm}{43}\RG{32.506}{0.004}{10.9091}{mi}{84}\RG{40.251}{0.004}{10.9091}{rm}{41}\RG{44.405}{0.004}{10.9091}{mi}{89}\RG{56.006}{0.004}{10.9091}{rm}{61}\RG{67.345}{0.004}{10.9091}{mi}{89}\RG{75.909}{0.004}{10.9091}{mi}{84}\RG{87.979}{0.004}{10.9091}{mi}{59}\RG{14.207}{-1.632}{7.9701}{mi}{110}%
}
\]
矩阵\allowbreak{}\hskip.1em \RB{11.714}{9.295}{0.501}{%
\RG{0.000}{-0.001}{10.9091}{mi}{84}\RG{7.744}{3.958}{7.9701}{rm}{50}%
}\hskip.1em\allowbreak{} 为上三角矩阵，且\allowbreak{}\hskip.1em \RB{38.192}{9.295}{4.681}{%
\RG{0.000}{-0.001}{10.9091}{mi}{117}\RG{6.080}{-1.637}{7.9701}{mi}{106}\RG{9.622}{-1.637}{7.9701}{mi}{106}\RG{16.653}{-0.001}{10.9091}{rm}{61}\RG{27.992}{-0.001}{10.9091}{mi}{21}\RG{34.223}{3.958}{7.9701}{rm}{50}\RG{34.223}{-2.697}{7.9701}{mi}{106}%
}\hskip.1em\allowbreak{} 。比较上式的第\allowbreak{}\hskip.1em \RB{21.192}{8.681}{3.239}{%
\RG{0.000}{-0.001}{10.9091}{rm}{40}\RG{4.154}{-0.001}{10.9091}{mi}{105}\RG{7.797}{-0.001}{10.9091}{mi}{59}\RG{12.576}{-0.001}{10.9091}{mi}{106}\RG{17.616}{-0.001}{10.9091}{rm}{41}%
}\hskip.1em\allowbreak{} 个元素，得到

\[
\RB{190.810}{17.763}{12.124}{%
\RG{86.291}{12.427}{7.9701}{mi}{106}\RG{89.835}{12.427}{7.9701}{sy}{0}\RG{95.846}{12.427}{7.9701}{rm}{49}\RG{149.479}{11.820}{7.9701}{mi}{110}\RG{86.434}{9.093}{9.0910}{ex}{88}\RG{145.257}{9.093}{9.0910}{ex}{88}\RG{51.533}{4.506}{7.9701}{rm}{50}\RG{0.000}{0.002}{10.9091}{rm}{40}\RG{4.156}{0.002}{10.9091}{rm}{50}\RG{11.924}{0.002}{10.9091}{rm}{43}\RG{22.652}{0.002}{10.9091}{mi}{21}\RG{34.393}{0.002}{10.9091}{sy}{0}\RG{45.302}{0.002}{10.9091}{mi}{21}\RG{55.851}{0.002}{10.9091}{rm}{41}\RG{60.005}{0.002}{10.9091}{mi}{121}\RG{74.873}{0.002}{10.9091}{rm}{61}\RG{101.606}{0.002}{10.9091}{mi}{121}\RG{113.601}{0.002}{10.9091}{mi}{117}\RG{129.806}{0.002}{10.9091}{sy}{0}\RG{164.750}{0.002}{10.9091}{mi}{116}\RG{175.394}{0.002}{10.9091}{mi}{121}\RG{188.302}{0.002}{10.9091}{mi}{58}\RG{28.883}{-1.634}{7.9701}{mi}{105}\RG{65.214}{-1.634}{7.9701}{mi}{105}\RG{67.846}{-1.634}{7.9701}{mi}{106}\RG{106.814}{-1.634}{7.9701}{mi}{105}\RG{109.447}{-1.634}{7.9701}{mi}{114}\RG{119.681}{-1.634}{7.9701}{mi}{114}\RG{123.382}{-1.634}{7.9701}{mi}{106}\RG{168.607}{-1.634}{7.9701}{mi}{105}\RG{171.240}{-1.634}{7.9701}{mi}{114}\RG{180.603}{-1.634}{7.9701}{mi}{114}\RG{184.305}{-1.634}{7.9701}{mi}{106}\RG{51.533}{-2.696}{7.9701}{mi}{106}\RG{86.212}{-10.746}{7.9701}{mi}{114}\RG{89.914}{-10.746}{7.9701}{rm}{61}\RG{95.921}{-10.746}{7.9701}{rm}{49}\RG{140.715}{-10.867}{7.9701}{mi}{114}\RG{144.417}{-10.867}{7.9701}{rm}{61}\RG{150.427}{-10.867}{7.9701}{mi}{105}\RG{153.058}{-10.867}{7.9701}{rm}{43}\RG{159.065}{-10.867}{7.9701}{rm}{49}%
}
\]
由于

\[
\RB{139.787}{9.844}{3.616}{%
\RG{132.842}{4.508}{7.9701}{rm}{50}\RG{0.000}{0.004}{10.9091}{sy}{106}\RG{3.030}{0.004}{10.9091}{rm}{50}\RG{10.798}{0.004}{10.9091}{rm}{43}\RG{21.529}{0.004}{10.9091}{mi}{21}\RG{30.845}{0.004}{10.9091}{sy}{106}\RG{36.899}{0.004}{10.9091}{sy}{21}\RG{48.421}{0.004}{10.9091}{rm}{50}\RG{56.186}{0.004}{10.9091}{sy}{0}\RG{67.095}{0.004}{10.9091}{sy}{106}\RG{70.126}{0.004}{10.9091}{mi}{21}\RG{79.442}{0.004}{10.9091}{sy}{106}\RG{85.503}{0.004}{10.9091}{mi}{62}\RG{96.841}{0.004}{10.9091}{rm}{49}\RG{105.213}{0.004}{10.9091}{mi}{62}\RG{116.551}{0.004}{10.9091}{sy}{106}\RG{119.581}{0.004}{10.9091}{mi}{21}\RG{129.811}{0.004}{10.9091}{sy}{106}\RG{137.160}{0.004}{10.9091}{mi}{59}\RG{27.760}{-1.633}{7.9701}{mi}{105}\RG{76.357}{-1.633}{7.9701}{mi}{105}\RG{125.813}{-1.633}{7.9701}{mi}{106}%
}
\]
系数\allowbreak{}\hskip.1em \RB{43.330}{9.294}{4.682}{%
\RG{0.000}{-0.001}{10.9091}{rm}{50}\RG{5.760}{-0.001}{10.9091}{rm}{43}\RG{14.489}{-0.001}{10.9091}{mi}{21}\RG{20.720}{-1.637}{7.9701}{mi}{105}\RG{24.225}{-0.001}{10.9091}{sy}{0}\RG{33.130}{-0.001}{10.9091}{mi}{21}\RG{39.361}{3.957}{7.9701}{rm}{50}\RG{39.361}{-2.698}{7.9701}{mi}{106}%
}\hskip.1em\allowbreak{} 对所有\allowbreak{}\hskip.1em \RB{13.187}{7.666}{2.727}{%
\RG{0.000}{-0.001}{10.9091}{mi}{105}\RG{3.644}{-0.001}{10.9091}{mi}{59}\RG{8.422}{-0.001}{10.9091}{mi}{106}%
}\hskip.1em\allowbreak{} 均非零。按\allowbreak{}\hskip.1em \RB{79.116}{7.754}{2.607}{%
\RG{0.000}{-0.001}{10.9091}{mi}{105}\RG{6.675}{-0.001}{10.9091}{rm}{61}\RG{18.013}{-0.001}{10.9091}{mi}{110}\RG{24.395}{-0.001}{10.9091}{mi}{59}\RG{29.173}{-0.001}{10.9091}{mi}{110}\RG{35.971}{-0.001}{10.9091}{sy}{0}\RG{44.876}{-0.001}{10.9091}{rm}{49}\RG{50.216}{-0.001}{10.9091}{mi}{59}\RG{54.995}{-0.001}{10.9091}{mi}{58}\RG{59.784}{-0.001}{10.9091}{mi}{58}\RG{64.573}{-0.001}{10.9091}{mi}{58}\RG{69.351}{-0.001}{10.9091}{mi}{59}\RG{74.143}{-0.001}{10.9091}{rm}{49}%
}\hskip.1em\allowbreak{} 的次序，并在每行内按\allowbreak{}\hskip.1em \RB{65.348}{7.754}{2.727}{%
\RG{0.000}{-0.001}{10.9091}{mi}{106}\RG{8.063}{-0.001}{10.9091}{rm}{61}\RG{19.398}{-0.001}{10.9091}{rm}{49}\RG{24.743}{-0.001}{10.9091}{mi}{59}\RG{29.528}{-0.001}{10.9091}{rm}{50}\RG{34.869}{-0.001}{10.9091}{mi}{59}\RG{39.647}{-0.001}{10.9091}{mi}{58}\RG{44.436}{-0.001}{10.9091}{mi}{58}\RG{49.226}{-0.001}{10.9091}{mi}{58}\RG{54.004}{-0.001}{10.9091}{mi}{59}\RG{58.793}{-0.001}{10.9091}{mi}{110}%
}\hskip.1em\allowbreak{}的次序归纳。计算\allowbreak{}\hskip.1em \RB{11.344}{5.311}{3.621}{%
\RG{0.000}{-0.000}{10.9091}{mi}{121}\RG{5.208}{-1.637}{7.9701}{mi}{105}\RG{7.840}{-1.637}{7.9701}{mi}{106}%
}\hskip.1em\allowbreak{} 时，第一项和式中的\allowbreak{}\hskip.1em \RB{11.701}{5.311}{2.726}{%
\RG{0.000}{-0.000}{10.9091}{mi}{121}\RG{5.209}{-1.637}{7.9701}{mi}{105}\RG{7.841}{-1.637}{7.9701}{mi}{114}%
}\hskip.1em\allowbreak{} 属于本行此前已计算的元素，第二项和式中的\allowbreak{}\hskip.1em \RB{12.413}{5.311}{3.620}{%
\RG{0.000}{-0.000}{10.9091}{mi}{121}\RG{5.208}{-1.636}{7.9701}{mi}{114}\RG{8.910}{-1.636}{7.9701}{mi}{106}%
}\hskip.1em\allowbreak{} 属于此前已计算的行；这些元素均为零，边界处的空和亦为零。故\allowbreak{}\hskip.1em \RB{32.020}{7.755}{3.620}{%
\RG{0.000}{-0.000}{10.9091}{mi}{121}\RG{5.209}{-1.636}{7.9701}{mi}{105}\RG{7.841}{-1.636}{7.9701}{mi}{106}\RG{15.080}{-0.000}{10.9091}{rm}{61}\RG{26.633}{-0.000}{10.9091}{rm}{48}%
}\hskip.1em\allowbreak{} 。归纳得\allowbreak{}\hskip.1em \RB{31.361}{8.181}{0.728}{%
\RG{0.000}{0.001}{10.9091}{mi}{89}\RG{11.602}{0.001}{10.9091}{rm}{61}\RG{22.941}{0.001}{10.9091}{mi}{79}%
}\hskip.1em\allowbreak{} ，从而\allowbreak{}\hskip.1em \RB{82.518}{9.297}{0.728}{%
\RG{0.000}{0.001}{10.9091}{mi}{88}\RG{12.753}{0.001}{10.9091}{rm}{61}\RG{24.092}{0.001}{10.9091}{mi}{80}\RG{32.449}{0.001}{10.9091}{mi}{89}\RG{41.023}{0.001}{10.9091}{mi}{80}\RG{49.401}{3.960}{7.9701}{sy}{0}\RG{55.411}{3.960}{7.9701}{rm}{49}\RG{62.760}{0.001}{10.9091}{rm}{61}\RG{74.098}{0.001}{10.9091}{mi}{79}%
}\hskip.1em\allowbreak{} 。零矩阵满足原方程，因此它是唯一解：\allowbreak{}\hskip.1em \RB{8.366}{7.995}{0.500}{%
\RG{0.000}{0.001}{10.9091}{am}{3}%
}\hskip.1em\allowbreak{}

% Original page 3.
\section*{2016S05}
设\allowbreak{}\hskip.1em \RB{42.345}{8.682}{3.620}{%
\RG{0.000}{-0.000}{10.9091}{mi}{65}\RG{11.039}{-0.000}{10.9091}{rm}{61}\RG{22.373}{-0.000}{10.9091}{rm}{40}\RG{26.531}{-0.000}{10.9091}{mi}{97}\RG{32.139}{-1.636}{7.9701}{mi}{105}\RG{34.771}{-1.636}{7.9701}{mi}{106}\RG{38.768}{-0.000}{10.9091}{rm}{41}%
}\hskip.1em\allowbreak{} 为\allowbreak{}\hskip.1em \RB{6.555}{5.311}{0.609}{%
\RG{0.000}{-0.000}{10.9091}{mi}{110}%
}\hskip.1em\allowbreak{} 阶复矩阵。证明：存在\allowbreak{}\hskip.1em \RB{24.896}{8.256}{0.805}{%
\RG{0.000}{-0.000}{10.9091}{mi}{14}\RG{8.171}{-0.000}{10.9091}{mi}{62}\RG{19.509}{-0.000}{10.9091}{rm}{48}%
}\hskip.1em\allowbreak{} ，使对每个\allowbreak{}\hskip.1em \RB{41.362}{8.682}{3.238}{%
\RG{0.000}{-0.000}{10.9091}{mi}{115}\RG{8.062}{-0.000}{10.9091}{sy}{50}\RG{18.365}{-0.000}{10.9091}{rm}{40}\RG{22.522}{-0.000}{10.9091}{rm}{48}\RG{27.860}{-0.000}{10.9091}{mi}{59}\RG{32.638}{-0.000}{10.9091}{mi}{14}\RG{37.785}{-0.000}{10.9091}{rm}{41}%
}\hskip.1em\allowbreak{} ，矩阵

\[
\RB{315.171}{33.222}{27.761}{%
\RG{152.554}{32.368}{9.0910}{ex}{48}\RG{309.371}{32.368}{9.0910}{ex}{49}\RG{160.508}{23.699}{10.9091}{mi}{97}\RG{176.722}{23.699}{10.9091}{rm}{43}\RG{187.455}{23.699}{10.9091}{mi}{115}\RG{212.819}{23.699}{10.9091}{mi}{97}\RG{246.964}{23.699}{10.9091}{sy}{1}\RG{251.808}{23.699}{10.9091}{sy}{1}\RG{256.663}{23.699}{10.9091}{sy}{1}\RG{282.679}{23.699}{10.9091}{mi}{97}\RG{166.116}{22.062}{7.9701}{rm}{49}\RG{169.978}{22.062}{7.9701}{rm}{49}\RG{218.437}{22.062}{7.9701}{rm}{49}\RG{222.299}{22.062}{7.9701}{rm}{50}\RG{288.287}{22.062}{7.9701}{rm}{49}\RG{292.150}{22.062}{7.9701}{mi}{110}\RG{152.554}{16.368}{9.0910}{ex}{66}\RG{309.371}{16.368}{9.0910}{ex}{67}\RG{233.555}{12.894}{7.9701}{rm}{50}\RG{152.554}{10.913}{9.0910}{ex}{66}\RG{309.371}{10.913}{9.0910}{ex}{67}\RG{169.602}{8.934}{10.9091}{mi}{97}\RG{201.577}{8.934}{10.9091}{mi}{97}\RG{217.790}{8.934}{10.9091}{rm}{43}\RG{228.523}{8.934}{10.9091}{mi}{115}\RG{246.964}{8.934}{10.9091}{sy}{1}\RG{251.808}{8.934}{10.9091}{sy}{1}\RG{256.663}{8.934}{10.9091}{sy}{1}\RG{282.679}{8.934}{10.9091}{mi}{97}\RG{175.210}{7.298}{7.9701}{rm}{50}\RG{179.072}{7.298}{7.9701}{rm}{49}\RG{207.184}{7.298}{7.9701}{rm}{50}\RG{211.046}{7.298}{7.9701}{rm}{50}\RG{288.287}{7.298}{7.9701}{rm}{50}\RG{292.150}{7.298}{7.9701}{mi}{110}\RG{152.554}{5.459}{9.0910}{ex}{66}\RG{309.371}{5.459}{9.0910}{ex}{67}\RG{95.611}{4.508}{7.9701}{rm}{50}\RG{128.888}{4.508}{7.9701}{mi}{110}\RG{0.000}{0.004}{10.9091}{mi}{65}\RG{8.007}{0.004}{10.9091}{rm}{40}\RG{12.162}{0.004}{10.9091}{mi}{115}\RG{17.194}{0.004}{10.9091}{rm}{41}\RG{24.372}{0.004}{10.9091}{rm}{61}\RG{35.716}{0.004}{10.9091}{mi}{65}\RG{46.148}{0.004}{10.9091}{rm}{43}\RG{56.882}{0.004}{10.9091}{rm}{100}\RG{62.817}{0.004}{10.9091}{rm}{105}\RG{65.784}{0.004}{10.9091}{rm}{97}\RG{71.130}{0.004}{10.9091}{rm}{103}\RG{76.617}{0.004}{10.9091}{rm}{40}\RG{80.762}{0.004}{10.9091}{mi}{115}\RG{85.791}{0.004}{10.9091}{mi}{59}\RG{90.569}{0.004}{10.9091}{mi}{115}\RG{99.929}{0.004}{10.9091}{mi}{59}\RG{104.708}{0.004}{10.9091}{mi}{58}\RG{109.497}{0.004}{10.9091}{mi}{58}\RG{114.275}{0.004}{10.9091}{mi}{58}\RG{119.064}{0.004}{10.9091}{mi}{59}\RG{123.853}{0.004}{10.9091}{mi}{115}\RG{134.031}{0.004}{10.9091}{rm}{41}\RG{141.209}{0.004}{10.9091}{rm}{61}\RG{152.554}{0.004}{9.0910}{ex}{66}\RG{309.371}{0.004}{9.0910}{ex}{67}\RT{174.981}{-2.105}{10.9091}{.}\RT{218.208}{-2.105}{10.9091}{.}\RT{288.470}{-2.105}{10.9091}{.}\RT{248.475}{-3.015}{10.9091}{.}\RG{152.554}{-5.451}{9.0910}{ex}{66}\RG{309.371}{-5.451}{9.0910}{ex}{67}\RT{174.981}{-5.742}{10.9091}{.}\RT{218.208}{-5.742}{10.9091}{.}\RT{252.720}{-5.742}{10.9091}{.}\RT{288.470}{-5.742}{10.9091}{.}\RT{256.965}{-8.469}{10.9091}{.}\RT{174.981}{-9.378}{10.9091}{.}\RT{218.208}{-9.378}{10.9091}{.}\RT{288.470}{-9.378}{10.9091}{.}\RG{152.554}{-11.270}{9.0910}{ex}{64}\RG{309.371}{-11.270}{9.0910}{ex}{65}\RG{304.228}{-20.183}{7.9701}{mi}{110}\RG{169.190}{-24.142}{10.9091}{mi}{97}\RG{212.417}{-24.142}{10.9091}{mi}{97}\RG{246.964}{-24.142}{10.9091}{sy}{1}\RG{251.808}{-24.142}{10.9091}{sy}{1}\RG{256.663}{-24.142}{10.9091}{sy}{1}\RG{270.600}{-24.142}{10.9091}{mi}{97}\RG{288.464}{-24.142}{10.9091}{rm}{43}\RG{299.197}{-24.142}{10.9091}{mi}{115}\RG{174.797}{-25.778}{7.9701}{mi}{110}\RG{179.486}{-25.778}{7.9701}{rm}{49}\RG{218.025}{-25.778}{7.9701}{mi}{110}\RG{222.714}{-25.778}{7.9701}{rm}{50}\RG{276.208}{-25.778}{7.9701}{mi}{110}\RG{280.899}{-25.778}{7.9701}{mi}{110}%
}
\]
均可对角化。

证明：当\allowbreak{}\hskip.1em \RB{25.714}{7.754}{0.610}{%
\RG{0.000}{-0.001}{10.9091}{mi}{110}\RG{9.406}{-0.001}{10.9091}{rm}{61}\RG{20.741}{-0.001}{10.9091}{rm}{49}%
}\hskip.1em\allowbreak{} 时，所有\allowbreak{}\hskip.1em \RB{20.770}{8.681}{3.239}{%
\RG{0.000}{-0.001}{10.9091}{mi}{65}\RG{8.007}{-0.001}{10.9091}{rm}{40}\RG{12.162}{-0.001}{10.9091}{mi}{115}\RG{17.194}{-0.001}{10.9091}{rm}{41}%
}\hskip.1em\allowbreak{} 都是对角矩阵，取\allowbreak{}\hskip.1em \RB{24.478}{8.256}{0.610}{%
\RG{0.000}{-0.001}{10.9091}{mi}{14}\RG{8.171}{-0.001}{10.9091}{rm}{61}\RG{19.505}{-0.001}{10.9091}{rm}{49}%
}\hskip.1em\allowbreak{} 即满足要求。以下设\allowbreak{}\hskip.1em \RB{26.221}{7.754}{1.995}{%
\RG{0.000}{-0.001}{10.9091}{mi}{110}\RG{9.406}{-0.001}{10.9091}{sy}{21}\RG{20.921}{-0.001}{10.9091}{rm}{50}%
}\hskip.1em\allowbreak{} 。令

\[
\RB{233.059}{53.122}{8.360}{%
\RG{77.951}{52.622}{9.0910}{ex}{88}\RG{50.190}{43.531}{10.9091}{mi}{114}\RG{61.120}{43.531}{10.9091}{rm}{61}\RG{98.394}{43.531}{10.9091}{sy}{106}\RG{101.424}{43.531}{10.9091}{mi}{97}\RG{113.662}{43.531}{10.9091}{sy}{106}\RG{127.601}{43.531}{10.9091}{rm}{40}\RG{131.758}{43.531}{10.9091}{rm}{49}\RG{140.126}{43.531}{10.9091}{sy}{20}\RG{151.641}{43.531}{10.9091}{mi}{105}\RG{158.315}{43.531}{10.9091}{sy}{20}\RG{169.831}{43.531}{10.9091}{mi}{110}\RG{176.207}{43.531}{10.9091}{rm}{41}\RG{180.362}{43.531}{10.9091}{mi}{59}\RG{55.004}{41.894}{7.9701}{mi}{105}\RG{107.031}{41.894}{7.9701}{mi}{105}\RG{109.664}{41.894}{7.9701}{mi}{106}\RG{72.459}{32.608}{7.9701}{rm}{49}\RG{76.323}{32.608}{7.9701}{sy}{20}\RG{82.332}{32.608}{7.9701}{mi}{106}\RG{85.877}{32.608}{7.9701}{sy}{20}\RG{91.887}{32.608}{7.9701}{mi}{110}\RG{78.425}{25.052}{7.9701}{mi}{106}\RG{81.969}{25.052}{7.9701}{sy}{54}\RG{81.969}{25.052}{7.9701}{rm}{61}\RG{87.979}{25.052}{7.9701}{mi}{105}\RG{55.553}{10.547}{9.0910}{ex}{16}\RG{147.876}{10.547}{9.0910}{ex}{17}\RG{0.000}{0.001}{10.9091}{mi}{67}\RG{11.457}{0.001}{10.9091}{rm}{61}\RG{29.141}{0.001}{10.9091}{rm}{109}\RG{38.043}{0.001}{10.9091}{rm}{97}\RG{43.388}{0.001}{10.9091}{rm}{120}\RG{60.982}{0.001}{10.9091}{sy}{106}\RG{64.013}{0.001}{10.9091}{mi}{97}\RG{77.761}{0.001}{10.9091}{sy}{0}\RG{88.670}{0.001}{10.9091}{mi}{97}\RG{101.820}{0.001}{10.9091}{sy}{106}\RG{107.275}{0.001}{10.9091}{rm}{43}\RG{118.008}{0.001}{10.9091}{mi}{114}\RG{128.331}{0.001}{10.9091}{rm}{43}\RG{139.063}{0.001}{10.9091}{mi}{114}\RG{153.306}{0.001}{10.9091}{mi}{59}\RG{179.903}{0.001}{10.9091}{mi}{115}\RG{192.289}{0.001}{10.9091}{rm}{61}\RG{203.628}{0.001}{10.9091}{mi}{67}\RG{214.478}{0.001}{10.9091}{rm}{43}\RG{225.213}{0.001}{10.9091}{rm}{50}\RG{230.552}{0.001}{10.9091}{mi}{58}\RG{69.620}{-1.635}{7.9701}{mi}{105}\RG{72.253}{-1.635}{7.9701}{mi}{105}\RG{94.277}{-1.635}{7.9701}{mi}{106}\RG{97.819}{-1.635}{7.9701}{mi}{106}\RG{122.822}{-1.635}{7.9701}{mi}{105}\RG{143.878}{-1.635}{7.9701}{mi}{106}\RG{184.941}{-1.635}{7.9701}{rm}{48}\RG{22.795}{-6.377}{7.9701}{rm}{49}\RG{26.658}{-6.377}{7.9701}{sy}{20}\RG{32.669}{-6.377}{7.9701}{mi}{105}\RG{35.302}{-6.377}{7.9701}{mi}{60}\RG{41.309}{-6.377}{7.9701}{mi}{106}\RG{44.855}{-6.377}{7.9701}{sy}{20}\RG{50.864}{-6.377}{7.9701}{mi}{110}%
}
\]
由\allowbreak{}\hskip.1em \RB{28.359}{8.177}{1.998}{%
\RG{0.000}{-0.004}{10.9091}{mi}{67}\RG{11.457}{-0.004}{10.9091}{sy}{21}\RG{22.972}{-0.004}{10.9091}{rm}{48}%
}\hskip.1em\allowbreak{} 得\allowbreak{}\hskip.1em \RB{29.195}{7.751}{2.292}{%
\RG{0.000}{-0.004}{10.9091}{mi}{115}\RG{5.032}{-1.640}{7.9701}{rm}{48}\RG{12.380}{-0.004}{10.9091}{sy}{21}\RG{23.895}{-0.004}{10.9091}{rm}{50}%
}\hskip.1em\allowbreak{} 。对于\allowbreak{}\hskip.1em \RB{64.032}{7.751}{2.729}{%
\RG{0.000}{-0.004}{10.9091}{rm}{49}\RG{8.371}{-0.004}{10.9091}{sy}{20}\RG{19.886}{-0.004}{10.9091}{mi}{105}\RG{26.552}{-0.004}{10.9091}{mi}{60}\RG{37.898}{-0.004}{10.9091}{mi}{106}\RG{45.962}{-0.004}{10.9091}{sy}{20}\RG{57.477}{-0.004}{10.9091}{mi}{110}%
}\hskip.1em\allowbreak{} ，有

\[
\RB{185.349}{10.489}{3.348}{%
\RG{5.032}{5.174}{7.9701}{mi}{106}\RG{64.937}{5.174}{7.9701}{mi}{106}\RG{68.481}{5.174}{7.9701}{sy}{0}\RG{74.492}{5.174}{7.9701}{mi}{105}\RG{27.715}{4.506}{7.9701}{mi}{105}\RG{51.434}{4.506}{7.9701}{mi}{105}\RG{0.000}{0.002}{10.9091}{mi}{115}\RG{11.774}{0.002}{10.9091}{sy}{0}\RG{22.683}{0.002}{10.9091}{mi}{115}\RG{35.063}{0.002}{10.9091}{rm}{61}\RG{46.402}{0.002}{10.9091}{mi}{115}\RG{55.752}{0.002}{10.9091}{rm}{40}\RG{59.906}{0.002}{10.9091}{mi}{115}\RG{80.002}{0.002}{10.9091}{sy}{0}\RG{90.911}{0.002}{10.9091}{rm}{49}\RG{96.256}{0.002}{10.9091}{rm}{41}\RG{103.436}{0.002}{10.9091}{sy}{21}\RG{114.952}{0.002}{10.9091}{mi}{115}\RG{124.302}{0.002}{10.9091}{rm}{40}\RG{128.455}{0.002}{10.9091}{mi}{115}\RG{140.229}{0.002}{10.9091}{sy}{0}\RG{151.139}{0.002}{10.9091}{rm}{49}\RG{156.484}{0.002}{10.9091}{rm}{41}\RG{163.664}{0.002}{10.9091}{mi}{62}\RG{174.998}{0.002}{10.9091}{mi}{67}\RG{182.842}{0.002}{10.9091}{mi}{58}\RG{119.984}{-1.634}{7.9701}{rm}{48}\RG{133.487}{-1.634}{7.9701}{rm}{48}\RG{5.032}{-2.382}{7.9701}{rm}{48}\RG{64.937}{-2.382}{7.9701}{rm}{48}\RG{27.715}{-2.696}{7.9701}{rm}{48}\RG{51.434}{-2.696}{7.9701}{rm}{48}%
}
\]
因而

\[
\RB{208.482}{11.603}{40.534}{%
\RG{0.000}{10.912}{9.0910}{ex}{12}\RG{102.467}{10.912}{9.0910}{ex}{12}\RG{0.000}{5.458}{9.0910}{ex}{12}\RG{102.467}{5.458}{9.0910}{ex}{12}\RG{38.523}{5.175}{7.9701}{mi}{106}\RG{125.074}{5.175}{7.9701}{mi}{106}\RG{93.994}{4.507}{7.9701}{mi}{105}\RG{147.758}{4.507}{7.9701}{mi}{105}\RG{0.000}{0.003}{9.0910}{ex}{12}\RG{3.030}{0.003}{10.9091}{rm}{40}\RG{7.184}{0.003}{10.9091}{mi}{97}\RG{22.759}{0.003}{10.9091}{rm}{43}\RG{33.491}{0.003}{10.9091}{mi}{115}\RG{42.841}{0.003}{10.9091}{rm}{41}\RG{49.419}{0.003}{10.9091}{sy}{0}\RG{60.328}{0.003}{10.9091}{rm}{40}\RG{64.483}{0.003}{10.9091}{mi}{97}\RG{78.231}{0.003}{10.9091}{rm}{43}\RG{88.963}{0.003}{10.9091}{mi}{115}\RG{98.313}{0.003}{10.9091}{rm}{41}\RG{102.467}{0.003}{9.0910}{ex}{12}\RG{108.527}{0.003}{10.9091}{sy}{21}\RG{120.043}{0.003}{10.9091}{mi}{115}\RG{131.817}{0.003}{10.9091}{sy}{0}\RG{142.726}{0.003}{10.9091}{mi}{115}\RG{154.500}{0.003}{10.9091}{sy}{0}\RG{165.410}{0.003}{10.9091}{sy}{106}\RG{168.440}{0.003}{10.9091}{mi}{97}\RG{184.015}{0.003}{10.9091}{sy}{0}\RG{194.924}{0.003}{10.9091}{mi}{97}\RG{206.247}{0.003}{10.9091}{sy}{106}\RG{12.792}{-1.633}{7.9701}{mi}{106}\RG{16.333}{-1.633}{7.9701}{mi}{106}\RG{70.090}{-1.633}{7.9701}{mi}{105}\RG{72.723}{-1.633}{7.9701}{mi}{105}\RG{174.047}{-1.633}{7.9701}{mi}{106}\RG{177.589}{-1.633}{7.9701}{mi}{106}\RG{200.531}{-1.633}{7.9701}{mi}{105}\RG{203.164}{-1.633}{7.9701}{mi}{105}\RG{38.523}{-2.381}{7.9701}{rm}{48}\RG{125.074}{-2.381}{7.9701}{rm}{48}\RG{93.994}{-2.694}{7.9701}{rm}{48}\RG{147.758}{-2.694}{7.9701}{rm}{48}\RG{108.527}{-19.423}{10.9091}{mi}{62}\RG{119.862}{-19.423}{10.9091}{mi}{67}\RG{130.716}{-19.423}{10.9091}{sy}{0}\RG{141.625}{-19.423}{10.9091}{sy}{106}\RG{144.655}{-19.423}{10.9091}{mi}{97}\RG{160.230}{-19.423}{10.9091}{sy}{0}\RG{171.139}{-19.423}{10.9091}{mi}{97}\RG{182.463}{-19.423}{10.9091}{sy}{106}\RG{150.263}{-21.059}{7.9701}{mi}{106}\RG{153.804}{-21.059}{7.9701}{mi}{106}\RG{176.746}{-21.059}{7.9701}{mi}{105}\RG{179.379}{-21.059}{7.9701}{mi}{105}\RG{108.527}{-36.914}{10.9091}{sy}{21}\RG{120.043}{-36.914}{10.9091}{mi}{114}\RG{130.366}{-36.914}{10.9091}{rm}{43}\RG{141.099}{-36.914}{10.9091}{mi}{114}\RG{149.911}{-36.914}{10.9091}{mi}{58}\RG{124.856}{-38.550}{7.9701}{mi}{105}\RG{145.913}{-38.550}{7.9701}{mi}{106}%
}
\]
故\allowbreak{}\hskip.1em \RB{25.088}{8.684}{3.236}{%
\RG{0.000}{0.002}{10.9091}{mi}{65}\RG{8.008}{0.002}{10.9091}{rm}{40}\RG{12.162}{0.002}{10.9091}{mi}{115}\RG{17.194}{-1.634}{7.9701}{rm}{48}\RG{21.512}{0.002}{10.9091}{rm}{41}%
}\hskip.1em\allowbreak{} 的\allowbreak{}\hskip.1em \RB{6.555}{5.313}{0.607}{%
\RG{0.000}{0.002}{10.9091}{mi}{110}%
}\hskip.1em\allowbreak{} 个Gershgorin圆盘

\[
\RB{225.439}{9.817}{3.351}{%
\RG{122.945}{4.502}{7.9701}{mi}{105}\RG{0.000}{-0.002}{10.9091}{mi}{68}\RG{14.984}{-0.002}{10.9091}{rm}{61}\RG{26.323}{-0.002}{10.9091}{sy}{102}\RG{31.778}{-0.002}{10.9091}{mi}{122}\RG{40.256}{-0.002}{10.9091}{sy}{50}\RG{50.560}{-0.002}{10.9091}{bb}{67}\RG{61.469}{-0.002}{10.9091}{rm}{58}\RG{67.467}{-0.002}{10.9091}{sy}{106}\RG{70.497}{-0.002}{10.9091}{mi}{122}\RG{78.369}{-0.002}{10.9091}{sy}{0}\RG{89.278}{-0.002}{10.9091}{rm}{40}\RG{93.432}{-0.002}{10.9091}{mi}{97}\RG{107.180}{-0.002}{10.9091}{rm}{43}\RG{117.913}{-0.002}{10.9091}{mi}{115}\RG{127.263}{-0.002}{10.9091}{rm}{41}\RG{131.416}{-0.002}{10.9091}{sy}{106}\RG{137.471}{-0.002}{10.9091}{sy}{20}\RG{148.993}{-0.002}{10.9091}{mi}{114}\RG{156.892}{-0.002}{10.9091}{sy}{103}\RG{173.256}{-0.002}{10.9091}{rm}{40}\RG{177.413}{-0.002}{10.9091}{rm}{49}\RG{185.781}{-0.002}{10.9091}{sy}{20}\RG{197.296}{-0.002}{10.9091}{mi}{105}\RG{203.971}{-0.002}{10.9091}{sy}{20}\RG{215.486}{-0.002}{10.9091}{mi}{110}\RG{221.863}{-0.002}{10.9091}{rm}{41}\RG{8.869}{-1.638}{7.9701}{mi}{105}\RG{99.040}{-1.638}{7.9701}{mi}{105}\RG{101.672}{-1.638}{7.9701}{mi}{105}\RG{153.807}{-1.638}{7.9701}{mi}{105}\RG{122.945}{-2.698}{7.9701}{rm}{48}%
}
\]
两两不交。由Gershgorin圆盘定理的计数结论，每个圆盘恰含一个特征值，按代数重数计。因此，\allowbreak{}\hskip.1em \RB{25.088}{8.686}{3.235}{%
\RG{0.000}{0.004}{10.9091}{mi}{65}\RG{8.007}{0.004}{10.9091}{rm}{40}\RG{12.161}{0.004}{10.9091}{mi}{115}\RG{17.193}{-1.632}{7.9701}{rm}{48}\RG{21.512}{0.004}{10.9091}{rm}{41}%
}\hskip.1em\allowbreak{} 的特征值两两不同。

记\allowbreak{}\hskip.1em \RB{9.662}{7.949}{2.212}{%
\RG{0.000}{-0.002}{10.9091}{mi}{73}\RG{4.711}{-1.639}{7.9701}{mi}{110}%
}\hskip.1em\allowbreak{} 为\allowbreak{}\hskip.1em \RB{6.555}{5.309}{0.611}{%
\RG{0.000}{-0.002}{10.9091}{mi}{110}%
}\hskip.1em\allowbreak{} 阶单位矩阵，\allowbreak{}\hskip.1em \RB{21.663}{8.680}{3.240}{%
\RG{0.000}{-0.002}{10.9091}{rm}{1}\RG{8.902}{-0.002}{10.9091}{rm}{40}\RG{13.055}{-0.002}{10.9091}{mi}{115}\RG{18.087}{-0.002}{10.9091}{rm}{41}%
}\hskip.1em\allowbreak{} 为首一多项式\allowbreak{}\hskip.1em \RB{71.824}{8.680}{3.240}{%
\RG{0.000}{-0.002}{10.9091}{rm}{100}\RG{5.935}{-0.002}{10.9091}{rm}{101}\RG{10.680}{-0.002}{10.9091}{rm}{116}\RG{14.837}{-0.002}{10.9091}{rm}{40}\RG{18.989}{-0.002}{10.9091}{mi}{122}\RG{24.433}{-0.002}{10.9091}{mi}{73}\RG{29.149}{-1.639}{7.9701}{mi}{110}\RG{36.353}{-0.002}{10.9091}{sy}{0}\RG{46.899}{-0.002}{10.9091}{mi}{65}\RG{54.906}{-0.002}{10.9091}{rm}{40}\RG{59.060}{-0.002}{10.9091}{mi}{115}\RG{64.091}{-0.002}{10.9091}{rm}{41}\RG{68.248}{-0.002}{10.9091}{rm}{41}%
}\hskip.1em\allowbreak{} 关于\allowbreak{}\hskip.1em \RB{5.518}{5.309}{0.611}{%
\RG{0.000}{-0.002}{10.9091}{mi}{122}%
}\hskip.1em\allowbreak{} 的判别式，即其全部根两两之差平方的乘积。由对称多项式基本定理，判别式是多项式系数的多项式。由于\allowbreak{}\hskip.1em \RB{20.770}{8.681}{3.239}{%
\RG{0.000}{-0.001}{10.9091}{mi}{65}\RG{8.007}{-0.001}{10.9091}{rm}{40}\RG{12.162}{-0.001}{10.9091}{mi}{115}\RG{17.194}{-0.001}{10.9091}{rm}{41}%
}\hskip.1em\allowbreak{} 的各元素属于\allowbreak{}\hskip.1em \RB{18.132}{8.681}{3.228}{%
\RG{0.000}{-0.001}{10.9091}{bb}{67}\RG{7.879}{-0.001}{10.9091}{rm}{91}\RG{10.909}{-0.001}{10.9091}{mi}{115}\RG{15.941}{-0.001}{10.9091}{rm}{93}%
}\hskip.1em\allowbreak{} ，特征多项式的各系数也属于\allowbreak{}\hskip.1em \RB{18.133}{8.681}{3.228}{%
\RG{0.000}{-0.001}{10.9091}{bb}{67}\RG{7.880}{-0.001}{10.9091}{rm}{91}\RG{10.910}{-0.001}{10.9091}{mi}{115}\RG{15.942}{-0.001}{10.9091}{rm}{93}%
}\hskip.1em\allowbreak{} ，所以\allowbreak{}\hskip.1em \RB{54.703}{8.681}{3.239}{%
\RG{0.000}{-0.001}{10.9091}{rm}{1}\RG{8.902}{-0.001}{10.9091}{rm}{40}\RG{13.056}{-0.001}{10.9091}{mi}{115}\RG{18.088}{-0.001}{10.9091}{rm}{41}\RG{25.770}{-0.001}{10.9091}{sy}{50}\RG{36.571}{-0.001}{10.9091}{bb}{67}\RG{44.450}{-0.001}{10.9091}{rm}{91}\RG{47.481}{-0.001}{10.9091}{mi}{115}\RG{52.512}{-0.001}{10.9091}{rm}{93}%
}\hskip.1em\allowbreak{} 。上述特征值互异性给出\allowbreak{}\hskip.1em \RB{46.312}{8.682}{3.238}{%
\RG{0.000}{0.000}{10.9091}{rm}{1}\RG{8.902}{0.000}{10.9091}{rm}{40}\RG{13.056}{0.000}{10.9091}{mi}{115}\RG{18.088}{-1.636}{7.9701}{rm}{48}\RG{22.406}{0.000}{10.9091}{rm}{41}\RG{29.590}{0.000}{10.9091}{sy}{54}\RG{29.590}{0.000}{10.9091}{rm}{61}\RG{40.925}{0.000}{10.9091}{rm}{48}%
}\hskip.1em\allowbreak{} ，故\allowbreak{}\hskip.1em \RB{8.889}{8.289}{0.500}{%
\RG{0.000}{0.000}{10.9091}{rm}{1}%
}\hskip.1em\allowbreak{} 不是零多项式。

设\allowbreak{}\hskip.1em \RB{152.442}{8.685}{3.235}{%
\RG{0.000}{0.003}{10.9091}{mi}{69}\RG{12.359}{0.003}{10.9091}{rm}{61}\RG{24.481}{0.003}{10.9091}{sy}{102}\RG{29.935}{0.003}{10.9091}{mi}{114}\RG{38.858}{0.003}{10.9091}{sy}{50}\RG{49.945}{0.003}{10.9091}{bb}{82}\RG{61.636}{0.003}{10.9091}{rm}{58}\RG{68.417}{0.003}{10.9091}{mi}{114}\RG{77.329}{0.003}{10.9091}{mi}{62}\RG{89.461}{0.003}{10.9091}{rm}{48}\RG{94.802}{0.003}{10.9091}{mi}{59}\RG{103.689}{0.003}{10.9091}{rm}{1}\RG{112.591}{0.003}{10.9091}{rm}{40}\RG{116.745}{0.003}{10.9091}{mi}{114}\RG{121.855}{0.003}{10.9091}{rm}{41}\RG{129.819}{0.003}{10.9091}{rm}{61}\RG{141.950}{0.003}{10.9091}{rm}{48}\RG{147.284}{0.003}{10.9091}{sy}{103}%
}\hskip.1em\allowbreak{} 。非零一元多项式只有有限个根，故\allowbreak{}\hskip.1em \RB{8.715}{7.932}{0.500}{%
\RG{0.000}{0.003}{10.9091}{mi}{69}%
}\hskip.1em\allowbreak{} 有限。若\allowbreak{}\hskip.1em \RB{32.212}{7.933}{0.518}{%
\RG{0.000}{0.004}{10.9091}{mi}{69}\RG{12.042}{0.004}{10.9091}{rm}{61}\RG{23.847}{0.004}{10.9091}{bb}{63}%
}\hskip.1em\allowbreak{} ，取\allowbreak{}\hskip.1em \RB{25.412}{8.260}{0.605}{%
\RG{0.000}{0.004}{10.9091}{mi}{14}\RG{8.635}{0.004}{10.9091}{rm}{61}\RG{20.439}{0.004}{10.9091}{rm}{49}%
}\hskip.1em\allowbreak{} ；否则取\allowbreak{}\hskip.1em \RB{48.829}{8.260}{0.605}{%
\RG{0.000}{0.004}{10.9091}{mi}{14}\RG{8.636}{0.004}{10.9091}{rm}{61}\RG{20.440}{0.004}{10.9091}{rm}{109}\RG{29.342}{0.004}{10.9091}{rm}{105}\RG{32.309}{0.004}{10.9091}{rm}{110}\RG{40.114}{0.004}{10.9091}{mi}{69}%
}\hskip.1em\allowbreak{} 。于是\allowbreak{}\hskip.1em \RB{25.827}{8.260}{0.802}{%
\RG{0.000}{0.004}{10.9091}{mi}{14}\RG{8.629}{0.004}{10.9091}{mi}{62}\RG{20.439}{0.004}{10.9091}{rm}{48}%
}\hskip.1em\allowbreak{} ，且对每个\allowbreak{}\hskip.1em \RB{42.293}{8.686}{3.235}{%
\RG{0.000}{0.004}{10.9091}{mi}{115}\RG{8.527}{0.004}{10.9091}{sy}{50}\RG{19.296}{0.004}{10.9091}{rm}{40}\RG{23.452}{0.004}{10.9091}{rm}{48}\RG{28.791}{0.004}{10.9091}{mi}{59}\RG{33.569}{0.004}{10.9091}{mi}{14}\RG{38.716}{0.004}{10.9091}{rm}{41}%
}\hskip.1em\allowbreak{} ，均有\allowbreak{}\hskip.1em \RB{42.929}{8.686}{3.235}{%
\RG{0.000}{0.004}{10.9091}{rm}{1}\RG{8.902}{0.004}{10.9091}{rm}{40}\RG{13.055}{0.004}{10.9091}{mi}{115}\RG{18.087}{0.004}{10.9091}{rm}{41}\RG{25.737}{0.004}{10.9091}{sy}{54}\RG{25.737}{0.004}{10.9091}{rm}{61}\RG{37.541}{0.004}{10.9091}{rm}{48}%
}\hskip.1em\allowbreak{} 。因此\allowbreak{}\hskip.1em \RB{20.769}{8.677}{3.243}{%
\RG{0.000}{-0.005}{10.9091}{mi}{65}\RG{8.007}{-0.005}{10.9091}{rm}{40}\RG{12.161}{-0.005}{10.9091}{mi}{115}\RG{17.193}{-0.005}{10.9091}{rm}{41}%
}\hskip.1em\allowbreak{} 有\allowbreak{}\hskip.1em \RB{6.555}{5.306}{0.614}{%
\RG{0.000}{-0.005}{10.9091}{mi}{110}%
}\hskip.1em\allowbreak{} 个互异特征值，从而可对角化。\allowbreak{}\hskip.1em \RB{8.366}{7.990}{0.505}{%
\RG{0.000}{-0.005}{10.9091}{am}{3}%
}\hskip.1em\allowbreak{}

% Original page 4.
\section*{2016S06}
（1）设\allowbreak{}\hskip.1em \RB{71.424}{8.681}{3.621}{%
\RG{0.000}{-0.001}{10.9091}{mi}{65}\RG{8.007}{-0.001}{10.9091}{rm}{40}\RG{12.161}{-0.001}{10.9091}{mi}{21}\RG{18.393}{-0.001}{10.9091}{rm}{41}\RG{25.571}{-0.001}{10.9091}{rm}{61}\RG{36.917}{-0.001}{10.9091}{rm}{40}\RG{41.069}{-0.001}{10.9091}{mi}{97}\RG{46.677}{-1.637}{7.9701}{mi}{105}\RG{49.310}{-1.637}{7.9701}{mi}{106}\RG{53.306}{-0.001}{10.9091}{rm}{40}\RG{57.460}{-0.001}{10.9091}{mi}{21}\RG{63.692}{-0.001}{10.9091}{rm}{41}\RG{67.848}{-0.001}{10.9091}{rm}{41}%
}\hskip.1em\allowbreak{} 是\allowbreak{}\hskip.1em \RB{6.555}{5.310}{0.610}{%
\RG{0.000}{-0.001}{10.9091}{mi}{110}%
}\hskip.1em\allowbreak{} 阶\allowbreak{}\hskip.1em \RB{6.336}{8.070}{0.621}{%
\RG{0.000}{-0.001}{10.9091}{mi}{21}%
}\hskip.1em\allowbreak{} –矩阵，其行列式定义为

\[
\RB{230.316}{10.454}{13.839}{%
\RG{58.885}{9.086}{9.0910}{ex}{88}\RG{110.060}{4.500}{7.9701}{mi}{78}\RG{116.968}{4.500}{7.9701}{rm}{40}\RG{119.973}{4.500}{7.9701}{mi}{105}\RG{126.392}{4.500}{7.9701}{mi}{59}\RG{128.538}{4.500}{7.9701}{mi}{58}\RG{130.683}{4.500}{7.9701}{mi}{58}\RG{132.829}{4.500}{7.9701}{mi}{58}\RG{134.974}{4.500}{7.9701}{mi}{59}\RG{137.120}{4.500}{7.9701}{mi}{105}\RG{144.408}{4.500}{7.9701}{rm}{41}\RG{139.755}{3.590}{5.9776}{mi}{110}\RG{122.604}{3.489}{5.9776}{rm}{49}\RG{0.000}{-0.005}{10.9091}{sy}{106}\RG{3.030}{-0.005}{10.9091}{mi}{65}\RG{11.038}{-0.005}{10.9091}{rm}{40}\RG{15.192}{-0.005}{10.9091}{mi}{21}\RG{21.423}{-0.005}{10.9091}{rm}{41}\RG{25.577}{-0.005}{10.9091}{sy}{106}\RG{31.638}{-0.005}{10.9091}{rm}{61}\RG{87.926}{-0.005}{10.9091}{rm}{40}\RG{92.080}{-0.005}{10.9091}{sy}{0}\RG{100.564}{-0.005}{10.9091}{rm}{49}\RG{105.910}{-0.005}{10.9091}{rm}{41}\RG{147.868}{-0.005}{10.9091}{mi}{97}\RG{164.212}{-0.005}{10.9091}{rm}{40}\RG{168.366}{-0.005}{10.9091}{mi}{21}\RG{174.597}{-0.005}{10.9091}{rm}{41}\RG{180.570}{-0.005}{10.9091}{sy}{1}\RG{185.414}{-0.005}{10.9091}{sy}{1}\RG{190.268}{-0.005}{10.9091}{sy}{1}\RG{195.115}{-0.005}{10.9091}{mi}{97}\RG{213.149}{-0.005}{10.9091}{rm}{40}\RG{217.303}{-0.005}{10.9091}{mi}{21}\RG{223.535}{-0.005}{10.9091}{rm}{41}\RG{227.688}{-0.005}{10.9091}{mi}{59}\RG{153.476}{-1.642}{7.9701}{mi}{105}\RG{159.894}{-1.642}{7.9701}{rm}{49}\RG{200.723}{-1.642}{7.9701}{mi}{105}\RG{208.006}{-1.642}{7.9701}{mi}{110}\RG{203.354}{-2.550}{5.9776}{mi}{110}\RG{156.106}{-2.652}{5.9776}{rm}{49}\RG{42.976}{-11.520}{7.9701}{rm}{40}\RG{45.982}{-11.520}{7.9701}{mi}{105}\RG{52.400}{-11.520}{7.9701}{mi}{59}\RG{54.546}{-11.520}{7.9701}{mi}{58}\RG{56.691}{-11.520}{7.9701}{mi}{58}\RG{58.837}{-11.520}{7.9701}{mi}{58}\RG{60.982}{-11.520}{7.9701}{mi}{59}\RG{63.128}{-11.520}{7.9701}{mi}{105}\RG{70.417}{-11.520}{7.9701}{rm}{41}\RG{73.422}{-11.520}{7.9701}{sy}{50}\RG{78.574}{-11.520}{7.9701}{mi}{83}\RG{65.764}{-12.429}{5.9776}{mi}{110}\RG{83.273}{-12.429}{5.9776}{mi}{110}\RG{48.612}{-12.531}{5.9776}{rm}{49}%
}
\]
其中\allowbreak{}\hskip.1em \RB{11.520}{8.184}{2.205}{%
\RG{0.000}{0.004}{10.9091}{mi}{83}\RG{6.569}{-1.632}{7.9701}{mi}{110}%
}\hskip.1em\allowbreak{} 是\allowbreak{}\hskip.1em \RB{35.820}{7.758}{2.602}{%
\RG{0.000}{0.004}{10.9091}{rm}{49}\RG{5.341}{0.004}{10.9091}{mi}{59}\RG{10.119}{0.004}{10.9091}{mi}{58}\RG{14.908}{0.004}{10.9091}{mi}{58}\RG{19.687}{0.004}{10.9091}{mi}{58}\RG{24.476}{0.004}{10.9091}{mi}{59}\RG{29.265}{0.004}{10.9091}{mi}{110}%
}\hskip.1em\allowbreak{} 的全排列的集合，\allowbreak{}\hskip.1em \RB{58.100}{8.686}{3.235}{%
\RG{0.000}{0.004}{10.9091}{mi}{78}\RG{9.694}{0.004}{10.9091}{rm}{40}\RG{13.847}{0.004}{10.9091}{mi}{105}\RG{17.491}{-1.632}{7.9701}{rm}{49}\RG{21.810}{0.004}{10.9091}{mi}{59}\RG{26.588}{0.004}{10.9091}{mi}{58}\RG{31.377}{0.004}{10.9091}{mi}{58}\RG{36.156}{0.004}{10.9091}{mi}{58}\RG{40.945}{0.004}{10.9091}{mi}{59}\RG{45.734}{0.004}{10.9091}{mi}{105}\RG{49.381}{-1.632}{7.9701}{mi}{110}\RG{54.524}{0.004}{10.9091}{rm}{41}%
}\hskip.1em\allowbreak{} 是排列的逆序数。利用上述定义证明：\allowbreak{}\hskip.1em \RB{6.555}{5.316}{0.605}{%
\RG{0.000}{0.005}{10.9091}{mi}{110}%
}\hskip.1em\allowbreak{} 阶\allowbreak{}\hskip.1em \RB{6.336}{8.076}{0.615}{%
\RG{0.000}{0.005}{10.9091}{mi}{21}%
}\hskip.1em\allowbreak{} –矩阵的行列式满足九条性质，其中前八条参照《高等代数学》（谢启鸿、姚慕生、吴泉水，第四版）的§1.3和§1.4，第九条参照该教材的定理1.4.1和定理1.4.2。（2）证明：\allowbreak{}\hskip.1em \RB{6.336}{8.070}{0.621}{%
\RG{0.000}{-0.001}{10.9091}{mi}{21}%
}\hskip.1em\allowbreak{} –矩阵的行列式满足Laplace定理和Cauchy–Binet公式。特别地，对\allowbreak{}\hskip.1em \RB{6.555}{5.310}{0.610}{%
\RG{0.000}{-0.001}{10.9091}{mi}{110}%
}\hskip.1em\allowbreak{} 阶\allowbreak{}\hskip.1em \RB{6.336}{8.071}{0.620}{%
\RG{0.000}{-0.000}{10.9091}{mi}{21}%
}\hskip.1em\allowbreak{} –矩阵\allowbreak{}\hskip.1em \RB{49.959}{8.682}{3.238}{%
\RG{0.000}{-0.000}{10.9091}{mi}{65}\RG{8.007}{-0.000}{10.9091}{rm}{40}\RG{12.161}{-0.000}{10.9091}{mi}{21}\RG{18.393}{-0.000}{10.9091}{rm}{41}\RG{22.546}{-0.000}{10.9091}{mi}{59}\RG{27.325}{-0.000}{10.9091}{mi}{66}\RG{35.997}{-0.000}{10.9091}{rm}{40}\RG{40.151}{-0.000}{10.9091}{mi}{21}\RG{46.382}{-0.000}{10.9091}{rm}{41}%
}\hskip.1em\allowbreak{} ，有

\[
\RB{128.379}{8.682}{3.238}{%
\RG{0.000}{-0.000}{10.9091}{sy}{106}\RG{3.030}{-0.000}{10.9091}{mi}{65}\RG{11.038}{-0.000}{10.9091}{rm}{40}\RG{15.192}{-0.000}{10.9091}{mi}{21}\RG{21.423}{-0.000}{10.9091}{rm}{41}\RG{25.577}{-0.000}{10.9091}{mi}{66}\RG{34.243}{-0.000}{10.9091}{rm}{40}\RG{38.397}{-0.000}{10.9091}{mi}{21}\RG{44.627}{-0.000}{10.9091}{rm}{41}\RG{48.782}{-0.000}{10.9091}{sy}{106}\RG{54.842}{-0.000}{10.9091}{rm}{61}\RG{66.181}{-0.000}{10.9091}{sy}{106}\RG{69.211}{-0.000}{10.9091}{mi}{65}\RG{77.219}{-0.000}{10.9091}{rm}{40}\RG{81.373}{-0.000}{10.9091}{mi}{21}\RG{87.604}{-0.000}{10.9091}{rm}{41}\RG{91.758}{-0.000}{10.9091}{sy}{106}\RG{96.601}{-0.000}{10.9091}{sy}{106}\RG{99.637}{-0.000}{10.9091}{mi}{66}\RG{108.302}{-0.000}{10.9091}{rm}{40}\RG{112.456}{-0.000}{10.9091}{mi}{21}\RG{118.688}{-0.000}{10.9091}{rm}{41}\RG{122.842}{-0.000}{10.9091}{sy}{106}\RG{125.872}{-0.000}{10.9091}{mi}{58}%
}
\]
（3）设\allowbreak{}\hskip.1em \RB{26.222}{7.753}{1.996}{%
\RG{0.000}{-0.002}{10.9091}{mi}{110}\RG{9.406}{-0.002}{10.9091}{sy}{21}\RG{20.922}{-0.002}{10.9091}{rm}{50}%
}\hskip.1em\allowbreak{} ，\allowbreak{}\hskip.1em \RB{26.399}{8.680}{3.240}{%
\RG{0.000}{-0.002}{10.9091}{mi}{65}\RG{8.007}{-0.002}{10.9091}{rm}{40}\RG{12.162}{-0.002}{10.9091}{mi}{21}\RG{18.393}{-0.002}{10.9091}{rm}{41}\RG{22.546}{3.957}{7.9701}{sy}{3}%
}\hskip.1em\allowbreak{} 表示\allowbreak{}\hskip.1em \RB{21.969}{8.680}{3.240}{%
\RG{0.000}{-0.002}{10.9091}{mi}{65}\RG{8.008}{-0.002}{10.9091}{rm}{40}\RG{12.162}{-0.002}{10.9091}{mi}{21}\RG{18.393}{-0.002}{10.9091}{rm}{41}%
}\hskip.1em\allowbreak{} 的伴随矩阵。证明

\[
\RB{210.868}{9.296}{55.709}{%
\RG{87.398}{4.508}{7.9701}{sy}{3}\RG{128.632}{4.508}{7.9701}{sy}{3}\RT{8.966}{0.003}{10.9091}{(}\RT{13.210}{0.003}{10.9091}{3}\RT{18.665}{0.003}{10.9091}{.}\RT{21.697}{0.003}{10.9091}{1}\RT{27.152}{0.003}{10.9091}{)}\RG{42.305}{0.003}{10.9091}{mi}{65}\RG{50.312}{0.003}{10.9091}{rm}{40}\RG{54.466}{0.003}{10.9091}{mi}{21}\RG{60.698}{0.003}{10.9091}{rm}{41}\RG{64.851}{0.003}{10.9091}{mi}{65}\RG{72.859}{0.003}{10.9091}{rm}{40}\RG{77.013}{0.003}{10.9091}{mi}{21}\RG{83.244}{0.003}{10.9091}{rm}{41}\RG{94.747}{0.003}{10.9091}{rm}{61}\RG{106.085}{0.003}{10.9091}{mi}{65}\RG{114.093}{0.003}{10.9091}{rm}{40}\RG{118.247}{0.003}{10.9091}{mi}{21}\RG{124.478}{0.003}{10.9091}{rm}{41}\RG{132.950}{0.003}{10.9091}{mi}{65}\RG{140.958}{0.003}{10.9091}{rm}{40}\RG{145.112}{0.003}{10.9091}{mi}{21}\RG{151.343}{0.003}{10.9091}{rm}{41}\RG{158.521}{0.003}{10.9091}{rm}{61}\RG{169.866}{0.003}{10.9091}{sy}{106}\RG{172.897}{0.003}{10.9091}{mi}{65}\RG{180.904}{0.003}{10.9091}{rm}{40}\RG{185.058}{0.003}{10.9091}{mi}{21}\RG{191.289}{0.003}{10.9091}{rm}{41}\RG{195.443}{0.003}{10.9091}{sy}{106}\RG{198.473}{0.003}{10.9091}{mi}{73}\RG{208.328}{0.003}{10.9091}{rm}{59}\RG{203.185}{-1.633}{7.9701}{mi}{110}\RG{87.398}{-12.983}{7.9701}{sy}{3}\RG{129.290}{-12.983}{7.9701}{sy}{3}\RG{156.155}{-12.983}{7.9701}{sy}{3}\RT{0.000}{-17.488}{10.9091}{(}\RT{4.244}{-17.488}{10.9091}{3}\RT{9.698}{-17.488}{10.9091}{.}\RT{12.731}{-17.488}{10.9091}{2}\RT{18.186}{-17.488}{10.9091}{)}\RG{33.339}{-17.488}{10.9091}{rm}{40}\RG{37.493}{-17.488}{10.9091}{mi}{65}\RG{45.501}{-17.488}{10.9091}{rm}{40}\RG{49.654}{-17.488}{10.9091}{mi}{21}\RG{55.886}{-17.488}{10.9091}{rm}{41}\RG{60.040}{-17.488}{10.9091}{mi}{66}\RG{68.705}{-17.488}{10.9091}{rm}{40}\RG{72.859}{-17.488}{10.9091}{mi}{21}\RG{79.090}{-17.488}{10.9091}{rm}{41}\RG{83.247}{-17.488}{10.9091}{rm}{41}\RG{94.747}{-17.488}{10.9091}{rm}{61}\RG{106.085}{-17.488}{10.9091}{mi}{66}\RG{114.751}{-17.488}{10.9091}{rm}{40}\RG{118.905}{-17.488}{10.9091}{mi}{21}\RG{125.136}{-17.488}{10.9091}{rm}{41}\RG{133.608}{-17.488}{10.9091}{mi}{65}\RG{141.616}{-17.488}{10.9091}{rm}{40}\RG{145.769}{-17.488}{10.9091}{mi}{21}\RG{152.001}{-17.488}{10.9091}{rm}{41}\RG{160.474}{-17.488}{10.9091}{rm}{59}\RG{84.368}{-30.474}{7.9701}{sy}{3}\RG{134.692}{-30.474}{7.9701}{mi}{110}\RG{139.382}{-30.474}{7.9701}{sy}{0}\RG{145.392}{-30.474}{7.9701}{rm}{49}\RT{25.452}{-34.979}{10.9091}{(}\RT{29.696}{-34.979}{10.9091}{3}\RT{35.150}{-34.979}{10.9091}{.}\RT{38.183}{-34.979}{10.9091}{3}\RT{43.638}{-34.979}{10.9091}{)}\RG{58.790}{-34.979}{10.9091}{sy}{106}\RG{61.821}{-34.979}{10.9091}{mi}{65}\RG{69.829}{-34.979}{10.9091}{rm}{40}\RG{73.983}{-34.979}{10.9091}{mi}{21}\RG{80.214}{-34.979}{10.9091}{rm}{41}\RG{88.686}{-34.979}{10.9091}{sy}{106}\RG{94.747}{-34.979}{10.9091}{rm}{61}\RG{106.085}{-34.979}{10.9091}{sy}{106}\RG{109.116}{-34.979}{10.9091}{mi}{65}\RG{117.124}{-34.979}{10.9091}{rm}{40}\RG{121.278}{-34.979}{10.9091}{mi}{21}\RG{127.508}{-34.979}{10.9091}{rm}{41}\RG{131.663}{-34.979}{10.9091}{sy}{106}\RG{149.710}{-34.979}{10.9091}{rm}{59}\RG{78.926}{-47.965}{7.9701}{sy}{3}\RG{87.398}{-47.965}{7.9701}{sy}{3}\RG{134.692}{-47.965}{7.9701}{mi}{110}\RG{139.382}{-47.965}{7.9701}{sy}{0}\RG{145.392}{-47.965}{7.9701}{rm}{50}\RT{18.886}{-52.470}{10.9091}{(}\RT{23.130}{-52.470}{10.9091}{3}\RT{28.585}{-52.470}{10.9091}{.}\RT{31.617}{-52.470}{10.9091}{4}\RT{37.072}{-52.470}{10.9091}{)}\RG{52.225}{-52.470}{10.9091}{rm}{40}\RG{56.380}{-52.470}{10.9091}{mi}{65}\RG{64.387}{-52.470}{10.9091}{rm}{40}\RG{68.541}{-52.470}{10.9091}{mi}{21}\RG{74.772}{-52.470}{10.9091}{rm}{41}\RG{83.244}{-52.470}{10.9091}{rm}{41}\RG{94.747}{-52.470}{10.9091}{rm}{61}\RG{106.085}{-52.470}{10.9091}{sy}{106}\RG{109.116}{-52.470}{10.9091}{mi}{65}\RG{117.124}{-52.470}{10.9091}{rm}{40}\RG{121.278}{-52.470}{10.9091}{mi}{21}\RG{127.508}{-52.470}{10.9091}{rm}{41}\RG{131.663}{-52.470}{10.9091}{sy}{106}\RG{149.710}{-52.470}{10.9091}{mi}{65}\RG{157.718}{-52.470}{10.9091}{rm}{40}\RG{161.872}{-52.470}{10.9091}{mi}{21}\RG{168.102}{-52.470}{10.9091}{rm}{41}\RG{172.257}{-52.470}{10.9091}{mi}{58}%
}
\]
（4）设\allowbreak{}\hskip.1em \RB{10.035}{7.948}{0.503}{%
\RG{0.000}{-0.003}{10.9091}{mi}{75}%
}\hskip.1em\allowbreak{} 为数域，\allowbreak{}\hskip.1em \RB{54.420}{8.679}{3.241}{%
\RG{0.000}{-0.003}{10.9091}{mi}{65}\RG{11.038}{-0.003}{10.9091}{sy}{50}\RG{21.341}{-0.003}{10.9091}{mi}{77}\RG{31.675}{-1.639}{7.9701}{mi}{110}\RG{36.818}{-0.003}{10.9091}{rm}{40}\RG{40.973}{-0.003}{10.9091}{mi}{75}\RG{50.843}{-0.003}{10.9091}{rm}{41}%
}\hskip.1em\allowbreak{} 的特征多项式为\allowbreak{}\hskip.1em \RB{76.755}{8.679}{3.241}{%
\RG{0.000}{-0.003}{10.9091}{mi}{102}\RG{6.430}{-0.003}{10.9091}{rm}{40}\RG{10.584}{-0.003}{10.9091}{mi}{21}\RG{16.815}{-0.003}{10.9091}{rm}{41}\RG{23.993}{-0.003}{10.9091}{rm}{61}\RG{35.338}{-0.003}{10.9091}{sy}{106}\RG{38.368}{-0.003}{10.9091}{mi}{21}\RG{44.598}{-0.003}{10.9091}{mi}{73}\RG{49.311}{-1.639}{7.9701}{mi}{110}\RG{56.241}{-0.003}{10.9091}{sy}{0}\RG{66.513}{-0.003}{10.9091}{mi}{65}\RG{74.520}{-0.003}{10.9091}{sy}{106}%
}\hskip.1em\allowbreak{} 。对特征矩阵\allowbreak{}\hskip.1em \RB{36.347}{8.286}{2.212}{%
\RG{0.000}{-0.003}{10.9091}{mi}{21}\RG{6.229}{-0.003}{10.9091}{mi}{73}\RG{10.943}{-1.639}{7.9701}{mi}{110}\RG{17.872}{-0.003}{10.9091}{sy}{0}\RG{28.145}{-0.003}{10.9091}{mi}{65}%
}\hskip.1em\allowbreak{}利用（3.1），证明Cayley–Hamilton定理\allowbreak{}\hskip.1em \RB{42.503}{8.680}{3.240}{%
\RG{0.000}{-0.002}{10.9091}{mi}{102}\RG{6.430}{-0.002}{10.9091}{rm}{40}\RG{10.584}{-0.002}{10.9091}{mi}{65}\RG{18.592}{-0.002}{10.9091}{rm}{41}\RG{25.770}{-0.002}{10.9091}{rm}{61}\RG{37.115}{-0.002}{10.9091}{rm}{48}%
}\hskip.1em\allowbreak{} 。（5）设\allowbreak{}\hskip.1em \RB{21.969}{8.684}{3.236}{%
\RG{0.000}{0.002}{10.9091}{mi}{65}\RG{8.007}{0.002}{10.9091}{rm}{40}\RG{12.162}{0.002}{10.9091}{mi}{21}\RG{18.393}{0.002}{10.9091}{rm}{41}%
}\hskip.1em\allowbreak{} 为\allowbreak{}\hskip.1em \RB{6.555}{5.313}{0.607}{%
\RG{0.000}{0.002}{10.9091}{mi}{110}%
}\hskip.1em\allowbreak{} 阶\allowbreak{}\hskip.1em \RB{6.336}{8.073}{0.618}{%
\RG{0.000}{0.002}{10.9091}{mi}{21}%
}\hskip.1em\allowbreak{} –矩阵，证明下列结论等价：（5.1）\allowbreak{}\hskip.1em \RB{21.969}{8.684}{3.236}{%
\RG{0.000}{0.002}{10.9091}{mi}{65}\RG{8.008}{0.002}{10.9091}{rm}{40}\RG{12.162}{0.002}{10.9091}{mi}{21}\RG{18.393}{0.002}{10.9091}{rm}{41}%
}\hskip.1em\allowbreak{} 是可逆\allowbreak{}\hskip.1em \RB{6.336}{8.073}{0.618}{%
\RG{0.000}{0.002}{10.9091}{mi}{21}%
}\hskip.1em\allowbreak{} –矩阵；（5.2）\allowbreak{}\hskip.1em \RB{27.811}{8.684}{3.236}{%
\RG{0.000}{0.002}{10.9091}{sy}{106}\RG{3.030}{0.002}{10.9091}{mi}{65}\RG{11.038}{0.002}{10.9091}{rm}{40}\RG{15.192}{0.002}{10.9091}{mi}{21}\RG{21.423}{0.002}{10.9091}{rm}{41}\RG{25.577}{0.002}{10.9091}{sy}{106}%
}\hskip.1em\allowbreak{} 是非零常数；（5.3）\allowbreak{}\hskip.1em \RB{21.969}{8.684}{3.236}{%
\RG{0.000}{0.002}{10.9091}{mi}{65}\RG{8.007}{0.002}{10.9091}{rm}{40}\RG{12.161}{0.002}{10.9091}{mi}{21}\RG{18.393}{0.002}{10.9091}{rm}{41}%
}\hskip.1em\allowbreak{} 的相抵标准型是\allowbreak{}\hskip.1em \RB{9.662}{7.953}{2.207}{%
\RG{0.000}{0.002}{10.9091}{mi}{73}\RG{4.711}{-1.634}{7.9701}{mi}{110}%
}\hskip.1em\allowbreak{} ；（5.4）\allowbreak{}\hskip.1em \RB{21.969}{8.684}{3.236}{%
\RG{0.000}{0.002}{10.9091}{mi}{65}\RG{8.007}{0.002}{10.9091}{rm}{40}\RG{12.161}{0.002}{10.9091}{mi}{21}\RG{18.393}{0.002}{10.9091}{rm}{41}%
}\hskip.1em\allowbreak{} 只通过初等行变换或只通过初等列变换就可变为\allowbreak{}\hskip.1em \RB{9.662}{7.954}{2.206}{%
\RG{0.000}{0.003}{10.9091}{mi}{73}\RG{4.711}{-1.633}{7.9701}{mi}{110}%
}\hskip.1em\allowbreak{} ；（5.5）\allowbreak{}\hskip.1em \RB{21.969}{8.685}{3.235}{%
\RG{0.000}{0.003}{10.9091}{mi}{65}\RG{8.007}{0.003}{10.9091}{rm}{40}\RG{12.161}{0.003}{10.9091}{mi}{21}\RG{18.393}{0.003}{10.9091}{rm}{41}%
}\hskip.1em\allowbreak{} 是有限个初等\allowbreak{}\hskip.1em \RB{6.336}{8.074}{0.617}{%
\RG{0.000}{0.003}{10.9091}{mi}{21}%
}\hskip.1em\allowbreak{} –矩阵的乘积。上述条件成立时，证明\allowbreak{}\hskip.1em \RB{107.406}{0.500}{27.787}{%
\RG{0.000}{-17.065}{10.9091}{mi}{65}\RG{8.007}{-17.065}{10.9091}{rm}{40}\RG{12.162}{-17.065}{10.9091}{mi}{21}\RG{18.393}{-17.065}{10.9091}{rm}{41}\RG{22.546}{-12.561}{7.9701}{sy}{0}\RG{28.557}{-12.561}{7.9701}{rm}{49}\RG{35.906}{-17.065}{10.9091}{rm}{61}\RG{48.335}{-24.548}{10.9091}{sy}{106}\RG{51.365}{-24.548}{10.9091}{mi}{65}\RG{59.374}{-24.548}{10.9091}{rm}{40}\RG{59.969}{-9.684}{10.9091}{rm}{49}\RG{63.527}{-24.548}{10.9091}{mi}{21}\RG{69.758}{-24.548}{10.9091}{rm}{41}\RG{73.913}{-24.548}{10.9091}{sy}{106}\RG{78.034}{-17.065}{10.9091}{mi}{65}\RG{86.041}{-17.065}{10.9091}{rm}{40}\RG{90.196}{-17.065}{10.9091}{mi}{21}\RG{96.426}{-17.065}{10.9091}{rm}{41}\RG{100.580}{-12.561}{7.9701}{sy}{3}\RG{104.899}{-17.065}{10.9091}{mi}{58}\RR{48.335}{-14.337}{28.608}{0.364}%
}\hskip.1em\allowbreak{}

证明：设\allowbreak{}\hskip.1em \RB{10.035}{7.949}{0.502}{%
\RG{0.000}{-0.002}{10.9091}{mi}{75}%
}\hskip.1em\allowbreak{} 为系数数域，\allowbreak{}\hskip.1em \RB{46.797}{8.680}{3.229}{%
\RG{0.000}{-0.002}{10.9091}{mi}{82}\RG{12.693}{-0.002}{10.9091}{rm}{61}\RG{25.474}{-0.002}{10.9091}{mi}{75}\RG{35.345}{-0.002}{10.9091}{rm}{91}\RG{38.376}{-0.002}{10.9091}{mi}{21}\RG{44.606}{-0.002}{10.9091}{rm}{93}%
}\hskip.1em\allowbreak{} ，\allowbreak{}\hskip.1em \RB{49.484}{8.680}{3.240}{%
\RG{0.000}{-0.002}{10.9091}{mi}{70}\RG{12.871}{-0.002}{10.9091}{rm}{61}\RG{25.652}{-0.002}{10.9091}{mi}{75}\RG{35.523}{-0.002}{10.9091}{rm}{40}\RG{39.677}{-0.002}{10.9091}{mi}{21}\RG{45.908}{-0.002}{10.9091}{rm}{41}%
}\hskip.1em\allowbreak{} 为\allowbreak{}\hskip.1em \RB{8.573}{7.949}{0.731}{%
\RG{0.000}{-0.002}{10.9091}{mi}{82}%
}\hskip.1em\allowbreak{} 的分式域。对非负整数\allowbreak{}\hskip.1em \RB{16.340}{5.309}{2.607}{%
\RG{0.000}{-0.002}{10.9091}{mi}{117}\RG{6.076}{-0.002}{10.9091}{mi}{59}\RG{10.855}{-0.002}{10.9091}{mi}{118}%
}\hskip.1em\allowbreak{} ，记\allowbreak{}\hskip.1em \RB{37.384}{8.680}{3.542}{%
\RG{0.000}{-0.002}{10.9091}{mi}{77}\RG{10.334}{-1.638}{7.9701}{mi}{117}\RG{14.807}{-1.638}{7.9701}{mi}{59}\RG{16.952}{-1.638}{7.9701}{mi}{118}\RG{21.433}{-0.002}{10.9091}{rm}{40}\RG{25.587}{-0.002}{10.9091}{mi}{82}\RG{33.808}{-0.002}{10.9091}{rm}{41}%
}\hskip.1em\allowbreak{} 为\allowbreak{}\hskip.1em \RB{8.573}{7.949}{0.731}{%
\RG{0.000}{-0.002}{10.9091}{mi}{82}%
}\hskip.1em\allowbreak{} 上的\allowbreak{}\hskip.1em \RB{24.865}{5.855}{0.611}{%
\RG{0.000}{-0.002}{10.9091}{mi}{117}\RG{8.486}{-0.002}{10.9091}{sy}{2}\RG{19.379}{-0.002}{10.9091}{mi}{118}%
}\hskip.1em\allowbreak{} 矩阵全体，\allowbreak{}\hskip.1em \RB{83.993}{8.680}{3.542}{%
\RG{0.000}{-0.002}{10.9091}{mi}{77}\RG{10.334}{-1.638}{7.9701}{mi}{117}\RG{15.263}{-0.002}{10.9091}{rm}{40}\RG{19.416}{-0.002}{10.9091}{mi}{82}\RG{27.637}{-0.002}{10.9091}{rm}{41}\RG{34.816}{-0.002}{10.9091}{rm}{61}\RG{46.159}{-0.002}{10.9091}{mi}{77}\RG{56.494}{-1.638}{7.9701}{mi}{117}\RG{60.966}{-1.638}{7.9701}{mi}{59}\RG{63.112}{-1.638}{7.9701}{mi}{117}\RG{68.043}{-0.002}{10.9091}{rm}{40}\RG{72.197}{-0.002}{10.9091}{mi}{82}\RG{80.417}{-0.002}{10.9091}{rm}{41}%
}\hskip.1em\allowbreak{} ，\allowbreak{}\hskip.1em \RB{9.458}{7.949}{2.211}{%
\RG{0.000}{-0.002}{10.9091}{mi}{73}\RG{4.710}{-1.638}{7.9701}{mi}{117}%
}\hskip.1em\allowbreak{} 为\allowbreak{}\hskip.1em \RB{6.260}{5.309}{0.611}{%
\RG{0.000}{-0.002}{10.9091}{mi}{117}%
}\hskip.1em\allowbreak{} 阶单位矩阵。空矩阵的行列式约定为\allowbreak{}\hskip.1em \RB{4.973}{7.754}{0.501}{%
\RG{0.000}{-0.001}{10.9091}{rm}{49}%
}\hskip.1em\allowbreak{} 。在（1）至（3）及（5）中略去矩阵元素的自变量\allowbreak{}\hskip.1em \RB{6.336}{8.070}{0.621}{%
\RG{0.000}{-0.001}{10.9091}{mi}{21}%
}\hskip.1em\allowbreak{} ，以\allowbreak{}\hskip.1em \RB{24.856}{8.288}{0.610}{%
\RG{0.000}{-0.001}{10.9091}{rm}{100}\RG{5.935}{-0.001}{10.9091}{rm}{101}\RG{10.680}{-0.001}{10.9091}{rm}{116}\RG{16.654}{-0.001}{10.9091}{mi}{65}%
}\hskip.1em\allowbreak{} 表示\allowbreak{}\hskip.1em \RB{27.812}{8.681}{3.239}{%
\RG{0.000}{-0.001}{10.9091}{sy}{106}\RG{3.030}{-0.001}{10.9091}{mi}{65}\RG{11.038}{-0.001}{10.9091}{rm}{40}\RG{15.192}{-0.001}{10.9091}{mi}{21}\RG{21.423}{-0.001}{10.9091}{rm}{41}\RG{25.578}{-0.001}{10.9091}{sy}{106}%
}\hskip.1em\allowbreak{} 。（1）对置换\allowbreak{}\hskip.1em \RB{34.536}{8.182}{2.207}{%
\RG{0.000}{0.002}{10.9091}{mi}{27}\RG{11.103}{0.002}{10.9091}{sy}{50}\RG{23.015}{0.002}{10.9091}{mi}{83}\RG{29.585}{-1.634}{7.9701}{mi}{110}%
}\hskip.1em\allowbreak{} ，记其符号为\allowbreak{}\hskip.1em \RB{76.829}{9.916}{3.236}{%
\RG{0.000}{0.002}{10.9091}{mi}{34}\RG{4.999}{0.002}{10.9091}{rm}{40}\RG{9.153}{0.002}{10.9091}{mi}{27}\RG{15.616}{0.002}{10.9091}{rm}{41}\RG{24.409}{0.002}{10.9091}{rm}{61}\RG{37.358}{0.002}{10.9091}{rm}{40}\RG{41.511}{0.002}{10.9091}{sy}{0}\RG{49.997}{0.002}{10.9091}{rm}{49}\RG{55.342}{0.002}{10.9091}{rm}{41}\RG{59.491}{3.961}{7.9701}{mi}{78}\RG{66.399}{3.961}{7.9701}{rm}{40}\RG{69.405}{3.961}{7.9701}{mi}{27}\RG{74.081}{3.961}{7.9701}{rm}{41}%
}\hskip.1em\allowbreak{} 。以\allowbreak{}\hskip.1em \RB{16.224}{9.298}{0.607}{%
\RG{0.000}{0.002}{10.9091}{mi}{27}\RG{6.462}{3.961}{7.9701}{sy}{0}\RG{12.473}{3.961}{7.9701}{rm}{49}%
}\hskip.1em\allowbreak{} 替换求和指标，并利用\allowbreak{}\hskip.1em \RB{63.659}{9.299}{3.235}{%
\RG{0.000}{0.003}{10.9091}{mi}{34}\RG{4.999}{0.003}{10.9091}{rm}{40}\RG{9.153}{0.003}{10.9091}{mi}{27}\RG{15.616}{3.962}{7.9701}{sy}{0}\RG{21.626}{3.962}{7.9701}{rm}{49}\RG{25.945}{0.003}{10.9091}{rm}{41}\RG{33.123}{0.003}{10.9091}{rm}{61}\RG{44.467}{0.003}{10.9091}{mi}{34}\RG{49.465}{0.003}{10.9091}{rm}{40}\RG{53.620}{0.003}{10.9091}{mi}{27}\RG{60.083}{0.003}{10.9091}{rm}{41}%
}\hskip.1em\allowbreak{} ，得到

\[
\RB{122.342}{15.529}{45.178}{%
\RG{85.982}{11.822}{7.9701}{mi}{110}\RG{42.056}{9.094}{9.0910}{ex}{88}\RG{82.519}{9.094}{9.0910}{ex}{89}\RG{0.000}{0.004}{10.9091}{rm}{100}\RG{5.935}{0.004}{10.9091}{rm}{101}\RG{10.680}{0.004}{10.9091}{rm}{116}\RG{16.654}{0.004}{10.9091}{mi}{65}\RG{27.692}{0.004}{10.9091}{rm}{61}\RG{60.031}{0.004}{10.9091}{mi}{34}\RG{65.029}{0.004}{10.9091}{rm}{40}\RG{69.183}{0.004}{10.9091}{mi}{27}\RG{75.645}{0.004}{10.9091}{rm}{41}\RG{96.854}{0.004}{10.9091}{mi}{97}\RG{102.461}{-1.694}{7.9701}{mi}{27}\RG{107.138}{-1.694}{7.9701}{rm}{40}\RG{110.143}{-1.694}{7.9701}{mi}{106}\RG{113.688}{-1.694}{7.9701}{rm}{41}\RG{116.693}{-1.694}{7.9701}{mi}{59}\RG{118.838}{-1.694}{7.9701}{mi}{106}\RG{81.618}{-10.865}{7.9701}{mi}{106}\RG{85.162}{-10.865}{7.9701}{rm}{61}\RG{91.169}{-10.865}{7.9701}{rm}{49}\RG{39.031}{-11.027}{7.9701}{mi}{27}\RG{43.708}{-11.027}{7.9701}{sy}{50}\RG{48.860}{-11.027}{7.9701}{mi}{83}\RG{53.559}{-11.937}{5.9776}{mi}{110}\RG{84.512}{-20.865}{7.9701}{mi}{110}\RG{41.851}{-23.593}{9.0910}{ex}{88}\RG{81.048}{-23.593}{9.0910}{ex}{89}\RG{27.692}{-32.684}{10.9091}{rm}{61}\RG{59.623}{-32.684}{10.9091}{mi}{34}\RG{64.621}{-32.684}{10.9091}{rm}{40}\RG{68.775}{-32.684}{10.9091}{mi}{28}\RG{74.633}{-32.684}{10.9091}{rm}{41}\RG{94.928}{-32.684}{10.9091}{mi}{97}\RG{118.677}{-32.684}{10.9091}{mi}{58}\RG{100.535}{-34.381}{7.9701}{mi}{105}\RG{103.168}{-34.381}{7.9701}{mi}{59}\RG{105.313}{-34.381}{7.9701}{mi}{28}\RG{109.581}{-34.381}{7.9701}{rm}{40}\RG{112.587}{-34.381}{7.9701}{mi}{105}\RG{115.217}{-34.381}{7.9701}{rm}{41}\RG{80.605}{-43.552}{7.9701}{mi}{105}\RG{83.236}{-43.552}{7.9701}{rm}{61}\RG{89.243}{-43.552}{7.9701}{rm}{49}\RG{39.031}{-43.715}{7.9701}{mi}{28}\RG{43.300}{-43.715}{7.9701}{sy}{50}\RG{48.452}{-43.715}{7.9701}{mi}{83}\RG{53.151}{-44.624}{5.9776}{mi}{110}%
}
\]
各行性质对应的列性质均可由转置得到。记\allowbreak{}\hskip.1em \RB{7.619}{5.307}{2.213}{%
\RG{0.000}{-0.004}{10.9091}{mi}{114}\RG{4.813}{-1.641}{7.9701}{mi}{105}%
}\hskip.1em\allowbreak{} 为\allowbreak{}\hskip.1em \RB{8.202}{8.286}{0.504}{%
\RG{0.000}{-0.004}{10.9091}{mi}{65}%
}\hskip.1em\allowbreak{} 的第\allowbreak{}\hskip.1em \RB{3.587}{7.664}{0.613}{%
\RG{0.000}{-0.004}{10.9091}{mi}{105}%
}\hskip.1em\allowbreak{} 行。对行向量\allowbreak{}\hskip.1em \RB{44.368}{9.291}{0.940}{%
\RG{0.000}{-0.004}{10.9091}{mi}{117}\RG{10.064}{-0.004}{10.9091}{sy}{50}\RG{21.323}{-0.004}{10.9091}{mi}{82}\RG{29.543}{3.955}{7.9701}{rm}{49}\RG{33.406}{3.955}{7.9701}{sy}{2}\RG{39.417}{3.955}{7.9701}{mi}{110}%
}\hskip.1em\allowbreak{} ，记\allowbreak{}\hskip.1em \RB{23.113}{8.286}{2.242}{%
\RG{0.000}{-0.004}{10.9091}{mi}{65}\RG{8.008}{-1.640}{7.9701}{mi}{105}\RG{10.639}{-1.640}{7.9701}{sy}{32}\RG{18.366}{-1.640}{7.9701}{mi}{117}%
}\hskip.1em\allowbreak{} 为用\allowbreak{}\hskip.1em \RB{6.260}{5.307}{0.613}{%
\RG{0.000}{-0.004}{10.9091}{mi}{117}%
}\hskip.1em\allowbreak{} 替换\allowbreak{}\hskip.1em \RB{8.202}{8.286}{0.504}{%
\RG{0.000}{-0.004}{10.9091}{mi}{65}%
}\hskip.1em\allowbreak{} 的第\allowbreak{}\hskip.1em \RB{3.587}{7.664}{0.613}{%
\RG{0.000}{-0.004}{10.9091}{mi}{105}%
}\hskip.1em\allowbreak{} 行所得的矩阵。

1.转置不改变行列式。将定义中的\allowbreak{}\hskip.1em \RB{11.744}{5.316}{3.615}{%
\RG{0.000}{0.005}{10.9091}{mi}{97}\RG{5.607}{-1.632}{7.9701}{mi}{105}\RG{8.240}{-1.632}{7.9701}{mi}{106}%
}\hskip.1em\allowbreak{} 换为\allowbreak{}\hskip.1em \RB{11.955}{5.316}{3.615}{%
\RG{0.000}{0.005}{10.9091}{mi}{97}\RG{5.607}{-1.632}{7.9701}{mi}{106}\RG{9.149}{-1.632}{7.9701}{mi}{105}%
}\hskip.1em\allowbreak{} ，所得求和式即为上述第二个表达式，故

\[
\RB{80.833}{9.452}{2.686}{%
\RG{24.662}{3.956}{7.9701}{rm}{84}\RG{0.000}{-0.004}{10.9091}{rm}{100}\RG{5.935}{-0.004}{10.9091}{rm}{101}\RG{10.680}{-0.004}{10.9091}{rm}{116}\RG{16.653}{-0.004}{10.9091}{mi}{65}\RG{33.424}{-0.004}{10.9091}{rm}{61}\RG{44.758}{-0.004}{10.9091}{rm}{100}\RG{50.693}{-0.004}{10.9091}{rm}{101}\RG{55.438}{-0.004}{10.9091}{rm}{116}\RG{61.417}{-0.004}{10.9091}{mi}{65}\RT{69.424}{-0.004}{10.9091}{。}%
}
\]
% Original page 5.
2.交换两行使行列式变号。设\allowbreak{}\hskip.1em \RB{5.856}{5.314}{2.843}{%
\RG{0.000}{0.003}{10.9091}{mi}{26}%
}\hskip.1em\allowbreak{} 是交换两个不同指标的对换，\allowbreak{}\hskip.1em \RB{8.606}{7.954}{0.500}{%
\RG{0.000}{0.003}{10.9091}{mi}{66}%
}\hskip.1em\allowbreak{} 是对\allowbreak{}\hskip.1em \RB{8.202}{8.292}{0.500}{%
\RG{0.000}{0.003}{10.9091}{mi}{65}%
}\hskip.1em\allowbreak{} 作相应行交换所得的矩阵。由\allowbreak{}\hskip.1em \RB{49.395}{8.074}{4.012}{%
\RG{0.000}{0.003}{10.9091}{mi}{98}\RG{4.542}{-1.633}{7.9701}{mi}{105}\RG{7.174}{-1.633}{7.9701}{mi}{106}\RG{14.201}{0.003}{10.9091}{rm}{61}\RG{25.540}{0.003}{10.9091}{mi}{97}\RG{31.147}{-1.694}{7.9701}{mi}{26}\RG{35.104}{-1.694}{7.9701}{rm}{40}\RG{38.109}{-1.694}{7.9701}{mi}{105}\RG{40.741}{-1.694}{7.9701}{rm}{41}\RG{43.746}{-1.694}{7.9701}{mi}{59}\RG{45.891}{-1.694}{7.9701}{mi}{106}%
}\hskip.1em\allowbreak{} 及置换符号的乘法公式，有

\[
\RB{184.218}{48.216}{12.850}{%
\RG{86.639}{44.509}{7.9701}{mi}{110}\RG{42.713}{41.781}{9.0910}{ex}{88}\RG{83.176}{41.781}{9.0910}{ex}{89}\RG{0.000}{32.690}{10.9091}{rm}{100}\RG{5.935}{32.690}{10.9091}{rm}{101}\RG{10.680}{32.690}{10.9091}{rm}{116}\RG{16.653}{32.690}{10.9091}{mi}{66}\RG{28.349}{32.690}{10.9091}{rm}{61}\RG{60.688}{32.690}{10.9091}{mi}{34}\RG{65.686}{32.690}{10.9091}{rm}{40}\RG{69.840}{32.690}{10.9091}{mi}{27}\RG{76.302}{32.690}{10.9091}{rm}{41}\RG{97.511}{32.690}{10.9091}{mi}{97}\RG{103.118}{30.993}{7.9701}{mi}{26}\RG{107.075}{30.993}{7.9701}{mi}{27}\RG{111.753}{30.993}{7.9701}{rm}{40}\RG{114.758}{30.993}{7.9701}{mi}{106}\RG{118.302}{30.993}{7.9701}{rm}{41}\RG{121.307}{30.993}{7.9701}{mi}{59}\RG{123.452}{30.993}{7.9701}{mi}{106}\RG{82.275}{21.821}{7.9701}{mi}{106}\RG{85.819}{21.821}{7.9701}{rm}{61}\RG{91.826}{21.821}{7.9701}{rm}{49}\RG{39.688}{21.659}{7.9701}{mi}{27}\RG{44.365}{21.659}{7.9701}{sy}{50}\RG{49.517}{21.659}{7.9701}{mi}{83}\RG{54.216}{20.751}{5.9776}{mi}{110}\RG{95.930}{11.821}{7.9701}{mi}{110}\RG{52.812}{9.094}{9.0910}{ex}{88}\RG{92.466}{9.094}{9.0910}{ex}{89}\RG{28.349}{0.003}{10.9091}{rm}{61}\RG{39.688}{0.003}{10.9091}{sy}{0}\RG{70.583}{0.003}{10.9091}{mi}{34}\RG{75.582}{0.003}{10.9091}{rm}{40}\RG{79.735}{0.003}{10.9091}{mi}{28}\RG{85.593}{0.003}{10.9091}{rm}{41}\RG{106.801}{0.003}{10.9091}{mi}{97}\RG{135.408}{0.003}{10.9091}{rm}{61}\RG{146.747}{0.003}{10.9091}{sy}{0}\RG{157.050}{0.003}{10.9091}{rm}{100}\RG{162.984}{0.003}{10.9091}{rm}{101}\RG{167.730}{0.003}{10.9091}{rm}{116}\RG{173.703}{0.003}{10.9091}{mi}{65}\RG{181.711}{0.003}{10.9091}{mi}{58}\RG{112.409}{-1.695}{7.9701}{mi}{28}\RG{116.678}{-1.695}{7.9701}{rm}{40}\RG{119.682}{-1.695}{7.9701}{mi}{106}\RG{123.227}{-1.695}{7.9701}{rm}{41}\RG{126.232}{-1.695}{7.9701}{mi}{59}\RG{128.378}{-1.695}{7.9701}{mi}{106}\RG{91.565}{-10.866}{7.9701}{mi}{106}\RG{95.109}{-10.866}{7.9701}{rm}{61}\RG{101.117}{-10.866}{7.9701}{rm}{49}\RG{49.991}{-11.028}{7.9701}{mi}{28}\RG{54.261}{-11.028}{7.9701}{sy}{50}\RG{59.412}{-11.028}{7.9701}{mi}{83}\RG{64.111}{-11.937}{5.9776}{mi}{110}%
}
\]
3.两行相同时，行列式为零。交换这两行不改变矩阵，而由上一条使其行列式变号，故\allowbreak{}\hskip.1em \RB{51.573}{8.286}{0.732}{%
\RG{0.000}{-0.003}{10.9091}{rm}{50}\RG{7.156}{-0.003}{10.9091}{rm}{100}\RG{13.091}{-0.003}{10.9091}{rm}{101}\RG{17.837}{-0.003}{10.9091}{rm}{116}\RG{23.813}{-0.003}{10.9091}{mi}{65}\RG{34.851}{-0.003}{10.9091}{rm}{61}\RG{46.185}{-0.003}{10.9091}{rm}{48}%
}\hskip.1em\allowbreak{} 。数域\allowbreak{}\hskip.1em \RB{10.035}{7.948}{0.503}{%
\RG{0.000}{-0.003}{10.9091}{mi}{75}%
}\hskip.1em\allowbreak{} 的特征为零，因而\allowbreak{}\hskip.1em \RB{44.414}{8.286}{0.732}{%
\RG{0.000}{-0.003}{10.9091}{rm}{100}\RG{5.935}{-0.003}{10.9091}{rm}{101}\RG{10.680}{-0.003}{10.9091}{rm}{116}\RG{16.653}{-0.003}{10.9091}{mi}{65}\RG{27.692}{-0.003}{10.9091}{rm}{61}\RG{39.027}{-0.003}{10.9091}{rm}{48}%
}\hskip.1em\allowbreak{} 。4.某行的公因子可以提出。对\allowbreak{}\hskip.1em \RB{28.055}{8.068}{0.939}{%
\RG{0.000}{-0.003}{10.9091}{mi}{104}\RG{9.179}{-0.003}{10.9091}{sy}{50}\RG{19.482}{-0.003}{10.9091}{mi}{82}%
}\hskip.1em\allowbreak{} ，定义式的每一项恰含第\allowbreak{}\hskip.1em \RB{3.587}{7.665}{0.612}{%
\RG{0.000}{-0.003}{10.9091}{mi}{105}%
}\hskip.1em\allowbreak{} 行的一个元素，故

\[
\RB{95.805}{8.286}{3.303}{%
\RG{0.000}{-0.003}{10.9091}{rm}{100}\RG{5.935}{-0.003}{10.9091}{rm}{101}\RG{10.680}{-0.003}{10.9091}{rm}{116}\RG{16.654}{-0.003}{10.9091}{mi}{65}\RG{49.330}{-0.003}{10.9091}{rm}{61}\RG{60.668}{-0.003}{10.9091}{mi}{104}\RG{68.636}{-0.003}{10.9091}{rm}{100}\RG{74.570}{-0.003}{10.9091}{rm}{101}\RG{79.316}{-0.003}{10.9091}{rm}{116}\RG{85.290}{-0.003}{10.9091}{mi}{65}\RG{93.297}{-0.003}{10.9091}{mi}{58}\RG{24.662}{-1.640}{7.9701}{mi}{105}\RG{27.292}{-1.640}{7.9701}{sy}{32}\RG{35.020}{-1.640}{7.9701}{mi}{104}\RG{39.471}{-1.640}{7.9701}{mi}{114}\RG{42.960}{-2.748}{5.9776}{mi}{105}%
}
\]
取\allowbreak{}\hskip.1em \RB{25.901}{8.067}{0.733}{%
\RG{0.000}{-0.004}{10.9091}{mi}{104}\RG{9.179}{-0.004}{10.9091}{rm}{61}\RG{20.514}{-0.004}{10.9091}{rm}{48}%
}\hskip.1em\allowbreak{} 可知，含零行的矩阵的行列式为零。5.行列式对每一行具有可加性。对\allowbreak{}\hskip.1em \RB{52.799}{9.292}{2.608}{%
\RG{0.000}{-0.003}{10.9091}{mi}{117}\RG{6.076}{-0.003}{10.9091}{mi}{59}\RG{10.855}{-0.003}{10.9091}{mi}{118}\RG{19.451}{-0.003}{10.9091}{sy}{50}\RG{29.754}{-0.003}{10.9091}{mi}{82}\RG{37.974}{3.956}{7.9701}{rm}{49}\RG{41.838}{3.956}{7.9701}{sy}{2}\RG{47.848}{3.956}{7.9701}{mi}{110}%
}\hskip.1em\allowbreak{} ，在定义式中对第\allowbreak{}\hskip.1em \RB{3.587}{7.665}{0.612}{%
\RG{0.000}{-0.003}{10.9091}{mi}{105}%
}\hskip.1em\allowbreak{} 行的元素使用分配律，得到

\[
\RB{159.462}{8.287}{2.894}{%
\RG{0.000}{-0.002}{10.9091}{rm}{100}\RG{5.935}{-0.002}{10.9091}{rm}{101}\RG{10.680}{-0.002}{10.9091}{rm}{116}\RG{16.653}{-0.002}{10.9091}{mi}{65}\RG{53.012}{-0.002}{10.9091}{rm}{61}\RG{64.347}{-0.002}{10.9091}{rm}{100}\RG{70.282}{-0.002}{10.9091}{rm}{101}\RG{75.027}{-0.002}{10.9091}{rm}{116}\RG{81.005}{-0.002}{10.9091}{mi}{65}\RG{106.724}{-0.002}{10.9091}{rm}{43}\RG{117.459}{-0.002}{10.9091}{rm}{100}\RG{123.393}{-0.002}{10.9091}{rm}{101}\RG{128.139}{-0.002}{10.9091}{rm}{116}\RG{134.110}{-0.002}{10.9091}{mi}{65}\RG{156.955}{-0.002}{10.9091}{mi}{58}\RG{24.662}{-1.638}{7.9701}{mi}{105}\RG{27.292}{-1.638}{7.9701}{sy}{32}\RG{35.019}{-1.638}{7.9701}{mi}{117}\RG{39.494}{-1.638}{7.9701}{rm}{43}\RG{45.503}{-1.638}{7.9701}{mi}{118}\RG{89.013}{-1.638}{7.9701}{mi}{105}\RG{91.644}{-1.638}{7.9701}{sy}{32}\RG{99.371}{-1.638}{7.9701}{mi}{117}\RG{142.118}{-1.638}{7.9701}{mi}{105}\RG{144.749}{-1.638}{7.9701}{sy}{32}\RG{152.476}{-1.638}{7.9701}{mi}{118}%
}
\]
结合上一条，行列式对每一行均为\allowbreak{}\hskip.1em \RB{8.573}{7.954}{0.726}{%
\RG{0.000}{0.003}{10.9091}{mi}{82}%
}\hskip.1em\allowbreak{} –线性。6.两行成比例时，行列式为零。设\allowbreak{}\hskip.1em \RB{23.271}{8.315}{2.853}{%
\RG{0.000}{0.004}{10.9091}{mi}{105}\RG{6.921}{0.004}{10.9091}{sy}{54}\RG{6.921}{0.004}{10.9091}{rm}{61}\RG{18.506}{0.004}{10.9091}{mi}{106}%
}\hskip.1em\allowbreak{} ，并有\allowbreak{}\hskip.1em \RB{39.254}{8.075}{2.602}{%
\RG{0.000}{0.004}{10.9091}{mi}{104}\RG{6.153}{0.004}{10.9091}{mi}{59}\RG{10.931}{0.004}{10.9091}{mi}{107}\RG{20.133}{0.004}{10.9091}{sy}{50}\RG{30.681}{0.004}{10.9091}{mi}{82}%
}\hskip.1em\allowbreak{} 满足\allowbreak{}\hskip.1em \RB{26.398}{8.315}{2.853}{%
\RG{0.000}{0.004}{10.9091}{mi}{104}\RG{9.425}{0.004}{10.9091}{sy}{54}\RG{9.425}{0.004}{10.9091}{rm}{61}\RG{21.011}{0.004}{10.9091}{rm}{48}%
}\hskip.1em\allowbreak{} 及\allowbreak{}\hskip.1em \RB{43.149}{8.075}{3.616}{%
\RG{0.000}{0.004}{10.9091}{mi}{104}\RG{6.153}{0.004}{10.9091}{mi}{114}\RG{10.963}{-1.632}{7.9701}{mi}{105}\RG{17.325}{0.004}{10.9091}{rm}{61}\RG{28.910}{0.004}{10.9091}{mi}{107}\RG{34.823}{0.004}{10.9091}{mi}{114}\RG{39.645}{-1.632}{7.9701}{mi}{106}%
}\hskip.1em\allowbreak{} 。由提出公因子及相同行性质，有

\[
\RB{180.146}{8.293}{4.354}{%
\RG{0.000}{0.004}{10.9091}{mi}{104}\RG{7.967}{0.004}{10.9091}{rm}{100}\RG{13.902}{0.004}{10.9091}{rm}{101}\RG{18.647}{0.004}{10.9091}{rm}{116}\RG{24.621}{0.004}{10.9091}{mi}{65}\RG{35.659}{0.004}{10.9091}{rm}{61}\RG{46.994}{0.004}{10.9091}{rm}{100}\RG{52.929}{0.004}{10.9091}{rm}{101}\RG{57.674}{0.004}{10.9091}{rm}{116}\RG{63.651}{0.004}{10.9091}{mi}{65}\RG{96.328}{0.004}{10.9091}{rm}{61}\RG{107.666}{0.004}{10.9091}{mi}{107}\RG{115.406}{0.004}{10.9091}{rm}{100}\RG{121.341}{0.004}{10.9091}{rm}{101}\RG{126.087}{0.004}{10.9091}{rm}{116}\RG{132.061}{0.004}{10.9091}{mi}{65}\RG{160.960}{0.004}{10.9091}{rm}{61}\RG{172.295}{0.004}{10.9091}{rm}{48}\RG{177.639}{0.004}{10.9091}{mi}{58}\RG{71.660}{-1.632}{7.9701}{mi}{105}\RG{74.290}{-1.632}{7.9701}{sy}{32}\RG{82.018}{-1.632}{7.9701}{mi}{104}\RG{86.468}{-1.632}{7.9701}{mi}{114}\RG{140.068}{-1.632}{7.9701}{mi}{105}\RG{142.699}{-1.632}{7.9701}{sy}{32}\RG{150.427}{-1.632}{7.9701}{mi}{114}\RG{89.958}{-2.741}{5.9776}{mi}{105}\RG{153.914}{-2.741}{5.9776}{mi}{106}%
}
\]
由于\allowbreak{}\hskip.1em \RB{8.573}{7.954}{0.726}{%
\RG{0.000}{0.003}{10.9091}{mi}{82}%
}\hskip.1em\allowbreak{} 是整环，故\allowbreak{}\hskip.1em \RB{44.414}{8.292}{0.726}{%
\RG{0.000}{0.003}{10.9091}{rm}{100}\RG{5.935}{0.003}{10.9091}{rm}{101}\RG{10.680}{0.003}{10.9091}{rm}{116}\RG{16.654}{0.003}{10.9091}{mi}{65}\RG{27.692}{0.003}{10.9091}{rm}{61}\RG{39.027}{0.003}{10.9091}{rm}{48}%
}\hskip.1em\allowbreak{} 。若两行中有一行为零，结论已由第四条证明：7.将一行的多项式倍数加于另一行，不改变行列式。对\allowbreak{}\hskip.1em \RB{22.779}{8.315}{2.853}{%
\RG{0.000}{0.004}{10.9091}{mi}{105}\RG{6.674}{0.004}{10.9091}{sy}{54}\RG{6.674}{0.004}{10.9091}{rm}{61}\RG{18.013}{0.004}{10.9091}{mi}{106}%
}\hskip.1em\allowbreak{} 及\allowbreak{}\hskip.1em \RB{28.055}{8.075}{0.933}{%
\RG{0.000}{0.004}{10.9091}{mi}{104}\RG{9.179}{0.004}{10.9091}{sy}{50}\RG{19.482}{0.004}{10.9091}{mi}{82}%
}\hskip.1em\allowbreak{} ，由线性性与相同行性质，有

\[
\RB{203.574}{8.286}{4.362}{%
\RG{0.000}{-0.004}{10.9091}{rm}{100}\RG{5.935}{-0.004}{10.9091}{rm}{101}\RG{10.680}{-0.004}{10.9091}{rm}{116}\RG{16.654}{-0.004}{10.9091}{mi}{65}\RG{62.388}{-0.004}{10.9091}{rm}{61}\RG{73.723}{-0.004}{10.9091}{rm}{100}\RG{79.658}{-0.004}{10.9091}{rm}{101}\RG{84.403}{-0.004}{10.9091}{rm}{116}\RG{90.382}{-0.004}{10.9091}{mi}{65}\RG{100.814}{-0.004}{10.9091}{rm}{43}\RG{111.546}{-0.004}{10.9091}{mi}{104}\RG{119.513}{-0.004}{10.9091}{rm}{100}\RG{125.447}{-0.004}{10.9091}{rm}{101}\RG{130.193}{-0.004}{10.9091}{rm}{116}\RG{136.167}{-0.004}{10.9091}{mi}{65}\RG{165.066}{-0.004}{10.9091}{rm}{61}\RG{176.401}{-0.004}{10.9091}{rm}{100}\RG{182.336}{-0.004}{10.9091}{rm}{101}\RG{187.081}{-0.004}{10.9091}{rm}{116}\RG{193.059}{-0.004}{10.9091}{mi}{65}\RG{201.067}{-0.004}{10.9091}{mi}{58}\RG{24.662}{-1.641}{7.9701}{mi}{105}\RG{27.293}{-1.641}{7.9701}{sy}{32}\RG{35.020}{-1.641}{7.9701}{mi}{114}\RG{41.392}{-1.641}{7.9701}{rm}{43}\RG{47.402}{-1.641}{7.9701}{mi}{104}\RG{51.853}{-1.641}{7.9701}{mi}{114}\RG{144.175}{-1.641}{7.9701}{mi}{105}\RG{146.806}{-1.641}{7.9701}{sy}{32}\RG{154.533}{-1.641}{7.9701}{mi}{114}\RG{38.507}{-2.749}{5.9776}{mi}{105}\RG{55.343}{-2.749}{5.9776}{mi}{106}\RG{158.020}{-2.749}{5.9776}{mi}{106}%
}
\]
8.三角矩阵的行列式等于其对角元之积；块三角矩阵的行列式等于其对角块行列式之积。

\noindent \RB{408.074}{9.228}{18.032}{%
\RT{0.000}{0.001}{10.9091}{若}\RG{14.263}{0.001}{10.9091}{mi}{65}\RT{25.623}{0.001}{10.9091}{为}\RT{36.533}{0.001}{10.9091}{上}\RT{47.442}{0.001}{10.9091}{三}\RT{58.351}{0.001}{10.9091}{角}\RT{69.260}{0.001}{10.9091}{矩}\RT{80.170}{0.001}{10.9091}{阵}\RT{91.079}{0.001}{10.9091}{，}\RT{100.734}{0.001}{10.9091}{则}\RT{111.643}{0.001}{10.9091}{行}\RT{122.552}{0.001}{10.9091}{表}\RT{133.462}{0.001}{10.9091}{达}\RT{144.371}{0.001}{10.9091}{式}\RT{155.280}{0.001}{10.9091}{中}\RT{166.189}{0.001}{10.9091}{非}\RT{177.099}{0.001}{10.9091}{零}\RT{188.008}{0.001}{10.9091}{项}\RT{198.917}{0.001}{10.9091}{对}\RT{209.826}{0.001}{10.9091}{应}\RT{220.736}{0.001}{10.9091}{的}\RT{231.645}{0.001}{10.9091}{置}\RT{242.554}{0.001}{10.9091}{换}\RG{256.813}{0.001}{10.9091}{mi}{27}\RT{266.629}{0.001}{10.9091}{必}\RT{277.538}{0.001}{10.9091}{须}\RT{288.447}{0.001}{10.9091}{满}\RT{299.357}{0.001}{10.9091}{足}\RG{313.619}{0.001}{10.9091}{mi}{105}\RG{320.294}{0.001}{10.9091}{sy}{20}\RG{331.809}{0.001}{10.9091}{mi}{27}\RG{338.271}{0.001}{10.9091}{rm}{40}\RG{342.426}{0.001}{10.9091}{mi}{105}\RG{346.070}{0.001}{10.9091}{rm}{41}\RT{353.577}{0.001}{10.9091}{对}\RT{364.486}{0.001}{10.9091}{每}\RT{375.395}{0.001}{10.9091}{个}\RG{389.658}{0.001}{10.9091}{mi}{105}\RT{396.656}{0.001}{10.9091}{成}\RG{47.086}{-7.490}{9.0910}{ex}{80}\RG{79.598}{-7.490}{9.0910}{ex}{80}\RG{307.572}{-7.490}{9.0910}{ex}{81}\RG{316.158}{-10.368}{7.9701}{mi}{110}\RT{0.000}{-14.763}{10.9091}{立}\RT{10.909}{-14.763}{10.9091}{。}\RT{21.677}{-14.763}{10.9091}{又}\RT{32.586}{-14.763}{10.9091}{因}\RG{61.585}{-14.763}{10.9091}{mi}{105}\RG{68.260}{-14.763}{10.9091}{rm}{61}\RG{94.098}{-14.763}{10.9091}{mi}{27}\RG{100.561}{-14.763}{10.9091}{rm}{40}\RG{104.715}{-14.763}{10.9091}{mi}{105}\RG{108.359}{-14.763}{10.9091}{rm}{41}\RT{112.513}{-14.763}{10.9091}{，}\RT{123.280}{-14.763}{10.9091}{这}\RT{134.190}{-14.763}{10.9091}{些}\RT{145.099}{-14.763}{10.9091}{不}\RT{156.008}{-14.763}{10.9091}{等}\RT{166.917}{-14.763}{10.9091}{式}\RT{177.827}{-14.763}{10.9091}{只}\RT{188.736}{-14.763}{10.9091}{能}\RT{199.645}{-14.763}{10.9091}{全}\RT{210.554}{-14.763}{10.9091}{为}\RT{221.464}{-14.763}{10.9091}{等}\RT{232.373}{-14.763}{10.9091}{式}\RT{243.282}{-14.763}{10.9091}{，}\RT{254.039}{-14.763}{10.9091}{故}\RG{268.542}{-14.763}{10.9091}{rm}{100}\RG{274.476}{-14.763}{10.9091}{rm}{101}\RG{279.222}{-14.763}{10.9091}{rm}{116}\RG{285.196}{-14.763}{10.9091}{mi}{65}\RG{296.234}{-14.763}{10.9091}{rm}{61}\RG{330.936}{-14.763}{10.9091}{mi}{97}\RT{342.260}{-14.763}{10.9091}{。}\RT{353.027}{-14.763}{10.9091}{下}\RT{363.937}{-14.763}{10.9091}{三}\RT{374.846}{-14.763}{10.9091}{角}\RT{385.755}{-14.763}{10.9091}{情}\RT{396.664}{-14.763}{10.9091}{形}\RG{336.543}{-16.399}{7.9701}{mi}{105}\RG{339.176}{-16.399}{7.9701}{mi}{105}\RG{56.682}{-17.036}{7.9701}{mi}{105}\RG{89.194}{-17.036}{7.9701}{mi}{105}\RG{316.158}{-17.459}{7.9701}{mi}{105}\RG{318.789}{-17.459}{7.9701}{rm}{61}\RG{324.796}{-17.459}{7.9701}{rm}{49}%
}

由转置性质得到，特别地\allowbreak{}\hskip.1em \RB{45.846}{8.074}{2.206}{%
\RG{0.000}{0.003}{10.9091}{rm}{100}\RG{5.935}{0.003}{10.9091}{rm}{101}\RG{10.680}{0.003}{10.9091}{rm}{116}\RG{16.654}{0.003}{10.9091}{mi}{73}\RG{21.366}{-1.633}{7.9701}{mi}{110}\RG{29.539}{0.003}{10.9091}{rm}{61}\RG{40.874}{0.003}{10.9091}{rm}{49}%
}\hskip.1em\allowbreak{} 。更一般地，设\allowbreak{}\hskip.1em \RB{46.097}{7.758}{1.991}{%
\RG{0.000}{0.004}{10.9091}{rm}{48}\RG{8.370}{0.004}{10.9091}{sy}{20}\RG{19.886}{0.004}{10.9091}{mi}{114}\RG{28.027}{0.004}{10.9091}{sy}{20}\RG{39.542}{0.004}{10.9091}{mi}{110}%
}\hskip.1em\allowbreak{} ，

\[
\RB{227.551}{16.864}{11.399}{%
\RG{22.377}{16.001}{9.0910}{ex}{32}\RG{56.495}{16.001}{9.0910}{ex}{33}\RG{29.574}{7.157}{10.9091}{mi}{66}\RG{47.694}{7.157}{10.9091}{mi}{67}\RG{0.000}{0.001}{10.9091}{mi}{65}\RG{11.038}{0.001}{10.9091}{rm}{61}\RG{65.511}{0.001}{10.9091}{mi}{59}\RG{92.108}{0.001}{10.9091}{mi}{66}\RG{103.811}{0.001}{10.9091}{sy}{50}\RG{114.114}{0.001}{10.9091}{mi}{77}\RG{128.604}{0.001}{10.9091}{rm}{40}\RG{132.758}{0.001}{10.9091}{mi}{82}\RG{140.979}{0.001}{10.9091}{rm}{41}\RG{145.132}{0.001}{10.9091}{mi}{59}\RG{160.820}{0.001}{10.9091}{mi}{68}\RG{173.023}{0.001}{10.9091}{sy}{50}\RG{183.327}{0.001}{10.9091}{mi}{77}\RG{208.516}{0.001}{10.9091}{rm}{40}\RG{212.669}{0.001}{10.9091}{mi}{82}\RG{220.890}{0.001}{10.9091}{rm}{41}\RG{225.044}{0.001}{10.9091}{mi}{58}\RG{124.448}{-1.636}{7.9701}{mi}{114}\RG{193.661}{-1.636}{7.9701}{mi}{110}\RG{198.349}{-1.636}{7.9701}{sy}{0}\RG{204.359}{-1.636}{7.9701}{mi}{114}\RG{31.236}{-7.607}{10.9091}{rm}{48}\RG{47.330}{-7.607}{10.9091}{mi}{68}%
}
\]
定义式中的非零项所对应的置换必须把\allowbreak{}\hskip.1em \RB{44.990}{8.681}{3.228}{%
\RG{0.000}{-0.001}{10.9091}{sy}{102}\RG{5.454}{-0.001}{10.9091}{rm}{49}\RG{10.795}{-0.001}{10.9091}{mi}{59}\RG{15.573}{-0.001}{10.9091}{mi}{58}\RG{20.362}{-0.001}{10.9091}{mi}{58}\RG{25.152}{-0.001}{10.9091}{mi}{58}\RG{29.930}{-0.001}{10.9091}{mi}{59}\RG{34.719}{-0.001}{10.9091}{mi}{114}\RG{39.832}{-0.001}{10.9091}{sy}{103}%
}\hskip.1em\allowbreak{} 映至自身。这样的置换由前后两个指标块上的置换唯一确定，其符号为这两个置换符号之积。分别对两块中的置换求和，得到\allowbreak{}\hskip.1em \RB{91.930}{8.290}{0.608}{%
\RG{0.000}{0.001}{10.9091}{rm}{100}\RG{5.935}{0.001}{10.9091}{rm}{101}\RG{10.680}{0.001}{10.9091}{rm}{116}\RG{16.654}{0.001}{10.9091}{mi}{65}\RG{27.692}{0.001}{10.9091}{rm}{61}\RG{39.027}{0.001}{10.9091}{rm}{100}\RG{44.961}{0.001}{10.9091}{rm}{101}\RG{49.707}{0.001}{10.9091}{rm}{116}\RG{55.684}{0.001}{10.9091}{mi}{66}\RG{66.168}{0.001}{10.9091}{rm}{100}\RG{72.103}{0.001}{10.9091}{rm}{101}\RG{76.848}{0.001}{10.9091}{rm}{116}\RG{82.822}{0.001}{10.9091}{mi}{68}%
}\hskip.1em\allowbreak{} 。块下三角情形仍由转置性质得到。9.同行（列）展开及异行（列）代数余子式恒等式成立。记\allowbreak{}\hskip.1em \RB{26.691}{10.503}{3.226}{%
\RG{0.000}{0.002}{10.9091}{mi}{65}\RG{8.007}{0.002}{10.9091}{rm}{91}\RG{10.334}{3.239}{9.0910}{ex}{98}\RG{11.038}{0.002}{10.9091}{mi}{105}\RG{14.681}{0.002}{10.9091}{mi}{59}\RG{19.458}{3.239}{9.0910}{ex}{98}\RG{19.467}{0.002}{10.9091}{mi}{106}\RG{24.500}{0.002}{10.9091}{rm}{93}%
}\hskip.1em\allowbreak{} 为删去第\allowbreak{}\hskip.1em \RB{3.587}{7.669}{0.607}{%
\RG{0.000}{0.002}{10.9091}{mi}{105}%
}\hskip.1em\allowbreak{} 行、第\allowbreak{}\hskip.1em \RB{4.766}{7.669}{2.724}{%
\RG{0.000}{0.002}{10.9091}{mi}{106}%
}\hskip.1em\allowbreak{}列所得的矩阵，记\allowbreak{}\hskip.1em \RB{108.578}{10.504}{3.618}{%
\RG{0.000}{0.002}{10.9091}{mi}{67}\RG{7.642}{-1.634}{7.9701}{mi}{105}\RG{10.275}{-1.634}{7.9701}{mi}{106}\RG{17.302}{0.002}{10.9091}{rm}{61}\RG{28.637}{0.002}{10.9091}{rm}{40}\RG{32.795}{0.002}{10.9091}{sy}{0}\RG{41.280}{0.002}{10.9091}{rm}{49}\RG{46.626}{0.002}{10.9091}{rm}{41}\RG{50.775}{3.961}{7.9701}{mi}{105}\RG{53.406}{3.961}{7.9701}{rm}{43}\RG{59.416}{3.961}{7.9701}{mi}{106}\RG{65.233}{0.002}{10.9091}{rm}{100}\RG{71.167}{0.002}{10.9091}{rm}{101}\RG{75.913}{0.002}{10.9091}{rm}{116}\RG{81.887}{0.002}{10.9091}{mi}{65}\RG{89.894}{0.002}{10.9091}{rm}{91}\RG{92.222}{3.240}{9.0910}{ex}{98}\RG{92.925}{0.002}{10.9091}{mi}{105}\RG{96.569}{0.002}{10.9091}{mi}{59}\RG{101.345}{3.240}{9.0910}{ex}{98}\RG{101.355}{0.002}{10.9091}{mi}{106}\RG{106.387}{0.002}{10.9091}{rm}{93}%
}\hskip.1em\allowbreak{} 为\allowbreak{}\hskip.1em \RB{11.744}{5.313}{3.618}{%
\RG{0.000}{0.002}{10.9091}{mi}{97}\RG{5.607}{-1.634}{7.9701}{mi}{105}\RG{8.240}{-1.634}{7.9701}{mi}{106}%
}\hskip.1em\allowbreak{} 的代数余子式。在行表达式中，将置换按\allowbreak{}\hskip.1em \RB{37.549}{8.685}{3.236}{%
\RG{0.000}{0.003}{10.9091}{mi}{27}\RG{6.462}{0.003}{10.9091}{rm}{40}\RG{10.616}{0.003}{10.9091}{mi}{105}\RG{14.261}{0.003}{10.9091}{rm}{41}\RG{21.439}{0.003}{10.9091}{rm}{61}\RG{32.784}{0.003}{10.9091}{mi}{106}%
}\hskip.1em\allowbreak{} 分组。把第\allowbreak{}\hskip.1em \RB{3.587}{7.670}{0.606}{%
\RG{0.000}{0.003}{10.9091}{mi}{105}%
}\hskip.1em\allowbreak{} 行及第\allowbreak{}\hskip.1em \RB{4.766}{7.670}{2.723}{%
\RG{0.000}{0.003}{10.9091}{mi}{106}%
}\hskip.1em\allowbreak{} 列分别移至首位，共作\allowbreak{}\hskip.1em \RB{39.641}{7.757}{2.723}{%
\RG{0.000}{0.003}{10.9091}{mi}{105}\RG{5.862}{0.003}{10.9091}{rm}{43}\RG{16.388}{0.003}{10.9091}{mi}{106}\RG{23.638}{0.003}{10.9091}{sy}{0}\RG{34.341}{0.003}{10.9091}{rm}{50}%
}\hskip.1em\allowbreak{} 次相邻交换；删除这两个首位指标后，余下置换的符号因而与原置换相差因子\allowbreak{}\hskip.1em \RB{34.279}{9.277}{3.235}{%
\RG{0.000}{0.004}{10.9091}{rm}{40}\RG{4.155}{0.004}{10.9091}{sy}{0}\RG{12.639}{0.004}{10.9091}{rm}{49}\RG{17.985}{0.004}{10.9091}{rm}{41}\RG{22.134}{3.962}{7.9701}{mi}{105}\RG{24.766}{3.962}{7.9701}{rm}{43}\RG{30.775}{3.962}{7.9701}{mi}{106}%
}\hskip.1em\allowbreak{} 。对每一组求和，即得

\[
\RB{83.283}{15.527}{12.851}{%
\RG{43.395}{11.820}{7.9701}{mi}{110}\RG{39.174}{9.092}{9.0910}{ex}{88}\RG{0.000}{0.002}{10.9091}{rm}{100}\RG{5.935}{0.002}{10.9091}{rm}{101}\RG{10.680}{0.002}{10.9091}{rm}{116}\RG{16.654}{0.002}{10.9091}{mi}{65}\RG{27.692}{0.002}{10.9091}{rm}{61}\RG{54.267}{0.002}{10.9091}{mi}{97}\RG{66.504}{0.002}{10.9091}{mi}{67}\RG{80.776}{0.002}{10.9091}{mi}{58}\RG{59.875}{-1.635}{7.9701}{mi}{105}\RG{62.507}{-1.635}{7.9701}{mi}{106}\RG{74.146}{-1.635}{7.9701}{mi}{105}\RG{76.779}{-1.635}{7.9701}{mi}{106}\RG{39.031}{-10.867}{7.9701}{mi}{106}\RG{42.575}{-10.867}{7.9701}{rm}{61}\RG{48.582}{-10.867}{7.9701}{rm}{49}%
}
\]
% Original page 6.
对\allowbreak{}\hskip.1em \RB{56.261}{7.757}{2.614}{%
\RG{0.000}{0.003}{10.9091}{rm}{49}\RG{8.371}{0.003}{10.9091}{sy}{20}\RG{19.886}{0.003}{10.9091}{mi}{112}\RG{25.243}{0.003}{10.9091}{mi}{59}\RG{30.021}{0.003}{10.9091}{mi}{113}\RG{38.190}{0.003}{10.9091}{sy}{20}\RG{49.706}{0.003}{10.9091}{mi}{110}%
}\hskip.1em\allowbreak{} ，定义\allowbreak{}\hskip.1em \RB{12.599}{8.259}{3.545}{%
\RG{0.000}{0.003}{10.9091}{mi}{14}\RG{4.731}{-1.634}{7.9701}{mi}{112}\RG{8.622}{-1.634}{7.9701}{mi}{113}%
}\hskip.1em\allowbreak{} 在\allowbreak{}\hskip.1em \RB{25.041}{5.314}{2.614}{%
\RG{0.000}{0.003}{10.9091}{mi}{112}\RG{8.391}{0.003}{10.9091}{rm}{61}\RG{19.730}{0.003}{10.9091}{mi}{113}%
}\hskip.1em\allowbreak{} 时为\allowbreak{}\hskip.1em \RB{4.973}{7.757}{0.500}{%
\RG{0.000}{0.003}{10.9091}{rm}{49}%
}\hskip.1em\allowbreak{} ，在\allowbreak{}\hskip.1em \RB{25.041}{8.314}{2.854}{%
\RG{0.000}{0.003}{10.9091}{mi}{112}\RG{8.391}{0.003}{10.9091}{sy}{54}\RG{8.391}{0.003}{10.9091}{rm}{61}\RG{19.730}{0.003}{10.9091}{mi}{113}%
}\hskip.1em\allowbreak{} 时为\allowbreak{}\hskip.1em \RB{5.387}{7.757}{0.726}{%
\RG{0.000}{0.003}{10.9091}{rm}{48}%
}\hskip.1em\allowbreak{} 。当\allowbreak{}\hskip.1em \RB{25.040}{8.314}{2.854}{%
\RG{0.000}{0.003}{10.9091}{mi}{112}\RG{8.391}{0.003}{10.9091}{sy}{54}\RG{8.391}{0.003}{10.9091}{rm}{61}\RG{19.729}{0.003}{10.9091}{mi}{113}%
}\hskip.1em\allowbreak{} 时，以第\allowbreak{}\hskip.1em \RB{5.715}{5.314}{2.614}{%
\RG{0.000}{0.003}{10.9091}{mi}{112}%
}\hskip.1em\allowbreak{} 行替换第\allowbreak{}\hskip.1em \RB{5.311}{5.314}{2.614}{%
\RG{0.000}{0.003}{10.9091}{mi}{113}%
}\hskip.1em\allowbreak{} 行，所得矩阵含两行相同，其行列式为零；其第\allowbreak{}\hskip.1em \RB{5.311}{5.314}{2.614}{%
\RG{0.000}{0.003}{10.9091}{mi}{113}%
}\hskip.1em\allowbreak{} 行的代数余子式仍为原矩阵的\allowbreak{}\hskip.1em \RB{14.848}{8.184}{3.616}{%
\RG{0.000}{0.004}{10.9091}{mi}{67}\RG{7.642}{-1.632}{7.9701}{mi}{113}\RG{11.344}{-1.632}{7.9701}{mi}{106}%
}\hskip.1em\allowbreak{} 。再结合\allowbreak{}\hskip.1em \RB{25.040}{5.315}{2.613}{%
\RG{0.000}{0.004}{10.9091}{mi}{112}\RG{8.391}{0.004}{10.9091}{rm}{61}\RG{19.729}{0.004}{10.9091}{mi}{113}%
}\hskip.1em\allowbreak{} 时的展开式及转置性质，得到

\[
\RB{100.347}{15.523}{44.131}{%
\RG{4.365}{11.816}{7.9701}{mi}{110}\RG{0.143}{9.089}{9.0910}{ex}{88}\RG{15.236}{-0.002}{10.9091}{mi}{97}\RG{28.732}{-0.002}{10.9091}{mi}{67}\RG{47.115}{-0.002}{10.9091}{rm}{61}\RG{58.453}{-0.002}{10.9091}{mi}{14}\RG{73.058}{-0.002}{10.9091}{rm}{100}\RG{78.992}{-0.002}{10.9091}{rm}{101}\RG{83.738}{-0.002}{10.9091}{rm}{116}\RG{89.712}{-0.002}{10.9091}{mi}{65}\RG{97.719}{-0.002}{10.9091}{mi}{59}\RG{20.844}{-1.639}{7.9701}{mi}{112}\RG{24.735}{-1.639}{7.9701}{mi}{106}\RG{36.375}{-1.639}{7.9701}{mi}{113}\RG{40.076}{-1.639}{7.9701}{mi}{106}\RG{63.184}{-1.639}{7.9701}{mi}{112}\RG{67.075}{-1.639}{7.9701}{mi}{113}\RG{0.000}{-10.871}{7.9701}{mi}{106}\RG{3.544}{-10.871}{7.9701}{rm}{61}\RG{9.552}{-10.871}{7.9701}{rm}{49}\RG{6.335}{-20.871}{7.9701}{mi}{110}\RG{2.113}{-23.598}{9.0910}{ex}{88}\RG{17.063}{-32.689}{10.9091}{mi}{97}\RG{29.646}{-32.689}{10.9091}{mi}{67}\RG{47.115}{-32.689}{10.9091}{rm}{61}\RG{58.453}{-32.689}{10.9091}{mi}{14}\RG{73.058}{-32.689}{10.9091}{rm}{100}\RG{78.992}{-32.689}{10.9091}{rm}{101}\RG{83.738}{-32.689}{10.9091}{rm}{116}\RG{89.712}{-32.689}{10.9091}{mi}{65}\RG{97.719}{-32.689}{10.9091}{mi}{58}\RG{22.670}{-34.326}{7.9701}{mi}{105}\RG{25.303}{-34.326}{7.9701}{mi}{112}\RG{37.288}{-34.326}{7.9701}{mi}{105}\RG{39.921}{-34.326}{7.9701}{mi}{113}\RG{63.184}{-34.326}{7.9701}{mi}{112}\RG{67.075}{-34.326}{7.9701}{mi}{113}\RG{2.426}{-43.558}{7.9701}{mi}{105}\RG{5.057}{-43.558}{7.9701}{rm}{61}\RG{11.065}{-43.558}{7.9701}{rm}{49}%
}
\]
（2）对非负整数\allowbreak{}\hskip.1em \RB{5.311}{5.312}{2.615}{%
\RG{0.000}{0.001}{10.9091}{mi}{113}%
}\hskip.1em\allowbreak{} ，记\allowbreak{}\hskip.1em \RB{70.502}{8.683}{3.226}{%
\RG{0.000}{0.001}{10.9091}{rm}{91}\RG{3.030}{0.001}{10.9091}{mi}{113}\RG{8.158}{0.001}{10.9091}{rm}{93}\RG{14.147}{0.001}{10.9091}{rm}{61}\RG{25.494}{0.001}{10.9091}{sy}{102}\RG{30.949}{0.001}{10.9091}{rm}{49}\RG{36.290}{0.001}{10.9091}{mi}{59}\RG{41.068}{0.001}{10.9091}{mi}{58}\RG{45.857}{0.001}{10.9091}{mi}{58}\RG{50.635}{0.001}{10.9091}{mi}{58}\RG{55.425}{0.001}{10.9091}{mi}{59}\RG{60.214}{0.001}{10.9091}{mi}{113}\RG{65.344}{0.001}{10.9091}{sy}{103}%
}\hskip.1em\allowbreak{} ，其中\allowbreak{}\hskip.1em \RB{34.073}{8.683}{3.226}{%
\RG{0.000}{0.001}{10.9091}{rm}{91}\RG{3.022}{0.001}{10.9091}{rm}{48}\RG{8.367}{0.001}{10.9091}{rm}{93}\RG{14.368}{0.001}{10.9091}{rm}{61}\RG{25.707}{0.001}{10.9091}{bb}{63}%
}\hskip.1em\allowbreak{} ；约定\allowbreak{}\hskip.1em \RB{10.627}{8.181}{2.288}{%
\RG{0.000}{0.001}{10.9091}{mi}{83}\RG{6.570}{-1.635}{7.9701}{rm}{48}%
}\hskip.1em\allowbreak{} 只含一个空置换，其符号及空乘积均为\allowbreak{}\hskip.1em \RB{4.973}{7.756}{0.500}{%
\RG{0.000}{0.002}{10.9091}{rm}{49}%
}\hskip.1em\allowbreak{} 。对行指标集\allowbreak{}\hskip.1em \RB{5.758}{7.953}{0.500}{%
\RG{0.000}{0.002}{10.9091}{mi}{73}%
}\hskip.1em\allowbreak{} 和列指标集\allowbreak{}\hskip.1em \RB{7.253}{7.953}{0.727}{%
\RG{0.000}{0.002}{10.9091}{mi}{74}%
}\hskip.1em\allowbreak{} ，记\allowbreak{}\hskip.1em \RB{30.497}{8.684}{3.225}{%
\RG{0.000}{0.002}{10.9091}{mi}{65}\RG{8.007}{0.002}{10.9091}{rm}{91}\RG{11.038}{0.002}{10.9091}{mi}{73}\RG{16.601}{0.002}{10.9091}{mi}{59}\RG{21.391}{0.002}{10.9091}{mi}{74}\RG{28.306}{0.002}{10.9091}{rm}{93}%
}\hskip.1em\allowbreak{} 为按指标递增次序排列的子矩阵。在\allowbreak{}\hskip.1em \RB{43.991}{8.685}{3.225}{%
\RG{0.000}{0.003}{10.9091}{mi}{73}\RG{5.564}{0.003}{10.9091}{mi}{59}\RG{10.353}{0.003}{10.9091}{mi}{74}\RG{20.588}{0.003}{10.9091}{sy}{18}\RG{32.393}{0.003}{10.9091}{rm}{91}\RG{35.424}{0.003}{10.9091}{mi}{110}\RG{41.800}{0.003}{10.9091}{rm}{93}%
}\hskip.1em\allowbreak{} 时，记\allowbreak{}\hskip.1em \RB{52.986}{8.685}{3.225}{%
\RG{0.000}{0.003}{10.9091}{mi}{73}\RG{5.569}{3.961}{7.9701}{mi}{99}\RG{12.691}{0.003}{10.9091}{rm}{61}\RG{24.321}{0.003}{10.9091}{rm}{91}\RG{27.350}{0.003}{10.9091}{mi}{110}\RG{33.726}{0.003}{10.9091}{rm}{93}\RG{39.233}{0.003}{10.9091}{sy}{110}\RG{47.228}{0.003}{10.9091}{mi}{73}%
}\hskip.1em\allowbreak{} 、\allowbreak{}\hskip.1em \RB{55.826}{8.685}{3.225}{%
\RG{0.000}{0.003}{10.9091}{mi}{74}\RG{6.914}{3.961}{7.9701}{mi}{99}\RG{14.035}{0.003}{10.9091}{rm}{61}\RG{25.665}{0.003}{10.9091}{rm}{91}\RG{28.694}{0.003}{10.9091}{mi}{110}\RG{35.071}{0.003}{10.9091}{rm}{93}\RG{40.579}{0.003}{10.9091}{sy}{110}\RG{48.573}{0.003}{10.9091}{mi}{74}%
}\hskip.1em\allowbreak{} ，并记\allowbreak{}\hskip.1em \RB{61.014}{8.685}{3.236}{%
\RG{0.000}{0.003}{10.9091}{mi}{115}\RG{5.032}{0.003}{10.9091}{rm}{40}\RG{9.186}{0.003}{10.9091}{mi}{73}\RG{14.754}{0.003}{10.9091}{rm}{41}\RG{22.227}{0.003}{10.9091}{rm}{61}\RG{33.857}{7.275}{9.0910}{ex}{80}\RG{43.453}{-2.270}{7.9701}{mi}{105}\RG{46.084}{-2.270}{7.9701}{sy}{50}\RG{51.236}{-2.270}{7.9701}{mi}{73}\RG{57.426}{0.003}{10.9091}{mi}{105}%
}\hskip.1em\allowbreak{} 。固定满足\allowbreak{}\hskip.1em \RB{31.725}{8.685}{3.225}{%
\RG{0.000}{0.003}{10.9091}{sy}{106}\RG{3.030}{0.003}{10.9091}{mi}{73}\RG{8.600}{0.003}{10.9091}{sy}{106}\RG{14.950}{0.003}{10.9091}{rm}{61}\RG{26.578}{0.003}{10.9091}{mi}{114}%
}\hskip.1em\allowbreak{}的集合\allowbreak{}\hskip.1em \RB{31.712}{8.686}{3.224}{%
\RG{0.000}{0.004}{10.9091}{mi}{73}\RG{8.599}{0.004}{10.9091}{sy}{18}\RG{20.115}{0.004}{10.9091}{rm}{91}\RG{23.145}{0.004}{10.9091}{mi}{110}\RG{29.521}{0.004}{10.9091}{rm}{93}%
}\hskip.1em\allowbreak{} ，其中\allowbreak{}\hskip.1em \RB{46.097}{7.758}{1.991}{%
\RG{0.000}{0.004}{10.9091}{rm}{48}\RG{8.371}{0.004}{10.9091}{sy}{20}\RG{19.886}{0.004}{10.9091}{mi}{114}\RG{28.027}{0.004}{10.9091}{sy}{20}\RG{39.543}{0.004}{10.9091}{mi}{110}%
}\hskip.1em\allowbreak{} 。Laplace展开式为

\[
\RB{225.144}{10.456}{22.200}{%
\RG{42.426}{9.088}{9.0910}{ex}{88}\RG{81.088}{4.501}{7.9701}{mi}{115}\RG{84.661}{4.501}{7.9701}{rm}{40}\RG{87.666}{4.501}{7.9701}{mi}{73}\RG{91.584}{4.501}{7.9701}{rm}{41}\RG{94.588}{4.501}{7.9701}{rm}{43}\RG{100.599}{4.501}{7.9701}{mi}{115}\RG{104.172}{4.501}{7.9701}{rm}{40}\RG{107.177}{4.501}{7.9701}{mi}{74}\RG{112.082}{4.501}{7.9701}{rm}{41}\RG{200.367}{4.501}{7.9701}{mi}{99}\RG{215.868}{4.501}{7.9701}{mi}{99}\RG{0.000}{-0.003}{10.9091}{rm}{100}\RG{5.935}{-0.003}{10.9091}{rm}{101}\RG{10.680}{-0.003}{10.9091}{rm}{116}\RG{16.654}{-0.003}{10.9091}{mi}{65}\RG{27.692}{-0.003}{10.9091}{rm}{61}\RG{58.954}{-0.003}{10.9091}{rm}{40}\RG{63.108}{-0.003}{10.9091}{sy}{0}\RG{71.592}{-0.003}{10.9091}{rm}{49}\RG{76.938}{-0.003}{10.9091}{rm}{41}\RG{117.360}{-0.003}{10.9091}{rm}{100}\RG{123.295}{-0.003}{10.9091}{rm}{101}\RG{128.040}{-0.003}{10.9091}{rm}{116}\RG{134.014}{-0.003}{10.9091}{mi}{65}\RG{142.022}{-0.003}{10.9091}{rm}{91}\RG{145.052}{-0.003}{10.9091}{mi}{73}\RG{150.616}{-0.003}{10.9091}{mi}{59}\RG{155.405}{-0.003}{10.9091}{mi}{74}\RG{162.321}{-0.003}{10.9091}{rm}{93}\RG{167.099}{-0.003}{10.9091}{rm}{100}\RG{173.033}{-0.003}{10.9091}{rm}{101}\RG{177.779}{-0.003}{10.9091}{rm}{116}\RG{183.760}{-0.003}{10.9091}{mi}{65}\RG{191.767}{-0.003}{10.9091}{rm}{91}\RG{194.798}{-0.003}{10.9091}{mi}{73}\RG{204.168}{-0.003}{10.9091}{mi}{59}\RG{208.947}{-0.003}{10.9091}{mi}{74}\RG{219.669}{-0.003}{10.9091}{rm}{93}\RG{222.637}{-0.003}{10.9091}{mi}{58}\RG{39.031}{-11.519}{7.9701}{mi}{74}\RG{43.935}{-11.519}{7.9701}{sy}{18}\RG{49.945}{-11.519}{7.9701}{rm}{91}\RG{52.118}{-11.519}{7.9701}{mi}{110}\RG{56.807}{-11.519}{7.9701}{rm}{93}\RG{39.537}{-19.882}{7.9701}{sy}{106}\RG{41.684}{-19.882}{7.9701}{mi}{74}\RG{46.588}{-19.882}{7.9701}{sy}{106}\RG{48.735}{-19.882}{7.9701}{rm}{61}\RG{54.745}{-19.882}{7.9701}{mi}{114}%
}
\]
为证明此式，在行表达式中按照\allowbreak{}\hskip.1em \RB{41.961}{8.678}{3.242}{%
\RG{0.000}{-0.004}{10.9091}{mi}{27}\RG{6.462}{-0.004}{10.9091}{rm}{40}\RG{10.616}{-0.004}{10.9091}{mi}{73}\RG{16.185}{-0.004}{10.9091}{rm}{41}\RG{23.364}{-0.004}{10.9091}{rm}{61}\RG{34.708}{-0.004}{10.9091}{mi}{74}%
}\hskip.1em\allowbreak{} 分组。记\allowbreak{}\hskip.1em \RB{35.307}{8.176}{2.213}{%
\RG{0.000}{-0.004}{10.9091}{mi}{117}\RG{6.080}{-1.640}{7.9701}{mi}{73}\RG{13.482}{-0.004}{10.9091}{sy}{50}\RG{23.786}{-0.004}{10.9091}{mi}{83}\RG{30.356}{-1.640}{7.9701}{mi}{110}%
}\hskip.1em\allowbreak{} 为将前\allowbreak{}\hskip.1em \RB{5.147}{5.307}{0.613}{%
\RG{0.000}{-0.004}{10.9091}{mi}{114}%
}\hskip.1em\allowbreak{} 个指标依次映至\allowbreak{}\hskip.1em \RB{5.758}{7.948}{0.503}{%
\RG{0.000}{-0.003}{10.9091}{mi}{73}%
}\hskip.1em\allowbreak{} 、其余指标依次映至\allowbreak{}\hskip.1em \RB{9.422}{7.948}{0.503}{%
\RG{0.000}{-0.003}{10.9091}{mi}{73}\RG{5.569}{3.956}{7.9701}{mi}{99}%
}\hskip.1em\allowbreak{} 的置换。若\allowbreak{}\hskip.1em \RB{88.371}{8.679}{3.230}{%
\RG{0.000}{-0.003}{10.9091}{mi}{73}\RG{8.698}{-0.003}{10.9091}{rm}{61}\RG{20.135}{-0.003}{10.9091}{sy}{102}\RG{25.590}{-0.003}{10.9091}{mi}{105}\RG{29.234}{-1.639}{7.9701}{rm}{49}\RG{36.681}{-0.003}{10.9091}{mi}{60}\RG{48.119}{-0.003}{10.9091}{sy}{1}\RG{52.962}{-0.003}{10.9091}{sy}{1}\RG{57.817}{-0.003}{10.9091}{sy}{1}\RG{63.975}{-0.003}{10.9091}{mi}{60}\RG{75.408}{-0.003}{10.9091}{mi}{105}\RG{79.057}{-1.639}{7.9701}{mi}{114}\RG{83.213}{-0.003}{10.9091}{sy}{103}%
}\hskip.1em\allowbreak{} ，则\allowbreak{}\hskip.1em \RB{10.281}{5.308}{2.139}{%
\RG{0.000}{-0.003}{10.9091}{mi}{117}\RG{6.079}{-1.639}{7.9701}{mi}{73}%
}\hskip.1em\allowbreak{} 的逆序数为\allowbreak{}\hskip.1em \RB{59.523}{8.679}{3.272}{%
\RG{0.000}{7.270}{9.0910}{ex}{80}\RG{9.596}{4.391}{7.9701}{mi}{114}\RG{9.596}{-2.699}{7.9701}{mi}{107}\RG{13.814}{-2.700}{7.9701}{rm}{61}\RG{19.821}{-2.700}{7.9701}{rm}{49}\RG{24.142}{-0.003}{10.9091}{rm}{40}\RG{28.296}{-0.003}{10.9091}{mi}{105}\RG{31.940}{-1.639}{7.9701}{mi}{107}\RG{39.076}{-0.003}{10.9091}{sy}{0}\RG{50.025}{-0.003}{10.9091}{mi}{107}\RG{55.946}{-0.003}{10.9091}{rm}{41}%
}\hskip.1em\allowbreak{} ，从而\allowbreak{}\hskip.1em \RB{110.894}{9.911}{3.240}{%
\RG{0.000}{-0.002}{10.9091}{mi}{34}\RG{4.999}{-0.002}{10.9091}{rm}{40}\RG{9.153}{-0.002}{10.9091}{mi}{117}\RG{15.232}{-1.639}{7.9701}{mi}{73}\RG{19.604}{-0.002}{10.9091}{rm}{41}\RG{26.783}{-0.002}{10.9091}{rm}{61}\RG{38.128}{-0.002}{10.9091}{rm}{40}\RG{42.281}{-0.002}{10.9091}{sy}{0}\RG{50.767}{-0.002}{10.9091}{rm}{49}\RG{56.112}{-0.002}{10.9091}{rm}{41}\RG{60.261}{3.957}{7.9701}{mi}{115}\RG{63.835}{3.957}{7.9701}{rm}{40}\RG{66.840}{3.957}{7.9701}{mi}{73}\RG{70.758}{3.957}{7.9701}{rm}{41}\RG{73.763}{3.957}{7.9701}{sy}{0}\RG{79.774}{3.957}{7.9701}{mi}{114}\RG{83.475}{3.957}{7.9701}{rm}{40}\RG{86.481}{3.957}{7.9701}{mi}{114}\RG{90.182}{3.957}{7.9701}{rm}{43}\RG{96.189}{3.957}{7.9701}{rm}{49}\RG{100.051}{3.957}{7.9701}{rm}{41}\RG{103.061}{3.957}{7.9701}{mi}{61}\RG{106.925}{3.957}{7.9701}{rm}{50}%
}\hskip.1em\allowbreak{} 。以同样方式定义\allowbreak{}\hskip.1em \RB{11.336}{5.309}{2.292}{%
\RG{0.000}{-0.002}{10.9091}{mi}{117}\RG{6.079}{-1.639}{7.9701}{mi}{74}%
}\hskip.1em\allowbreak{} 。每个满足\allowbreak{}\hskip.1em \RB{41.961}{8.680}{3.240}{%
\RG{0.000}{-0.002}{10.9091}{mi}{27}\RG{6.462}{-0.002}{10.9091}{rm}{40}\RG{10.616}{-0.002}{10.9091}{mi}{73}\RG{16.185}{-0.002}{10.9091}{rm}{41}\RG{23.364}{-0.002}{10.9091}{rm}{61}\RG{34.708}{-0.002}{10.9091}{mi}{74}%
}\hskip.1em\allowbreak{} 的置换唯一写成

\[
\RB{201.061}{9.841}{3.237}{%
\RG{73.490}{4.505}{7.9701}{sy}{0}\RG{79.500}{4.505}{7.9701}{rm}{49}\RG{0.000}{0.001}{10.9091}{mi}{27}\RG{9.493}{0.001}{10.9091}{rm}{61}\RG{20.832}{0.001}{10.9091}{mi}{117}\RG{32.270}{0.001}{10.9091}{rm}{40}\RG{36.424}{0.001}{10.9091}{mi}{11}\RG{45.712}{0.001}{10.9091}{sy}{8}\RG{56.621}{0.001}{10.9091}{mi}{12}\RG{63.256}{0.001}{10.9091}{rm}{41}\RG{67.410}{0.001}{10.9091}{mi}{117}\RG{83.819}{0.001}{10.9091}{mi}{59}\RG{110.416}{0.001}{10.9091}{mi}{11}\RG{120.317}{0.001}{10.9091}{sy}{50}\RG{130.620}{0.001}{10.9091}{mi}{83}\RG{141.345}{0.001}{10.9091}{mi}{59}\RG{157.033}{0.001}{10.9091}{mi}{12}\RG{166.705}{0.001}{10.9091}{sy}{50}\RG{177.008}{0.001}{10.9091}{mi}{83}\RG{198.433}{0.001}{10.9091}{mi}{59}\RG{26.912}{-1.636}{7.9701}{mi}{74}\RG{137.189}{-1.636}{7.9701}{mi}{114}\RG{183.577}{-1.636}{7.9701}{mi}{110}\RG{188.267}{-1.636}{7.9701}{sy}{0}\RG{194.277}{-1.636}{7.9701}{mi}{114}\RG{73.490}{-2.696}{7.9701}{mi}{73}%
}
\]
其中\allowbreak{}\hskip.1em \RB{26.860}{8.195}{2.613}{%
\RG{0.000}{0.004}{10.9091}{mi}{11}\RG{9.288}{0.004}{10.9091}{sy}{8}\RG{20.197}{0.004}{10.9091}{mi}{12}%
}\hskip.1em\allowbreak{} 分别在前后两个指标块上按\allowbreak{}\hskip.1em \RB{6.926}{5.315}{0.605}{%
\RG{0.000}{0.004}{10.9091}{mi}{11}%
}\hskip.1em\allowbreak{} 、\allowbreak{}\hskip.1em \RB{6.664}{8.195}{2.613}{%
\RG{0.000}{0.004}{10.9091}{mi}{12}%
}\hskip.1em\allowbreak{} 作用。由于\allowbreak{}\hskip.1em \RB{36.456}{8.686}{3.235}{%
\RG{0.000}{0.004}{10.9091}{mi}{114}\RG{5.110}{0.004}{10.9091}{rm}{40}\RG{9.264}{0.004}{10.9091}{mi}{114}\RG{16.799}{0.004}{10.9091}{rm}{43}\RG{27.534}{0.004}{10.9091}{rm}{49}\RG{32.880}{0.004}{10.9091}{rm}{41}%
}\hskip.1em\allowbreak{} 为偶数，有

\[
\RB{133.347}{10.456}{3.241}{%
\RG{56.273}{4.501}{7.9701}{mi}{115}\RG{59.846}{4.501}{7.9701}{rm}{40}\RG{62.851}{4.501}{7.9701}{mi}{73}\RG{66.769}{4.501}{7.9701}{rm}{41}\RG{69.773}{4.501}{7.9701}{rm}{43}\RG{75.784}{4.501}{7.9701}{mi}{115}\RG{79.358}{4.501}{7.9701}{rm}{40}\RG{82.363}{4.501}{7.9701}{mi}{74}\RG{87.267}{4.501}{7.9701}{rm}{41}\RG{0.000}{-0.003}{10.9091}{mi}{34}\RG{4.999}{-0.003}{10.9091}{rm}{40}\RG{9.153}{-0.003}{10.9091}{mi}{27}\RG{15.616}{-0.003}{10.9091}{rm}{41}\RG{22.794}{-0.003}{10.9091}{rm}{61}\RG{34.140}{-0.003}{10.9091}{rm}{40}\RG{38.293}{-0.003}{10.9091}{sy}{0}\RG{46.778}{-0.003}{10.9091}{rm}{49}\RG{52.124}{-0.003}{10.9091}{rm}{41}\RG{90.727}{-0.003}{10.9091}{mi}{34}\RG{95.726}{-0.003}{10.9091}{rm}{40}\RG{99.880}{-0.003}{10.9091}{mi}{11}\RG{106.744}{-0.003}{10.9091}{rm}{41}\RG{110.898}{-0.003}{10.9091}{mi}{34}\RG{115.896}{-0.003}{10.9091}{rm}{40}\RG{120.050}{-0.003}{10.9091}{mi}{12}\RG{126.685}{-0.003}{10.9091}{rm}{41}\RG{130.840}{-0.003}{10.9091}{mi}{58}%
}
\]
对应乘积分解为两个子矩阵中的乘积。先分别对\allowbreak{}\hskip.1em \RB{18.315}{8.191}{2.616}{%
\RG{0.000}{-0.000}{10.9091}{mi}{11}\RG{6.862}{-0.000}{10.9091}{mi}{59}\RG{11.651}{-0.000}{10.9091}{mi}{12}%
}\hskip.1em\allowbreak{} 求和，再对\allowbreak{}\hskip.1em \RB{7.253}{7.951}{0.729}{%
\RG{0.000}{-0.000}{10.9091}{mi}{74}%
}\hskip.1em\allowbreak{} 求和，便得到所述Laplace展开式。转置后即得按指定若干列展开的形式。

\noindent \RB{437.080}{9.231}{19.249}{%
\RT{21.819}{0.004}{10.9091}{设}\RG{36.549}{0.004}{10.9091}{mi}{85}\RG{48.287}{0.004}{10.9091}{sy}{50}\RG{58.904}{0.004}{10.9091}{mi}{77}\RG{82.563}{0.004}{10.9091}{rm}{40}\RG{86.717}{0.004}{10.9091}{mi}{82}\RG{94.938}{0.004}{10.9091}{rm}{41}\RT{99.092}{0.004}{10.9091}{、}\RG{110.001}{0.004}{10.9091}{mi}{86}\RG{121.946}{0.004}{10.9091}{rm}{61}\RG{133.597}{0.004}{10.9091}{rm}{40}\RG{137.752}{0.004}{10.9091}{mi}{118}\RG{151.132}{0.004}{10.9091}{rm}{41}\RG{158.629}{0.004}{10.9091}{sy}{50}\RG{169.247}{0.004}{10.9091}{mi}{77}\RG{192.906}{0.004}{10.9091}{rm}{40}\RG{197.061}{0.004}{10.9091}{mi}{82}\RG{205.281}{0.004}{10.9091}{rm}{41}\RT{209.434}{0.004}{10.9091}{，}\RT{220.464}{0.004}{10.9091}{其}\RT{231.493}{0.004}{10.9091}{中}\RG{246.228}{0.004}{10.9091}{mi}{109}\RG{255.577}{0.004}{10.9091}{mi}{59}\RG{260.355}{0.004}{10.9091}{mi}{112}\RT{269.538}{0.004}{10.9091}{为}\RT{280.567}{0.004}{10.9091}{正}\RT{291.597}{0.004}{10.9091}{整}\RT{302.626}{0.004}{10.9091}{数}\RT{313.535}{0.004}{10.9091}{。}\RT{324.564}{0.004}{10.9091}{记}\RG{339.303}{0.004}{10.9091}{mi}{117}\RT{353.202}{0.004}{10.9091}{为}\RG{367.932}{0.004}{10.9091}{mi}{85}\RT{380.146}{0.004}{10.9091}{的}\RT{391.176}{0.004}{10.9091}{第}\RG{405.908}{0.004}{10.9091}{mi}{106}\RT{414.762}{0.004}{10.9091}{列}\RT{425.671}{0.004}{10.9091}{。}\RG{69.237}{-1.633}{7.9701}{mi}{109}\RG{76.073}{-1.633}{7.9701}{mi}{59}\RG{78.219}{-1.633}{7.9701}{mi}{112}\RG{142.916}{-1.633}{7.9701}{mi}{106}\RG{146.458}{-1.633}{7.9701}{mi}{107}\RG{179.581}{-1.633}{7.9701}{mi}{112}\RG{183.472}{-1.633}{7.9701}{mi}{59}\RG{185.617}{-1.633}{7.9701}{mi}{109}\RG{345.382}{-1.633}{7.9701}{mi}{106}\RG{81.085}{-7.487}{9.0910}{ex}{80}\RG{90.682}{-9.589}{7.9701}{mi}{112}\RG{0.000}{-14.761}{10.9091}{mi}{85}\RG{8.389}{-14.761}{10.9091}{mi}{86}\RT{20.628}{-14.761}{10.9091}{的}\RT{31.538}{-14.761}{10.9091}{第}\RG{46.079}{-14.761}{10.9091}{mi}{107}\RT{55.634}{-14.761}{10.9091}{列}\RT{66.543}{-14.761}{10.9091}{为}\RG{106.373}{-14.761}{10.9091}{mi}{118}\RG{119.753}{-14.761}{10.9091}{mi}{117}\RT{129.830}{-14.761}{10.9091}{，}\RT{140.740}{-14.761}{10.9091}{故}\RT{151.649}{-14.761}{10.9091}{由}\RT{162.558}{-14.761}{10.9091}{各}\RT{173.468}{-14.761}{10.9091}{列}\RT{184.377}{-14.761}{10.9091}{的}\RT{195.286}{-14.761}{10.9091}{线}\RT{206.195}{-14.761}{10.9091}{性}\RT{217.105}{-14.761}{10.9091}{性}\RT{228.014}{-14.761}{10.9091}{，}\RG{111.536}{-16.397}{7.9701}{mi}{106}\RG{115.078}{-16.397}{7.9701}{mi}{107}\RG{125.832}{-16.397}{7.9701}{mi}{106}\RG{90.682}{-17.266}{7.9701}{mi}{106}\RG{94.226}{-17.266}{7.9701}{rm}{61}\RG{100.233}{-17.266}{7.9701}{rm}{49}%
}

\[
\RB{221.751}{16.866}{12.851}{%
\RG{93.391}{16.002}{9.0910}{ex}{32}\RG{133.631}{16.003}{9.0910}{ex}{33}\RG{71.096}{12.427}{7.9701}{mi}{112}\RG{104.216}{11.821}{7.9701}{mi}{109}\RG{66.475}{9.093}{9.0910}{ex}{88}\RG{101.826}{9.093}{9.0910}{ex}{89}\RG{0.000}{0.002}{10.9091}{rm}{100}\RG{5.935}{0.002}{10.9091}{rm}{101}\RG{10.680}{0.002}{10.9091}{rm}{116}\RG{14.837}{0.002}{10.9091}{rm}{40}\RG{18.989}{0.002}{10.9091}{mi}{85}\RG{27.379}{0.002}{10.9091}{mi}{86}\RG{35.985}{0.002}{10.9091}{rm}{41}\RG{43.163}{0.002}{10.9091}{rm}{61}\RG{116.498}{0.002}{10.9091}{mi}{118}\RG{142.647}{0.002}{10.9091}{rm}{100}\RG{148.581}{0.002}{10.9091}{rm}{101}\RG{153.327}{0.002}{10.9091}{rm}{116}\RG{157.483}{0.002}{10.9091}{rm}{40}\RG{161.637}{0.002}{10.9091}{mi}{117}\RG{175.084}{0.002}{10.9091}{mi}{59}\RG{179.862}{0.002}{10.9091}{mi}{58}\RG{184.651}{0.002}{10.9091}{mi}{58}\RG{189.429}{0.002}{10.9091}{mi}{58}\RG{194.219}{0.002}{10.9091}{mi}{59}\RG{199.008}{0.002}{10.9091}{mi}{117}\RG{215.090}{0.002}{10.9091}{rm}{41}\RG{219.244}{0.002}{10.9091}{mi}{58}\RG{121.662}{-1.634}{7.9701}{mi}{106}\RG{128.960}{-1.634}{7.9701}{mi}{107}\RG{167.716}{-1.634}{7.9701}{mi}{106}\RG{205.090}{-1.634}{7.9701}{mi}{106}\RG{208.215}{-2.543}{5.9776}{mi}{109}\RG{170.842}{-2.644}{5.9776}{rm}{49}\RG{124.786}{-2.917}{5.9776}{mi}{107}\RG{54.508}{-10.867}{7.9701}{mi}{106}\RG{61.421}{-10.867}{7.9701}{mi}{59}\RG{63.567}{-10.867}{7.9701}{mi}{58}\RG{65.712}{-10.867}{7.9701}{mi}{58}\RG{67.858}{-10.867}{7.9701}{mi}{58}\RG{70.003}{-10.867}{7.9701}{mi}{59}\RG{72.149}{-10.867}{7.9701}{mi}{106}\RG{81.699}{-10.867}{7.9701}{rm}{61}\RG{87.706}{-10.867}{7.9701}{rm}{49}\RG{100.588}{-11.110}{7.9701}{mi}{107}\RG{104.806}{-11.110}{7.9701}{rm}{61}\RG{110.813}{-11.110}{7.9701}{rm}{49}\RG{75.279}{-11.776}{5.9776}{mi}{109}\RG{57.633}{-11.877}{5.9776}{rm}{49}%
}
\]
若\allowbreak{}\hskip.1em \RB{44.179}{7.666}{2.727}{%
\RG{0.000}{-0.002}{10.9091}{mi}{106}\RG{4.418}{-1.638}{7.9701}{rm}{49}\RG{8.737}{-0.002}{10.9091}{mi}{59}\RG{13.516}{-0.002}{10.9091}{mi}{58}\RG{18.305}{-0.002}{10.9091}{mi}{58}\RG{23.083}{-0.002}{10.9091}{mi}{58}\RG{27.872}{-0.002}{10.9091}{mi}{59}\RG{32.661}{-0.002}{10.9091}{mi}{106}\RG{37.083}{-1.638}{7.9701}{mi}{109}%
}\hskip.1em\allowbreak{} 中有相同指标，相应行列式为零。其余各项按集合\allowbreak{}\hskip.1em \RB{121.372}{8.680}{3.229}{%
\RG{0.000}{-0.002}{10.9091}{mi}{74}\RG{10.059}{-0.002}{10.9091}{rm}{61}\RG{21.512}{-0.002}{10.9091}{sy}{102}\RG{26.967}{-0.002}{10.9091}{mi}{115}\RG{31.998}{-1.638}{7.9701}{rm}{49}\RG{39.461}{-0.002}{10.9091}{mi}{60}\RG{50.914}{-0.002}{10.9091}{sy}{1}\RG{55.757}{-0.002}{10.9091}{sy}{1}\RG{60.612}{-0.002}{10.9091}{sy}{1}\RG{66.786}{-0.002}{10.9091}{mi}{60}\RG{78.241}{-0.002}{10.9091}{mi}{115}\RG{83.270}{-1.638}{7.9701}{mi}{109}\RG{90.560}{-0.002}{10.9091}{sy}{103}\RG{99.156}{-0.002}{10.9091}{sy}{18}\RG{110.789}{-0.002}{10.9091}{rm}{91}\RG{113.819}{-0.002}{10.9091}{mi}{112}\RG{119.181}{-0.002}{10.9091}{rm}{93}%
}\hskip.1em\allowbreak{}分组，并唯一写成\allowbreak{}\hskip.1em \RB{42.730}{7.666}{4.016}{%
\RG{0.000}{-0.001}{10.9091}{mi}{106}\RG{4.419}{-1.637}{7.9701}{mi}{107}\RG{12.121}{-0.001}{10.9091}{rm}{61}\RG{23.460}{-0.001}{10.9091}{mi}{115}\RG{28.491}{-1.698}{7.9701}{mi}{28}\RG{32.760}{-1.698}{7.9701}{rm}{40}\RG{35.765}{-1.698}{7.9701}{mi}{107}\RG{39.983}{-1.698}{7.9701}{rm}{41}%
}\hskip.1em\allowbreak{} ，其中\allowbreak{}\hskip.1em \RB{32.857}{8.179}{2.210}{%
\RG{0.000}{-0.001}{10.9091}{mi}{28}\RG{8.888}{-0.001}{10.9091}{sy}{50}\RG{19.191}{-0.001}{10.9091}{mi}{83}\RG{25.760}{-1.637}{7.9701}{mi}{109}%
}\hskip.1em\allowbreak{} 。按换列变号性质，每组的贡献为

\[
\RB{153.889}{42.348}{3.226}{%
\RG{106.874}{38.640}{7.9701}{mi}{109}\RG{63.609}{35.913}{9.0910}{ex}{88}\RG{104.484}{35.913}{9.0910}{ex}{89}\RG{0.000}{26.821}{10.9091}{rm}{100}\RG{5.935}{26.821}{10.9091}{rm}{101}\RG{10.680}{26.821}{10.9091}{rm}{116}\RG{16.654}{26.821}{10.9091}{mi}{85}\RG{25.048}{26.821}{10.9091}{rm}{91}\RG{28.070}{26.821}{10.9091}{rm}{91}\RG{31.109}{26.821}{10.9091}{mi}{109}\RG{40.453}{26.821}{10.9091}{rm}{93}\RG{43.420}{26.821}{10.9091}{mi}{59}\RG{48.198}{26.821}{10.9091}{mi}{74}\RG{55.119}{26.821}{10.9091}{rm}{93}\RG{82.264}{26.821}{10.9091}{mi}{34}\RG{87.263}{26.821}{10.9091}{rm}{40}\RG{91.416}{26.821}{10.9091}{mi}{28}\RG{97.274}{26.821}{10.9091}{rm}{41}\RG{119.156}{26.821}{10.9091}{mi}{118}\RG{124.320}{25.185}{7.9701}{mi}{115}\RG{141.125}{25.185}{7.9701}{mi}{59}\RG{143.270}{25.185}{7.9701}{mi}{107}\RG{127.893}{23.599}{5.9776}{mi}{28}\RG{131.699}{23.599}{5.9776}{rm}{40}\RG{134.326}{23.599}{5.9776}{mi}{107}\RG{138.044}{23.599}{5.9776}{rm}{41}\RG{59.905}{15.791}{7.9701}{mi}{28}\RG{64.174}{15.791}{7.9701}{sy}{50}\RG{69.325}{15.791}{7.9701}{mi}{83}\RG{103.246}{15.711}{7.9701}{mi}{107}\RG{107.464}{15.711}{7.9701}{rm}{61}\RG{113.471}{15.711}{7.9701}{rm}{49}\RG{74.026}{14.882}{5.9776}{mi}{109}\RG{23.031}{0.002}{10.9091}{rm}{61}\RG{34.366}{0.002}{10.9091}{rm}{100}\RG{40.300}{0.002}{10.9091}{rm}{101}\RG{45.046}{0.002}{10.9091}{rm}{116}\RG{51.023}{0.002}{10.9091}{mi}{85}\RG{59.417}{0.002}{10.9091}{rm}{91}\RG{62.439}{0.002}{10.9091}{rm}{91}\RG{65.478}{0.002}{10.9091}{mi}{109}\RG{74.821}{0.002}{10.9091}{rm}{93}\RG{77.789}{0.002}{10.9091}{mi}{59}\RG{82.567}{0.002}{10.9091}{mi}{74}\RG{89.488}{0.002}{10.9091}{rm}{93}\RG{94.267}{0.002}{10.9091}{rm}{100}\RG{100.201}{0.002}{10.9091}{rm}{101}\RG{104.947}{0.002}{10.9091}{rm}{116}\RG{110.929}{0.002}{10.9091}{mi}{86}\RG{119.529}{0.002}{10.9091}{rm}{91}\RG{122.560}{0.002}{10.9091}{mi}{74}\RG{128.287}{0.002}{10.9091}{mi}{59}\RG{133.073}{0.002}{10.9091}{rm}{91}\RG{136.103}{0.002}{10.9091}{mi}{109}\RG{145.446}{0.002}{10.9091}{rm}{93}\RG{148.414}{0.002}{10.9091}{rm}{93}\RG{151.381}{0.002}{10.9091}{mi}{58}%
}
\]
因此Cauchy–Binet公式为

\[
\RB{197.887}{9.589}{22.199}{%
\RG{58.964}{9.089}{9.0910}{ex}{88}\RG{0.000}{-0.002}{10.9091}{rm}{100}\RG{5.935}{-0.002}{10.9091}{rm}{101}\RG{10.680}{-0.002}{10.9091}{rm}{116}\RG{14.837}{-0.002}{10.9091}{rm}{40}\RG{18.989}{-0.002}{10.9091}{mi}{85}\RG{27.379}{-0.002}{10.9091}{mi}{86}\RG{35.985}{-0.002}{10.9091}{rm}{41}\RG{43.163}{-0.002}{10.9091}{rm}{61}\RG{78.368}{-0.002}{10.9091}{rm}{100}\RG{84.303}{-0.002}{10.9091}{rm}{101}\RG{89.049}{-0.002}{10.9091}{rm}{116}\RG{95.023}{-0.002}{10.9091}{mi}{85}\RG{103.417}{-0.002}{10.9091}{rm}{91}\RG{106.439}{-0.002}{10.9091}{rm}{91}\RG{109.478}{-0.002}{10.9091}{mi}{109}\RG{118.821}{-0.002}{10.9091}{rm}{93}\RG{121.789}{-0.002}{10.9091}{mi}{59}\RG{126.567}{-0.002}{10.9091}{mi}{74}\RG{133.488}{-0.002}{10.9091}{rm}{93}\RG{138.266}{-0.002}{10.9091}{rm}{100}\RG{144.201}{-0.002}{10.9091}{rm}{101}\RG{148.946}{-0.002}{10.9091}{rm}{116}\RG{154.927}{-0.002}{10.9091}{mi}{86}\RG{163.529}{-0.002}{10.9091}{rm}{91}\RG{166.559}{-0.002}{10.9091}{mi}{74}\RG{172.287}{-0.002}{10.9091}{mi}{59}\RG{177.072}{-0.002}{10.9091}{rm}{91}\RG{180.103}{-0.002}{10.9091}{mi}{109}\RG{189.446}{-0.002}{10.9091}{rm}{93}\RG{192.413}{-0.002}{10.9091}{rm}{93}\RG{195.380}{-0.002}{10.9091}{mi}{58}\RG{55.967}{-11.517}{7.9701}{mi}{74}\RG{60.872}{-11.517}{7.9701}{sy}{18}\RG{66.882}{-11.517}{7.9701}{rm}{91}\RG{69.054}{-11.517}{7.9701}{mi}{112}\RG{72.944}{-11.517}{7.9701}{rm}{93}\RG{54.508}{-19.881}{7.9701}{sy}{106}\RG{56.654}{-19.881}{7.9701}{mi}{74}\RG{61.559}{-19.881}{7.9701}{sy}{106}\RG{63.705}{-19.881}{7.9701}{rm}{61}\RG{69.715}{-19.881}{7.9701}{mi}{109}%
}
\]
当\allowbreak{}\hskip.1em \RB{29.410}{6.251}{2.615}{%
\RG{0.000}{0.002}{10.9091}{mi}{112}\RG{8.466}{0.002}{10.9091}{mi}{60}\RG{19.888}{0.002}{10.9091}{mi}{109}%
}\hskip.1em\allowbreak{} 时，右端为空和，左端展开的每一项均含重复列，故两端同为零。当\allowbreak{}\hskip.1em \RB{50.318}{5.313}{2.615}{%
\RG{0.000}{0.002}{10.9091}{mi}{109}\RG{12.454}{0.002}{10.9091}{rm}{61}\RG{23.872}{0.002}{10.9091}{mi}{112}\RG{32.344}{0.002}{10.9091}{rm}{61}\RG{43.763}{0.002}{10.9091}{mi}{110}%
}\hskip.1em\allowbreak{}时，求和仅含\allowbreak{}\hskip.1em \RB{32.881}{8.684}{3.226}{%
\RG{0.000}{0.002}{10.9091}{mi}{74}\RG{9.945}{0.002}{10.9091}{rm}{61}\RG{21.279}{0.002}{10.9091}{rm}{91}\RG{24.314}{0.002}{10.9091}{mi}{110}\RG{30.690}{0.002}{10.9091}{rm}{93}%
}\hskip.1em\allowbreak{} 一项，因而\allowbreak{}\hskip.1em \RB{105.926}{8.684}{3.237}{%
\RG{0.000}{0.002}{10.9091}{rm}{100}\RG{5.935}{0.002}{10.9091}{rm}{101}\RG{10.680}{0.002}{10.9091}{rm}{116}\RG{14.837}{0.002}{10.9091}{rm}{40}\RG{18.990}{0.002}{10.9091}{mi}{65}\RG{26.998}{0.002}{10.9091}{mi}{66}\RG{35.663}{0.002}{10.9091}{rm}{41}\RG{42.841}{0.002}{10.9091}{rm}{61}\RG{54.187}{0.002}{10.9091}{rm}{100}\RG{60.121}{0.002}{10.9091}{rm}{101}\RG{64.867}{0.002}{10.9091}{rm}{116}\RG{70.840}{0.002}{10.9091}{mi}{65}\RG{80.666}{0.002}{10.9091}{rm}{100}\RG{86.601}{0.002}{10.9091}{rm}{101}\RG{91.346}{0.002}{10.9091}{rm}{116}\RG{97.320}{0.002}{10.9091}{mi}{66}%
}\hskip.1em\allowbreak{} 。

% Original page 7.
（3.1）依伴随矩阵的定义，\allowbreak{}\hskip.1em \RB{56.332}{8.685}{3.618}{%
\RG{0.000}{0.003}{10.9091}{rm}{40}\RG{4.154}{0.003}{10.9091}{mi}{65}\RG{12.162}{3.961}{7.9701}{sy}{3}\RG{16.480}{0.003}{10.9091}{rm}{41}\RG{20.634}{-1.634}{7.9701}{mi}{105}\RG{23.266}{-1.634}{7.9701}{mi}{106}\RG{30.649}{0.003}{10.9091}{rm}{61}\RG{42.342}{0.003}{10.9091}{mi}{67}\RG{49.985}{-1.634}{7.9701}{mi}{106}\RG{53.527}{-1.634}{7.9701}{mi}{105}%
}\hskip.1em\allowbreak{} 。由（1）中的第九条性质，矩阵\allowbreak{}\hskip.1em \RB{19.868}{8.292}{0.500}{%
\RG{0.000}{0.003}{10.9091}{mi}{65}\RG{8.007}{0.003}{10.9091}{mi}{65}\RG{16.016}{3.961}{7.9701}{sy}{3}%
}\hskip.1em\allowbreak{} 与\allowbreak{}\hskip.1em \RB{20.528}{8.292}{0.500}{%
\RG{0.000}{0.003}{10.9091}{mi}{65}\RG{8.008}{3.961}{7.9701}{sy}{3}\RG{12.326}{0.003}{10.9091}{mi}{65}%
}\hskip.1em\allowbreak{}的第\allowbreak{}\hskip.1em \RB{21.192}{8.685}{3.235}{%
\RG{0.000}{0.003}{10.9091}{rm}{40}\RG{4.154}{0.003}{10.9091}{mi}{105}\RG{7.797}{0.003}{10.9091}{mi}{59}\RG{12.576}{0.003}{10.9091}{mi}{106}\RG{17.616}{0.003}{10.9091}{rm}{41}%
}\hskip.1em\allowbreak{} 个元素均等于\allowbreak{}\hskip.1em \RB{38.036}{8.292}{3.617}{%
\RG{0.000}{0.003}{10.9091}{mi}{14}\RG{4.731}{-1.633}{7.9701}{mi}{105}\RG{7.364}{-1.633}{7.9701}{mi}{106}\RG{13.179}{0.003}{10.9091}{rm}{100}\RG{19.114}{0.003}{10.9091}{rm}{101}\RG{23.860}{0.003}{10.9091}{rm}{116}\RG{29.834}{0.003}{10.9091}{mi}{65}%
}\hskip.1em\allowbreak{} ，故\allowbreak{}\hskip.1em \RB{112.037}{8.685}{3.235}{%
\RG{0.000}{0.003}{10.9091}{mi}{65}\RG{8.007}{0.003}{10.9091}{mi}{65}\RG{16.015}{3.962}{7.9701}{sy}{3}\RG{23.364}{0.003}{10.9091}{rm}{61}\RG{34.703}{0.003}{10.9091}{mi}{65}\RG{42.710}{3.962}{7.9701}{sy}{3}\RG{47.029}{0.003}{10.9091}{mi}{65}\RG{58.067}{0.003}{10.9091}{rm}{61}\RG{69.402}{0.003}{10.9091}{rm}{40}\RG{73.558}{0.003}{10.9091}{rm}{100}\RG{79.493}{0.003}{10.9091}{rm}{101}\RG{84.238}{0.003}{10.9091}{rm}{116}\RG{90.214}{0.003}{10.9091}{mi}{65}\RG{98.221}{0.003}{10.9091}{rm}{41}\RG{102.375}{0.003}{10.9091}{mi}{73}\RG{107.086}{-1.633}{7.9701}{mi}{110}%
}\hskip.1em\allowbreak{} 。（3.2）以下以\allowbreak{}\hskip.1em \RB{5.600}{10.499}{0.612}{%
\RG{0.000}{3.235}{9.0910}{ex}{98}\RG{0.704}{-0.003}{10.9091}{mi}{105}%
}\hskip.1em\allowbreak{} 简记指标集\allowbreak{}\hskip.1em \RB{38.183}{8.679}{3.230}{%
\RG{0.000}{-0.003}{10.9091}{rm}{91}\RG{3.030}{-0.003}{10.9091}{mi}{110}\RG{9.406}{-0.003}{10.9091}{rm}{93}\RG{15.422}{-0.003}{10.9091}{sy}{110}\RG{23.921}{-0.003}{10.9091}{sy}{102}\RG{29.380}{-0.003}{10.9091}{mi}{105}\RG{33.025}{-0.003}{10.9091}{sy}{103}%
}\hskip.1em\allowbreak{} 。对\allowbreak{}\hskip.1em \RB{60.695}{7.752}{2.728}{%
\RG{0.000}{-0.003}{10.9091}{rm}{49}\RG{9.933}{-0.003}{10.9091}{sy}{20}\RG{23.010}{-0.003}{10.9091}{mi}{105}\RG{26.654}{-0.003}{10.9091}{mi}{59}\RG{31.432}{-0.003}{10.9091}{mi}{106}\RG{41.064}{-0.003}{10.9091}{sy}{20}\RG{54.140}{-0.003}{10.9091}{mi}{110}%
}\hskip.1em\allowbreak{} ，将Cauchy–Binet公式用于\allowbreak{}\hskip.1em \RB{70.888}{10.499}{3.230}{%
\RG{0.000}{-0.003}{10.9091}{mi}{65}\RG{8.008}{-0.003}{10.9091}{rm}{91}\RG{11.030}{3.236}{9.0910}{ex}{98}\RG{11.038}{-0.003}{10.9091}{mi}{106}\RG{16.067}{-0.003}{10.9091}{mi}{59}\RG{20.856}{-0.003}{10.9091}{rm}{91}\RG{23.887}{-0.003}{10.9091}{mi}{110}\RG{30.263}{-0.003}{10.9091}{rm}{93}\RG{33.231}{-0.003}{10.9091}{rm}{93}\RG{36.198}{-0.003}{10.9091}{mi}{66}\RG{44.863}{-0.003}{10.9091}{rm}{91}\RG{47.885}{-0.003}{10.9091}{rm}{91}\RG{50.924}{-0.003}{10.9091}{mi}{110}\RG{57.300}{-0.003}{10.9091}{rm}{93}\RG{60.268}{-0.003}{10.9091}{mi}{59}\RG{64.349}{3.236}{9.0910}{ex}{98}\RG{65.053}{-0.003}{10.9091}{mi}{105}\RG{68.697}{-0.003}{10.9091}{rm}{93}%
}\hskip.1em\allowbreak{} ，得到

\[
\RB{202.398}{11.412}{68.147}{%
\RG{110.034}{10.549}{9.0910}{ex}{16}\RG{159.904}{10.549}{9.0910}{ex}{17}\RG{29.135}{4.507}{7.9701}{sy}{3}\RG{80.739}{4.507}{7.9701}{mi}{105}\RG{83.371}{4.507}{7.9701}{rm}{43}\RG{89.381}{4.507}{7.9701}{mi}{106}\RG{143.465}{3.240}{9.0910}{ex}{98}\RG{152.589}{3.240}{9.0910}{ex}{98}\RG{0.000}{0.003}{10.9091}{rm}{40}\RG{4.156}{0.003}{10.9091}{rm}{40}\RG{8.308}{0.003}{10.9091}{mi}{65}\RG{16.315}{0.003}{10.9091}{mi}{66}\RG{24.981}{0.003}{10.9091}{rm}{41}\RG{33.453}{0.003}{10.9091}{rm}{41}\RG{47.267}{0.003}{10.9091}{rm}{61}\RG{58.602}{0.003}{10.9091}{rm}{40}\RG{62.760}{0.003}{10.9091}{sy}{0}\RG{71.244}{0.003}{10.9091}{rm}{49}\RG{76.590}{0.003}{10.9091}{rm}{41}\RG{95.198}{0.003}{10.9091}{rm}{100}\RG{101.133}{0.003}{10.9091}{rm}{101}\RG{105.878}{0.003}{10.9091}{rm}{116}\RG{115.463}{0.003}{10.9091}{rm}{40}\RG{119.617}{0.003}{10.9091}{mi}{65}\RG{127.624}{0.003}{10.9091}{mi}{66}\RG{136.290}{0.003}{10.9091}{rm}{41}\RG{140.447}{0.003}{10.9091}{rm}{91}\RG{143.474}{0.003}{10.9091}{mi}{106}\RG{148.504}{0.003}{10.9091}{mi}{59}\RG{153.292}{0.003}{10.9091}{mi}{105}\RG{156.937}{0.003}{10.9091}{rm}{93}\RG{37.607}{-1.633}{7.9701}{mi}{105}\RG{40.240}{-1.633}{7.9701}{mi}{106}\RG{63.306}{-13.129}{7.9701}{mi}{110}\RG{59.085}{-15.856}{9.0910}{ex}{88}\RG{94.830}{-20.443}{7.9701}{mi}{105}\RG{97.462}{-20.443}{7.9701}{rm}{43}\RG{103.472}{-20.443}{7.9701}{mi}{106}\RG{147.235}{-21.285}{9.0910}{ex}{98}\RG{186.291}{-21.285}{9.0910}{ex}{98}\RG{136.972}{-21.710}{9.0910}{ex}{98}\RG{195.860}{-21.710}{9.0910}{ex}{98}\RG{47.267}{-24.947}{10.9091}{rm}{61}\RG{72.696}{-24.947}{10.9091}{rm}{40}\RG{76.851}{-24.947}{10.9091}{sy}{0}\RG{85.335}{-24.947}{10.9091}{rm}{49}\RG{90.681}{-24.947}{10.9091}{rm}{41}\RG{109.289}{-24.947}{10.9091}{rm}{100}\RG{115.224}{-24.947}{10.9091}{rm}{101}\RG{119.969}{-24.947}{10.9091}{rm}{116}\RG{125.942}{-24.947}{10.9091}{mi}{65}\RG{133.950}{-24.947}{10.9091}{rm}{91}\RG{136.981}{-24.947}{10.9091}{mi}{106}\RG{142.010}{-24.947}{10.9091}{mi}{59}\RG{146.799}{-24.947}{10.9091}{mi}{107}\RG{152.721}{-24.947}{10.9091}{rm}{93}\RG{157.499}{-24.947}{10.9091}{rm}{100}\RG{163.434}{-24.947}{10.9091}{rm}{101}\RG{168.179}{-24.947}{10.9091}{rm}{116}\RG{174.160}{-24.947}{10.9091}{mi}{66}\RG{182.826}{-24.947}{10.9091}{rm}{91}\RG{185.856}{-24.947}{10.9091}{mi}{107}\RG{191.769}{-24.947}{10.9091}{mi}{59}\RG{196.563}{-24.947}{10.9091}{mi}{105}\RG{200.207}{-24.947}{10.9091}{rm}{93}\RG{58.605}{-36.059}{7.9701}{mi}{107}\RG{62.823}{-36.059}{7.9701}{rm}{61}\RG{68.830}{-36.059}{7.9701}{rm}{49}\RG{63.306}{-44.645}{7.9701}{mi}{110}\RG{59.085}{-47.372}{9.0910}{ex}{88}\RG{85.516}{-51.958}{7.9701}{sy}{3}\RG{113.453}{-51.958}{7.9701}{sy}{3}\RG{157.329}{-51.958}{7.9701}{sy}{3}\RG{169.655}{-51.958}{7.9701}{sy}{3}\RG{47.267}{-56.463}{10.9091}{rm}{61}\RG{72.696}{-56.463}{10.9091}{rm}{40}\RG{76.851}{-56.463}{10.9091}{mi}{66}\RG{89.834}{-56.463}{10.9091}{rm}{41}\RG{101.291}{-56.463}{10.9091}{rm}{40}\RG{105.445}{-56.463}{10.9091}{mi}{65}\RG{117.771}{-56.463}{10.9091}{rm}{41}\RG{133.171}{-56.463}{10.9091}{rm}{61}\RG{144.506}{-56.463}{10.9091}{rm}{40}\RG{148.664}{-56.463}{10.9091}{mi}{66}\RG{161.648}{-56.463}{10.9091}{mi}{65}\RG{173.973}{-56.463}{10.9091}{rm}{41}\RG{184.757}{-56.463}{10.9091}{mi}{58}\RG{93.988}{-58.100}{7.9701}{mi}{105}\RG{96.621}{-58.100}{7.9701}{mi}{107}\RG{121.925}{-58.100}{7.9701}{mi}{107}\RG{126.136}{-58.100}{7.9701}{mi}{106}\RG{178.128}{-58.100}{7.9701}{mi}{105}\RG{180.761}{-58.100}{7.9701}{mi}{106}\RG{58.605}{-67.574}{7.9701}{mi}{107}\RG{62.823}{-67.574}{7.9701}{rm}{61}\RG{68.830}{-67.574}{7.9701}{rm}{49}%
}
\]
故\allowbreak{}\hskip.1em \RB{68.513}{8.681}{3.239}{%
\RG{0.000}{-0.001}{10.9091}{rm}{40}\RG{4.155}{-0.001}{10.9091}{mi}{65}\RG{12.162}{-0.001}{10.9091}{mi}{66}\RG{20.827}{-0.001}{10.9091}{rm}{41}\RG{24.982}{3.958}{7.9701}{sy}{3}\RG{32.330}{-0.001}{10.9091}{rm}{61}\RG{43.668}{-0.001}{10.9091}{mi}{66}\RG{52.334}{3.958}{7.9701}{sy}{3}\RG{56.652}{-0.001}{10.9091}{mi}{65}\RG{64.660}{3.958}{7.9701}{sy}{3}%
}\hskip.1em\allowbreak{} ，此证明不要求\allowbreak{}\hskip.1em \RB{8.202}{8.288}{0.501}{%
\RG{0.000}{-0.001}{10.9091}{mi}{65}%
}\hskip.1em\allowbreak{} 或\allowbreak{}\hskip.1em \RB{8.606}{7.950}{0.501}{%
\RG{0.000}{-0.001}{10.9091}{mi}{66}%
}\hskip.1em\allowbreak{} 可逆。（3.3）与（3.4）引入独立于\allowbreak{}\hskip.1em \RB{6.336}{8.073}{0.618}{%
\RG{0.000}{0.002}{10.9091}{mi}{21}%
}\hskip.1em\allowbreak{} 的不定元\allowbreak{}\hskip.1em \RB{4.013}{7.320}{0.607}{%
\RG{0.000}{0.002}{10.9091}{mi}{116}%
}\hskip.1em\allowbreak{} ，令\allowbreak{}\hskip.1em \RB{113.581}{8.684}{3.236}{%
\RG{0.000}{0.002}{10.9091}{mi}{72}\RG{12.729}{0.002}{10.9091}{rm}{61}\RG{24.073}{0.002}{10.9091}{mi}{65}\RG{34.507}{0.002}{10.9091}{rm}{43}\RG{45.241}{0.002}{10.9091}{mi}{116}\RG{49.103}{0.002}{10.9091}{mi}{73}\RG{53.810}{-1.634}{7.9701}{mi}{110}\RG{61.989}{0.002}{10.9091}{sy}{50}\RG{72.298}{0.002}{10.9091}{mi}{77}\RG{82.632}{-1.634}{7.9701}{mi}{110}\RG{87.775}{0.002}{10.9091}{rm}{40}\RG{91.928}{0.002}{10.9091}{mi}{82}\RG{100.149}{0.002}{10.9091}{rm}{91}\RG{103.180}{0.002}{10.9091}{mi}{116}\RG{107.037}{0.002}{10.9091}{rm}{93}\RG{110.004}{0.002}{10.9091}{rm}{41}%
}\hskip.1em\allowbreak{} 、\allowbreak{}\hskip.1em \RB{58.610}{8.684}{3.236}{%
\RG{0.000}{0.002}{10.9091}{mi}{100}\RG{5.551}{0.002}{10.9091}{rm}{40}\RG{9.705}{0.002}{10.9091}{mi}{116}\RG{13.562}{0.002}{10.9091}{rm}{41}\RG{20.751}{0.002}{10.9091}{rm}{61}\RG{32.097}{0.002}{10.9091}{rm}{100}\RG{38.031}{0.002}{10.9091}{rm}{101}\RG{42.777}{0.002}{10.9091}{rm}{116}\RG{48.750}{0.002}{10.9091}{mi}{72}%
}\hskip.1em\allowbreak{} 。在行列式的定义式中，只有主对角置换的乘积含\allowbreak{}\hskip.1em \RB{8.808}{7.669}{0.606}{%
\RG{0.000}{0.003}{10.9091}{mi}{116}\RG{3.857}{3.961}{7.9701}{mi}{110}%
}\hskip.1em\allowbreak{} 项，且其系数为\allowbreak{}\hskip.1em \RB{4.973}{7.757}{0.500}{%
\RG{0.000}{0.003}{10.9091}{rm}{49}%
}\hskip.1em\allowbreak{} 。因此多项式\allowbreak{}\hskip.1em \RB{17.138}{8.685}{3.235}{%
\RG{0.000}{0.003}{10.9091}{mi}{100}\RG{5.551}{0.003}{10.9091}{rm}{40}\RG{9.705}{0.003}{10.9091}{mi}{116}\RG{13.562}{0.003}{10.9091}{rm}{41}%
}\hskip.1em\allowbreak{} 关于\allowbreak{}\hskip.1em \RB{4.013}{7.322}{0.605}{%
\RG{0.000}{0.004}{10.9091}{mi}{116}%
}\hskip.1em\allowbreak{} 首一且次数为\allowbreak{}\hskip.1em \RB{6.555}{5.315}{0.605}{%
\RG{0.000}{0.004}{10.9091}{mi}{110}%
}\hskip.1em\allowbreak{} ，故它在域\allowbreak{}\hskip.1em \RB{19.986}{8.686}{3.235}{%
\RG{0.000}{0.004}{10.9091}{mi}{70}\RG{8.399}{0.004}{10.9091}{rm}{40}\RG{12.553}{0.004}{10.9091}{mi}{116}\RG{16.410}{0.004}{10.9091}{rm}{41}%
}\hskip.1em\allowbreak{} 中非零。\allowbreak{}\hskip.1em \RB{17.299}{8.686}{3.224}{%
\RG{0.000}{0.004}{10.9091}{mi}{82}\RG{8.221}{0.004}{10.9091}{rm}{91}\RG{11.251}{0.004}{10.9091}{mi}{116}\RG{15.108}{0.004}{10.9091}{rm}{93}%
}\hskip.1em\allowbreak{} 仍为特征零整环，故（1）、（2）及（3.1）的上述证明对\allowbreak{}\hskip.1em \RB{17.299}{8.686}{3.224}{%
\RG{0.000}{0.004}{10.9091}{mi}{82}\RG{8.221}{0.004}{10.9091}{rm}{91}\RG{11.251}{0.004}{10.9091}{mi}{116}\RG{15.108}{0.004}{10.9091}{rm}{93}%
}\hskip.1em\allowbreak{} 中的矩阵仍成立。由\allowbreak{}\hskip.1em \RB{65.452}{8.686}{3.235}{%
\RG{0.000}{0.004}{10.9091}{mi}{72}\RG{9.687}{0.004}{10.9091}{mi}{72}\RG{19.386}{3.963}{7.9701}{sy}{3}\RG{26.736}{0.004}{10.9091}{rm}{61}\RG{38.074}{0.004}{10.9091}{mi}{100}\RG{43.625}{0.004}{10.9091}{rm}{40}\RG{47.778}{0.004}{10.9091}{mi}{116}\RG{51.635}{0.004}{10.9091}{rm}{41}\RG{55.790}{0.004}{10.9091}{mi}{73}\RG{60.501}{-1.632}{7.9701}{mi}{110}%
}\hskip.1em\allowbreak{} 及行列式乘法公式，有

\[
\RB{157.569}{11.410}{5.945}{%
\RG{79.404}{10.546}{9.0910}{ex}{16}\RG{112.404}{10.546}{9.0910}{ex}{17}\RG{45.881}{4.504}{7.9701}{sy}{3}\RG{149.919}{4.504}{7.9701}{mi}{110}\RG{0.000}{0.000}{10.9091}{mi}{100}\RG{5.551}{0.000}{10.9091}{rm}{40}\RG{9.705}{0.000}{10.9091}{mi}{116}\RG{13.562}{0.000}{10.9091}{rm}{41}\RG{19.529}{0.000}{10.9091}{rm}{100}\RG{25.464}{0.000}{10.9091}{rm}{101}\RG{30.209}{0.000}{10.9091}{rm}{116}\RG{36.188}{0.000}{10.9091}{mi}{72}\RG{53.230}{0.000}{10.9091}{rm}{61}\RG{64.564}{0.000}{10.9091}{rm}{100}\RG{70.499}{0.000}{10.9091}{rm}{101}\RG{75.245}{0.000}{10.9091}{rm}{116}\RG{84.834}{0.000}{10.9091}{mi}{100}\RG{90.384}{0.000}{10.9091}{rm}{40}\RG{94.538}{0.000}{10.9091}{mi}{116}\RG{98.395}{0.000}{10.9091}{rm}{41}\RG{102.550}{0.000}{10.9091}{mi}{73}\RG{120.864}{0.000}{10.9091}{rm}{61}\RG{132.202}{0.000}{10.9091}{mi}{100}\RG{137.753}{0.000}{10.9091}{rm}{40}\RG{141.907}{0.000}{10.9091}{mi}{116}\RG{145.764}{0.000}{10.9091}{rm}{41}\RG{155.061}{0.000}{10.9091}{mi}{58}\RG{107.261}{-1.636}{7.9701}{mi}{110}%
}
\]
在\allowbreak{}\hskip.1em \RB{19.987}{8.677}{3.243}{%
\RG{0.000}{-0.005}{10.9091}{mi}{70}\RG{8.399}{-0.005}{10.9091}{rm}{40}\RG{12.553}{-0.005}{10.9091}{mi}{116}\RG{16.411}{-0.005}{10.9091}{rm}{41}%
}\hskip.1em\allowbreak{} 中约去\allowbreak{}\hskip.1em \RB{17.138}{8.677}{3.243}{%
\RG{0.000}{-0.005}{10.9091}{mi}{100}\RG{5.550}{-0.005}{10.9091}{rm}{40}\RG{9.705}{-0.005}{10.9091}{mi}{116}\RG{13.562}{-0.005}{10.9091}{rm}{41}%
}\hskip.1em\allowbreak{} ，得到\allowbreak{}\hskip.1em \RB{80.626}{9.291}{3.243}{%
\RG{0.000}{-0.005}{10.9091}{rm}{100}\RG{5.935}{-0.005}{10.9091}{rm}{101}\RG{10.680}{-0.005}{10.9091}{rm}{116}\RG{16.654}{-0.005}{10.9091}{mi}{72}\RG{26.347}{3.955}{7.9701}{sy}{3}\RG{35.408}{-0.005}{10.9091}{rm}{61}\RG{48.460}{-0.005}{10.9091}{mi}{100}\RG{54.011}{-0.005}{10.9091}{rm}{40}\RG{58.164}{-0.005}{10.9091}{mi}{116}\RG{62.022}{-0.005}{10.9091}{rm}{41}\RG{66.175}{3.955}{7.9701}{mi}{110}\RG{70.864}{3.955}{7.9701}{sy}{0}\RG{76.875}{3.955}{7.9701}{rm}{49}%
}\hskip.1em\allowbreak{} 。又由（3.1）得到\allowbreak{}\hskip.1em \RB{79.406}{9.291}{3.243}{%
\RG{0.000}{-0.005}{10.9091}{mi}{72}\RG{9.694}{3.955}{7.9701}{sy}{0}\RG{15.703}{3.955}{7.9701}{rm}{49}\RG{24.765}{-0.005}{10.9091}{rm}{61}\RG{37.816}{-0.005}{10.9091}{mi}{100}\RG{43.366}{-0.005}{10.9091}{rm}{40}\RG{47.521}{-0.005}{10.9091}{mi}{116}\RG{51.378}{-0.005}{10.9091}{rm}{41}\RG{55.532}{3.955}{7.9701}{sy}{0}\RG{61.543}{3.955}{7.9701}{rm}{49}\RG{65.861}{-0.005}{10.9091}{mi}{72}\RG{75.553}{3.955}{7.9701}{sy}{3}%
}\hskip.1em\allowbreak{} ，从而\allowbreak{}\hskip.1em \RB{84.922}{9.292}{3.242}{%
\RG{0.000}{-0.004}{10.9091}{rm}{40}\RG{4.154}{-0.004}{10.9091}{mi}{72}\RG{13.847}{3.956}{7.9701}{sy}{3}\RG{18.165}{-0.004}{10.9091}{rm}{41}\RG{22.320}{3.956}{7.9701}{sy}{0}\RG{28.330}{3.956}{7.9701}{rm}{49}\RG{35.678}{-0.004}{10.9091}{rm}{61}\RG{47.017}{-0.004}{10.9091}{mi}{100}\RG{52.568}{-0.004}{10.9091}{rm}{40}\RG{56.722}{-0.004}{10.9091}{mi}{116}\RG{60.579}{-0.004}{10.9091}{rm}{41}\RG{64.733}{3.956}{7.9701}{sy}{0}\RG{70.743}{3.956}{7.9701}{rm}{49}\RG{75.061}{-0.004}{10.9091}{mi}{72}%
}\hskip.1em\allowbreak{} 。再次应用（3.1），有

\[
\RB{170.174}{27.335}{3.235}{%
\RG{13.847}{21.999}{7.9701}{sy}{3}\RG{22.320}{21.999}{7.9701}{sy}{3}\RG{71.508}{21.999}{7.9701}{sy}{3}\RG{93.828}{21.999}{7.9701}{sy}{3}\RG{102.300}{21.999}{7.9701}{sy}{0}\RG{108.311}{21.999}{7.9701}{rm}{49}\RG{0.000}{17.495}{10.9091}{rm}{40}\RG{4.154}{17.495}{10.9091}{mi}{72}\RG{18.165}{17.495}{10.9091}{rm}{41}\RG{29.669}{17.495}{10.9091}{rm}{61}\RG{41.003}{17.495}{10.9091}{rm}{40}\RG{45.160}{17.495}{10.9091}{rm}{100}\RG{51.094}{17.495}{10.9091}{rm}{101}\RG{55.840}{17.495}{10.9091}{rm}{116}\RG{61.815}{17.495}{10.9091}{mi}{72}\RG{75.826}{17.495}{10.9091}{rm}{41}\RG{79.983}{17.495}{10.9091}{rm}{40}\RG{84.135}{17.495}{10.9091}{mi}{72}\RG{98.146}{17.495}{10.9091}{rm}{41}\RG{58.723}{4.508}{7.9701}{mi}{110}\RG{63.411}{4.508}{7.9701}{sy}{0}\RG{69.422}{4.508}{7.9701}{rm}{49}\RG{91.456}{4.508}{7.9701}{sy}{0}\RG{97.466}{4.508}{7.9701}{rm}{49}\RG{143.562}{4.508}{7.9701}{mi}{110}\RG{148.251}{4.508}{7.9701}{sy}{0}\RG{154.261}{4.508}{7.9701}{rm}{50}\RG{29.668}{0.004}{10.9091}{rm}{61}\RG{41.006}{0.004}{10.9091}{mi}{100}\RG{46.557}{0.004}{10.9091}{rm}{40}\RG{50.711}{0.004}{10.9091}{mi}{116}\RG{54.569}{0.004}{10.9091}{rm}{41}\RG{73.740}{0.004}{10.9091}{mi}{100}\RG{79.291}{0.004}{10.9091}{rm}{40}\RG{83.445}{0.004}{10.9091}{mi}{116}\RG{87.302}{0.004}{10.9091}{rm}{41}\RG{101.784}{0.004}{10.9091}{mi}{72}\RG{114.508}{0.004}{10.9091}{rm}{61}\RG{125.846}{0.004}{10.9091}{mi}{100}\RG{131.397}{0.004}{10.9091}{rm}{40}\RG{135.551}{0.004}{10.9091}{mi}{116}\RG{139.408}{0.004}{10.9091}{rm}{41}\RG{158.579}{0.004}{10.9091}{mi}{72}\RG{167.667}{0.004}{10.9091}{mi}{58}%
}
\]
这些等式两端的各个标量或矩阵元素均属于\allowbreak{}\hskip.1em \RB{17.298}{8.681}{3.228}{%
\RG{0.000}{-0.001}{10.9091}{mi}{82}\RG{8.220}{-0.001}{10.9091}{rm}{91}\RG{11.250}{-0.001}{10.9091}{mi}{116}\RG{15.108}{-0.001}{10.9091}{rm}{93}%
}\hskip.1em\allowbreak{} ，而\allowbreak{}\hskip.1em \RB{17.299}{8.681}{3.228}{%
\RG{0.000}{-0.001}{10.9091}{mi}{82}\RG{8.220}{-0.001}{10.9091}{rm}{91}\RG{11.250}{-0.001}{10.9091}{mi}{116}\RG{15.108}{-0.001}{10.9091}{rm}{93}%
}\hskip.1em\allowbreak{} 嵌入\allowbreak{}\hskip.1em \RB{19.986}{8.681}{3.239}{%
\RG{0.000}{-0.001}{10.9091}{mi}{70}\RG{8.399}{-0.001}{10.9091}{rm}{40}\RG{12.553}{-0.001}{10.9091}{mi}{116}\RG{16.410}{-0.001}{10.9091}{rm}{41}%
}\hskip.1em\allowbreak{} ，故它们也是\allowbreak{}\hskip.1em \RB{17.299}{8.681}{3.228}{%
\RG{0.000}{-0.001}{10.9091}{mi}{82}\RG{8.221}{-0.001}{10.9091}{rm}{91}\RG{11.250}{-0.001}{10.9091}{mi}{116}\RG{15.108}{-0.001}{10.9091}{rm}{93}%
}\hskip.1em\allowbreak{} 上的标量或矩阵恒等式。对这些多项式恒等式取\allowbreak{}\hskip.1em \RB{23.610}{7.755}{0.729}{%
\RG{0.000}{-0.000}{10.9091}{mi}{116}\RG{6.888}{-0.000}{10.9091}{rm}{61}\RG{18.222}{-0.000}{10.9091}{rm}{48}%
}\hskip.1em\allowbreak{} ，得到

\[
\RB{101.851}{9.837}{20.733}{%
\RG{24.662}{4.501}{7.9701}{sy}{3}\RG{76.319}{4.501}{7.9701}{mi}{110}\RG{81.007}{4.501}{7.9701}{sy}{0}\RG{87.017}{4.501}{7.9701}{rm}{49}\RG{0.000}{-0.004}{10.9091}{rm}{100}\RG{5.935}{-0.004}{10.9091}{rm}{101}\RG{10.680}{-0.004}{10.9091}{rm}{116}\RG{16.654}{-0.004}{10.9091}{mi}{65}\RG{32.010}{-0.004}{10.9091}{rm}{61}\RG{43.345}{-0.004}{10.9091}{rm}{40}\RG{47.501}{-0.004}{10.9091}{rm}{100}\RG{53.436}{-0.004}{10.9091}{rm}{101}\RG{58.181}{-0.004}{10.9091}{rm}{116}\RG{64.157}{-0.004}{10.9091}{mi}{65}\RG{72.164}{-0.004}{10.9091}{rm}{41}\RG{91.336}{-0.004}{10.9091}{mi}{59}\RG{16.190}{-12.991}{7.9701}{sy}{3}\RG{24.662}{-12.991}{7.9701}{sy}{3}\RG{76.319}{-12.991}{7.9701}{mi}{110}\RG{81.007}{-12.991}{7.9701}{sy}{0}\RG{87.017}{-12.991}{7.9701}{rm}{50}\RG{4.028}{-17.495}{10.9091}{rm}{40}\RG{8.182}{-17.495}{10.9091}{mi}{65}\RG{20.508}{-17.495}{10.9091}{rm}{41}\RG{32.010}{-17.495}{10.9091}{rm}{61}\RG{43.345}{-17.495}{10.9091}{rm}{40}\RG{47.501}{-17.495}{10.9091}{rm}{100}\RG{53.436}{-17.495}{10.9091}{rm}{101}\RG{58.181}{-17.495}{10.9091}{rm}{116}\RG{64.157}{-17.495}{10.9091}{mi}{65}\RG{72.164}{-17.495}{10.9091}{rm}{41}\RG{91.336}{-17.495}{10.9091}{mi}{65}\RG{99.344}{-17.495}{10.9091}{mi}{58}%
}
\]
当\allowbreak{}\hskip.1em \RB{26.675}{7.757}{0.606}{%
\RG{0.000}{0.003}{10.9091}{mi}{110}\RG{9.724}{0.003}{10.9091}{rm}{61}\RG{21.375}{0.003}{10.9091}{rm}{50}%
}\hskip.1em\allowbreak{} 时，第二式中的零次幂是恒等于\allowbreak{}\hskip.1em \RB{4.973}{7.757}{0.500}{%
\RG{0.000}{0.003}{10.9091}{rm}{49}%
}\hskip.1em\allowbreak{} 的多项式，故该式为\allowbreak{}\hskip.1em \RB{48.158}{8.685}{3.236}{%
\RG{0.000}{0.003}{10.9091}{rm}{40}\RG{4.155}{0.003}{10.9091}{mi}{65}\RG{12.162}{3.962}{7.9701}{sy}{3}\RG{16.480}{0.003}{10.9091}{rm}{41}\RG{20.635}{3.962}{7.9701}{sy}{3}\RG{28.300}{0.003}{10.9091}{rm}{61}\RG{39.956}{0.003}{10.9091}{mi}{65}%
}\hskip.1em\allowbreak{} 。上述论证包括\allowbreak{}\hskip.1em \RB{44.414}{8.293}{0.725}{%
\RG{0.000}{0.004}{10.9091}{rm}{100}\RG{5.935}{0.004}{10.9091}{rm}{101}\RG{10.680}{0.004}{10.9091}{rm}{116}\RG{16.654}{0.004}{10.9091}{mi}{65}\RG{27.692}{0.004}{10.9091}{rm}{61}\RG{39.027}{0.004}{10.9091}{rm}{48}%
}\hskip.1em\allowbreak{} 的情形。

% Original page 8.
（4）此处\allowbreak{}\hskip.1em \RB{54.421}{8.685}{3.235}{%
\RG{0.000}{0.003}{10.9091}{mi}{65}\RG{11.039}{0.003}{10.9091}{sy}{50}\RG{21.342}{0.003}{10.9091}{mi}{77}\RG{31.675}{-1.634}{7.9701}{mi}{110}\RG{36.819}{0.003}{10.9091}{rm}{40}\RG{40.973}{0.003}{10.9091}{mi}{75}\RG{50.844}{0.003}{10.9091}{rm}{41}%
}\hskip.1em\allowbreak{} 是常矩阵。若\allowbreak{}\hskip.1em \RB{25.714}{7.757}{0.606}{%
\RG{0.000}{0.003}{10.9091}{mi}{110}\RG{9.406}{0.003}{10.9091}{rm}{61}\RG{20.741}{0.003}{10.9091}{rm}{49}%
}\hskip.1em\allowbreak{} ，写\allowbreak{}\hskip.1em \RB{35.714}{8.685}{3.235}{%
\RG{0.000}{0.003}{10.9091}{mi}{65}\RG{11.038}{0.003}{10.9091}{rm}{61}\RG{22.372}{0.003}{10.9091}{rm}{40}\RG{26.530}{0.003}{10.9091}{mi}{97}\RG{32.138}{0.003}{10.9091}{rm}{41}%
}\hskip.1em\allowbreak{} ，则\allowbreak{}\hskip.1em \RB{59.866}{8.685}{3.235}{%
\RG{0.000}{0.003}{10.9091}{mi}{102}\RG{6.430}{0.003}{10.9091}{rm}{40}\RG{10.584}{0.003}{10.9091}{mi}{21}\RG{16.816}{0.003}{10.9091}{rm}{41}\RG{23.994}{0.003}{10.9091}{rm}{61}\RG{35.338}{0.003}{10.9091}{mi}{21}\RG{43.580}{0.003}{10.9091}{sy}{0}\RG{54.075}{0.003}{10.9091}{mi}{97}%
}\hskip.1em\allowbreak{} ，故\allowbreak{}\hskip.1em \RB{42.503}{8.685}{3.235}{%
\RG{0.000}{0.003}{10.9091}{mi}{102}\RG{6.430}{0.003}{10.9091}{rm}{40}\RG{10.584}{0.003}{10.9091}{mi}{65}\RG{18.592}{0.003}{10.9091}{rm}{41}\RG{25.770}{0.003}{10.9091}{rm}{61}\RG{37.115}{0.003}{10.9091}{rm}{48}%
}\hskip.1em\allowbreak{} 。以下设\allowbreak{}\hskip.1em \RB{26.221}{7.757}{1.992}{%
\RG{0.000}{0.003}{10.9091}{mi}{110}\RG{9.406}{0.003}{10.9091}{sy}{21}\RG{20.921}{0.003}{10.9091}{rm}{50}%
}\hskip.1em\allowbreak{} 。记

\[
\RB{174.863}{15.531}{12.848}{%
\RG{39.703}{11.823}{7.9701}{mi}{110}\RG{35.481}{9.096}{9.0910}{ex}{88}\RG{65.401}{4.509}{7.9701}{mi}{106}\RG{0.000}{0.005}{10.9091}{mi}{102}\RG{6.430}{0.005}{10.9091}{rm}{40}\RG{10.583}{0.005}{10.9091}{mi}{21}\RG{16.815}{0.005}{10.9091}{rm}{41}\RG{23.993}{0.005}{10.9091}{rm}{61}\RG{50.574}{0.005}{10.9091}{mi}{99}\RG{59.170}{0.005}{10.9091}{mi}{21}\RG{69.399}{0.005}{10.9091}{mi}{59}\RG{95.996}{0.005}{10.9091}{mi}{99}\RG{107.630}{0.005}{10.9091}{sy}{50}\RG{117.933}{0.005}{10.9091}{mi}{75}\RG{127.206}{0.005}{10.9091}{mi}{59}\RG{142.905}{0.005}{10.9091}{mi}{99}\RG{155.677}{0.005}{10.9091}{rm}{61}\RG{167.011}{0.005}{10.9091}{rm}{49}\RG{172.355}{0.005}{10.9091}{mi}{58}\RG{55.171}{-1.632}{7.9701}{mi}{106}\RG{100.601}{-1.632}{7.9701}{mi}{106}\RG{147.502}{-1.632}{7.9701}{mi}{110}\RG{35.338}{-10.864}{7.9701}{mi}{106}\RG{38.882}{-10.864}{7.9701}{rm}{61}\RG{44.890}{-10.864}{7.9701}{rm}{48}%
}
\]
由于\allowbreak{}\hskip.1em \RB{37.620}{8.287}{2.212}{%
\RG{0.000}{-0.002}{10.9091}{mi}{21}\RG{6.229}{-0.002}{10.9091}{mi}{73}\RG{10.942}{-1.639}{7.9701}{mi}{110}\RG{18.509}{-0.002}{10.9091}{sy}{0}\RG{29.418}{-0.002}{10.9091}{mi}{65}%
}\hskip.1em\allowbreak{} 的每个\allowbreak{}\hskip.1em \RB{32.786}{8.680}{3.240}{%
\RG{0.000}{-0.002}{10.9091}{rm}{40}\RG{4.154}{-0.002}{10.9091}{mi}{110}\RG{12.955}{-0.002}{10.9091}{sy}{0}\RG{23.864}{-0.002}{10.9091}{rm}{49}\RG{29.210}{-0.002}{10.9091}{rm}{41}%
}\hskip.1em\allowbreak{} 阶子式的次数不超过\allowbreak{}\hskip.1em \RB{24.682}{7.753}{0.611}{%
\RG{0.000}{-0.002}{10.9091}{mi}{110}\RG{8.800}{-0.002}{10.9091}{sy}{0}\RG{19.709}{-0.002}{10.9091}{rm}{49}%
}\hskip.1em\allowbreak{} ，可写

\[
\RB{191.361}{17.152}{12.856}{%
\RG{64.422}{11.815}{7.9701}{mi}{110}\RG{69.110}{11.815}{7.9701}{sy}{0}\RG{75.121}{11.815}{7.9701}{rm}{49}\RG{65.137}{9.088}{9.0910}{ex}{88}\RG{45.734}{4.501}{7.9701}{sy}{3}\RG{99.147}{4.501}{7.9701}{mi}{106}\RG{0.000}{-0.004}{10.9091}{rm}{40}\RG{4.154}{-0.004}{10.9091}{mi}{21}\RG{10.383}{-0.004}{10.9091}{mi}{73}\RG{22.664}{-0.004}{10.9091}{sy}{0}\RG{33.573}{-0.004}{10.9091}{mi}{65}\RG{41.580}{-0.004}{10.9091}{rm}{41}\RG{53.083}{-0.004}{10.9091}{rm}{61}\RG{80.802}{-0.004}{10.9091}{mi}{66}\RG{92.917}{-0.004}{10.9091}{mi}{21}\RG{103.146}{-0.004}{10.9091}{mi}{59}\RG{129.743}{-0.004}{10.9091}{mi}{66}\RG{144.894}{-0.004}{10.9091}{sy}{50}\RG{155.197}{-0.004}{10.9091}{mi}{77}\RG{170.675}{-0.004}{10.9091}{rm}{40}\RG{174.828}{-0.004}{10.9091}{mi}{75}\RG{184.700}{-0.004}{10.9091}{rm}{41}\RG{188.854}{-0.004}{10.9091}{mi}{58}\RG{15.097}{-1.640}{7.9701}{mi}{110}\RG{88.917}{-1.640}{7.9701}{mi}{106}\RG{137.865}{-1.640}{7.9701}{mi}{106}\RG{165.531}{-1.640}{7.9701}{mi}{110}\RG{64.994}{-10.873}{7.9701}{mi}{106}\RG{68.538}{-10.873}{7.9701}{rm}{61}\RG{74.545}{-10.873}{7.9701}{rm}{48}%
}
\]
在（3.1）给出的恒等式\allowbreak{}\hskip.1em \RB{140.788}{8.679}{3.241}{%
\RG{0.000}{-0.003}{10.9091}{rm}{40}\RG{4.155}{-0.003}{10.9091}{mi}{21}\RG{10.384}{-0.003}{10.9091}{mi}{73}\RG{15.097}{-1.640}{7.9701}{mi}{110}\RG{22.664}{-0.003}{10.9091}{sy}{0}\RG{33.573}{-0.003}{10.9091}{mi}{65}\RG{41.581}{-0.003}{10.9091}{rm}{41}\RG{45.738}{-0.003}{10.9091}{rm}{40}\RG{49.889}{-0.003}{10.9091}{mi}{21}\RG{56.118}{-0.003}{10.9091}{mi}{73}\RG{60.831}{-1.640}{7.9701}{mi}{110}\RG{68.398}{-0.003}{10.9091}{sy}{0}\RG{79.308}{-0.003}{10.9091}{mi}{65}\RG{87.316}{-0.003}{10.9091}{rm}{41}\RG{91.469}{3.956}{7.9701}{sy}{3}\RG{98.818}{-0.003}{10.9091}{rm}{61}\RG{110.157}{-0.003}{10.9091}{mi}{102}\RG{116.587}{-0.003}{10.9091}{rm}{40}\RG{120.741}{-0.003}{10.9091}{mi}{21}\RG{126.972}{-0.003}{10.9091}{rm}{41}\RG{131.126}{-0.003}{10.9091}{mi}{73}\RG{135.837}{-1.640}{7.9701}{mi}{110}%
}\hskip.1em\allowbreak{} 中比较\allowbreak{}\hskip.1em \RB{6.336}{8.068}{0.623}{%
\RG{0.000}{-0.003}{10.9091}{mi}{21}%
}\hskip.1em\allowbreak{} 的各次幂系数，得到

\[
\RB{185.531}{25.439}{19.783}{%
\RG{32.310}{17.488}{10.9091}{mi}{66}\RG{58.473}{17.488}{10.9091}{rm}{61}\RG{69.811}{17.488}{10.9091}{mi}{73}\RG{79.666}{17.488}{10.9091}{mi}{59}\RG{40.425}{15.852}{7.9701}{mi}{110}\RG{45.114}{15.852}{7.9701}{sy}{0}\RG{51.124}{15.852}{7.9701}{rm}{49}\RG{74.522}{15.852}{7.9701}{mi}{110}\RG{0.000}{-0.003}{10.9091}{mi}{66}\RG{24.412}{-0.003}{10.9091}{sy}{0}\RG{35.321}{-0.003}{10.9091}{mi}{65}\RG{43.328}{-0.003}{10.9091}{mi}{66}\RG{58.473}{-0.003}{10.9091}{rm}{61}\RG{69.811}{-0.003}{10.9091}{mi}{99}\RG{78.408}{-0.003}{10.9091}{mi}{73}\RG{110.080}{-0.003}{10.9091}{rm}{40}\RG{114.237}{-0.003}{10.9091}{rm}{49}\RG{122.605}{-0.003}{10.9091}{sy}{20}\RG{134.121}{-0.003}{10.9091}{mi}{106}\RG{142.184}{-0.003}{10.9091}{sy}{20}\RG{153.699}{-0.003}{10.9091}{mi}{110}\RG{162.500}{-0.003}{10.9091}{sy}{0}\RG{173.409}{-0.003}{10.9091}{rm}{49}\RG{178.755}{-0.003}{10.9091}{rm}{41}\RG{182.904}{-0.003}{10.9091}{mi}{59}\RG{8.115}{-1.639}{7.9701}{mi}{106}\RG{11.659}{-1.639}{7.9701}{sy}{0}\RG{17.669}{-1.639}{7.9701}{rm}{49}\RG{51.444}{-1.639}{7.9701}{mi}{106}\RG{74.408}{-1.639}{7.9701}{mi}{106}\RG{83.118}{-1.639}{7.9701}{mi}{110}\RG{26.517}{-17.494}{10.9091}{sy}{0}\RG{35.001}{-17.494}{10.9091}{mi}{65}\RG{43.009}{-17.494}{10.9091}{mi}{66}\RG{58.473}{-17.494}{10.9091}{rm}{61}\RG{69.811}{-17.494}{10.9091}{mi}{99}\RG{78.727}{-17.494}{10.9091}{mi}{73}\RG{88.581}{-17.494}{10.9091}{mi}{58}\RG{51.124}{-19.130}{7.9701}{rm}{48}\RG{74.408}{-19.130}{7.9701}{rm}{48}\RG{83.438}{-19.130}{7.9701}{mi}{110}%
}
\]
将中间一组等式分别左乘\allowbreak{}\hskip.1em \RB{11.511}{9.271}{0.503}{%
\RG{0.000}{-0.003}{10.9091}{mi}{65}\RG{8.007}{3.957}{7.9701}{mi}{106}%
}\hskip.1em\allowbreak{} ，再利用首尾两个等式，有

\[
\RB{226.809}{78.565}{2.285}{%
\RG{63.422}{73.228}{7.9701}{mi}{110}\RG{68.111}{73.228}{7.9701}{sy}{0}\RG{74.121}{73.228}{7.9701}{rm}{49}\RG{64.138}{70.501}{9.0910}{ex}{88}\RG{45.122}{65.914}{7.9701}{mi}{110}\RG{96.406}{65.914}{7.9701}{mi}{106}\RG{0.000}{61.409}{10.9091}{mi}{102}\RG{6.430}{61.409}{10.9091}{rm}{40}\RG{10.584}{61.409}{10.9091}{mi}{65}\RG{18.592}{61.409}{10.9091}{rm}{41}\RG{25.770}{61.409}{10.9091}{rm}{61}\RG{37.115}{61.409}{10.9091}{mi}{65}\RG{52.689}{61.409}{10.9091}{rm}{43}\RG{79.803}{61.409}{10.9091}{mi}{99}\RG{88.399}{61.409}{10.9091}{mi}{65}\RG{102.829}{61.409}{10.9091}{rm}{43}\RG{113.562}{61.409}{10.9091}{mi}{99}\RG{122.477}{61.409}{10.9091}{mi}{73}\RG{84.400}{59.773}{7.9701}{mi}{106}\RG{118.159}{59.773}{7.9701}{rm}{48}\RG{127.189}{59.773}{7.9701}{mi}{110}\RG{63.995}{50.541}{7.9701}{mi}{106}\RG{67.539}{50.541}{7.9701}{rm}{61}\RG{73.546}{50.541}{7.9701}{rm}{49}\RG{86.554}{38.985}{7.9701}{mi}{110}\RG{91.243}{38.985}{7.9701}{sy}{0}\RG{97.253}{38.985}{7.9701}{rm}{49}\RG{101.116}{37.712}{9.0910}{ex}{16}\RG{187.867}{37.712}{9.0910}{ex}{17}\RG{87.270}{36.258}{9.0910}{ex}{88}\RG{45.122}{31.672}{7.9701}{mi}{110}\RG{114.554}{31.672}{7.9701}{mi}{106}\RG{161.882}{31.672}{7.9701}{mi}{106}\RG{165.426}{31.672}{7.9701}{rm}{43}\RG{171.433}{31.672}{7.9701}{rm}{49}\RG{25.776}{27.167}{10.9091}{rm}{61}\RG{37.115}{27.167}{10.9091}{mi}{65}\RG{50.265}{27.167}{10.9091}{mi}{66}\RG{75.822}{27.167}{10.9091}{rm}{43}\RG{106.546}{27.167}{10.9091}{mi}{65}\RG{118.553}{27.167}{10.9091}{mi}{66}\RG{142.964}{27.167}{10.9091}{sy}{0}\RG{153.873}{27.167}{10.9091}{mi}{65}\RG{175.754}{27.167}{10.9091}{mi}{66}\RG{195.721}{27.167}{10.9091}{sy}{0}\RG{206.630}{27.167}{10.9091}{mi}{65}\RG{214.638}{27.167}{10.9091}{mi}{66}\RG{58.380}{25.530}{7.9701}{mi}{110}\RG{63.069}{25.530}{7.9701}{sy}{0}\RG{69.079}{25.530}{7.9701}{rm}{49}\RG{126.668}{25.530}{7.9701}{mi}{106}\RG{130.212}{25.530}{7.9701}{sy}{0}\RG{136.222}{25.530}{7.9701}{rm}{49}\RG{183.869}{25.530}{7.9701}{mi}{106}\RG{222.753}{25.530}{7.9701}{rm}{48}\RG{87.127}{16.298}{7.9701}{mi}{106}\RG{90.671}{16.298}{7.9701}{rm}{61}\RG{96.678}{16.298}{7.9701}{rm}{49}\RG{45.122}{4.509}{7.9701}{mi}{110}\RG{128.336}{4.509}{7.9701}{mi}{110}\RG{25.776}{0.004}{10.9091}{rm}{61}\RG{37.115}{0.004}{10.9091}{mi}{65}\RG{50.265}{0.004}{10.9091}{mi}{66}\RG{75.822}{0.004}{10.9091}{rm}{43}\RG{86.554}{0.004}{10.9091}{mi}{65}\RG{94.562}{0.004}{10.9091}{mi}{66}\RG{109.419}{0.004}{10.9091}{sy}{0}\RG{120.328}{0.004}{10.9091}{mi}{65}\RG{133.480}{0.004}{10.9091}{mi}{66}\RG{159.035}{0.004}{10.9091}{sy}{0}\RG{169.944}{0.004}{10.9091}{mi}{65}\RG{177.952}{0.004}{10.9091}{mi}{66}\RG{193.416}{0.004}{10.9091}{rm}{61}\RG{204.751}{0.004}{10.9091}{rm}{48}\RG{210.096}{0.004}{10.9091}{mi}{58}\RG{58.380}{-1.632}{7.9701}{mi}{110}\RG{63.069}{-1.632}{7.9701}{sy}{0}\RG{69.079}{-1.632}{7.9701}{rm}{49}\RG{102.677}{-1.632}{7.9701}{rm}{48}\RG{141.595}{-1.632}{7.9701}{mi}{110}\RG{146.283}{-1.632}{7.9701}{sy}{0}\RG{152.293}{-1.632}{7.9701}{rm}{49}\RG{186.068}{-1.632}{7.9701}{rm}{48}%
}
\]
（5）记\allowbreak{}\hskip.1em \RB{68.578}{8.677}{3.232}{%
\RG{0.000}{-0.005}{10.9091}{mi}{75}\RG{9.872}{3.955}{7.9701}{sy}{2}\RG{19.990}{-0.005}{10.9091}{rm}{61}\RG{31.952}{-0.005}{10.9091}{mi}{75}\RG{44.496}{-0.005}{10.9091}{sy}{110}\RG{52.623}{-0.005}{10.9091}{sy}{102}\RG{58.079}{-0.005}{10.9091}{rm}{48}\RG{63.420}{-0.005}{10.9091}{sy}{103}%
}\hskip.1em\allowbreak{} 。此处可逆是指逆矩阵仍属于\allowbreak{}\hskip.1em \RB{31.428}{8.677}{3.243}{%
\RG{0.000}{-0.005}{10.9091}{mi}{77}\RG{10.334}{-1.641}{7.9701}{mi}{110}\RG{15.477}{-0.005}{10.9091}{rm}{40}\RG{19.631}{-0.005}{10.9091}{mi}{82}\RG{27.852}{-0.005}{10.9091}{rm}{41}%
}\hskip.1em\allowbreak{} 。初等行变换允许交换两行、将一行乘以\allowbreak{}\hskip.1em \RB{33.168}{8.142}{0.940}{%
\RG{0.000}{-0.004}{10.9091}{mi}{99}\RG{7.628}{-0.004}{10.9091}{sy}{50}\RG{17.931}{-0.004}{10.9091}{mi}{75}\RG{27.802}{3.955}{7.9701}{sy}{2}%
}\hskip.1em\allowbreak{} ，或将一行的\allowbreak{}\hskip.1em \RB{28.056}{8.067}{0.940}{%
\RG{0.000}{-0.004}{10.9091}{mi}{104}\RG{9.180}{-0.004}{10.9091}{sy}{50}\RG{19.483}{-0.004}{10.9091}{mi}{82}%
}\hskip.1em\allowbreak{} 倍加于另一行。初等\allowbreak{}\hskip.1em \RB{6.336}{8.067}{0.624}{%
\RG{0.000}{-0.004}{10.9091}{mi}{21}%
}\hskip.1em\allowbreak{} –矩阵是由\allowbreak{}\hskip.1em \RB{9.662}{7.947}{2.213}{%
\RG{0.000}{-0.004}{10.9091}{mi}{73}\RG{4.711}{-1.640}{7.9701}{mi}{110}%
}\hskip.1em\allowbreak{} 经一次这类行变换得到的矩阵。它们的逆矩阵分别对应原交换、乘以\allowbreak{}\hskip.1em \RB{14.358}{9.292}{0.612}{%
\RG{0.000}{-0.003}{10.9091}{mi}{99}\RG{4.597}{3.956}{7.9701}{sy}{0}\RG{10.607}{3.956}{7.9701}{rm}{49}%
}\hskip.1em\allowbreak{} 、加上\allowbreak{}\hskip.1em \RB{14.821}{8.068}{0.612}{%
\RG{0.000}{-0.003}{10.9091}{sy}{0}\RG{8.485}{-0.003}{10.9091}{mi}{104}%
}\hskip.1em\allowbreak{} 倍，仍为初等矩阵。由（1），它们的行列式均属于\allowbreak{}\hskip.1em \RB{15.236}{8.144}{0.502}{%
\RG{0.000}{-0.002}{10.9091}{mi}{75}\RG{9.871}{3.957}{7.9701}{sy}{2}%
}\hskip.1em\allowbreak{} 。

若（5.1）成立，则由行列式乘法公式，\allowbreak{}\hskip.1em \RB{80.809}{9.296}{0.608}{%
\RG{0.000}{0.001}{10.9091}{rm}{100}\RG{5.935}{0.001}{10.9091}{rm}{101}\RG{10.680}{0.001}{10.9091}{rm}{116}\RG{16.654}{0.001}{10.9091}{mi}{65}\RG{26.480}{0.001}{10.9091}{rm}{100}\RG{32.415}{0.001}{10.9091}{rm}{101}\RG{37.160}{0.001}{10.9091}{rm}{116}\RG{43.134}{0.001}{10.9091}{mi}{65}\RG{51.142}{3.959}{7.9701}{sy}{0}\RG{57.153}{3.959}{7.9701}{rm}{49}\RG{64.501}{0.001}{10.9091}{rm}{61}\RG{75.836}{0.001}{10.9091}{rm}{49}%
}\hskip.1em\allowbreak{} 。两个非零多项式的次数之和等于乘积的次数，故\allowbreak{}\hskip.1em \RB{63.767}{8.291}{2.735}{%
\RG{0.000}{0.002}{10.9091}{rm}{100}\RG{5.935}{0.002}{10.9091}{rm}{101}\RG{10.680}{0.002}{10.9091}{rm}{103}\RG{17.978}{0.002}{10.9091}{rm}{100}\RG{23.913}{0.002}{10.9091}{rm}{101}\RG{28.659}{0.002}{10.9091}{rm}{116}\RG{34.639}{0.002}{10.9091}{mi}{65}\RG{46.357}{0.002}{10.9091}{rm}{61}\RG{58.379}{0.002}{10.9091}{rm}{48}%
}\hskip.1em\allowbreak{} ，得到（5.2）。若（5.2）成立，则\allowbreak{}\hskip.1em \RB{68.028}{9.297}{3.236}{%
\RG{0.000}{0.002}{10.9091}{rm}{40}\RG{4.156}{0.002}{10.9091}{rm}{100}\RG{10.091}{0.002}{10.9091}{rm}{101}\RG{14.837}{0.002}{10.9091}{rm}{116}\RG{20.808}{0.002}{10.9091}{mi}{65}\RG{28.815}{0.002}{10.9091}{rm}{41}\RG{32.969}{3.960}{7.9701}{sy}{0}\RG{38.980}{3.960}{7.9701}{rm}{49}\RG{47.009}{0.002}{10.9091}{sy}{50}\RG{57.993}{0.002}{10.9091}{mi}{75}%
}\hskip.1em\allowbreak{} ，由（3.1），矩阵\allowbreak{}\hskip.1em \RB{101.375}{9.298}{3.236}{%
\RG{0.000}{0.002}{10.9091}{rm}{40}\RG{4.156}{0.002}{10.9091}{rm}{100}\RG{10.091}{0.002}{10.9091}{rm}{101}\RG{14.837}{0.002}{10.9091}{rm}{116}\RG{20.808}{0.002}{10.9091}{mi}{65}\RG{28.815}{0.002}{10.9091}{rm}{41}\RG{32.969}{3.961}{7.9701}{sy}{0}\RG{38.980}{3.961}{7.9701}{rm}{49}\RG{43.298}{0.002}{10.9091}{mi}{65}\RG{51.306}{3.961}{7.9701}{sy}{3}\RG{59.149}{0.002}{10.9091}{sy}{50}\RG{69.947}{0.002}{10.9091}{mi}{77}\RG{80.281}{-1.634}{7.9701}{mi}{110}\RG{85.424}{0.002}{10.9091}{rm}{40}\RG{89.578}{0.002}{10.9091}{mi}{82}\RG{97.799}{0.002}{10.9091}{rm}{41}%
}\hskip.1em\allowbreak{} 是\allowbreak{}\hskip.1em \RB{8.202}{8.291}{0.500}{%
\RG{0.000}{0.002}{10.9091}{mi}{65}%
}\hskip.1em\allowbreak{} 的双侧逆矩阵，得到（5.1）及所要求的逆矩阵公式。当\allowbreak{}\hskip.1em \RB{25.714}{7.757}{0.606}{%
\RG{0.000}{0.003}{10.9091}{mi}{110}\RG{9.406}{0.003}{10.9091}{rm}{61}\RG{20.741}{0.003}{10.9091}{rm}{49}%
}\hskip.1em\allowbreak{} 时，（3.1）同样成立，因为相应的空子式行列式为\allowbreak{}\hskip.1em \RB{4.973}{7.757}{0.500}{%
\RG{0.000}{0.003}{10.9091}{rm}{49}%
}\hskip.1em\allowbreak{} 。

由多项式矩阵的Smith标准形定理，存在有限个初等矩阵的乘积\allowbreak{}\hskip.1em \RB{20.377}{8.176}{2.620}{%
\RG{0.000}{-0.004}{10.9091}{mi}{80}\RG{7.178}{-0.004}{10.9091}{mi}{59}\RG{11.957}{-0.004}{10.9091}{mi}{81}%
}\hskip.1em\allowbreak{} ，使

\[
\RB{176.300}{8.686}{3.234}{%
\RG{0.000}{0.005}{10.9091}{mi}{80}\RG{8.367}{0.005}{10.9091}{mi}{65}\RG{16.375}{0.005}{10.9091}{mi}{81}\RG{27.867}{0.005}{10.9091}{rm}{61}\RG{39.206}{0.005}{10.9091}{mi}{68}\RG{51.402}{0.005}{10.9091}{rm}{61}\RG{62.737}{0.005}{10.9091}{rm}{100}\RG{68.671}{0.005}{10.9091}{rm}{105}\RG{71.639}{0.005}{10.9091}{rm}{97}\RG{76.984}{0.005}{10.9091}{rm}{103}\RG{82.471}{0.005}{10.9091}{rm}{40}\RG{86.622}{0.005}{10.9091}{mi}{100}\RG{96.491}{0.005}{10.9091}{mi}{59}\RG{101.269}{0.005}{10.9091}{mi}{58}\RG{106.058}{0.005}{10.9091}{mi}{58}\RG{110.837}{0.005}{10.9091}{mi}{58}\RG{115.626}{0.005}{10.9091}{mi}{59}\RG{120.415}{0.005}{10.9091}{mi}{100}\RG{130.124}{0.005}{10.9091}{mi}{59}\RG{134.910}{0.005}{10.9091}{rm}{48}\RG{140.250}{0.005}{10.9091}{mi}{59}\RG{145.029}{0.005}{10.9091}{mi}{58}\RG{149.818}{0.005}{10.9091}{mi}{58}\RG{154.607}{0.005}{10.9091}{mi}{58}\RG{159.385}{0.005}{10.9091}{mi}{59}\RG{164.177}{0.005}{10.9091}{rm}{48}\RG{169.523}{0.005}{10.9091}{rm}{41}\RG{173.673}{0.005}{10.9091}{mi}{59}\RG{92.173}{-1.633}{7.9701}{rm}{49}\RG{125.968}{-1.633}{7.9701}{mi}{114}%
}
\]
其中\allowbreak{}\hskip.1em \RB{46.638}{7.757}{1.992}{%
\RG{0.000}{0.003}{10.9091}{rm}{48}\RG{8.506}{0.003}{10.9091}{sy}{20}\RG{20.156}{0.003}{10.9091}{mi}{114}\RG{28.433}{0.003}{10.9091}{sy}{20}\RG{40.083}{0.003}{10.9091}{mi}{110}%
}\hskip.1em\allowbreak{} ，各\allowbreak{}\hskip.1em \RB{8.356}{8.074}{2.206}{%
\RG{0.000}{0.003}{10.9091}{mi}{100}\RG{5.551}{-1.633}{7.9701}{mi}{105}%
}\hskip.1em\allowbreak{} 为非零首一多项式，且\allowbreak{}\hskip.1em \RB{35.937}{8.685}{3.225}{%
\RG{0.000}{0.003}{10.9091}{mi}{100}\RG{5.551}{-1.633}{7.9701}{mi}{105}\RG{11.802}{0.003}{10.9091}{sy}{106}\RG{17.997}{0.003}{10.9091}{mi}{100}\RG{23.548}{-1.633}{7.9701}{mi}{105}\RG{26.179}{-1.633}{7.9701}{rm}{43}\RG{32.186}{-1.633}{7.9701}{rm}{49}%
}\hskip.1em\allowbreak{} （\allowbreak{}\hskip.1em \RB{43.589}{7.757}{1.992}{%
\RG{0.000}{0.003}{10.9091}{rm}{49}\RG{8.506}{0.003}{10.9091}{sy}{20}\RG{20.157}{0.003}{10.9091}{mi}{105}\RG{26.965}{0.003}{10.9091}{mi}{60}\RG{38.441}{0.003}{10.9091}{mi}{114}%
}\hskip.1em\allowbreak{} ）。这一定理适用，因为\allowbreak{}\hskip.1em \RB{43.913}{8.686}{3.224}{%
\RG{0.000}{0.004}{10.9091}{mi}{82}\RG{11.250}{0.004}{10.9091}{rm}{61}\RG{22.589}{0.004}{10.9091}{mi}{75}\RG{32.461}{0.004}{10.9091}{rm}{91}\RG{35.491}{0.004}{10.9091}{mi}{21}\RG{41.722}{0.004}{10.9091}{rm}{93}%
}\hskip.1em\allowbreak{} 是欧几里得整环。由初等矩阵的行列式性质，\allowbreak{}\hskip.1em \RB{82.306}{8.184}{2.613}{%
\RG{0.000}{0.004}{10.9091}{rm}{100}\RG{5.935}{0.004}{10.9091}{rm}{101}\RG{10.680}{0.004}{10.9091}{rm}{116}\RG{16.654}{0.004}{10.9091}{mi}{80}\RG{23.833}{0.004}{10.9091}{mi}{59}\RG{28.622}{0.004}{10.9091}{rm}{100}\RG{34.557}{0.004}{10.9091}{rm}{101}\RG{39.302}{0.004}{10.9091}{rm}{116}\RG{45.276}{0.004}{10.9091}{mi}{81}\RG{56.767}{0.004}{10.9091}{sy}{50}\RG{67.070}{0.004}{10.9091}{mi}{75}\RG{76.940}{3.962}{7.9701}{sy}{2}%
}\hskip.1em\allowbreak{} 。若（5.2）成立，则\allowbreak{}\hskip.1em \RB{54.390}{8.150}{0.932}{%
\RG{0.000}{0.004}{10.9091}{rm}{100}\RG{5.935}{0.004}{10.9091}{rm}{101}\RG{10.680}{0.004}{10.9091}{rm}{116}\RG{16.654}{0.004}{10.9091}{mi}{68}\RG{28.850}{0.004}{10.9091}{sy}{50}\RG{39.153}{0.004}{10.9091}{mi}{75}\RG{49.025}{3.963}{7.9701}{sy}{2}%
}\hskip.1em\allowbreak{} ，所以\allowbreak{}\hskip.1em \RB{26.034}{5.315}{0.605}{%
\RG{0.000}{0.004}{10.9091}{mi}{114}\RG{8.141}{0.004}{10.9091}{rm}{61}\RG{19.479}{0.004}{10.9091}{mi}{110}%
}\hskip.1em\allowbreak{} ，并且\allowbreak{}\hskip.1em \RB{70.746}{8.105}{3.266}{%
\RG{0.000}{7.277}{9.0910}{ex}{80}\RG{9.596}{4.398}{7.9701}{mi}{110}\RG{9.596}{-2.693}{7.9701}{mi}{105}\RG{12.227}{-2.693}{7.9701}{rm}{61}\RG{18.234}{-2.693}{7.9701}{rm}{49}\RG{24.373}{0.004}{10.9091}{rm}{100}\RG{30.308}{0.004}{10.9091}{rm}{101}\RG{35.053}{0.004}{10.9091}{rm}{103}\RG{42.358}{0.004}{10.9091}{mi}{100}\RG{47.909}{-1.633}{7.9701}{mi}{105}\RG{54.024}{0.004}{10.9091}{rm}{61}\RG{65.359}{0.004}{10.9091}{rm}{48}%
}\hskip.1em\allowbreak{} 。各\allowbreak{}\hskip.1em \RB{8.355}{8.075}{2.205}{%
\RG{0.000}{0.004}{10.9091}{mi}{100}\RG{5.550}{-1.633}{7.9701}{mi}{105}%
}\hskip.1em\allowbreak{} 首一，因而均等于\allowbreak{}\hskip.1em \RB{4.973}{7.759}{0.500}{%
\RG{0.000}{0.004}{10.9091}{rm}{49}%
}\hskip.1em\allowbreak{} ，得到（5.3）。反之，若\allowbreak{}\hskip.1em \RB{33.197}{7.955}{2.204}{%
\RG{0.000}{0.004}{10.9091}{mi}{68}\RG{12.196}{0.004}{10.9091}{rm}{61}\RG{23.535}{0.004}{10.9091}{mi}{73}\RG{28.246}{-1.632}{7.9701}{mi}{110}%
}\hskip.1em\allowbreak{} ，则\allowbreak{}\hskip.1em \RB{97.773}{8.294}{2.612}{%
\RG{0.000}{0.004}{10.9091}{rm}{100}\RG{5.935}{0.004}{10.9091}{rm}{101}\RG{10.680}{0.004}{10.9091}{rm}{116}\RG{16.654}{0.004}{10.9091}{mi}{80}\RG{26.842}{0.004}{10.9091}{rm}{100}\RG{32.776}{0.004}{10.9091}{rm}{101}\RG{37.522}{0.004}{10.9091}{rm}{116}\RG{43.496}{0.004}{10.9091}{mi}{65}\RG{53.322}{0.004}{10.9091}{rm}{100}\RG{59.256}{0.004}{10.9091}{rm}{101}\RG{64.002}{0.004}{10.9091}{rm}{116}\RG{69.975}{0.004}{10.9091}{mi}{81}\RG{81.465}{0.004}{10.9091}{rm}{61}\RG{92.800}{0.004}{10.9091}{rm}{49}%
}\hskip.1em\allowbreak{} ，故\allowbreak{}\hskip.1em \RB{53.231}{8.294}{0.932}{%
\RG{0.000}{0.004}{10.9091}{rm}{100}\RG{5.935}{0.004}{10.9091}{rm}{101}\RG{10.680}{0.004}{10.9091}{rm}{116}\RG{16.653}{0.004}{10.9091}{mi}{65}\RG{27.692}{0.004}{10.9091}{sy}{50}\RG{37.995}{0.004}{10.9091}{mi}{75}\RG{47.866}{3.964}{7.9701}{sy}{2}%
}\hskip.1em\allowbreak{} 。因此（5.2）与（5.3）等价。

现由（5.2）证明只用初等行变换即可把\allowbreak{}\hskip.1em \RB{8.202}{8.288}{0.501}{%
\RG{0.000}{-0.001}{10.9091}{mi}{65}%
}\hskip.1em\allowbreak{} 化为\allowbreak{}\hskip.1em \RB{9.662}{7.950}{2.211}{%
\RG{0.000}{-0.001}{10.9091}{mi}{73}\RG{4.711}{-1.638}{7.9701}{mi}{110}%
}\hskip.1em\allowbreak{} 。对阶数\allowbreak{}\hskip.1em \RB{6.555}{5.310}{0.610}{%
\RG{0.000}{-0.001}{10.9091}{mi}{110}%
}\hskip.1em\allowbreak{} 作归纳。当\allowbreak{}\hskip.1em \RB{26.976}{7.754}{0.610}{%
\RG{0.000}{-0.001}{10.9091}{mi}{110}\RG{10.036}{-0.001}{10.9091}{rm}{61}\RG{22.003}{-0.001}{10.9091}{rm}{49}%
}\hskip.1em\allowbreak{} 时，将唯一的一行乘以非零常数\allowbreak{}\hskip.1em \RB{15.368}{9.851}{2.885}{%
\RG{0.000}{-0.001}{10.9091}{mi}{97}\RG{5.607}{4.515}{7.9701}{sy}{0}\RG{5.607}{-2.385}{7.9701}{rm}{49}\RG{9.469}{-2.385}{7.9701}{rm}{49}\RG{11.617}{4.515}{7.9701}{rm}{49}%
}\hskip.1em\allowbreak{} 即可。设\allowbreak{}\hskip.1em \RB{27.944}{7.754}{0.807}{%
\RG{0.000}{-0.001}{10.9091}{mi}{110}\RG{10.517}{-0.001}{10.9091}{mi}{62}\RG{22.972}{-0.001}{10.9091}{rm}{49}%
}\hskip.1em\allowbreak{} ，并假设结论对\allowbreak{}\hskip.1em \RB{25.574}{7.754}{0.610}{%
\RG{0.000}{-0.001}{10.9091}{mi}{110}\RG{9.246}{-0.001}{10.9091}{sy}{0}\RG{20.601}{-0.001}{10.9091}{rm}{49}%
}\hskip.1em\allowbreak{} 阶矩阵成立。因\allowbreak{}\hskip.1em \RB{44.847}{8.311}{2.857}{%
\RG{0.000}{-0.000}{10.9091}{rm}{100}\RG{5.935}{-0.000}{10.9091}{rm}{101}\RG{10.680}{-0.000}{10.9091}{rm}{116}\RG{16.654}{-0.000}{10.9091}{mi}{65}\RG{27.906}{-0.000}{10.9091}{sy}{54}\RG{27.906}{-0.000}{10.9091}{rm}{61}\RG{39.459}{-0.000}{10.9091}{rm}{48}%
}\hskip.1em\allowbreak{} ，第一列至少有一个非零元素。对该列中的非零元素逐对实行多项式欧几里得算法，可仅用初等行变换使其变为\allowbreak{}\hskip.1em \RB{58.749}{9.456}{3.237}{%
\RG{0.000}{0.001}{10.9091}{rm}{40}\RG{4.154}{0.001}{10.9091}{mi}{100}\RG{9.707}{0.001}{10.9091}{mi}{59}\RG{14.490}{0.001}{10.9091}{rm}{48}\RG{19.831}{0.001}{10.9091}{mi}{59}\RG{24.609}{0.001}{10.9091}{mi}{58}\RG{29.398}{0.001}{10.9091}{mi}{58}\RG{34.176}{0.001}{10.9091}{mi}{58}\RG{38.966}{0.001}{10.9091}{mi}{59}\RG{43.758}{0.001}{10.9091}{rm}{48}\RG{49.103}{0.001}{10.9091}{rm}{41}\RG{53.252}{3.959}{7.9701}{rm}{84}%
}\hskip.1em\allowbreak{} ，其中\allowbreak{}\hskip.1em \RB{27.456}{8.072}{0.936}{%
\RG{0.000}{0.001}{10.9091}{mi}{100}\RG{8.580}{0.001}{10.9091}{sy}{50}\RG{18.884}{0.001}{10.9091}{mi}{82}%
}\hskip.1em\allowbreak{} 非零。具体地，对同列中两个非零元素\allowbreak{}\hskip.1em \RB{16.558}{5.313}{2.724}{%
\RG{0.000}{0.002}{10.9091}{mi}{120}\RG{6.066}{0.002}{10.9091}{mi}{59}\RG{10.844}{0.002}{10.9091}{mi}{121}%
}\hskip.1em\allowbreak{} ，由多项式除法定理，存在\allowbreak{}\hskip.1em \RB{37.437}{7.953}{2.615}{%
\RG{0.000}{0.002}{10.9091}{mi}{113}\RG{5.116}{0.002}{10.9091}{mi}{59}\RG{9.906}{0.002}{10.9091}{mi}{115}\RG{18.268}{0.002}{10.9091}{sy}{50}\RG{28.864}{0.002}{10.9091}{mi}{82}%
}\hskip.1em\allowbreak{} 使\allowbreak{}\hskip.1em \RB{50.148}{6.785}{2.724}{%
\RG{0.000}{0.002}{10.9091}{mi}{120}\RG{9.393}{0.002}{10.9091}{rm}{61}\RG{21.024}{0.002}{10.9091}{mi}{113}\RG{26.141}{0.002}{10.9091}{mi}{121}\RG{34.293}{0.002}{10.9091}{rm}{43}\RG{45.143}{0.002}{10.9091}{mi}{115}%
}\hskip.1em\allowbreak{} ，且\allowbreak{}\hskip.1em \RB{25.371}{7.756}{0.727}{%
\RG{0.000}{0.002}{10.9091}{mi}{115}\RG{8.354}{0.002}{10.9091}{rm}{61}\RG{19.983}{0.002}{10.9091}{rm}{48}%
}\hskip.1em\allowbreak{} 或\allowbreak{}\hskip.1em \RB{61.671}{8.073}{2.735}{%
\RG{0.000}{0.002}{10.9091}{rm}{100}\RG{5.935}{0.002}{10.9091}{rm}{101}\RG{10.680}{0.002}{10.9091}{rm}{103}\RG{17.985}{0.002}{10.9091}{mi}{115}\RG{26.330}{0.002}{10.9091}{mi}{60}\RG{37.971}{0.002}{10.9091}{rm}{100}\RG{43.905}{0.002}{10.9091}{rm}{101}\RG{48.651}{0.002}{10.9091}{rm}{103}\RG{55.956}{0.002}{10.9091}{mi}{121}%
}\hskip.1em\allowbreak{} ；以含\allowbreak{}\hskip.1em \RB{6.107}{5.313}{0.607}{%
\RG{0.000}{0.002}{10.9091}{mi}{120}%
}\hskip.1em\allowbreak{} 的行减去含\allowbreak{}\hskip.1em \RB{5.715}{5.313}{2.724}{%
\RG{0.000}{0.002}{10.9091}{mi}{121}%
}\hskip.1em\allowbreak{} 的行的\allowbreak{}\hskip.1em \RB{5.311}{5.313}{2.615}{%
\RG{0.000}{0.002}{10.9091}{mi}{113}%
}\hskip.1em\allowbreak{} 倍，再交换这两行，便完成一次辗转相除。非零余式的次数

% Original page 9.
严格下降，故每对元素的处理在有限次后终止；依次并入该列其余元素，便得到上述列形状。

设所得矩阵为

\[
\RB{230.659}{16.865}{11.398}{%
\RG{25.937}{16.002}{9.0910}{ex}{32}\RG{56.202}{16.002}{9.0910}{ex}{33}\RG{33.134}{7.158}{10.9091}{mi}{100}\RG{49.716}{7.158}{10.9091}{mi}{98}\RG{117.917}{4.506}{7.9701}{rm}{49}\RG{121.780}{4.506}{7.9701}{sy}{2}\RG{127.790}{4.506}{7.9701}{rm}{40}\RG{130.796}{4.506}{7.9701}{mi}{110}\RG{135.484}{4.506}{7.9701}{sy}{0}\RG{141.495}{4.506}{7.9701}{rm}{49}\RG{145.357}{4.506}{7.9701}{rm}{41}\RG{0.000}{0.001}{10.9091}{mi}{87}\RG{14.599}{0.001}{10.9091}{rm}{61}\RG{65.217}{0.001}{10.9091}{mi}{59}\RG{91.814}{0.001}{10.9091}{mi}{98}\RG{99.393}{0.001}{10.9091}{sy}{50}\RG{109.696}{0.001}{10.9091}{mi}{82}\RG{148.818}{0.001}{10.9091}{mi}{59}\RG{164.506}{0.001}{10.9091}{mi}{67}\RG{175.969}{0.001}{10.9091}{sy}{50}\RG{186.272}{0.001}{10.9091}{mi}{77}\RG{211.624}{0.001}{10.9091}{rm}{40}\RG{215.777}{0.001}{10.9091}{mi}{82}\RG{223.997}{0.001}{10.9091}{rm}{41}\RG{228.152}{0.001}{10.9091}{mi}{58}\RG{196.606}{-1.635}{7.9701}{mi}{110}\RG{201.295}{-1.635}{7.9701}{sy}{0}\RG{207.305}{-1.635}{7.9701}{rm}{49}\RG{33.239}{-7.606}{10.9091}{rm}{48}\RG{47.776}{-7.606}{10.9091}{mi}{67}%
}
\]
初等行变换只使行列式乘以非零常数，故\allowbreak{}\hskip.1em \RB{57.947}{8.143}{0.939}{%
\RG{0.000}{-0.003}{10.9091}{rm}{100}\RG{5.935}{-0.003}{10.9091}{rm}{101}\RG{10.680}{-0.003}{10.9091}{rm}{116}\RG{16.654}{-0.003}{10.9091}{mi}{87}\RG{31.830}{-0.003}{10.9091}{sy}{50}\RG{42.711}{-0.003}{10.9091}{mi}{75}\RG{52.582}{3.956}{7.9701}{sy}{2}%
}\hskip.1em\allowbreak{} 。由块三角行列式公式，\allowbreak{}\hskip.1em \RB{44.270}{8.177}{0.732}{%
\RG{0.000}{-0.003}{10.9091}{mi}{100}\RG{7.369}{-0.003}{10.9091}{rm}{100}\RG{13.304}{-0.003}{10.9091}{rm}{101}\RG{18.050}{-0.003}{10.9091}{rm}{116}\RG{24.023}{-0.003}{10.9091}{mi}{67}\RG{36.057}{-0.003}{10.9091}{rm}{61}%
}\hskip.1em\allowbreak{}\allowbreak{}\hskip.1em \RB{56.793}{8.144}{0.938}{%
\RG{0.000}{-0.002}{10.9091}{rm}{100}\RG{5.935}{-0.002}{10.9091}{rm}{101}\RG{10.680}{-0.002}{10.9091}{rm}{116}\RG{16.654}{-0.002}{10.9091}{mi}{87}\RG{31.253}{-0.002}{10.9091}{sy}{50}\RG{41.556}{-0.002}{10.9091}{mi}{75}\RG{51.427}{3.957}{7.9701}{sy}{2}%
}\hskip.1em\allowbreak{} ，从而\allowbreak{}\hskip.1em \RB{34.121}{8.144}{0.938}{%
\RG{0.000}{-0.002}{10.9091}{mi}{100}\RG{8.580}{-0.002}{10.9091}{sy}{50}\RG{18.884}{-0.002}{10.9091}{mi}{75}\RG{28.755}{3.957}{7.9701}{sy}{2}%
}\hskip.1em\allowbreak{} 且\allowbreak{}\hskip.1em \RB{53.650}{8.178}{0.938}{%
\RG{0.000}{-0.002}{10.9091}{rm}{100}\RG{5.935}{-0.002}{10.9091}{rm}{101}\RG{10.680}{-0.002}{10.9091}{rm}{116}\RG{16.654}{-0.002}{10.9091}{mi}{67}\RG{28.110}{-0.002}{10.9091}{sy}{50}\RG{38.413}{-0.002}{10.9091}{mi}{75}\RG{48.285}{3.957}{7.9701}{sy}{2}%
}\hskip.1em\allowbreak{} 。将第一行乘以\allowbreak{}\hskip.1em \RB{15.312}{9.293}{0.611}{%
\RG{0.000}{-0.002}{10.9091}{mi}{100}\RG{5.551}{3.957}{7.9701}{sy}{0}\RG{11.561}{3.957}{7.9701}{rm}{49}%
}\hskip.1em\allowbreak{} ，再由归纳假设只在后\allowbreak{}\hskip.1em \RB{23.497}{7.753}{0.611}{%
\RG{0.000}{-0.002}{10.9091}{mi}{110}\RG{8.208}{-0.002}{10.9091}{sy}{0}\RG{18.524}{-0.002}{10.9091}{rm}{49}%
}\hskip.1em\allowbreak{}行上作初等行变换，得到

\[
\RB{214.053}{16.865}{11.398}{%
\RG{0.000}{16.002}{9.0910}{ex}{32}\RG{41.357}{16.002}{9.0910}{ex}{33}\RG{32.488}{11.117}{7.9701}{sy}{48}\RG{7.197}{7.158}{10.9091}{rm}{49}\RG{27.947}{7.158}{10.9091}{mi}{98}\RG{81.518}{4.506}{7.9701}{sy}{48}\RG{107.133}{4.506}{7.9701}{sy}{48}\RG{139.920}{4.506}{7.9701}{sy}{48}\RG{180.645}{4.506}{7.9701}{rm}{49}\RG{184.508}{4.506}{7.9701}{sy}{2}\RG{190.519}{4.506}{7.9701}{rm}{40}\RG{193.524}{4.506}{7.9701}{mi}{110}\RG{198.212}{4.506}{7.9701}{sy}{0}\RG{204.223}{4.506}{7.9701}{rm}{49}\RG{208.085}{4.506}{7.9701}{rm}{41}\RG{50.373}{0.002}{10.9091}{mi}{59}\RG{76.969}{0.002}{10.9091}{mi}{98}\RG{87.098}{0.002}{10.9091}{rm}{61}\RG{98.433}{0.002}{10.9091}{rm}{40}\RG{102.591}{0.002}{10.9091}{mi}{98}\RG{111.452}{0.002}{10.9091}{mi}{59}\RG{116.230}{0.002}{10.9091}{mi}{58}\RG{121.019}{0.002}{10.9091}{mi}{58}\RG{125.797}{0.002}{10.9091}{mi}{58}\RG{130.586}{0.002}{10.9091}{mi}{59}\RG{135.376}{0.002}{10.9091}{mi}{98}\RG{154.937}{0.002}{10.9091}{rm}{41}\RG{162.122}{0.002}{10.9091}{sy}{50}\RG{172.425}{0.002}{10.9091}{mi}{82}\RG{211.546}{0.002}{10.9091}{mi}{58}\RG{107.133}{-2.696}{7.9701}{rm}{49}\RG{139.920}{-2.696}{7.9701}{mi}{110}\RG{144.609}{-2.696}{7.9701}{sy}{0}\RG{150.618}{-2.696}{7.9701}{rm}{49}\RG{7.197}{-7.606}{10.9091}{rm}{48}\RG{21.628}{-7.606}{10.9091}{mi}{73}\RG{26.340}{-9.242}{7.9701}{mi}{110}\RG{31.029}{-9.242}{7.9701}{sy}{0}\RG{37.039}{-9.242}{7.9701}{rm}{49}%
}
\]
对每个\allowbreak{}\hskip.1em \RB{73.700}{7.751}{2.729}{%
\RG{0.000}{-0.004}{10.9091}{mi}{106}\RG{8.188}{-0.004}{10.9091}{rm}{61}\RG{19.654}{-0.004}{10.9091}{rm}{49}\RG{24.991}{-0.004}{10.9091}{mi}{59}\RG{29.769}{-0.004}{10.9091}{mi}{58}\RG{34.558}{-0.004}{10.9091}{mi}{58}\RG{39.348}{-0.004}{10.9091}{mi}{58}\RG{44.126}{-0.004}{10.9091}{mi}{59}\RG{48.915}{-0.004}{10.9091}{mi}{110}\RG{57.768}{-0.004}{10.9091}{sy}{0}\RG{68.727}{-0.004}{10.9091}{rm}{49}%
}\hskip.1em\allowbreak{} ，将第\allowbreak{}\hskip.1em \RB{23.258}{7.751}{2.729}{%
\RG{0.000}{-0.004}{10.9091}{mi}{106}\RG{7.507}{-0.004}{10.9091}{rm}{43}\RG{18.286}{-0.004}{10.9091}{rm}{49}%
}\hskip.1em\allowbreak{} 行的\allowbreak{}\hskip.1em \RB{8.045}{8.521}{4.684}{%
\RG{0.000}{-0.004}{10.9091}{mi}{98}\RG{4.542}{3.956}{7.9701}{sy}{48}\RG{4.542}{-2.700}{7.9701}{mi}{106}%
}\hskip.1em\allowbreak{} 倍从第一行减去，即得\allowbreak{}\hskip.1em \RB{9.662}{7.947}{2.213}{%
\RG{0.000}{-0.004}{10.9091}{mi}{73}\RG{4.711}{-1.640}{7.9701}{mi}{110}%
}\hskip.1em\allowbreak{} 。归纳完成。对\allowbreak{}\hskip.1em \RB{13.505}{9.452}{0.504}{%
\RG{0.000}{-0.004}{10.9091}{mi}{65}\RG{8.008}{3.956}{7.9701}{rm}{84}%
}\hskip.1em\allowbreak{}应用此结论，并将所得等式转置，便得到仅用初等列变换把\allowbreak{}\hskip.1em \RB{8.202}{8.286}{0.503}{%
\RG{0.000}{-0.003}{10.9091}{mi}{65}%
}\hskip.1em\allowbreak{} 化为\allowbreak{}\hskip.1em \RB{9.662}{7.948}{2.212}{%
\RG{0.000}{-0.003}{10.9091}{mi}{73}\RG{4.711}{-1.639}{7.9701}{mi}{110}%
}\hskip.1em\allowbreak{} 的过程。因此（5.2）推出（5.4），且两种单侧变换均可实现。

若初等行变换对应的矩阵依次为\allowbreak{}\hskip.1em \RB{47.753}{7.930}{2.605}{%
\RG{0.000}{0.001}{10.9091}{mi}{69}\RG{7.908}{-1.635}{7.9701}{rm}{49}\RG{12.226}{0.001}{10.9091}{mi}{59}\RG{17.004}{0.001}{10.9091}{mi}{58}\RG{21.793}{0.001}{10.9091}{mi}{58}\RG{26.572}{0.001}{10.9091}{mi}{58}\RG{31.361}{0.001}{10.9091}{mi}{59}\RG{36.150}{0.001}{10.9091}{mi}{69}\RG{44.060}{-1.635}{7.9701}{mi}{115}%
}\hskip.1em\allowbreak{} ，且\allowbreak{}\hskip.1em \RB{72.564}{8.290}{2.208}{%
\RG{0.000}{0.001}{10.9091}{mi}{69}\RG{7.908}{-1.635}{7.9701}{mi}{115}\RG{13.753}{0.001}{10.9091}{sy}{1}\RG{18.597}{0.001}{10.9091}{sy}{1}\RG{23.452}{0.001}{10.9091}{sy}{1}\RG{28.299}{0.001}{10.9091}{mi}{69}\RG{36.207}{-1.635}{7.9701}{rm}{49}\RG{40.525}{0.001}{10.9091}{mi}{65}\RG{51.562}{0.001}{10.9091}{rm}{61}\RG{62.901}{0.001}{10.9091}{mi}{73}\RG{67.613}{-1.635}{7.9701}{mi}{110}%
}\hskip.1em\allowbreak{} ，则\allowbreak{}\hskip.1em \RB{75.923}{9.852}{3.268}{%
\RG{0.000}{0.001}{10.9091}{mi}{65}\RG{11.038}{0.001}{10.9091}{rm}{61}\RG{22.377}{0.001}{10.9091}{mi}{69}\RG{30.284}{-2.383}{7.9701}{rm}{49}\RG{30.923}{4.516}{7.9701}{sy}{0}\RG{36.934}{4.516}{7.9701}{rm}{49}\RG{43.070}{0.001}{10.9091}{sy}{1}\RG{47.913}{0.001}{10.9091}{sy}{1}\RG{52.768}{0.001}{10.9091}{sy}{1}\RG{57.616}{0.001}{10.9091}{mi}{69}\RG{65.523}{-2.696}{7.9701}{mi}{115}\RG{66.163}{3.960}{7.9701}{sy}{0}\RG{72.172}{3.960}{7.9701}{rm}{49}%
}\hskip.1em\allowbreak{} 。若初等列变换对应的矩阵依次为\allowbreak{}\hskip.1em \RB{45.767}{7.931}{2.604}{%
\RG{0.000}{0.002}{10.9091}{mi}{70}\RG{6.914}{-1.634}{7.9701}{rm}{49}\RG{11.233}{0.002}{10.9091}{mi}{59}\RG{16.011}{0.002}{10.9091}{mi}{58}\RG{20.800}{0.002}{10.9091}{mi}{58}\RG{25.579}{0.002}{10.9091}{mi}{58}\RG{30.368}{0.002}{10.9091}{mi}{59}\RG{35.157}{0.002}{10.9091}{mi}{70}\RG{42.074}{-1.634}{7.9701}{mi}{115}%
}\hskip.1em\allowbreak{} ，且\allowbreak{}\hskip.1em \RB{70.578}{8.291}{2.207}{%
\RG{0.000}{0.002}{10.9091}{mi}{65}\RG{8.007}{0.002}{10.9091}{mi}{70}\RG{14.922}{-1.634}{7.9701}{rm}{49}\RG{21.059}{0.002}{10.9091}{sy}{1}\RG{25.903}{0.002}{10.9091}{sy}{1}\RG{30.757}{0.002}{10.9091}{sy}{1}\RG{35.604}{0.002}{10.9091}{mi}{70}\RG{42.519}{-1.634}{7.9701}{mi}{115}\RG{49.578}{0.002}{10.9091}{rm}{61}\RG{60.916}{0.002}{10.9091}{mi}{73}\RG{65.627}{-1.634}{7.9701}{mi}{110}%
}\hskip.1em\allowbreak{} ，则\allowbreak{}\hskip.1em \RB{75.627}{9.853}{3.267}{%
\RG{0.000}{0.002}{10.9091}{mi}{65}\RG{11.039}{0.002}{10.9091}{rm}{61}\RG{22.377}{0.002}{10.9091}{mi}{70}\RG{29.292}{-2.695}{7.9701}{mi}{115}\RG{30.775}{3.960}{7.9701}{sy}{0}\RG{36.785}{3.960}{7.9701}{rm}{49}\RG{42.922}{0.002}{10.9091}{sy}{1}\RG{47.766}{0.002}{10.9091}{sy}{1}\RG{52.620}{0.002}{10.9091}{sy}{1}\RG{57.467}{0.002}{10.9091}{mi}{70}\RG{64.382}{-2.382}{7.9701}{rm}{49}\RG{65.866}{4.517}{7.9701}{sy}{0}\RG{71.876}{4.517}{7.9701}{rm}{49}%
}\hskip.1em\allowbreak{} 。各因子的逆仍为初等矩阵，故（5.4）推出（5.5）；空乘积约定为\allowbreak{}\hskip.1em \RB{9.662}{7.954}{2.207}{%
\RG{0.000}{0.003}{10.9091}{mi}{73}\RG{4.711}{-1.634}{7.9701}{mi}{110}%
}\hskip.1em\allowbreak{} 。最后，由行列式乘法公式，有限个初等矩阵的乘积的行列式是非零常数，故（5.5）推出（5.2）。上述五个条件因而等价。若\allowbreak{}\hskip.1em \RB{22.100}{8.184}{2.602}{%
\RG{0.000}{0.004}{10.9091}{mi}{66}\RG{8.662}{0.004}{10.9091}{mi}{59}\RG{13.451}{0.004}{10.9091}{mi}{67}%
}\hskip.1em\allowbreak{} 都是\allowbreak{}\hskip.1em \RB{8.202}{8.293}{0.500}{%
\RG{0.000}{0.004}{10.9091}{mi}{65}%
}\hskip.1em\allowbreak{} 的双侧逆矩阵，则\allowbreak{}\hskip.1em \RB{127.236}{8.686}{3.235}{%
\RG{0.000}{0.004}{10.9091}{mi}{66}\RG{11.696}{0.004}{10.9091}{rm}{61}\RG{23.035}{0.004}{10.9091}{mi}{66}\RG{31.700}{0.004}{10.9091}{rm}{40}\RG{35.854}{0.004}{10.9091}{mi}{65}\RG{43.862}{0.004}{10.9091}{mi}{67}\RG{52.288}{0.004}{10.9091}{rm}{41}\RG{59.466}{0.004}{10.9091}{rm}{61}\RG{70.812}{0.004}{10.9091}{rm}{40}\RG{74.965}{0.004}{10.9091}{mi}{66}\RG{83.627}{0.004}{10.9091}{mi}{65}\RG{91.637}{0.004}{10.9091}{rm}{41}\RG{95.791}{0.004}{10.9091}{mi}{67}\RG{107.248}{0.004}{10.9091}{rm}{61}\RG{118.587}{0.004}{10.9091}{mi}{67}%
}\hskip.1em\allowbreak{} ，故所得逆矩阵唯一。\allowbreak{}\hskip.1em \RB{8.366}{7.998}{0.500}{%
\RG{0.000}{0.004}{10.9091}{am}{3}%
}\hskip.1em\allowbreak{}

% Original page 10.
\section*{2016S07}
设\allowbreak{}\hskip.1em \RB{8.202}{8.289}{0.500}{%
\RG{0.000}{-0.000}{10.9091}{mi}{65}%
}\hskip.1em\allowbreak{} 为三阶实矩阵，满足\allowbreak{}\hskip.1em \RB{55.102}{9.455}{2.289}{%
\RG{0.000}{-0.000}{10.9091}{mi}{65}\RG{8.007}{-0.000}{10.9091}{mi}{65}\RG{16.016}{3.958}{7.9701}{rm}{84}\RG{24.777}{-0.000}{10.9091}{rm}{61}\RG{36.116}{-0.000}{10.9091}{mi}{107}\RG{42.038}{3.958}{7.9701}{rm}{50}\RG{46.357}{-0.000}{10.9091}{mi}{73}\RG{51.068}{-1.636}{7.9701}{rm}{51}%
}\hskip.1em\allowbreak{} 且\allowbreak{}\hskip.1em \RB{48.987}{9.295}{0.609}{%
\RG{0.000}{-0.000}{10.9091}{rm}{100}\RG{5.935}{-0.000}{10.9091}{rm}{101}\RG{10.680}{-0.000}{10.9091}{rm}{116}\RG{16.654}{-0.000}{10.9091}{mi}{65}\RG{27.692}{-0.000}{10.9091}{rm}{61}\RG{39.031}{-0.000}{10.9091}{mi}{107}\RG{44.952}{3.958}{7.9701}{rm}{51}%
}\hskip.1em\allowbreak{} ，其中\allowbreak{}\hskip.1em \RB{25.854}{8.071}{1.995}{%
\RG{0.000}{-0.000}{10.9091}{mi}{107}\RG{8.952}{-0.000}{10.9091}{sy}{21}\RG{20.467}{-0.000}{10.9091}{rm}{48}%
}\hskip.1em\allowbreak{} ，\allowbreak{}\hskip.1em \RB{13.505}{9.455}{0.500}{%
\RG{0.000}{-0.000}{10.9091}{mi}{65}\RG{8.008}{3.958}{7.9701}{rm}{84}%
}\hskip.1em\allowbreak{} 表示\allowbreak{}\hskip.1em \RB{8.202}{8.289}{0.500}{%
\RG{0.000}{-0.000}{10.9091}{mi}{65}%
}\hskip.1em\allowbreak{} 的转置，\allowbreak{}\hskip.1em \RB{8.746}{7.951}{2.289}{%
\RG{0.000}{-0.000}{10.9091}{mi}{73}\RG{4.711}{-1.636}{7.9701}{rm}{51}%
}\hskip.1em\allowbreak{}为三阶单位矩阵。证明：存在\allowbreak{}\hskip.1em \RB{46.369}{8.683}{3.226}{%
\RG{0.000}{0.001}{10.9091}{mi}{116}\RG{6.888}{0.001}{10.9091}{sy}{50}\RG{17.191}{0.001}{10.9091}{rm}{91}\RG{20.221}{0.001}{10.9091}{sy}{0}\RG{28.707}{0.001}{10.9091}{rm}{49}\RG{34.048}{0.001}{10.9091}{mi}{59}\RG{38.833}{0.001}{10.9091}{rm}{51}\RG{44.178}{0.001}{10.9091}{rm}{93}%
}\hskip.1em\allowbreak{} ，使

\[
\RB{140.959}{9.837}{2.293}{%
\RG{8.007}{4.501}{7.9701}{rm}{51}\RG{43.446}{4.501}{7.9701}{rm}{50}\RG{70.700}{4.501}{7.9701}{rm}{50}\RG{102.281}{4.501}{7.9701}{rm}{51}\RG{0.000}{-0.004}{10.9091}{mi}{65}\RG{14.750}{-0.004}{10.9091}{sy}{0}\RG{25.659}{-0.004}{10.9091}{mi}{116}\RG{29.521}{-0.004}{10.9091}{mi}{107}\RG{35.434}{-0.004}{10.9091}{mi}{65}\RG{50.188}{-0.004}{10.9091}{rm}{43}\RG{60.921}{-0.004}{10.9091}{mi}{116}\RG{64.783}{-0.004}{10.9091}{mi}{107}\RG{75.018}{-0.004}{10.9091}{mi}{65}\RG{85.450}{-0.004}{10.9091}{sy}{0}\RG{96.360}{-0.004}{10.9091}{mi}{107}\RG{106.599}{-0.004}{10.9091}{mi}{73}\RG{118.659}{-0.004}{10.9091}{rm}{61}\RG{129.997}{-0.004}{10.9091}{mi}{79}\RG{138.452}{-0.004}{10.9091}{mi}{58}\RG{111.310}{-1.640}{7.9701}{rm}{51}%
}
\]
证明：当\allowbreak{}\hskip.1em \RB{25.674}{8.068}{0.732}{%
\RG{0.000}{-0.003}{10.9091}{mi}{107}\RG{8.952}{-0.003}{10.9091}{rm}{61}\RG{20.286}{-0.003}{10.9091}{rm}{48}%
}\hskip.1em\allowbreak{} 时，写\allowbreak{}\hskip.1em \RB{42.344}{8.679}{3.623}{%
\RG{0.000}{-0.003}{10.9091}{mi}{65}\RG{11.038}{-0.003}{10.9091}{rm}{61}\RG{22.372}{-0.003}{10.9091}{rm}{40}\RG{26.530}{-0.003}{10.9091}{mi}{97}\RG{32.138}{-1.639}{7.9701}{mi}{105}\RG{34.771}{-1.639}{7.9701}{mi}{106}\RG{38.767}{-0.003}{10.9091}{rm}{41}%
}\hskip.1em\allowbreak{} 。由\allowbreak{}\hskip.1em \RB{44.536}{9.452}{0.732}{%
\RG{0.000}{-0.003}{10.9091}{mi}{65}\RG{8.007}{-0.003}{10.9091}{mi}{65}\RG{16.015}{3.956}{7.9701}{rm}{84}\RG{24.777}{-0.003}{10.9091}{rm}{61}\RG{36.116}{-0.003}{10.9091}{mi}{79}%
}\hskip.1em\allowbreak{} 得

\[
\RB{117.373}{17.153}{12.854}{%
\RG{4.634}{11.817}{7.9701}{rm}{51}\RG{19.727}{11.817}{7.9701}{rm}{51}\RG{0.000}{9.090}{9.0910}{ex}{88}\RG{15.093}{9.090}{9.0910}{ex}{88}\RG{35.793}{4.503}{7.9701}{rm}{50}\RG{85.269}{4.503}{7.9701}{rm}{84}\RG{30.186}{-0.002}{10.9091}{mi}{97}\RG{45.453}{-0.002}{10.9091}{rm}{61}\RG{56.788}{-0.002}{10.9091}{rm}{116}\RG{60.944}{-0.002}{10.9091}{rm}{114}\RG{65.101}{-0.002}{10.9091}{rm}{40}\RG{69.254}{-0.002}{10.9091}{mi}{65}\RG{77.262}{-0.002}{10.9091}{mi}{65}\RG{91.001}{-0.002}{10.9091}{rm}{41}\RG{98.180}{-0.002}{10.9091}{rm}{61}\RG{109.525}{-0.002}{10.9091}{rm}{48}\RG{114.865}{-0.002}{10.9091}{mi}{58}\RG{35.793}{-2.698}{7.9701}{mi}{105}\RG{38.426}{-2.698}{7.9701}{mi}{106}\RG{0.313}{-10.870}{7.9701}{mi}{105}\RG{2.945}{-10.870}{7.9701}{rm}{61}\RG{8.952}{-10.870}{7.9701}{rm}{49}\RG{14.950}{-10.870}{7.9701}{mi}{106}\RG{18.494}{-10.870}{7.9701}{rm}{61}\RG{24.501}{-10.870}{7.9701}{rm}{49}%
}
\]
各\allowbreak{}\hskip.1em \RB{11.744}{5.312}{3.619}{%
\RG{0.000}{0.001}{10.9091}{mi}{97}\RG{5.607}{-1.635}{7.9701}{mi}{105}\RG{8.240}{-1.635}{7.9701}{mi}{106}%
}\hskip.1em\allowbreak{} 都是实数，故每个\allowbreak{}\hskip.1em \RB{31.989}{7.756}{3.619}{%
\RG{0.000}{0.001}{10.9091}{mi}{97}\RG{5.607}{-1.635}{7.9701}{mi}{105}\RG{8.240}{-1.635}{7.9701}{mi}{106}\RG{15.267}{0.001}{10.9091}{rm}{61}\RG{26.602}{0.001}{10.9091}{rm}{48}%
}\hskip.1em\allowbreak{} ，即\allowbreak{}\hskip.1em \RB{30.797}{8.290}{0.728}{%
\RG{0.000}{0.001}{10.9091}{mi}{65}\RG{11.039}{0.001}{10.9091}{rm}{61}\RG{22.377}{0.001}{10.9091}{mi}{79}%
}\hskip.1em\allowbreak{} 。此时取\allowbreak{}\hskip.1em \RB{23.610}{7.756}{0.728}{%
\RG{0.000}{0.001}{10.9091}{mi}{116}\RG{6.888}{0.001}{10.9091}{rm}{61}\RG{18.222}{0.001}{10.9091}{rm}{48}%
}\hskip.1em\allowbreak{} 满足结论。

以下设\allowbreak{}\hskip.1em \RB{25.678}{8.066}{0.810}{%
\RG{0.000}{-0.005}{10.9091}{mi}{107}\RG{8.946}{-0.005}{10.9091}{mi}{62}\RG{20.291}{-0.005}{10.9091}{rm}{48}%
}\hskip.1em\allowbreak{} ，令\allowbreak{}\hskip.1em \RB{47.281}{9.290}{2.621}{%
\RG{0.000}{-0.005}{10.9091}{mi}{81}\RG{11.490}{-0.005}{10.9091}{rm}{61}\RG{22.829}{-0.005}{10.9091}{mi}{107}\RG{28.751}{3.954}{7.9701}{sy}{0}\RG{34.761}{3.954}{7.9701}{rm}{49}\RG{39.079}{-0.005}{10.9091}{mi}{65}%
}\hskip.1em\allowbreak{} 。则\allowbreak{}\hskip.1em \RB{45.767}{9.450}{2.621}{%
\RG{0.000}{-0.005}{10.9091}{mi}{81}\RG{8.466}{-0.005}{10.9091}{mi}{81}\RG{16.921}{3.954}{7.9701}{rm}{84}\RG{25.683}{-0.005}{10.9091}{rm}{61}\RG{37.022}{-0.005}{10.9091}{mi}{73}\RG{41.733}{-1.641}{7.9701}{rm}{51}%
}\hskip.1em\allowbreak{} 且\allowbreak{}\hskip.1em \RB{44.452}{8.176}{2.621}{%
\RG{0.000}{-0.005}{10.9091}{rm}{100}\RG{5.935}{-0.005}{10.9091}{rm}{101}\RG{10.680}{-0.005}{10.9091}{rm}{116}\RG{16.654}{-0.005}{10.9091}{mi}{81}\RG{28.145}{-0.005}{10.9091}{rm}{61}\RG{39.479}{-0.005}{10.9091}{rm}{49}%
}\hskip.1em\allowbreak{} ，从而\allowbreak{}\hskip.1em \RB{47.113}{9.450}{2.621}{%
\RG{0.000}{-0.005}{10.9091}{mi}{81}\RG{8.460}{3.954}{7.9701}{sy}{0}\RG{14.470}{3.954}{7.9701}{rm}{49}\RG{21.819}{-0.005}{10.9091}{rm}{61}\RG{33.157}{-0.005}{10.9091}{mi}{81}\RG{41.617}{3.954}{7.9701}{rm}{84}%
}\hskip.1em\allowbreak{} 。记\allowbreak{}\hskip.1em \RB{30.730}{8.677}{3.243}{%
\RG{0.000}{-0.005}{10.9091}{rm}{97}\RG{5.346}{-0.005}{10.9091}{rm}{100}\RG{11.280}{-0.005}{10.9091}{rm}{106}\RG{14.542}{-0.005}{10.9091}{rm}{40}\RG{18.694}{-0.005}{10.9091}{mi}{81}\RG{27.154}{-0.005}{10.9091}{rm}{41}%
}\hskip.1em\allowbreak{}为\allowbreak{}\hskip.1em \RB{8.420}{8.176}{2.620}{%
\RG{0.000}{-0.004}{10.9091}{mi}{81}%
}\hskip.1em\allowbreak{} 的伴随矩阵，由伴随矩阵公式得

\[
\RB{127.138}{10.002}{3.237}{%
\RG{85.741}{4.505}{7.9701}{sy}{0}\RG{91.751}{4.505}{7.9701}{rm}{49}\RG{118.899}{4.505}{7.9701}{rm}{84}\RG{0.000}{0.001}{10.9091}{rm}{97}\RG{5.346}{0.001}{10.9091}{rm}{100}\RG{11.280}{0.001}{10.9091}{rm}{106}\RG{14.542}{0.001}{10.9091}{rm}{40}\RG{18.694}{0.001}{10.9091}{mi}{81}\RG{27.154}{0.001}{10.9091}{rm}{41}\RG{34.332}{0.001}{10.9091}{rm}{61}\RG{45.678}{0.001}{10.9091}{rm}{100}\RG{51.612}{0.001}{10.9091}{rm}{101}\RG{56.358}{0.001}{10.9091}{rm}{116}\RG{60.514}{0.001}{10.9091}{rm}{40}\RG{64.667}{0.001}{10.9091}{mi}{81}\RG{73.127}{0.001}{10.9091}{rm}{41}\RG{77.281}{0.001}{10.9091}{mi}{81}\RG{99.100}{0.001}{10.9091}{rm}{61}\RG{110.439}{0.001}{10.9091}{mi}{81}\RG{124.631}{0.001}{10.9091}{mi}{58}%
}
\]
定义实数\allowbreak{}\hskip.1em \RB{42.725}{8.679}{3.241}{%
\RG{0.000}{-0.003}{10.9091}{mi}{116}\RG{6.888}{-0.003}{10.9091}{rm}{61}\RG{18.222}{-0.003}{10.9091}{rm}{116}\RG{22.379}{-0.003}{10.9091}{rm}{114}\RG{26.535}{-0.003}{10.9091}{rm}{40}\RG{30.688}{-0.003}{10.9091}{mi}{81}\RG{39.149}{-0.003}{10.9091}{rm}{41}%
}\hskip.1em\allowbreak{} ，并记\allowbreak{}\hskip.1em \RB{98.857}{8.679}{3.551}{%
\RG{0.000}{-0.003}{10.9091}{mi}{112}\RG{5.362}{-1.640}{7.9701}{mi}{81}\RG{11.894}{-0.003}{10.9091}{rm}{40}\RG{16.048}{-0.003}{10.9091}{mi}{122}\RG{21.496}{-0.003}{10.9091}{rm}{41}\RG{28.674}{-0.003}{10.9091}{rm}{61}\RG{40.020}{-0.003}{10.9091}{rm}{100}\RG{45.955}{-0.003}{10.9091}{rm}{101}\RG{50.700}{-0.003}{10.9091}{rm}{116}\RG{54.857}{-0.003}{10.9091}{rm}{40}\RG{59.009}{-0.003}{10.9091}{mi}{122}\RG{64.453}{-0.003}{10.9091}{mi}{73}\RG{69.168}{-1.640}{7.9701}{rm}{51}\RG{75.911}{-0.003}{10.9091}{sy}{0}\RG{86.820}{-0.003}{10.9091}{mi}{81}\RG{95.281}{-0.003}{10.9091}{rm}{41}%
}\hskip.1em\allowbreak{} 。三阶矩阵特征多项式的系数公式给出

\[
\RB{202.898}{9.842}{18.829}{%
\RG{45.467}{4.506}{7.9701}{rm}{51}\RG{93.643}{4.506}{7.9701}{rm}{50}\RG{0.000}{0.002}{10.9091}{mi}{112}\RG{11.893}{0.002}{10.9091}{rm}{40}\RG{16.047}{0.002}{10.9091}{mi}{122}\RG{21.495}{0.002}{10.9091}{rm}{41}\RG{28.673}{0.002}{10.9091}{rm}{61}\RG{40.018}{0.002}{10.9091}{mi}{122}\RG{52.209}{0.002}{10.9091}{sy}{0}\RG{63.118}{0.002}{10.9091}{rm}{116}\RG{67.275}{0.002}{10.9091}{rm}{114}\RG{71.431}{0.002}{10.9091}{rm}{40}\RG{75.581}{0.002}{10.9091}{mi}{81}\RG{84.041}{0.002}{10.9091}{rm}{41}\RG{88.194}{0.002}{10.9091}{mi}{122}\RG{100.385}{0.002}{10.9091}{rm}{43}\RG{111.120}{0.002}{10.9091}{rm}{116}\RG{115.277}{0.002}{10.9091}{rm}{114}\RG{119.433}{0.002}{10.9091}{rm}{40}\RG{123.590}{0.002}{10.9091}{rm}{97}\RG{128.935}{0.002}{10.9091}{rm}{100}\RG{134.870}{0.002}{10.9091}{rm}{106}\RG{138.132}{0.002}{10.9091}{rm}{40}\RG{142.273}{0.002}{10.9091}{mi}{81}\RG{150.733}{0.002}{10.9091}{rm}{41}\RG{154.890}{0.002}{10.9091}{rm}{41}\RG{159.041}{0.002}{10.9091}{mi}{122}\RG{166.914}{0.002}{10.9091}{sy}{0}\RG{177.823}{0.002}{10.9091}{rm}{100}\RG{183.758}{0.002}{10.9091}{rm}{101}\RG{188.503}{0.002}{10.9091}{rm}{116}\RG{194.477}{0.002}{10.9091}{mi}{81}\RG{5.361}{-1.634}{7.9701}{mi}{81}\RG{45.467}{-12.985}{7.9701}{rm}{51}\RG{72.424}{-12.985}{7.9701}{rm}{50}\RG{28.679}{-17.489}{10.9091}{rm}{61}\RG{40.018}{-17.489}{10.9091}{mi}{122}\RG{52.209}{-17.489}{10.9091}{sy}{0}\RG{63.118}{-17.489}{10.9091}{mi}{116}\RG{66.980}{-17.489}{10.9091}{mi}{122}\RG{79.167}{-17.489}{10.9091}{rm}{43}\RG{89.899}{-17.489}{10.9091}{mi}{116}\RG{93.761}{-17.489}{10.9091}{mi}{122}\RG{101.628}{-17.489}{10.9091}{sy}{0}\RG{112.538}{-17.489}{10.9091}{rm}{49}\RG{117.879}{-17.489}{10.9091}{mi}{58}%
}
\]
特别地，\allowbreak{}\hskip.1em \RB{45.307}{8.682}{3.547}{%
\RG{0.000}{0.000}{10.9091}{mi}{112}\RG{5.361}{-1.636}{7.9701}{mi}{81}\RG{11.894}{0.000}{10.9091}{rm}{40}\RG{16.050}{0.000}{10.9091}{rm}{49}\RG{21.396}{0.000}{10.9091}{rm}{41}\RG{28.574}{0.000}{10.9091}{rm}{61}\RG{39.920}{0.000}{10.9091}{rm}{48}%
}\hskip.1em\allowbreak{} ，所以\allowbreak{}\hskip.1em \RB{4.973}{7.755}{0.500}{%
\RG{0.000}{0.000}{10.9091}{rm}{49}%
}\hskip.1em\allowbreak{} 是\allowbreak{}\hskip.1em \RB{8.420}{8.180}{2.616}{%
\RG{0.000}{0.000}{10.9091}{mi}{81}%
}\hskip.1em\allowbreak{} 的一个特征值。

任取\allowbreak{}\hskip.1em \RB{8.420}{8.183}{2.614}{%
\RG{0.000}{0.003}{10.9091}{mi}{81}%
}\hskip.1em\allowbreak{} 的复特征值\allowbreak{}\hskip.1em \RB{6.336}{8.074}{0.617}{%
\RG{0.000}{0.003}{10.9091}{mi}{21}%
}\hskip.1em\allowbreak{} 及相应的非零特征向量\allowbreak{}\hskip.1em \RB{31.196}{9.299}{0.934}{%
\RG{0.000}{0.003}{10.9091}{mi}{118}\RG{8.783}{0.003}{10.9091}{sy}{50}\RG{19.282}{0.003}{10.9091}{bb}{67}\RG{27.161}{3.962}{7.9701}{rm}{51}%
}\hskip.1em\allowbreak{} 。对\allowbreak{}\hskip.1em \RB{97.303}{9.459}{3.235}{%
\RG{0.000}{0.003}{10.9091}{mi}{122}\RG{8.675}{0.003}{10.9091}{rm}{61}\RG{20.206}{0.003}{10.9091}{rm}{40}\RG{24.364}{0.003}{10.9091}{mi}{122}\RG{29.326}{-1.633}{7.9701}{rm}{49}\RG{33.644}{0.003}{10.9091}{mi}{59}\RG{38.423}{0.003}{10.9091}{mi}{122}\RG{43.392}{-1.633}{7.9701}{rm}{50}\RG{47.710}{0.003}{10.9091}{mi}{59}\RG{52.488}{0.003}{10.9091}{mi}{122}\RG{57.457}{-1.633}{7.9701}{rm}{51}\RG{61.776}{0.003}{10.9091}{rm}{41}\RG{65.930}{3.962}{7.9701}{rm}{84}\RG{74.889}{0.003}{10.9091}{sy}{50}\RG{85.388}{0.003}{10.9091}{bb}{67}\RG{93.268}{3.962}{7.9701}{rm}{51}%
}\hskip.1em\allowbreak{} ，定义标准Hermite范数\allowbreak{}\hskip.1em \RB{93.530}{9.917}{3.266}{%
\RG{0.000}{0.004}{10.9091}{sy}{107}\RG{5.454}{0.004}{10.9091}{mi}{122}\RG{10.903}{0.004}{10.9091}{sy}{107}\RG{19.387}{0.004}{10.9091}{rm}{61}\RG{30.722}{0.004}{10.9091}{rm}{40}\RG{34.880}{7.276}{9.0910}{ex}{80}\RG{44.476}{4.397}{7.9701}{rm}{51}\RG{44.476}{-2.694}{7.9701}{mi}{105}\RG{47.108}{-2.694}{7.9701}{rm}{61}\RG{53.115}{-2.694}{7.9701}{rm}{49}\RG{59.254}{0.004}{10.9091}{sy}{106}\RG{62.284}{0.004}{10.9091}{mi}{122}\RG{67.246}{-1.632}{7.9701}{mi}{105}\RG{70.332}{0.004}{10.9091}{sy}{106}\RG{73.362}{3.962}{7.9701}{rm}{50}\RG{77.680}{0.004}{10.9091}{rm}{41}\RG{81.834}{3.962}{7.9701}{rm}{49}\RG{85.698}{3.962}{7.9701}{mi}{61}\RG{89.561}{3.962}{7.9701}{rm}{50}%
}\hskip.1em\allowbreak{} 。矩阵\allowbreak{}\hskip.1em \RB{8.420}{8.184}{2.613}{%
\RG{0.000}{0.004}{10.9091}{mi}{81}%
}\hskip.1em\allowbreak{} 为实矩阵且\allowbreak{}\hskip.1em \RB{45.767}{9.459}{2.613}{%
\RG{0.000}{0.004}{10.9091}{mi}{81}\RG{8.460}{3.962}{7.9701}{rm}{84}\RG{14.192}{0.004}{10.9091}{mi}{81}\RG{25.683}{0.004}{10.9091}{rm}{61}\RG{37.022}{0.004}{10.9091}{mi}{73}\RG{41.733}{-1.632}{7.9701}{rm}{51}%
}\hskip.1em\allowbreak{} ，故它保持该范数，因而

\[
\RB{103.210}{8.677}{3.232}{%
\RG{0.000}{-0.005}{10.9091}{sy}{106}\RG{3.030}{-0.005}{10.9091}{mi}{21}\RG{9.262}{-0.005}{10.9091}{sy}{106}\RG{14.106}{-0.005}{10.9091}{sy}{107}\RG{19.565}{-0.005}{10.9091}{mi}{118}\RG{25.121}{-0.005}{10.9091}{sy}{107}\RG{33.605}{-0.005}{10.9091}{rm}{61}\RG{44.944}{-0.005}{10.9091}{sy}{107}\RG{50.399}{-0.005}{10.9091}{mi}{81}\RG{58.865}{-0.005}{10.9091}{mi}{118}\RG{64.414}{-0.005}{10.9091}{sy}{107}\RG{72.900}{-0.005}{10.9091}{rm}{61}\RG{84.238}{-0.005}{10.9091}{sy}{107}\RG{89.693}{-0.005}{10.9091}{mi}{118}\RG{95.248}{-0.005}{10.9091}{sy}{107}\RG{100.703}{-0.005}{10.9091}{mi}{58}%
}
\]
由\allowbreak{}\hskip.1em \RB{25.308}{8.311}{2.856}{%
\RG{0.000}{-0.000}{10.9091}{mi}{118}\RG{8.586}{-0.000}{10.9091}{sy}{54}\RG{8.586}{-0.000}{10.9091}{rm}{61}\RG{19.921}{-0.000}{10.9091}{rm}{48}%
}\hskip.1em\allowbreak{} 得\allowbreak{}\hskip.1em \RB{31.630}{8.682}{3.227}{%
\RG{0.000}{-0.000}{10.9091}{sy}{106}\RG{3.030}{-0.000}{10.9091}{mi}{21}\RG{9.262}{-0.000}{10.9091}{sy}{106}\RG{15.322}{-0.000}{10.9091}{rm}{61}\RG{26.657}{-0.000}{10.9091}{rm}{49}%
}\hskip.1em\allowbreak{} 。记\allowbreak{}\hskip.1em \RB{8.420}{8.180}{2.616}{%
\RG{0.000}{-0.000}{10.9091}{mi}{81}%
}\hskip.1em\allowbreak{} 其余两个特征值为\allowbreak{}\hskip.1em \RB{17.311}{5.311}{2.846}{%
\RG{0.000}{-0.000}{10.9091}{mi}{22}\RG{6.426}{-0.000}{10.9091}{mi}{59}\RG{11.204}{-0.000}{10.9091}{mi}{23}%
}\hskip.1em\allowbreak{} ，按代数重数计。由特征多项式的二次项系数得\allowbreak{}\hskip.1em \RB{62.415}{7.756}{2.845}{%
\RG{0.000}{0.001}{10.9091}{mi}{116}\RG{6.888}{0.001}{10.9091}{rm}{61}\RG{18.222}{0.001}{10.9091}{rm}{49}\RG{25.990}{0.001}{10.9091}{rm}{43}\RG{36.724}{0.001}{10.9091}{mi}{22}\RG{45.575}{0.001}{10.9091}{rm}{43}\RG{56.308}{0.001}{10.9091}{mi}{23}%
}\hskip.1em\allowbreak{} ，从而

\[
\RB{149.006}{8.679}{3.230}{%
\RG{0.000}{-0.003}{10.9091}{sy}{106}\RG{3.030}{-0.003}{10.9091}{mi}{116}\RG{9.312}{-0.003}{10.9091}{sy}{0}\RG{20.221}{-0.003}{10.9091}{rm}{49}\RG{25.562}{-0.003}{10.9091}{sy}{106}\RG{31.623}{-0.003}{10.9091}{rm}{61}\RG{42.962}{-0.003}{10.9091}{sy}{106}\RG{45.992}{-0.003}{10.9091}{mi}{22}\RG{54.843}{-0.003}{10.9091}{rm}{43}\RG{65.576}{-0.003}{10.9091}{mi}{23}\RG{71.542}{-0.003}{10.9091}{sy}{106}\RG{77.597}{-0.003}{10.9091}{sy}{20}\RG{89.117}{-0.003}{10.9091}{sy}{106}\RG{92.148}{-0.003}{10.9091}{mi}{22}\RG{98.574}{-0.003}{10.9091}{sy}{106}\RG{104.029}{-0.003}{10.9091}{rm}{43}\RG{114.762}{-0.003}{10.9091}{sy}{106}\RG{117.792}{-0.003}{10.9091}{mi}{23}\RG{123.758}{-0.003}{10.9091}{sy}{106}\RG{129.819}{-0.003}{10.9091}{rm}{61}\RG{141.153}{-0.003}{10.9091}{rm}{50}\RG{146.498}{-0.003}{10.9091}{mi}{58}%
}
\]
由于\allowbreak{}\hskip.1em \RB{25.360}{7.975}{0.935}{%
\RG{0.000}{0.002}{10.9091}{mi}{116}\RG{6.888}{0.002}{10.9091}{sy}{50}\RG{17.191}{0.002}{10.9091}{bb}{82}%
}\hskip.1em\allowbreak{} ，这给出\allowbreak{}\hskip.1em \RB{52.162}{7.757}{1.993}{%
\RG{0.000}{0.002}{10.9091}{sy}{0}\RG{8.485}{0.002}{10.9091}{rm}{49}\RG{16.856}{0.002}{10.9091}{sy}{20}\RG{28.371}{0.002}{10.9091}{mi}{116}\RG{35.259}{0.002}{10.9091}{sy}{20}\RG{46.774}{0.002}{10.9091}{rm}{51}%
}\hskip.1em\allowbreak{} 。

由Cayley–Hamilton定理，\allowbreak{}\hskip.1em \RB{51.451}{8.678}{3.552}{%
\RG{0.000}{-0.004}{10.9091}{mi}{112}\RG{5.361}{-1.641}{7.9701}{mi}{81}\RG{11.894}{-0.004}{10.9091}{rm}{40}\RG{16.047}{-0.004}{10.9091}{mi}{81}\RG{24.507}{-0.004}{10.9091}{rm}{41}\RG{31.686}{-0.004}{10.9091}{rm}{61}\RG{43.030}{-0.004}{10.9091}{mi}{79}%
}\hskip.1em\allowbreak{} 。利用\allowbreak{}\hskip.1em \RB{36.710}{8.286}{2.620}{%
\RG{0.000}{-0.004}{10.9091}{mi}{65}\RG{11.038}{-0.004}{10.9091}{rm}{61}\RG{22.377}{-0.004}{10.9091}{mi}{107}\RG{28.289}{-0.004}{10.9091}{mi}{81}%
}\hskip.1em\allowbreak{} ，得到

\[
\RB{238.553}{9.842}{18.219}{%
\RG{8.008}{4.505}{7.9701}{rm}{51}\RG{43.447}{4.505}{7.9701}{rm}{50}\RG{70.701}{4.505}{7.9701}{rm}{50}\RG{102.282}{4.505}{7.9701}{rm}{51}\RG{135.920}{4.505}{7.9701}{rm}{51}\RG{152.852}{4.505}{7.9701}{rm}{51}\RG{182.821}{4.505}{7.9701}{rm}{50}\RG{0.000}{0.001}{10.9091}{mi}{65}\RG{14.750}{0.001}{10.9091}{sy}{0}\RG{25.660}{0.001}{10.9091}{mi}{116}\RG{29.522}{0.001}{10.9091}{mi}{107}\RG{35.435}{0.001}{10.9091}{mi}{65}\RG{50.189}{0.001}{10.9091}{rm}{43}\RG{60.921}{0.001}{10.9091}{mi}{116}\RG{64.783}{0.001}{10.9091}{mi}{107}\RG{75.019}{0.001}{10.9091}{mi}{65}\RG{85.450}{0.001}{10.9091}{sy}{0}\RG{96.360}{0.001}{10.9091}{mi}{107}\RG{106.600}{0.001}{10.9091}{mi}{73}\RG{118.659}{0.001}{10.9091}{rm}{61}\RG{129.998}{0.001}{10.9091}{mi}{107}\RG{140.238}{0.001}{10.9091}{rm}{40}\RG{144.392}{0.001}{10.9091}{mi}{81}\RG{159.595}{0.001}{10.9091}{sy}{0}\RG{170.504}{0.001}{10.9091}{mi}{116}\RG{174.366}{0.001}{10.9091}{mi}{81}\RG{189.564}{0.001}{10.9091}{rm}{43}\RG{200.297}{0.001}{10.9091}{mi}{116}\RG{204.158}{0.001}{10.9091}{mi}{81}\RG{215.038}{0.001}{10.9091}{sy}{0}\RG{225.947}{0.001}{10.9091}{mi}{73}\RG{234.977}{0.001}{10.9091}{rm}{41}\RG{111.311}{-1.635}{7.9701}{rm}{51}\RG{230.659}{-1.635}{7.9701}{rm}{51}\RG{118.659}{-17.490}{10.9091}{rm}{61}\RG{129.998}{-17.490}{10.9091}{mi}{79}\RG{138.453}{-17.490}{10.9091}{mi}{58}%
}
\]
两种\allowbreak{}\hskip.1em \RB{5.966}{8.070}{0.610}{%
\RG{0.000}{-0.001}{10.9091}{mi}{107}%
}\hskip.1em\allowbreak{} 的情形均已证明：\allowbreak{}\hskip.1em \RB{8.366}{7.994}{0.501}{%
\RG{0.000}{-0.001}{10.9091}{am}{3}%
}\hskip.1em\allowbreak{}

% Original page 11.
\section*{2016S08}
设\allowbreak{}\hskip.1em \RB{35.609}{7.949}{1.996}{%
\RG{0.000}{-0.002}{10.9091}{mi}{70}\RG{12.744}{-0.002}{10.9091}{sy}{18}\RG{25.574}{-0.002}{10.9091}{mi}{75}%
}\hskip.1em\allowbreak{} 是两个数域，\allowbreak{}\hskip.1em \RB{21.391}{8.287}{2.607}{%
\RG{0.000}{-0.002}{10.9091}{mi}{65}\RG{8.007}{-0.002}{10.9091}{mi}{59}\RG{12.786}{-0.002}{10.9091}{mi}{66}%
}\hskip.1em\allowbreak{} 是\allowbreak{}\hskip.1em \RB{8.562}{7.927}{0.502}{%
\RG{0.000}{-0.002}{10.9091}{mi}{70}%
}\hskip.1em\allowbreak{} 上的方阵。分别用线性空间理论及多项式理论证明：\allowbreak{}\hskip.1em \RB{21.391}{8.288}{2.606}{%
\RG{0.000}{-0.001}{10.9091}{mi}{65}\RG{8.007}{-0.001}{10.9091}{mi}{59}\RG{12.786}{-0.001}{10.9091}{mi}{66}%
}\hskip.1em\allowbreak{} 在\allowbreak{}\hskip.1em \RB{8.562}{7.928}{0.501}{%
\RG{0.000}{-0.001}{10.9091}{mi}{70}%
}\hskip.1em\allowbreak{} 上相似，当且仅当它们在\allowbreak{}\hskip.1em \RB{10.035}{7.950}{0.501}{%
\RG{0.000}{-0.001}{10.9091}{mi}{75}%
}\hskip.1em\allowbreak{} 上相似。

证明：若\allowbreak{}\hskip.1em \RB{21.391}{8.285}{2.610}{%
\RG{0.000}{-0.005}{10.9091}{mi}{65}\RG{8.007}{-0.005}{10.9091}{mi}{59}\RG{12.786}{-0.005}{10.9091}{mi}{66}%
}\hskip.1em\allowbreak{} 的阶数不同，则它们在两个数域上均不相似。以下设二者均为\allowbreak{}\hskip.1em \RB{6.555}{5.306}{0.614}{%
\RG{0.000}{-0.005}{10.9091}{mi}{110}%
}\hskip.1em\allowbreak{} 阶方阵。若\allowbreak{}\hskip.1em \RB{53.308}{8.677}{3.243}{%
\RG{0.000}{-0.005}{10.9091}{mi}{80}\RG{11.399}{-0.005}{10.9091}{sy}{50}\RG{21.702}{-0.005}{10.9091}{mi}{77}\RG{32.036}{-1.641}{7.9701}{mi}{110}\RG{37.179}{-0.005}{10.9091}{rm}{40}\RG{41.334}{-0.005}{10.9091}{mi}{70}\RG{49.732}{-0.005}{10.9091}{rm}{41}%
}\hskip.1em\allowbreak{} 可逆且\allowbreak{}\hskip.1em \RB{47.707}{8.285}{0.505}{%
\RG{0.000}{-0.005}{10.9091}{mi}{65}\RG{8.007}{-0.005}{10.9091}{mi}{80}\RG{19.406}{-0.005}{10.9091}{rm}{61}\RG{30.745}{-0.005}{10.9091}{mi}{80}\RG{39.102}{-0.005}{10.9091}{mi}{66}%
}\hskip.1em\allowbreak{} ，则\allowbreak{}\hskip.1em \RB{44.775}{8.306}{2.861}{%
\RG{0.000}{-0.005}{10.9091}{rm}{100}\RG{5.935}{-0.005}{10.9091}{rm}{101}\RG{10.680}{-0.005}{10.9091}{rm}{116}\RG{16.654}{-0.005}{10.9091}{mi}{80}\RG{28.053}{-0.005}{10.9091}{sy}{54}\RG{28.053}{-0.005}{10.9091}{rm}{61}\RG{39.388}{-0.005}{10.9091}{rm}{48}%
}\hskip.1em\allowbreak{} 在\allowbreak{}\hskip.1em \RB{10.035}{7.946}{0.505}{%
\RG{0.000}{-0.005}{10.9091}{mi}{75}%
}\hskip.1em\allowbreak{} 中仍成立，故\allowbreak{}\hskip.1em \RB{21.391}{8.285}{2.610}{%
\RG{0.000}{-0.005}{10.9091}{mi}{65}\RG{8.007}{-0.005}{10.9091}{mi}{59}\RG{12.786}{-0.005}{10.9091}{mi}{66}%
}\hskip.1em\allowbreak{} 在\allowbreak{}\hskip.1em \RB{10.035}{7.946}{0.505}{%
\RG{0.000}{-0.005}{10.9091}{mi}{75}%
}\hskip.1em\allowbreak{} 上相似。只须

证明：逆命题。

先用线性空间理论证明：设\allowbreak{}\hskip.1em \RB{54.816}{8.681}{3.239}{%
\RG{0.000}{-0.001}{10.9091}{mi}{80}\RG{11.416}{-0.001}{10.9091}{sy}{50}\RG{21.736}{-0.001}{10.9091}{mi}{77}\RG{32.070}{-1.637}{7.9701}{mi}{110}\RG{37.214}{-0.001}{10.9091}{rm}{40}\RG{41.368}{-0.001}{10.9091}{mi}{75}\RG{51.239}{-0.001}{10.9091}{rm}{41}%
}\hskip.1em\allowbreak{} 可逆且\allowbreak{}\hskip.1em \RB{47.752}{8.288}{0.501}{%
\RG{0.000}{-0.001}{10.9091}{mi}{65}\RG{8.007}{-0.001}{10.9091}{mi}{80}\RG{19.424}{-0.001}{10.9091}{rm}{61}\RG{30.779}{-0.001}{10.9091}{mi}{80}\RG{39.146}{-0.001}{10.9091}{mi}{66}%
}\hskip.1em\allowbreak{} 。令\allowbreak{}\hskip.1em \RB{8.715}{7.928}{0.501}{%
\RG{0.000}{-0.001}{10.9091}{mi}{69}%
}\hskip.1em\allowbreak{} 为\allowbreak{}\hskip.1em \RB{8.573}{7.950}{0.501}{%
\RG{0.000}{-0.001}{10.9091}{mi}{80}%
}\hskip.1em\allowbreak{} 的全部元素在\allowbreak{}\hskip.1em \RB{10.035}{7.951}{0.500}{%
\RG{0.000}{0.000}{10.9091}{mi}{75}%
}\hskip.1em\allowbreak{} 中张成的\allowbreak{}\hskip.1em \RB{8.562}{7.929}{0.500}{%
\RG{0.000}{0.000}{10.9091}{mi}{70}%
}\hskip.1em\allowbreak{} -线性子空间，取\allowbreak{}\hskip.1em \RB{8.715}{7.929}{0.500}{%
\RG{0.000}{0.000}{10.9091}{mi}{69}%
}\hskip.1em\allowbreak{} 的一组基\allowbreak{}\hskip.1em \RB{45.693}{5.311}{2.605}{%
\RG{0.000}{0.000}{10.9091}{mi}{11}\RG{6.794}{-1.636}{7.9701}{rm}{49}\RG{11.113}{0.000}{10.9091}{mi}{59}\RG{15.891}{0.000}{10.9091}{mi}{58}\RG{20.680}{0.000}{10.9091}{mi}{58}\RG{25.458}{0.000}{10.9091}{mi}{58}\RG{30.247}{0.000}{10.9091}{mi}{59}\RG{35.037}{0.000}{10.9091}{mi}{11}\RG{41.833}{-1.636}{7.9701}{mi}{114}%
}\hskip.1em\allowbreak{} 。逐项按此基展开，得到唯一确定的矩阵\allowbreak{}\hskip.1em \RB{91.111}{8.683}{3.237}{%
\RG{0.000}{0.001}{10.9091}{mi}{80}\RG{6.885}{-1.635}{7.9701}{rm}{49}\RG{11.203}{0.001}{10.9091}{mi}{59}\RG{15.981}{0.001}{10.9091}{mi}{58}\RG{20.770}{0.001}{10.9091}{mi}{58}\RG{25.549}{0.001}{10.9091}{mi}{58}\RG{30.338}{0.001}{10.9091}{mi}{59}\RG{35.127}{0.001}{10.9091}{mi}{80}\RG{42.015}{-1.635}{7.9701}{mi}{114}\RG{49.202}{0.001}{10.9091}{sy}{50}\RG{59.505}{0.001}{10.9091}{mi}{77}\RG{69.838}{-1.635}{7.9701}{mi}{110}\RG{74.982}{0.001}{10.9091}{rm}{40}\RG{79.136}{0.001}{10.9091}{mi}{70}\RG{87.535}{0.001}{10.9091}{rm}{41}%
}\hskip.1em\allowbreak{} ，使得

\[
\RB{169.442}{46.799}{11.442}{%
\RG{27.452}{43.092}{7.9701}{mi}{114}\RG{22.737}{40.364}{9.0910}{ex}{88}\RG{0.000}{31.273}{10.9091}{mi}{80}\RG{11.399}{31.273}{10.9091}{rm}{61}\RG{37.687}{31.273}{10.9091}{mi}{11}\RG{47.567}{31.273}{10.9091}{mi}{80}\RG{57.538}{31.273}{10.9091}{mi}{59}\RG{44.481}{29.637}{7.9701}{mi}{105}\RG{54.451}{29.637}{7.9701}{mi}{105}\RG{23.051}{20.405}{7.9701}{mi}{105}\RG{25.682}{20.405}{7.9701}{rm}{61}\RG{31.689}{20.405}{7.9701}{rm}{49}\RG{88.565}{11.818}{7.9701}{mi}{114}\RG{83.850}{9.091}{9.0910}{ex}{88}\RG{3.028}{0.000}{10.9091}{rm}{48}\RG{11.395}{0.000}{10.9091}{rm}{61}\RG{22.737}{0.000}{10.9091}{mi}{65}\RG{30.744}{0.000}{10.9091}{mi}{80}\RG{41.538}{0.000}{10.9091}{sy}{0}\RG{52.447}{0.000}{10.9091}{mi}{80}\RG{60.814}{0.000}{10.9091}{mi}{66}\RG{72.512}{0.000}{10.9091}{rm}{61}\RG{98.800}{-0.000}{10.9091}{mi}{11}\RG{108.679}{-0.000}{10.9091}{rm}{40}\RG{112.833}{-0.000}{10.9091}{mi}{65}\RG{120.841}{-0.000}{10.9091}{mi}{80}\RG{133.236}{-0.000}{10.9091}{sy}{0}\RG{144.145}{-0.000}{10.9091}{mi}{80}\RG{154.116}{-0.000}{10.9091}{mi}{66}\RG{162.781}{-0.000}{10.9091}{rm}{41}\RG{166.935}{-0.000}{10.9091}{mi}{58}\RG{105.594}{-1.636}{7.9701}{mi}{105}\RG{127.726}{-1.636}{7.9701}{mi}{105}\RG{151.030}{-1.636}{7.9701}{mi}{105}\RG{84.164}{-10.869}{7.9701}{mi}{105}\RG{86.795}{-10.869}{7.9701}{rm}{61}\RG{92.802}{-10.869}{7.9701}{rm}{49}%
}
\]
由于\allowbreak{}\hskip.1em \RB{45.694}{5.314}{2.603}{%
\RG{0.000}{0.003}{10.9091}{mi}{11}\RG{6.795}{-1.634}{7.9701}{rm}{49}\RG{11.113}{0.003}{10.9091}{mi}{59}\RG{15.891}{0.003}{10.9091}{mi}{58}\RG{20.680}{0.003}{10.9091}{mi}{58}\RG{25.469}{0.003}{10.9091}{mi}{58}\RG{30.247}{0.003}{10.9091}{mi}{59}\RG{35.037}{0.003}{10.9091}{mi}{11}\RG{41.834}{-1.634}{7.9701}{mi}{114}%
}\hskip.1em\allowbreak{} 在\allowbreak{}\hskip.1em \RB{8.562}{7.932}{0.500}{%
\RG{0.000}{0.003}{10.9091}{mi}{70}%
}\hskip.1em\allowbreak{} 上线性无关，对最后一个矩阵等式的每个元素比较系数，得到\allowbreak{}\hskip.1em \RB{50.923}{8.292}{2.206}{%
\RG{0.000}{0.003}{10.9091}{mi}{65}\RG{8.007}{0.003}{10.9091}{mi}{80}\RG{14.893}{-1.633}{7.9701}{mi}{105}\RG{21.009}{0.003}{10.9091}{rm}{61}\RG{32.348}{0.003}{10.9091}{mi}{80}\RG{39.233}{-1.633}{7.9701}{mi}{105}\RG{42.318}{0.003}{10.9091}{mi}{66}%
}\hskip.1em\allowbreak{} ，其中\allowbreak{}\hskip.1em \RB{43.224}{7.757}{1.992}{%
\RG{0.000}{0.003}{10.9091}{rm}{49}\RG{8.371}{0.003}{10.9091}{sy}{20}\RG{19.886}{0.003}{10.9091}{mi}{105}\RG{26.561}{0.003}{10.9091}{sy}{20}\RG{38.076}{0.003}{10.9091}{mi}{114}%
}\hskip.1em\allowbreak{} 。

令\allowbreak{}\hskip.1em \RB{39.819}{7.316}{2.608}{%
\RG{0.000}{-0.003}{10.9091}{mi}{116}\RG{3.857}{-1.640}{7.9701}{rm}{49}\RG{8.175}{-0.003}{10.9091}{mi}{59}\RG{12.953}{-0.003}{10.9091}{mi}{58}\RG{17.743}{-0.003}{10.9091}{mi}{58}\RG{22.532}{-0.003}{10.9091}{mi}{58}\RG{27.310}{-0.003}{10.9091}{mi}{59}\RG{32.099}{-0.003}{10.9091}{mi}{116}\RG{35.959}{-1.640}{7.9701}{mi}{114}%
}\hskip.1em\allowbreak{} 为独立未定元，定义

\[
\RB{203.932}{16.862}{11.443}{%
\RG{87.322}{15.998}{9.0910}{ex}{32}\RG{126.382}{15.998}{9.0910}{ex}{33}\RG{99.234}{11.816}{7.9701}{mi}{114}\RG{94.519}{9.090}{9.0910}{ex}{88}\RG{0.000}{-0.002}{10.9091}{mi}{72}\RG{9.693}{-0.002}{10.9091}{rm}{40}\RG{13.847}{-0.002}{10.9091}{mi}{116}\RG{22.023}{-0.002}{10.9091}{mi}{59}\RG{26.801}{-0.002}{10.9091}{mi}{58}\RG{31.590}{-0.002}{10.9091}{mi}{58}\RG{36.368}{-0.002}{10.9091}{mi}{58}\RG{41.157}{-0.002}{10.9091}{mi}{59}\RG{45.947}{-0.002}{10.9091}{mi}{116}\RG{53.963}{-0.002}{10.9091}{rm}{41}\RG{61.142}{-0.002}{10.9091}{rm}{61}\RG{72.487}{-0.002}{10.9091}{rm}{100}\RG{78.422}{-0.002}{10.9091}{rm}{101}\RG{83.167}{-0.002}{10.9091}{rm}{116}\RG{109.469}{-0.002}{10.9091}{mi}{116}\RG{116.411}{-0.002}{10.9091}{mi}{80}\RG{136.610}{-0.002}{10.9091}{sy}{50}\RG{146.913}{-0.002}{10.9091}{mi}{70}\RG{155.311}{-0.002}{10.9091}{rm}{91}\RG{158.341}{-0.002}{10.9091}{mi}{116}\RG{166.517}{-0.002}{10.9091}{mi}{59}\RG{171.295}{-0.002}{10.9091}{mi}{58}\RG{176.085}{-0.002}{10.9091}{mi}{58}\RG{180.863}{-0.002}{10.9091}{mi}{58}\RG{185.652}{-0.002}{10.9091}{mi}{59}\RG{190.441}{-0.002}{10.9091}{mi}{116}\RG{198.457}{-0.002}{10.9091}{rm}{93}\RG{201.424}{-0.002}{10.9091}{mi}{58}\RG{17.705}{-1.638}{7.9701}{rm}{49}\RG{49.807}{-1.638}{7.9701}{mi}{114}\RG{113.326}{-1.638}{7.9701}{mi}{105}\RG{123.297}{-1.638}{7.9701}{mi}{105}\RG{162.198}{-1.638}{7.9701}{rm}{49}\RG{194.301}{-1.638}{7.9701}{mi}{114}\RG{94.832}{-10.871}{7.9701}{mi}{105}\RG{97.464}{-10.871}{7.9701}{rm}{61}\RG{103.471}{-10.871}{7.9701}{rm}{49}%
}
\]
因为\allowbreak{}\hskip.1em \RB{123.135}{8.684}{3.236}{%
\RG{0.000}{0.002}{10.9091}{mi}{72}\RG{9.693}{0.002}{10.9091}{rm}{40}\RG{13.846}{0.002}{10.9091}{mi}{11}\RG{20.641}{-1.634}{7.9701}{rm}{49}\RG{24.959}{0.002}{10.9091}{mi}{59}\RG{29.737}{0.002}{10.9091}{mi}{58}\RG{34.526}{0.002}{10.9091}{mi}{58}\RG{39.316}{0.002}{10.9091}{mi}{58}\RG{44.094}{0.002}{10.9091}{mi}{59}\RG{48.883}{0.002}{10.9091}{mi}{11}\RG{55.681}{-1.634}{7.9701}{mi}{114}\RG{59.836}{0.002}{10.9091}{rm}{41}\RG{67.014}{0.002}{10.9091}{rm}{61}\RG{78.360}{0.002}{10.9091}{rm}{100}\RG{84.295}{0.002}{10.9091}{rm}{101}\RG{89.040}{0.002}{10.9091}{rm}{116}\RG{95.014}{0.002}{10.9091}{mi}{80}\RG{106.413}{0.002}{10.9091}{sy}{54}\RG{106.413}{0.002}{10.9091}{rm}{61}\RG{117.747}{0.002}{10.9091}{rm}{48}%
}\hskip.1em\allowbreak{} ，\allowbreak{}\hskip.1em \RB{9.860}{7.953}{0.500}{%
\RG{0.000}{0.002}{10.9091}{mi}{72}%
}\hskip.1em\allowbreak{} 不是零多项式。

数域\allowbreak{}\hskip.1em \RB{8.562}{7.925}{0.505}{%
\RG{0.000}{-0.005}{10.9091}{mi}{70}%
}\hskip.1em\allowbreak{} 包含\allowbreak{}\hskip.1em \RB{8.627}{8.176}{2.457}{%
\RG{0.000}{-0.005}{10.9091}{bb}{81}%
}\hskip.1em\allowbreak{} ，因而是无限域。在无限域上，任一非零多元多项式均在某个域内点取非零值：对变量个数归纳，一个变量时由非零多项式的根数不超过次数得到结论；多个变量时，将多项式视为最后一个变量的多项式，先由归纳假设选取其余变量，使一个非零系数保持非零，再由一个变量的结论选取最后一个变量。

因此可取\allowbreak{}\hskip.1em \RB{63.591}{7.930}{2.605}{%
\RG{0.000}{0.001}{10.9091}{mi}{99}\RG{4.597}{-1.635}{7.9701}{rm}{49}\RG{8.915}{0.001}{10.9091}{mi}{59}\RG{13.694}{0.001}{10.9091}{mi}{58}\RG{18.483}{0.001}{10.9091}{mi}{58}\RG{23.272}{0.001}{10.9091}{mi}{58}\RG{28.050}{0.001}{10.9091}{mi}{59}\RG{32.839}{0.001}{10.9091}{mi}{99}\RG{37.439}{-1.635}{7.9701}{mi}{114}\RG{44.676}{0.001}{10.9091}{sy}{50}\RG{55.029}{0.001}{10.9091}{mi}{70}%
}\hskip.1em\allowbreak{} ，使\allowbreak{}\hskip.1em \RB{79.454}{8.683}{3.237}{%
\RG{0.000}{0.001}{10.9091}{mi}{72}\RG{9.694}{0.001}{10.9091}{rm}{40}\RG{13.847}{0.001}{10.9091}{mi}{99}\RG{18.445}{-1.635}{7.9701}{rm}{49}\RG{22.763}{0.001}{10.9091}{mi}{59}\RG{27.541}{0.001}{10.9091}{mi}{58}\RG{32.330}{0.001}{10.9091}{mi}{58}\RG{37.119}{0.001}{10.9091}{mi}{58}\RG{41.898}{0.001}{10.9091}{mi}{59}\RG{46.687}{0.001}{10.9091}{mi}{99}\RG{51.287}{-1.635}{7.9701}{mi}{114}\RG{55.443}{0.001}{10.9091}{rm}{41}\RG{62.678}{0.001}{10.9091}{sy}{54}\RG{62.678}{0.001}{10.9091}{rm}{61}\RG{74.067}{0.001}{10.9091}{rm}{48}%
}\hskip.1em\allowbreak{} 。令\allowbreak{}\hskip.1em \RB{64.676}{8.181}{3.269}{%
\RG{0.000}{0.001}{10.9091}{mi}{81}\RG{11.541}{0.001}{10.9091}{rm}{61}\RG{22.929}{7.274}{9.0910}{ex}{80}\RG{32.525}{4.395}{7.9701}{mi}{114}\RG{32.525}{-2.696}{7.9701}{mi}{105}\RG{35.155}{-2.696}{7.9701}{rm}{61}\RG{41.163}{-2.696}{7.9701}{rm}{49}\RG{47.302}{0.001}{10.9091}{mi}{99}\RG{51.899}{-1.635}{7.9701}{mi}{105}\RG{54.985}{0.001}{10.9091}{mi}{80}\RG{61.870}{-1.635}{7.9701}{mi}{105}%
}\hskip.1em\allowbreak{} ，则\allowbreak{}\hskip.1em \RB{53.499}{8.683}{3.237}{%
\RG{0.000}{0.001}{10.9091}{mi}{81}\RG{11.541}{0.001}{10.9091}{sy}{50}\RG{21.893}{0.001}{10.9091}{mi}{77}\RG{32.227}{-1.635}{7.9701}{mi}{110}\RG{37.371}{0.001}{10.9091}{rm}{40}\RG{41.525}{0.001}{10.9091}{mi}{70}\RG{49.923}{0.001}{10.9091}{rm}{41}%
}\hskip.1em\allowbreak{} ，并且

\[
\RB{169.735}{8.679}{35.369}{%
\RG{0.000}{-0.003}{10.9091}{rm}{100}\RG{5.935}{-0.003}{10.9091}{rm}{101}\RG{10.680}{-0.003}{10.9091}{rm}{116}\RG{16.654}{-0.003}{10.9091}{mi}{81}\RG{28.145}{-0.003}{10.9091}{rm}{61}\RG{39.483}{-0.003}{10.9091}{mi}{72}\RG{49.176}{-0.003}{10.9091}{rm}{40}\RG{53.330}{-0.003}{10.9091}{mi}{99}\RG{62.245}{-0.003}{10.9091}{mi}{59}\RG{67.023}{-0.003}{10.9091}{mi}{58}\RG{71.813}{-0.003}{10.9091}{mi}{58}\RG{76.602}{-0.003}{10.9091}{mi}{58}\RG{81.380}{-0.003}{10.9091}{mi}{59}\RG{86.169}{-0.003}{10.9091}{mi}{99}\RG{94.926}{-0.003}{10.9091}{rm}{41}\RG{102.110}{-0.003}{10.9091}{sy}{54}\RG{102.110}{-0.003}{10.9091}{rm}{61}\RG{113.445}{-0.003}{10.9091}{rm}{48}\RG{118.790}{-0.003}{10.9091}{mi}{59}\RG{57.927}{-1.639}{7.9701}{rm}{49}\RG{90.769}{-1.639}{7.9701}{mi}{114}\RG{44.198}{-12.108}{7.9701}{mi}{114}\RG{99.177}{-12.108}{7.9701}{mi}{114}\RG{39.483}{-14.836}{9.0910}{ex}{88}\RG{94.462}{-14.836}{9.0910}{ex}{88}\RG{8.646}{-23.927}{10.9091}{mi}{65}\RG{16.653}{-23.927}{10.9091}{mi}{81}\RG{28.145}{-23.927}{10.9091}{rm}{61}\RG{54.432}{-23.927}{10.9091}{mi}{99}\RG{62.116}{-23.927}{10.9091}{mi}{65}\RG{70.123}{-23.927}{10.9091}{mi}{80}\RG{83.123}{-23.927}{10.9091}{rm}{61}\RG{109.412}{-23.927}{10.9091}{mi}{99}\RG{117.095}{-23.927}{10.9091}{mi}{80}\RG{127.066}{-23.927}{10.9091}{mi}{66}\RG{138.761}{-23.927}{10.9091}{rm}{61}\RG{150.100}{-23.927}{10.9091}{mi}{81}\RG{158.566}{-23.927}{10.9091}{mi}{66}\RG{167.228}{-23.927}{10.9091}{mi}{58}\RG{59.030}{-25.563}{7.9701}{mi}{105}\RG{77.008}{-25.563}{7.9701}{mi}{105}\RG{114.009}{-25.563}{7.9701}{mi}{105}\RG{123.979}{-25.563}{7.9701}{mi}{105}\RG{39.796}{-34.796}{7.9701}{mi}{105}\RG{42.427}{-34.796}{7.9701}{rm}{61}\RG{48.435}{-34.796}{7.9701}{rm}{49}\RG{94.775}{-34.796}{7.9701}{mi}{105}\RG{97.407}{-34.796}{7.9701}{rm}{61}\RG{103.414}{-34.796}{7.9701}{rm}{49}%
}
\]
于是\allowbreak{}\hskip.1em \RB{58.232}{9.291}{2.620}{%
\RG{0.000}{-0.004}{10.9091}{mi}{81}\RG{8.461}{3.955}{7.9701}{sy}{0}\RG{14.471}{3.955}{7.9701}{rm}{49}\RG{18.789}{-0.004}{10.9091}{mi}{65}\RG{26.796}{-0.004}{10.9091}{mi}{81}\RG{38.287}{-0.004}{10.9091}{rm}{61}\RG{49.626}{-0.004}{10.9091}{mi}{66}%
}\hskip.1em\allowbreak{} ，即\allowbreak{}\hskip.1em \RB{21.391}{8.286}{2.609}{%
\RG{0.000}{-0.004}{10.9091}{mi}{65}\RG{8.007}{-0.004}{10.9091}{mi}{59}\RG{12.786}{-0.004}{10.9091}{mi}{66}%
}\hskip.1em\allowbreak{} 在\allowbreak{}\hskip.1em \RB{8.562}{7.926}{0.504}{%
\RG{0.000}{-0.004}{10.9091}{mi}{70}%
}\hskip.1em\allowbreak{} 上相似。

% Original page 12.
再用多项式理论证明逆命题。对\allowbreak{}\hskip.1em \RB{53.149}{8.685}{3.225}{%
\RG{0.000}{0.003}{10.9091}{mi}{84}\RG{10.774}{0.003}{10.9091}{sy}{50}\RG{21.072}{0.003}{10.9091}{sy}{102}\RG{26.532}{0.003}{10.9091}{mi}{65}\RG{34.540}{0.003}{10.9091}{mi}{59}\RG{39.318}{0.003}{10.9091}{mi}{66}\RG{47.991}{0.003}{10.9091}{sy}{103}%
}\hskip.1em\allowbreak{} ，记

\[
\RB{84.460}{8.680}{3.240}{%
\RG{0.000}{-0.002}{10.9091}{mi}{77}\RG{16.360}{-0.002}{10.9091}{rm}{40}\RG{20.514}{-0.002}{10.9091}{mi}{21}\RG{26.746}{-0.002}{10.9091}{rm}{41}\RG{33.924}{-0.002}{10.9091}{rm}{61}\RG{45.268}{-0.002}{10.9091}{mi}{21}\RG{51.497}{-0.002}{10.9091}{mi}{73}\RG{63.778}{-0.002}{10.9091}{sy}{0}\RG{74.687}{-0.002}{10.9091}{mi}{84}\RG{81.833}{-0.002}{10.9091}{mi}{59}\RG{10.333}{-1.638}{7.9701}{mi}{84}\RG{56.211}{-1.638}{7.9701}{mi}{110}%
}
\]
并令\allowbreak{}\hskip.1em \RB{31.473}{8.686}{3.536}{%
\RG{0.000}{0.004}{10.9091}{mi}{100}\RG{5.550}{-1.632}{7.9701}{mi}{84}\RG{10.692}{-1.632}{7.9701}{mi}{59}\RG{12.837}{-1.632}{7.9701}{mi}{107}\RG{17.511}{0.004}{10.9091}{rm}{40}\RG{21.666}{0.004}{10.9091}{mi}{21}\RG{27.896}{0.004}{10.9091}{rm}{41}%
}\hskip.1em\allowbreak{} 为\allowbreak{}\hskip.1em \RB{30.322}{8.686}{3.235}{%
\RG{0.000}{0.004}{10.9091}{mi}{77}\RG{10.334}{-1.632}{7.9701}{mi}{84}\RG{16.360}{0.004}{10.9091}{rm}{40}\RG{20.514}{0.004}{10.9091}{mi}{21}\RG{26.746}{0.004}{10.9091}{rm}{41}%
}\hskip.1em\allowbreak{} 的全部\allowbreak{}\hskip.1em \RB{5.966}{8.075}{0.605}{%
\RG{0.000}{0.004}{10.9091}{mi}{107}%
}\hskip.1em\allowbreak{} 阶子式在\allowbreak{}\hskip.1em \RB{19.851}{8.686}{3.224}{%
\RG{0.000}{0.004}{10.9091}{mi}{70}\RG{8.399}{0.004}{10.9091}{rm}{91}\RG{11.429}{0.004}{10.9091}{mi}{21}\RG{17.660}{0.004}{10.9091}{rm}{93}%
}\hskip.1em\allowbreak{} 中的首一最大公因式，其中\allowbreak{}\hskip.1em \RB{46.909}{8.075}{1.991}{%
\RG{0.000}{0.004}{10.9091}{rm}{49}\RG{8.371}{0.004}{10.9091}{sy}{20}\RG{19.886}{0.004}{10.9091}{mi}{107}\RG{28.839}{0.004}{10.9091}{sy}{20}\RG{40.354}{0.004}{10.9091}{mi}{110}%
}\hskip.1em\allowbreak{} ；另置\allowbreak{}\hskip.1em \RB{36.496}{8.066}{3.545}{%
\RG{0.000}{-0.005}{10.9091}{mi}{100}\RG{5.551}{-1.642}{7.9701}{mi}{84}\RG{10.693}{-1.642}{7.9701}{mi}{59}\RG{12.840}{-1.642}{7.9701}{rm}{48}\RG{20.188}{-0.005}{10.9091}{rm}{61}\RG{31.523}{-0.005}{10.9091}{rm}{49}%
}\hskip.1em\allowbreak{} 。每个\allowbreak{}\hskip.1em \RB{5.966}{8.066}{0.614}{%
\RG{0.000}{-0.005}{10.9091}{mi}{107}%
}\hskip.1em\allowbreak{} 阶主子式的最高次项为\allowbreak{}\hskip.1em \RB{10.608}{9.501}{0.625}{%
\RG{0.000}{-0.005}{10.9091}{mi}{21}\RG{6.232}{3.954}{7.9701}{mi}{107}%
}\hskip.1em\allowbreak{} ，故这些最大公因式均非零。

最大公因式在扩域后保持不变。事实上，设\allowbreak{}\hskip.1em \RB{80.382}{8.680}{3.229}{%
\RG{0.000}{-0.002}{10.9091}{mi}{112}\RG{5.361}{-1.638}{7.9701}{rm}{49}\RG{9.680}{-0.002}{10.9091}{mi}{59}\RG{14.458}{-0.002}{10.9091}{mi}{58}\RG{19.247}{-0.002}{10.9091}{mi}{58}\RG{24.026}{-0.002}{10.9091}{mi}{58}\RG{28.815}{-0.002}{10.9091}{mi}{59}\RG{33.604}{-0.002}{10.9091}{mi}{112}\RG{38.968}{-1.638}{7.9701}{mi}{78}\RG{49.794}{-0.002}{10.9091}{sy}{50}\RG{60.532}{-0.002}{10.9091}{mi}{70}\RG{68.929}{-0.002}{10.9091}{rm}{91}\RG{71.960}{-0.002}{10.9091}{mi}{21}\RG{78.191}{-0.002}{10.9091}{rm}{93}%
}\hskip.1em\allowbreak{} 不全为零，\allowbreak{}\hskip.1em \RB{6.020}{8.069}{0.611}{%
\RG{0.000}{-0.002}{10.9091}{mi}{100}%
}\hskip.1em\allowbreak{} 为它们在\allowbreak{}\hskip.1em \RB{19.851}{8.681}{3.228}{%
\RG{0.000}{-0.001}{10.9091}{mi}{70}\RG{8.399}{-0.001}{10.9091}{rm}{91}\RG{11.429}{-0.001}{10.9091}{mi}{21}\RG{17.660}{-0.001}{10.9091}{rm}{93}%
}\hskip.1em\allowbreak{} 中的首一最大公因式。由多项式带余除法及Bézout等式，存在\allowbreak{}\hskip.1em \RB{80.951}{8.681}{3.228}{%
\RG{0.000}{-0.001}{10.9091}{mi}{117}\RG{6.079}{-1.638}{7.9701}{rm}{49}\RG{10.398}{-0.001}{10.9091}{mi}{59}\RG{15.176}{-0.001}{10.9091}{mi}{58}\RG{19.965}{-0.001}{10.9091}{mi}{58}\RG{24.744}{-0.001}{10.9091}{mi}{58}\RG{29.533}{-0.001}{10.9091}{mi}{59}\RG{34.322}{-0.001}{10.9091}{mi}{117}\RG{40.404}{-1.638}{7.9701}{mi}{78}\RG{50.797}{-0.001}{10.9091}{sy}{50}\RG{61.100}{-0.001}{10.9091}{mi}{70}\RG{69.499}{-0.001}{10.9091}{rm}{91}\RG{72.529}{-0.001}{10.9091}{mi}{21}\RG{78.760}{-0.001}{10.9091}{rm}{93}%
}\hskip.1em\allowbreak{} ，使得

\[
\RB{176.082}{17.288}{12.850}{%
\RG{142.155}{11.821}{7.9701}{mi}{78}\RG{139.043}{9.093}{9.0910}{ex}{88}\RG{0.000}{0.002}{10.9091}{mi}{100}\RG{8.580}{0.002}{10.9091}{sy}{106}\RG{14.641}{0.002}{10.9091}{mi}{112}\RG{34.911}{0.002}{10.9091}{rm}{40}\RG{39.067}{0.002}{10.9091}{rm}{49}\RG{47.436}{0.002}{10.9091}{sy}{20}\RG{58.951}{0.002}{10.9091}{mi}{106}\RG{67.014}{0.002}{10.9091}{sy}{20}\RG{78.530}{0.002}{10.9091}{mi}{78}\RG{88.223}{0.002}{10.9091}{rm}{41}\RG{92.377}{0.002}{10.9091}{mi}{59}\RG{118.974}{0.002}{10.9091}{mi}{100}\RG{127.562}{0.002}{10.9091}{rm}{61}\RG{154.136}{0.002}{10.9091}{mi}{117}\RG{164.215}{0.002}{10.9091}{mi}{112}\RG{173.574}{0.002}{10.9091}{mi}{58}\RG{20.002}{-1.634}{7.9701}{mi}{106}\RG{160.216}{-1.634}{7.9701}{mi}{106}\RG{169.576}{-1.634}{7.9701}{mi}{106}\RG{138.900}{-10.866}{7.9701}{mi}{106}\RG{142.444}{-10.866}{7.9701}{rm}{61}\RG{148.451}{-10.866}{7.9701}{rm}{49}%
}
\]
第一组整除关系在\allowbreak{}\hskip.1em \RB{21.324}{8.687}{3.223}{%
\RG{0.000}{0.005}{10.9091}{mi}{75}\RG{9.872}{0.005}{10.9091}{rm}{91}\RG{12.902}{0.005}{10.9091}{mi}{21}\RG{19.133}{0.005}{10.9091}{rm}{93}%
}\hskip.1em\allowbreak{} 中仍成立；第二个等式表明，\allowbreak{}\hskip.1em \RB{46.129}{5.316}{2.612}{%
\RG{0.000}{0.005}{10.9091}{mi}{112}\RG{5.361}{-1.632}{7.9701}{rm}{49}\RG{9.680}{0.005}{10.9091}{mi}{59}\RG{14.458}{0.005}{10.9091}{mi}{58}\RG{19.247}{0.005}{10.9091}{mi}{58}\RG{24.026}{0.005}{10.9091}{mi}{58}\RG{28.815}{0.005}{10.9091}{mi}{59}\RG{33.604}{0.005}{10.9091}{mi}{112}\RG{38.967}{-1.632}{7.9701}{mi}{78}%
}\hskip.1em\allowbreak{} 在\allowbreak{}\hskip.1em \RB{21.324}{8.687}{3.223}{%
\RG{0.000}{0.005}{10.9091}{mi}{75}\RG{9.872}{0.005}{10.9091}{rm}{91}\RG{12.901}{0.005}{10.9091}{mi}{21}\RG{19.133}{0.005}{10.9091}{rm}{93}%
}\hskip.1em\allowbreak{} 中的任一公因式均整除\allowbreak{}\hskip.1em \RB{6.020}{8.066}{0.614}{%
\RG{0.000}{-0.005}{10.9091}{mi}{100}%
}\hskip.1em\allowbreak{} 。因此\allowbreak{}\hskip.1em \RB{6.020}{8.066}{0.614}{%
\RG{0.000}{-0.005}{10.9091}{mi}{100}%
}\hskip.1em\allowbreak{} 也是它们在\allowbreak{}\hskip.1em \RB{21.324}{8.677}{3.232}{%
\RG{0.000}{-0.005}{10.9091}{mi}{75}\RG{9.871}{-0.005}{10.9091}{rm}{91}\RG{12.901}{-0.005}{10.9091}{mi}{21}\RG{19.133}{-0.005}{10.9091}{rm}{93}%
}\hskip.1em\allowbreak{} 中的首一最大公因式。

由\allowbreak{}\hskip.1em \RB{21.391}{8.288}{2.606}{%
\RG{0.000}{-0.001}{10.9091}{mi}{65}\RG{8.007}{-0.001}{10.9091}{mi}{59}\RG{12.786}{-0.001}{10.9091}{mi}{66}%
}\hskip.1em\allowbreak{} 在\allowbreak{}\hskip.1em \RB{10.035}{7.950}{0.501}{%
\RG{0.000}{-0.001}{10.9091}{mi}{75}%
}\hskip.1em\allowbreak{} 上相似，可取\allowbreak{}\hskip.1em \RB{54.781}{8.681}{3.239}{%
\RG{0.000}{-0.001}{10.9091}{mi}{80}\RG{11.399}{-0.001}{10.9091}{sy}{50}\RG{21.702}{-0.001}{10.9091}{mi}{77}\RG{32.036}{-1.638}{7.9701}{mi}{110}\RG{37.179}{-0.001}{10.9091}{rm}{40}\RG{41.333}{-0.001}{10.9091}{mi}{75}\RG{51.205}{-0.001}{10.9091}{rm}{41}%
}\hskip.1em\allowbreak{} 可逆，使

\[
\RB{105.382}{9.844}{3.235}{%
\RG{54.253}{4.508}{7.9701}{sy}{0}\RG{60.262}{4.508}{7.9701}{rm}{49}\RG{0.000}{0.004}{10.9091}{mi}{77}\RG{16.975}{0.004}{10.9091}{rm}{40}\RG{21.129}{0.004}{10.9091}{mi}{21}\RG{27.361}{0.004}{10.9091}{rm}{41}\RG{34.539}{0.004}{10.9091}{rm}{61}\RG{45.884}{0.004}{10.9091}{mi}{80}\RG{64.580}{0.004}{10.9091}{mi}{77}\RG{81.157}{0.004}{10.9091}{rm}{40}\RG{85.312}{0.004}{10.9091}{mi}{21}\RG{91.542}{0.004}{10.9091}{rm}{41}\RG{95.696}{0.004}{10.9091}{mi}{80}\RG{102.874}{0.004}{10.9091}{mi}{58}\RG{10.334}{-1.632}{7.9701}{mi}{66}\RG{74.915}{-1.632}{7.9701}{mi}{65}%
}
\]
对每个\allowbreak{}\hskip.1em \RB{5.966}{8.071}{0.609}{%
\RG{0.000}{0.000}{10.9091}{mi}{107}%
}\hskip.1em\allowbreak{} ，Cauchy–Binet公式表明，\allowbreak{}\hskip.1em \RB{16.616}{7.951}{2.136}{%
\RG{0.000}{0.000}{10.9091}{mi}{77}\RG{10.334}{-1.636}{7.9701}{mi}{66}%
}\hskip.1em\allowbreak{} 的每个\allowbreak{}\hskip.1em \RB{5.966}{8.071}{0.609}{%
\RG{0.000}{0.000}{10.9091}{mi}{107}%
}\hskip.1em\allowbreak{} 阶子式都是\allowbreak{}\hskip.1em \RB{16.390}{7.951}{2.136}{%
\RG{0.000}{0.000}{10.9091}{mi}{77}\RG{10.333}{-1.636}{7.9701}{mi}{65}%
}\hskip.1em\allowbreak{} 的\allowbreak{}\hskip.1em \RB{5.966}{8.071}{0.609}{%
\RG{0.000}{0.000}{10.9091}{mi}{107}%
}\hskip.1em\allowbreak{} 阶子式的\allowbreak{}\hskip.1em \RB{21.324}{8.682}{3.227}{%
\RG{0.000}{0.000}{10.9091}{mi}{75}\RG{9.872}{0.000}{10.9091}{rm}{91}\RG{12.902}{0.000}{10.9091}{mi}{21}\RG{19.133}{0.000}{10.9091}{rm}{93}%
}\hskip.1em\allowbreak{} -线性组合；用逆变换得到反向包含。因此，两组子式在\allowbreak{}\hskip.1em \RB{21.324}{8.683}{3.226}{%
\RG{0.000}{0.001}{10.9091}{mi}{75}\RG{9.872}{0.001}{10.9091}{rm}{91}\RG{12.902}{0.001}{10.9091}{mi}{21}\RG{19.133}{0.001}{10.9091}{rm}{93}%
}\hskip.1em\allowbreak{} 中生成同一理想，其首一最大公因式相同。结合扩域前后最大公因式不变，得到

\[
\RB{130.445}{8.679}{3.543}{%
\RG{0.000}{-0.003}{10.9091}{mi}{100}\RG{21.187}{-0.003}{10.9091}{rm}{61}\RG{32.526}{-0.003}{10.9091}{mi}{100}\RG{72.900}{-0.003}{10.9091}{rm}{40}\RG{77.056}{-0.003}{10.9091}{rm}{48}\RG{85.425}{-0.003}{10.9091}{sy}{20}\RG{96.941}{-0.003}{10.9091}{mi}{107}\RG{105.893}{-0.003}{10.9091}{sy}{20}\RG{117.408}{-0.003}{10.9091}{mi}{110}\RG{123.784}{-0.003}{10.9091}{rm}{41}\RG{127.938}{-0.003}{10.9091}{mi}{58}\RG{5.550}{-1.639}{7.9701}{mi}{65}\RG{11.339}{-1.639}{7.9701}{mi}{59}\RG{13.485}{-1.639}{7.9701}{mi}{107}\RG{38.076}{-1.639}{7.9701}{mi}{66}\RG{44.258}{-1.639}{7.9701}{mi}{59}\RG{46.404}{-1.639}{7.9701}{mi}{107}%
}
\]
特征矩阵的不变因子分别为\allowbreak{}\hskip.1em \RB{50.961}{8.684}{3.538}{%
\RG{0.000}{0.003}{10.9091}{mi}{100}\RG{5.550}{-1.634}{7.9701}{mi}{65}\RG{11.339}{-1.634}{7.9701}{mi}{59}\RG{13.485}{-1.634}{7.9701}{mi}{107}\RG{18.156}{0.002}{10.9091}{mi}{61}\RG{23.502}{0.002}{10.9091}{mi}{100}\RG{29.048}{-1.634}{7.9701}{mi}{65}\RG{34.837}{-1.634}{7.9701}{mi}{59}\RG{36.983}{-1.634}{7.9701}{mi}{107}\RG{41.200}{-1.634}{7.9701}{sy}{0}\RG{47.210}{-1.634}{7.9701}{rm}{49}%
}\hskip.1em\allowbreak{} 与\allowbreak{}\hskip.1em \RB{51.759}{8.684}{3.538}{%
\RG{0.000}{0.002}{10.9091}{mi}{100}\RG{5.550}{-1.634}{7.9701}{mi}{66}\RG{11.732}{-1.634}{7.9701}{mi}{59}\RG{13.877}{-1.634}{7.9701}{mi}{107}\RG{18.556}{0.002}{10.9091}{mi}{61}\RG{23.901}{0.002}{10.9091}{mi}{100}\RG{29.447}{-1.634}{7.9701}{mi}{66}\RG{35.629}{-1.634}{7.9701}{mi}{59}\RG{37.774}{-1.634}{7.9701}{mi}{107}\RG{41.998}{-1.634}{7.9701}{sy}{0}\RG{48.008}{-1.634}{7.9701}{rm}{49}%
}\hskip.1em\allowbreak{} ，故两组不变因子逐项相同。由有理标准形定理，\allowbreak{}\hskip.1em \RB{21.391}{8.292}{2.603}{%
\RG{0.000}{0.003}{10.9091}{mi}{65}\RG{8.007}{0.003}{10.9091}{mi}{59}\RG{12.786}{0.003}{10.9091}{mi}{66}%
}\hskip.1em\allowbreak{} 在\allowbreak{}\hskip.1em \RB{8.562}{7.932}{0.500}{%
\RG{0.000}{0.003}{10.9091}{mi}{70}%
}\hskip.1em\allowbreak{} 上相似。\allowbreak{}\hskip.1em \RB{8.366}{7.997}{0.500}{%
\RG{0.000}{0.003}{10.9091}{am}{3}%
}\hskip.1em\allowbreak{}

% Original page 13.
\section*{2016S09}
设\allowbreak{}\hskip.1em \RB{8.715}{7.949}{0.731}{%
\RG{0.000}{-0.002}{10.9091}{mi}{86}%
}\hskip.1em\allowbreak{} 是数域\allowbreak{}\hskip.1em \RB{10.035}{7.949}{0.502}{%
\RG{0.000}{-0.002}{10.9091}{mi}{75}%
}\hskip.1em\allowbreak{} 上的\allowbreak{}\hskip.1em \RB{6.555}{5.309}{0.611}{%
\RG{0.000}{-0.002}{10.9091}{mi}{110}%
}\hskip.1em\allowbreak{} 维线性空间，\allowbreak{}\hskip.1em \RB{7.122}{5.309}{2.858}{%
\RG{0.000}{-0.002}{10.9091}{mi}{39}%
}\hskip.1em\allowbreak{} 是\allowbreak{}\hskip.1em \RB{8.715}{7.949}{0.731}{%
\RG{0.000}{-0.002}{10.9091}{mi}{86}%
}\hskip.1em\allowbreak{} 上的线性变换。设\allowbreak{}\hskip.1em \RB{52.553}{8.309}{2.858}{%
\RG{0.000}{-0.002}{10.9091}{rm}{48}\RG{9.681}{-0.002}{10.9091}{sy}{54}\RG{9.681}{-0.002}{10.9091}{rm}{61}\RG{22.329}{-0.002}{10.9091}{mi}{118}\RG{32.225}{-0.002}{10.9091}{sy}{50}\RG{43.838}{-0.002}{10.9091}{mi}{86}%
}\hskip.1em\allowbreak{} 。若多项式\allowbreak{}\hskip.1em \RB{54.702}{8.681}{3.239}{%
\RG{0.000}{-0.001}{10.9091}{mi}{103}\RG{5.506}{-0.001}{10.9091}{rm}{40}\RG{9.661}{-0.001}{10.9091}{mi}{21}\RG{15.891}{-0.001}{10.9091}{rm}{41}\RG{23.076}{-0.001}{10.9091}{sy}{50}\RG{33.379}{-0.001}{10.9091}{mi}{75}\RG{43.250}{-0.001}{10.9091}{rm}{91}\RG{46.281}{-0.001}{10.9091}{mi}{21}\RG{52.511}{-0.001}{10.9091}{rm}{93}%
}\hskip.1em\allowbreak{} 满足\allowbreak{}\hskip.1em \RB{54.432}{8.681}{3.239}{%
\RG{0.000}{-0.001}{10.9091}{mi}{103}\RG{5.506}{-0.001}{10.9091}{rm}{40}\RG{9.660}{-0.001}{10.9091}{mi}{39}\RG{16.658}{-0.001}{10.9091}{rm}{41}\RG{20.814}{-0.001}{10.9091}{rm}{40}\RG{24.965}{-0.001}{10.9091}{mi}{118}\RG{30.521}{-0.001}{10.9091}{rm}{41}\RG{37.699}{-0.001}{10.9091}{rm}{61}\RG{49.045}{-0.001}{10.9091}{rm}{48}%
}\hskip.1em\allowbreak{} ，则称\allowbreak{}\hskip.1em \RB{19.468}{8.681}{3.239}{%
\RG{0.000}{-0.001}{10.9091}{mi}{103}\RG{5.507}{-0.001}{10.9091}{rm}{40}\RG{9.661}{-0.001}{10.9091}{mi}{21}\RG{15.891}{-0.001}{10.9091}{rm}{41}%
}\hskip.1em\allowbreak{} 为\allowbreak{}\hskip.1em \RB{7.122}{5.310}{2.857}{%
\RG{0.000}{-0.001}{10.9091}{mi}{39}%
}\hskip.1em\allowbreak{} 在\allowbreak{}\hskip.1em \RB{5.486}{5.310}{0.610}{%
\RG{0.000}{-0.001}{10.9091}{mi}{118}%
}\hskip.1em\allowbreak{} 处的零化多项式。若首一多项式\allowbreak{}\hskip.1em \RB{27.783}{8.681}{3.239}{%
\RG{0.000}{-0.001}{10.9091}{mi}{109}\RG{9.343}{-1.637}{7.9701}{mi}{118}\RG{13.822}{-0.001}{10.9091}{rm}{40}\RG{17.976}{-0.001}{10.9091}{mi}{21}\RG{24.206}{-0.001}{10.9091}{rm}{41}%
}\hskip.1em\allowbreak{}是其中所有非零零化多项式中次数最小者，则称其为\allowbreak{}\hskip.1em \RB{7.122}{5.311}{2.856}{%
\RG{0.000}{-0.000}{10.9091}{mi}{39}%
}\hskip.1em\allowbreak{} 在\allowbreak{}\hskip.1em \RB{5.486}{5.311}{0.609}{%
\RG{0.000}{-0.000}{10.9091}{mi}{118}%
}\hskip.1em\allowbreak{} 处的极小多项式。（1）证明：对任意\allowbreak{}\hskip.1em \RB{47.314}{8.314}{2.854}{%
\RG{0.000}{0.003}{10.9091}{rm}{48}\RG{8.371}{0.003}{10.9091}{sy}{54}\RG{8.371}{0.003}{10.9091}{rm}{61}\RG{19.709}{0.003}{10.9091}{mi}{118}\RG{28.296}{0.003}{10.9091}{sy}{50}\RG{38.599}{0.003}{10.9091}{mi}{86}%
}\hskip.1em\allowbreak{} ，\allowbreak{}\hskip.1em \RB{27.784}{8.685}{3.235}{%
\RG{0.000}{0.003}{10.9091}{mi}{109}\RG{9.343}{-1.633}{7.9701}{mi}{118}\RG{13.822}{0.003}{10.9091}{rm}{40}\RG{17.976}{0.003}{10.9091}{mi}{21}\RG{24.207}{0.003}{10.9091}{rm}{41}%
}\hskip.1em\allowbreak{} 存在且唯一。先证明其整除任意零化多项式的基本性质。（2）设\allowbreak{}\hskip.1em \RB{47.313}{8.308}{2.859}{%
\RG{0.000}{-0.003}{10.9091}{rm}{48}\RG{8.371}{-0.003}{10.9091}{sy}{54}\RG{8.371}{-0.003}{10.9091}{rm}{61}\RG{19.709}{-0.003}{10.9091}{mi}{118}\RG{28.295}{-0.003}{10.9091}{sy}{50}\RG{38.598}{-0.003}{10.9091}{mi}{86}%
}\hskip.1em\allowbreak{} ，记

\[
\RB{169.309}{9.839}{3.240}{%
\RG{125.842}{4.502}{7.9701}{rm}{50}\RG{0.000}{-0.002}{10.9091}{mi}{67}\RG{8.426}{-0.002}{10.9091}{rm}{40}\RG{12.580}{-0.002}{10.9091}{mi}{39}\RG{19.573}{-0.002}{10.9091}{mi}{59}\RG{24.351}{-0.002}{10.9091}{mi}{118}\RG{29.919}{-0.002}{10.9091}{rm}{41}\RG{37.097}{-0.002}{10.9091}{rm}{61}\RG{48.443}{-0.002}{10.9091}{rm}{115}\RG{52.653}{-0.002}{10.9091}{rm}{112}\RG{58.588}{-0.002}{10.9091}{rm}{97}\RG{63.934}{-0.002}{10.9091}{rm}{110}\RG{77.402}{-0.002}{10.9091}{sy}{102}\RG{82.857}{-0.002}{10.9091}{mi}{118}\RG{88.399}{-0.002}{10.9091}{mi}{59}\RG{93.188}{-0.002}{10.9091}{mi}{39}\RG{100.196}{-0.002}{10.9091}{rm}{40}\RG{104.349}{-0.002}{10.9091}{mi}{118}\RG{109.905}{-0.002}{10.9091}{rm}{41}\RG{114.059}{-0.002}{10.9091}{mi}{59}\RG{118.837}{-0.002}{10.9091}{mi}{39}\RG{130.160}{-0.002}{10.9091}{rm}{40}\RG{134.314}{-0.002}{10.9091}{mi}{118}\RG{139.870}{-0.002}{10.9091}{rm}{41}\RG{144.024}{-0.002}{10.9091}{mi}{59}\RG{148.802}{-0.002}{10.9091}{mi}{58}\RG{153.591}{-0.002}{10.9091}{mi}{58}\RG{158.369}{-0.002}{10.9091}{mi}{58}\RG{161.347}{-0.002}{10.9091}{sy}{103}\RG{166.802}{-0.002}{10.9091}{mi}{58}\RG{69.918}{-2.578}{7.9701}{mi}{75}%
}
\]
称其为由\allowbreak{}\hskip.1em \RB{5.486}{5.315}{0.605}{%
\RG{0.000}{0.004}{10.9091}{mi}{118}%
}\hskip.1em\allowbreak{} 生成的循环子空间，\allowbreak{}\hskip.1em \RB{5.486}{5.315}{0.605}{%
\RG{0.000}{0.004}{10.9091}{mi}{118}%
}\hskip.1em\allowbreak{} 称为其循环向量。它是包含\allowbreak{}\hskip.1em \RB{5.486}{5.315}{0.605}{%
\RG{0.000}{0.004}{10.9091}{mi}{118}%
}\hskip.1em\allowbreak{} 的最小的\allowbreak{}\hskip.1em \RB{7.122}{5.315}{2.853}{%
\RG{0.000}{0.004}{10.9091}{mi}{39}%
}\hskip.1em\allowbreak{} –不变子空间。设\allowbreak{}\hskip.1em \RB{74.081}{8.687}{3.234}{%
\RG{0.000}{0.005}{10.9091}{rm}{100}\RG{5.935}{0.005}{10.9091}{rm}{105}\RG{8.902}{0.005}{10.9091}{rm}{109}\RG{19.675}{0.005}{10.9091}{mi}{67}\RG{28.100}{0.005}{10.9091}{rm}{40}\RG{32.255}{0.005}{10.9091}{mi}{39}\RG{39.247}{0.005}{10.9091}{mi}{59}\RG{44.026}{0.005}{10.9091}{mi}{118}\RG{49.592}{0.005}{10.9091}{rm}{41}\RG{56.771}{0.005}{10.9091}{rm}{61}\RG{68.115}{0.005}{10.9091}{mi}{107}%
}\hskip.1em\allowbreak{} ，证明\allowbreak{}\hskip.1em \RB{101.149}{9.510}{3.234}{%
\RG{0.000}{0.005}{10.9091}{sy}{102}\RG{5.455}{0.005}{10.9091}{mi}{118}\RG{10.997}{0.005}{10.9091}{mi}{59}\RG{15.786}{0.005}{10.9091}{mi}{39}\RG{22.794}{0.005}{10.9091}{rm}{40}\RG{26.947}{0.005}{10.9091}{mi}{118}\RG{32.503}{0.005}{10.9091}{rm}{41}\RG{36.657}{0.005}{10.9091}{mi}{59}\RG{41.435}{0.005}{10.9091}{mi}{58}\RG{46.224}{0.005}{10.9091}{mi}{58}\RG{51.013}{0.005}{10.9091}{mi}{58}\RG{55.791}{0.005}{10.9091}{mi}{59}\RG{60.581}{0.005}{10.9091}{mi}{39}\RG{67.582}{3.963}{7.9701}{mi}{107}\RG{71.798}{3.963}{7.9701}{sy}{0}\RG{77.809}{3.963}{7.9701}{rm}{49}\RG{82.127}{0.005}{10.9091}{rm}{40}\RG{86.281}{0.005}{10.9091}{mi}{118}\RG{91.836}{0.005}{10.9091}{rm}{41}\RG{95.991}{0.005}{10.9091}{sy}{103}%
}\hskip.1em\allowbreak{} 是其一组基。若

\[
\RB{205.570}{10.048}{3.242}{%
\RG{6.998}{4.501}{7.9701}{mi}{107}\RG{174.534}{4.501}{7.9701}{mi}{107}\RG{178.750}{4.501}{7.9701}{sy}{0}\RG{184.761}{4.501}{7.9701}{rm}{49}\RG{0.000}{-0.004}{10.9091}{mi}{39}\RG{11.669}{-0.004}{10.9091}{rm}{40}\RG{15.824}{-0.004}{10.9091}{mi}{118}\RG{21.379}{-0.004}{10.9091}{rm}{41}\RG{28.558}{-0.004}{10.9091}{rm}{61}\RG{39.902}{-0.004}{10.9091}{sy}{0}\RG{48.387}{-0.004}{10.9091}{mi}{97}\RG{58.312}{-0.004}{10.9091}{mi}{118}\RG{66.292}{-0.004}{10.9091}{sy}{0}\RG{77.201}{-0.004}{10.9091}{mi}{97}\RG{87.128}{-0.004}{10.9091}{mi}{39}\RG{94.125}{-0.004}{10.9091}{rm}{40}\RG{98.279}{-0.004}{10.9091}{mi}{118}\RG{103.834}{-0.004}{10.9091}{rm}{41}\RG{110.412}{-0.004}{10.9091}{sy}{0}\RG{121.322}{-0.004}{10.9091}{sy}{1}\RG{126.165}{-0.004}{10.9091}{sy}{1}\RG{131.020}{-0.004}{10.9091}{sy}{1}\RG{136.474}{-0.004}{10.9091}{sy}{0}\RG{147.383}{-0.004}{10.9091}{mi}{97}\RG{167.536}{-0.004}{10.9091}{mi}{39}\RG{189.079}{-0.004}{10.9091}{rm}{40}\RG{193.233}{-0.004}{10.9091}{mi}{118}\RG{198.789}{-0.004}{10.9091}{rm}{41}\RG{202.943}{-0.004}{10.9091}{mi}{59}\RG{53.994}{-1.640}{7.9701}{rm}{48}\RG{82.809}{-1.640}{7.9701}{rm}{49}\RG{152.990}{-1.640}{7.9701}{mi}{107}\RG{157.208}{-1.640}{7.9701}{sy}{0}\RG{163.218}{-1.640}{7.9701}{rm}{49}%
}
\]
证明：\allowbreak{}\hskip.1em \RB{185.737}{9.504}{3.240}{%
\RG{0.000}{-0.002}{10.9091}{mi}{109}\RG{9.343}{-1.638}{7.9701}{mi}{118}\RG{13.821}{-0.002}{10.9091}{rm}{40}\RG{17.976}{-0.002}{10.9091}{mi}{21}\RG{24.206}{-0.002}{10.9091}{rm}{41}\RG{31.385}{-0.002}{10.9091}{rm}{61}\RG{42.729}{-0.002}{10.9091}{mi}{21}\RG{48.960}{3.957}{7.9701}{mi}{107}\RG{56.057}{-0.002}{10.9091}{rm}{43}\RG{66.789}{-0.002}{10.9091}{mi}{97}\RG{72.396}{-1.638}{7.9701}{mi}{107}\RG{76.614}{-1.638}{7.9701}{sy}{0}\RG{82.623}{-1.638}{7.9701}{rm}{49}\RG{86.942}{-0.002}{10.9091}{mi}{21}\RG{93.173}{3.957}{7.9701}{mi}{107}\RG{97.391}{3.957}{7.9701}{sy}{0}\RG{103.400}{3.957}{7.9701}{rm}{49}\RG{110.143}{-0.002}{10.9091}{rm}{43}\RG{120.876}{-0.002}{10.9091}{sy}{1}\RG{125.720}{-0.002}{10.9091}{sy}{1}\RG{130.574}{-0.002}{10.9091}{sy}{1}\RG{136.027}{-0.002}{10.9091}{rm}{43}\RG{146.760}{-0.002}{10.9091}{mi}{97}\RG{152.368}{-1.638}{7.9701}{rm}{49}\RG{156.686}{-0.002}{10.9091}{mi}{21}\RG{165.341}{-0.002}{10.9091}{rm}{43}\RG{176.074}{-0.002}{10.9091}{mi}{97}\RG{181.680}{-1.638}{7.9701}{rm}{48}%
}\hskip.1em\allowbreak{} 。（3）沿用（2）的记号和假设，证明下列结论。（3.1）\allowbreak{}\hskip.1em \RB{33.494}{8.686}{3.234}{%
\RG{0.000}{0.005}{10.9091}{mi}{67}\RG{8.425}{0.005}{10.9091}{rm}{40}\RG{12.580}{0.005}{10.9091}{mi}{39}\RG{19.573}{0.005}{10.9091}{mi}{59}\RG{24.351}{0.005}{10.9091}{mi}{118}\RG{29.918}{0.005}{10.9091}{rm}{41}%
}\hskip.1em\allowbreak{} 中任一向量均可写成\allowbreak{}\hskip.1em \RB{34.098}{8.686}{3.234}{%
\RG{0.000}{0.005}{10.9091}{mi}{103}\RG{5.506}{0.005}{10.9091}{rm}{40}\RG{9.661}{0.005}{10.9091}{mi}{39}\RG{16.658}{0.005}{10.9091}{rm}{41}\RG{20.814}{0.005}{10.9091}{rm}{40}\RG{24.966}{0.005}{10.9091}{mi}{118}\RG{30.522}{0.005}{10.9091}{rm}{41}%
}\hskip.1em\allowbreak{} ，其中\allowbreak{}\hskip.1em \RB{40.163}{8.686}{3.223}{%
\RG{0.000}{0.005}{10.9091}{mi}{103}\RG{8.537}{0.005}{10.9091}{sy}{50}\RG{18.840}{0.005}{10.9091}{mi}{75}\RG{28.710}{0.005}{10.9091}{rm}{91}\RG{31.741}{0.005}{10.9091}{mi}{21}\RG{37.972}{0.005}{10.9091}{rm}{93}%
}\hskip.1em\allowbreak{} ，且\allowbreak{}\hskip.1em \RB{25.259}{7.759}{2.721}{%
\RG{0.000}{0.005}{10.9091}{mi}{103}\RG{8.537}{0.005}{10.9091}{rm}{61}\RG{19.871}{0.005}{10.9091}{rm}{48}%
}\hskip.1em\allowbreak{} 或\allowbreak{}\hskip.1em \RB{43.828}{8.076}{2.732}{%
\RG{0.000}{0.005}{10.9091}{rm}{100}\RG{5.935}{0.005}{10.9091}{rm}{101}\RG{10.680}{0.005}{10.9091}{rm}{103}\RG{17.986}{0.005}{10.9091}{mi}{103}\RG{26.517}{0.005}{10.9091}{mi}{60}\RG{37.862}{0.005}{10.9091}{mi}{107}%
}\hskip.1em\allowbreak{} 。（3.2）\allowbreak{}\hskip.1em \RB{34.097}{8.680}{3.240}{%
\RG{0.000}{-0.002}{10.9091}{mi}{103}\RG{5.506}{-0.002}{10.9091}{rm}{40}\RG{9.660}{-0.002}{10.9091}{mi}{39}\RG{16.658}{-0.002}{10.9091}{rm}{41}\RG{20.814}{-0.002}{10.9091}{rm}{40}\RG{24.965}{-0.002}{10.9091}{mi}{118}\RG{30.521}{-0.002}{10.9091}{rm}{41}%
}\hskip.1em\allowbreak{} 是\allowbreak{}\hskip.1em \RB{33.495}{8.680}{3.240}{%
\RG{0.000}{-0.002}{10.9091}{mi}{67}\RG{8.426}{-0.002}{10.9091}{rm}{40}\RG{12.580}{-0.002}{10.9091}{mi}{39}\RG{19.573}{-0.002}{10.9091}{mi}{59}\RG{24.351}{-0.002}{10.9091}{mi}{118}\RG{29.919}{-0.002}{10.9091}{rm}{41}%
}\hskip.1em\allowbreak{} 的循环向量，当且仅当\allowbreak{}\hskip.1em \RB{51.764}{8.680}{3.240}{%
\RG{0.000}{-0.002}{10.9091}{rm}{40}\RG{4.154}{-0.002}{10.9091}{mi}{103}\RG{9.652}{-0.002}{10.9091}{mi}{59}\RG{14.441}{-0.002}{10.9091}{mi}{109}\RG{23.789}{-1.638}{7.9701}{mi}{118}\RG{28.267}{-0.002}{10.9091}{rm}{41}\RG{35.445}{-0.002}{10.9091}{rm}{61}\RG{46.791}{-0.002}{10.9091}{rm}{49}%
}\hskip.1em\allowbreak{} 。（3.3）对\allowbreak{}\hskip.1em \RB{13.458}{5.312}{2.208}{%
\RG{0.000}{0.001}{10.9091}{mi}{109}\RG{9.343}{-1.635}{7.9701}{mi}{118}%
}\hskip.1em\allowbreak{} 的任一非常数首一因式\allowbreak{}\hskip.1em \RB{6.336}{8.072}{0.608}{%
\RG{0.000}{0.001}{10.9091}{mi}{104}%
}\hskip.1em\allowbreak{} ，存在\allowbreak{}\hskip.1em \RB{74.467}{8.683}{3.237}{%
\RG{0.000}{0.001}{10.9091}{rm}{48}\RG{8.371}{0.001}{10.9091}{sy}{54}\RG{8.371}{0.001}{10.9091}{rm}{61}\RG{19.709}{0.001}{10.9091}{mi}{119}\RG{30.670}{0.001}{10.9091}{sy}{50}\RG{40.973}{0.001}{10.9091}{mi}{67}\RG{49.398}{0.001}{10.9091}{rm}{40}\RG{53.552}{0.001}{10.9091}{mi}{39}\RG{60.545}{0.001}{10.9091}{mi}{59}\RG{65.323}{0.001}{10.9091}{mi}{118}\RG{70.890}{0.001}{10.9091}{rm}{41}%
}\hskip.1em\allowbreak{} ，使得\allowbreak{}\hskip.1em \RB{36.244}{8.072}{2.208}{%
\RG{0.000}{0.001}{10.9091}{mi}{109}\RG{9.343}{-1.635}{7.9701}{mi}{119}\RG{18.570}{0.001}{10.9091}{rm}{61}\RG{29.908}{0.001}{10.9091}{mi}{104}%
}\hskip.1em\allowbreak{} 。（3.4）\allowbreak{}\hskip.1em \RB{33.494}{8.686}{3.235}{%
\RG{0.000}{0.004}{10.9091}{mi}{67}\RG{8.425}{0.004}{10.9091}{rm}{40}\RG{12.580}{0.004}{10.9091}{mi}{39}\RG{19.573}{0.004}{10.9091}{mi}{59}\RG{24.351}{0.004}{10.9091}{mi}{118}\RG{29.918}{0.004}{10.9091}{rm}{41}%
}\hskip.1em\allowbreak{} 只有有限个\allowbreak{}\hskip.1em \RB{7.122}{5.315}{2.853}{%
\RG{0.000}{0.004}{10.9091}{mi}{39}%
}\hskip.1em\allowbreak{} –不变子空间，它们恰为

\allowbreak{}\hskip.1em \RB{83.875}{11.406}{5.949}{%
\RG{0.000}{10.542}{9.0910}{ex}{110}\RG{6.061}{-0.004}{10.9091}{mi}{67}\RG{14.486}{-0.004}{10.9091}{rm}{40}\RG{18.641}{-0.004}{10.9091}{mi}{39}\RG{25.634}{-0.004}{10.9091}{mi}{59}\RG{30.412}{-0.004}{10.9091}{mi}{103}\RG{35.930}{-0.004}{10.9091}{rm}{40}\RG{40.084}{-0.004}{10.9091}{mi}{39}\RG{47.081}{-0.004}{10.9091}{rm}{41}\RG{51.238}{-0.004}{10.9091}{rm}{40}\RG{55.390}{-0.004}{10.9091}{mi}{118}\RG{60.945}{-0.004}{10.9091}{rm}{41}\RG{65.101}{-0.004}{10.9091}{rm}{41}\RG{72.280}{-0.004}{10.9091}{rm}{58}\RG{78.281}{-0.004}{10.9091}{mi}{103}%
}\hskip.1em\allowbreak{} 是\allowbreak{}\hskip.1em \RB{13.459}{5.307}{2.213}{%
\RG{0.000}{-0.004}{10.9091}{mi}{109}\RG{9.344}{-1.640}{7.9701}{mi}{118}%
}\hskip.1em\allowbreak{} 的首一因式\allowbreak{}\hskip.1em \RB{8.568}{11.406}{5.949}{%
\RG{0.000}{10.542}{9.0910}{ex}{111}\RG{6.061}{-0.004}{10.9091}{mi}{58}%
}\hskip.1em\allowbreak{}

（4）设\allowbreak{}\hskip.1em \RB{58.178}{8.311}{2.856}{%
\RG{0.000}{-0.000}{10.9091}{rm}{48}\RG{8.371}{-0.000}{10.9091}{sy}{54}\RG{8.371}{-0.000}{10.9091}{rm}{61}\RG{19.709}{-0.000}{10.9091}{mi}{117}\RG{25.786}{-0.000}{10.9091}{mi}{59}\RG{30.564}{-0.000}{10.9091}{mi}{118}\RG{39.160}{-0.000}{10.9091}{sy}{50}\RG{49.464}{-0.000}{10.9091}{mi}{86}%
}\hskip.1em\allowbreak{} 的极小多项式分别为\allowbreak{}\hskip.1em \RB{32.516}{5.311}{2.606}{%
\RG{0.000}{-0.000}{10.9091}{mi}{109}\RG{9.344}{-1.636}{7.9701}{mi}{117}\RG{14.273}{-0.000}{10.9091}{mi}{59}\RG{19.051}{-0.000}{10.9091}{mi}{109}\RG{28.401}{-1.636}{7.9701}{mi}{118}%
}\hskip.1em\allowbreak{} ，证明下列结论。（4.1）若\allowbreak{}\hskip.1em \RB{60.530}{8.686}{3.235}{%
\RG{0.000}{0.004}{10.9091}{rm}{40}\RG{4.154}{0.004}{10.9091}{mi}{109}\RG{13.497}{-1.633}{7.9701}{mi}{117}\RG{18.426}{0.004}{10.9091}{mi}{59}\RG{23.205}{0.004}{10.9091}{mi}{109}\RG{32.555}{-1.633}{7.9701}{mi}{118}\RG{37.033}{0.004}{10.9091}{rm}{41}\RG{44.212}{0.004}{10.9091}{rm}{61}\RG{55.557}{0.004}{10.9091}{rm}{49}%
}\hskip.1em\allowbreak{} ，则

\[
\RB{180.823}{8.686}{3.234}{%
\RG{0.000}{0.005}{10.9091}{mi}{67}\RG{8.425}{0.005}{10.9091}{rm}{40}\RG{12.580}{0.005}{10.9091}{mi}{39}\RG{19.573}{0.005}{10.9091}{mi}{59}\RG{24.351}{0.005}{10.9091}{mi}{117}\RG{30.442}{0.005}{10.9091}{rm}{41}\RG{37.021}{0.005}{10.9091}{rm}{43}\RG{47.753}{0.005}{10.9091}{mi}{67}\RG{56.179}{0.005}{10.9091}{rm}{40}\RG{60.333}{0.005}{10.9091}{mi}{39}\RG{67.325}{0.005}{10.9091}{mi}{59}\RG{72.104}{0.005}{10.9091}{mi}{118}\RG{77.671}{0.005}{10.9091}{rm}{41}\RG{84.850}{0.005}{10.9091}{rm}{61}\RG{96.193}{0.005}{10.9091}{mi}{67}\RG{104.620}{0.005}{10.9091}{rm}{40}\RG{108.773}{0.005}{10.9091}{mi}{39}\RG{115.766}{0.005}{10.9091}{mi}{59}\RG{120.544}{0.005}{10.9091}{mi}{117}\RG{126.636}{0.005}{10.9091}{rm}{41}\RG{133.214}{0.005}{10.9091}{sy}{8}\RG{144.123}{0.005}{10.9091}{mi}{67}\RG{152.549}{0.005}{10.9091}{rm}{40}\RG{156.703}{0.005}{10.9091}{mi}{39}\RG{163.696}{0.005}{10.9091}{mi}{59}\RG{168.474}{0.005}{10.9091}{mi}{118}\RG{174.042}{0.005}{10.9091}{rm}{41}\RG{178.196}{0.005}{10.9091}{mi}{59}%
}
\]
且\allowbreak{}\hskip.1em \RB{66.405}{5.308}{2.895}{%
\RG{0.000}{-0.003}{10.9091}{mi}{109}\RG{9.344}{-1.639}{7.9701}{mi}{117}\RG{13.818}{-1.639}{7.9701}{rm}{43}\RG{19.828}{-1.639}{7.9701}{mi}{118}\RG{27.336}{-0.003}{10.9091}{rm}{61}\RG{38.675}{-0.003}{10.9091}{mi}{109}\RG{48.018}{-1.639}{7.9701}{mi}{117}\RG{52.947}{-0.003}{10.9091}{mi}{109}\RG{62.291}{-1.639}{7.9701}{mi}{118}%
}\hskip.1em\allowbreak{} 。（4.2）存在\allowbreak{}\hskip.1em \RB{122.220}{8.682}{3.238}{%
\RG{0.000}{-0.000}{10.9091}{rm}{48}\RG{8.371}{-0.000}{10.9091}{sy}{54}\RG{8.371}{-0.000}{10.9091}{rm}{61}\RG{19.710}{-0.000}{10.9091}{mi}{119}\RG{30.670}{-0.000}{10.9091}{sy}{50}\RG{40.973}{-0.000}{10.9091}{mi}{67}\RG{49.399}{-0.000}{10.9091}{rm}{40}\RG{53.553}{-0.000}{10.9091}{mi}{39}\RG{60.545}{-0.000}{10.9091}{mi}{59}\RG{65.324}{-0.000}{10.9091}{mi}{117}\RG{71.415}{-0.000}{10.9091}{rm}{41}\RG{77.993}{-0.000}{10.9091}{rm}{43}\RG{88.725}{-0.000}{10.9091}{mi}{67}\RG{97.152}{-0.000}{10.9091}{rm}{40}\RG{101.305}{-0.000}{10.9091}{mi}{39}\RG{108.298}{-0.000}{10.9091}{mi}{59}\RG{113.076}{-0.000}{10.9091}{mi}{118}\RG{118.644}{-0.000}{10.9091}{rm}{41}%
}\hskip.1em\allowbreak{} ，使得\allowbreak{}\hskip.1em \RB{68.008}{8.682}{3.227}{%
\RG{0.000}{-0.000}{10.9091}{mi}{109}\RG{9.343}{-1.636}{7.9701}{mi}{119}\RG{18.569}{-0.000}{10.9091}{rm}{61}\RG{29.903}{-0.000}{10.9091}{rm}{91}\RG{32.938}{-0.000}{10.9091}{mi}{109}\RG{42.281}{-1.636}{7.9701}{mi}{117}\RG{47.210}{-0.000}{10.9091}{mi}{59}\RG{51.988}{-0.000}{10.9091}{mi}{109}\RG{61.339}{-1.636}{7.9701}{mi}{118}\RG{65.817}{-0.000}{10.9091}{rm}{93}%
}\hskip.1em\allowbreak{} 。（5）设\allowbreak{}\hskip.1em \RB{54.329}{8.686}{3.224}{%
\RG{0.000}{0.004}{10.9091}{sy}{102}\RG{5.454}{0.004}{10.9091}{mi}{118}\RG{10.619}{-1.633}{7.9701}{rm}{49}\RG{14.937}{0.004}{10.9091}{mi}{59}\RG{19.715}{0.004}{10.9091}{mi}{58}\RG{24.504}{0.004}{10.9091}{mi}{58}\RG{29.283}{0.004}{10.9091}{mi}{58}\RG{34.072}{0.004}{10.9091}{mi}{59}\RG{38.861}{0.004}{10.9091}{mi}{118}\RG{44.028}{-1.633}{7.9701}{mi}{110}\RG{49.171}{0.004}{10.9091}{sy}{103}%
}\hskip.1em\allowbreak{} 是\allowbreak{}\hskip.1em \RB{8.715}{7.955}{0.726}{%
\RG{0.000}{0.004}{10.9091}{mi}{86}%
}\hskip.1em\allowbreak{} 的一组基，记\allowbreak{}\hskip.1em \RB{46.410}{5.315}{3.297}{%
\RG{0.000}{0.004}{10.9091}{mi}{109}\RG{9.344}{-1.633}{7.9701}{mi}{105}\RG{17.415}{0.004}{10.9091}{rm}{61}\RG{30.710}{0.004}{10.9091}{mi}{109}\RG{40.053}{-1.633}{7.9701}{mi}{118}\RG{43.815}{-2.742}{5.9776}{mi}{105}%
}\hskip.1em\allowbreak{} ，\allowbreak{}\hskip.1em \RB{9.522}{5.315}{0.606}{%
\RG{0.000}{0.004}{10.9091}{mi}{109}%
}\hskip.1em\allowbreak{} 为\allowbreak{}\hskip.1em \RB{7.122}{5.315}{2.853}{%
\RG{0.000}{0.004}{10.9091}{mi}{39}%
}\hskip.1em\allowbreak{} 的极小多项式。证明\allowbreak{}\hskip.1em \RB{22.543}{5.315}{0.606}{%
\RG{0.000}{0.004}{10.9091}{mi}{109}\RG{14.330}{0.004}{10.9091}{rm}{61}%
}\hskip.1em\allowbreak{}

\[
\RB{69.482}{9.231}{3.224}{%
\RG{0.000}{0.004}{10.9091}{rm}{91}\RG{3.030}{0.004}{10.9091}{mi}{109}\RG{16.693}{0.004}{10.9091}{mi}{59}\RG{21.471}{0.004}{10.9091}{mi}{58}\RG{26.260}{0.004}{10.9091}{mi}{58}\RG{31.038}{0.004}{10.9091}{mi}{58}\RG{35.827}{0.004}{10.9091}{mi}{59}\RG{40.617}{0.004}{10.9091}{mi}{109}\RG{55.106}{0.004}{10.9091}{rm}{93}\RT{58.073}{0.004}{10.9091}{。}\RG{12.374}{-1.633}{7.9701}{rm}{49}\RG{49.963}{-1.633}{7.9701}{mi}{110}%
}
\]
（6）设\allowbreak{}\hskip.1em \RB{9.522}{5.309}{0.611}{%
\RG{0.000}{-0.002}{10.9091}{mi}{109}%
}\hskip.1em\allowbreak{} 是\allowbreak{}\hskip.1em \RB{7.122}{5.309}{2.858}{%
\RG{0.000}{-0.002}{10.9091}{mi}{39}%
}\hskip.1em\allowbreak{} 的极小多项式，证明存在\allowbreak{}\hskip.1em \RB{47.313}{8.309}{2.858}{%
\RG{0.000}{-0.002}{10.9091}{rm}{48}\RG{8.371}{-0.002}{10.9091}{sy}{54}\RG{8.371}{-0.002}{10.9091}{rm}{61}\RG{19.709}{-0.002}{10.9091}{mi}{118}\RG{28.295}{-0.002}{10.9091}{sy}{50}\RG{38.598}{-0.002}{10.9091}{mi}{86}%
}\hskip.1em\allowbreak{} ，使得\allowbreak{}\hskip.1em \RB{37.712}{5.309}{2.212}{%
\RG{0.000}{-0.002}{10.9091}{mi}{109}\RG{9.343}{-1.639}{7.9701}{mi}{118}\RG{16.851}{-0.002}{10.9091}{rm}{61}\RG{28.190}{-0.002}{10.9091}{mi}{109}%
}\hskip.1em\allowbreak{} 。（7）设\allowbreak{}\hskip.1em \RB{7.122}{5.312}{2.856}{%
\RG{0.000}{0.001}{10.9091}{mi}{39}%
}\hskip.1em\allowbreak{} 在\allowbreak{}\hskip.1em \RB{10.035}{7.952}{0.500}{%
\RG{0.000}{0.001}{10.9091}{mi}{75}%
}\hskip.1em\allowbreak{} 中有\allowbreak{}\hskip.1em \RB{6.555}{5.312}{0.608}{%
\RG{0.000}{0.001}{10.9091}{mi}{110}%
}\hskip.1em\allowbreak{} 个互异特征值\allowbreak{}\hskip.1em \RB{45.658}{8.072}{2.605}{%
\RG{0.000}{0.001}{10.9091}{mi}{21}\RG{6.232}{-1.635}{7.9701}{rm}{49}\RG{10.550}{0.001}{10.9091}{mi}{59}\RG{15.328}{0.001}{10.9091}{mi}{58}\RG{20.117}{0.001}{10.9091}{mi}{58}\RG{24.895}{0.001}{10.9091}{mi}{58}\RG{29.684}{0.001}{10.9091}{mi}{59}\RG{34.474}{0.001}{10.9091}{mi}{21}\RG{40.707}{-1.635}{7.9701}{mi}{110}%
}\hskip.1em\allowbreak{} ，相应的特征向量为\allowbreak{}\hskip.1em \RB{43.524}{5.312}{2.605}{%
\RG{0.000}{0.001}{10.9091}{mi}{118}\RG{5.164}{-1.635}{7.9701}{rm}{49}\RG{9.482}{0.001}{10.9091}{mi}{59}\RG{14.260}{0.001}{10.9091}{mi}{58}\RG{19.049}{0.001}{10.9091}{mi}{58}\RG{23.838}{0.001}{10.9091}{mi}{58}\RG{28.617}{0.001}{10.9091}{mi}{59}\RG{33.406}{0.001}{10.9091}{mi}{118}\RG{38.573}{-1.635}{7.9701}{mi}{110}%
}\hskip.1em\allowbreak{} 。证明\allowbreak{}\hskip.1em \RB{8.715}{7.952}{0.728}{%
\RG{0.000}{0.001}{10.9091}{mi}{86}%
}\hskip.1em\allowbreak{} 是循环空间，并求其全部循环向量。

证明：以下\allowbreak{}\hskip.1em \RB{27.516}{7.759}{1.990}{%
\RG{0.000}{0.005}{10.9091}{mi}{110}\RG{10.217}{0.005}{10.9091}{sy}{21}\RG{22.544}{0.005}{10.9091}{rm}{49}%
}\hskip.1em\allowbreak{} ，\allowbreak{}\hskip.1em \RB{21.324}{8.686}{3.223}{%
\RG{0.000}{0.005}{10.9091}{mi}{75}\RG{9.872}{0.005}{10.9091}{rm}{91}\RG{12.902}{0.005}{10.9091}{mi}{21}\RG{19.133}{0.005}{10.9091}{rm}{93}%
}\hskip.1em\allowbreak{} 表示\allowbreak{}\hskip.1em \RB{10.035}{7.956}{0.500}{%
\RG{0.000}{0.005}{10.9091}{mi}{75}%
}\hskip.1em\allowbreak{} 上的多项式环。对不同时为零的\allowbreak{}\hskip.1em \RB{52.407}{8.686}{3.223}{%
\RG{0.000}{0.005}{10.9091}{mi}{102}\RG{5.836}{0.005}{10.9091}{mi}{59}\RG{10.626}{0.005}{10.9091}{mi}{103}\RG{19.969}{0.005}{10.9091}{sy}{50}\RG{31.084}{0.005}{10.9091}{mi}{75}\RG{40.955}{0.005}{10.9091}{rm}{91}\RG{43.986}{0.005}{10.9091}{mi}{21}\RG{50.216}{0.005}{10.9091}{rm}{93}%
}\hskip.1em\allowbreak{} ，以\allowbreak{}\hskip.1em \RB{23.859}{8.686}{3.234}{%
\RG{0.000}{0.005}{10.9091}{rm}{40}\RG{4.155}{0.005}{10.9091}{mi}{102}\RG{9.991}{0.005}{10.9091}{mi}{59}\RG{14.780}{0.005}{10.9091}{mi}{103}\RG{20.282}{0.005}{10.9091}{rm}{41}%
}\hskip.1em\allowbreak{}表示其首一最大公因式；对任意正整数\allowbreak{}\hskip.1em \RB{5.006}{5.307}{0.613}{%
\RG{0.000}{-0.004}{10.9091}{mi}{115}%
}\hskip.1em\allowbreak{} 及非零多项式\allowbreak{}\hskip.1em \RB{42.468}{8.176}{2.729}{%
\RG{0.000}{-0.004}{10.9091}{mi}{102}\RG{5.265}{-1.641}{7.9701}{rm}{49}\RG{9.583}{-0.004}{10.9091}{mi}{59}\RG{14.361}{-0.004}{10.9091}{mi}{58}\RG{19.151}{-0.004}{10.9091}{mi}{58}\RG{23.940}{-0.004}{10.9091}{mi}{58}\RG{28.718}{-0.004}{10.9091}{mi}{59}\RG{33.507}{-0.004}{10.9091}{mi}{102}\RG{38.775}{-1.641}{7.9701}{mi}{115}%
}\hskip.1em\allowbreak{} ，以\allowbreak{}\hskip.1em \RB{48.025}{8.678}{3.231}{%
\RG{0.000}{-0.004}{10.9091}{rm}{91}\RG{3.030}{-0.004}{10.9091}{mi}{102}\RG{8.296}{-1.641}{7.9701}{rm}{49}\RG{12.615}{-0.004}{10.9091}{mi}{59}\RG{17.393}{-0.004}{10.9091}{mi}{58}\RG{22.182}{-0.004}{10.9091}{mi}{58}\RG{26.960}{-0.004}{10.9091}{mi}{58}\RG{31.749}{-0.004}{10.9091}{mi}{59}\RG{36.539}{-0.004}{10.9091}{mi}{102}\RG{41.806}{-1.641}{7.9701}{mi}{115}\RG{45.834}{-0.004}{10.9091}{rm}{93}%
}\hskip.1em\allowbreak{} 表示

\noindent \RB{429.879}{24.057}{4.683}{%
\RG{266.114}{18.720}{7.9701}{rm}{48}\RT{0.000}{14.761}{10.9091}{其}\RT{11.346}{14.761}{10.9091}{首}\RT{22.691}{14.761}{10.9091}{一}\RT{34.037}{14.761}{10.9091}{最}\RT{45.383}{14.761}{10.9091}{小}\RT{56.728}{14.761}{10.9091}{公}\RT{68.074}{14.761}{10.9091}{倍}\RT{79.430}{14.761}{10.9091}{式}\RT{90.340}{14.761}{10.9091}{。}\RT{101.685}{14.761}{10.9091}{记}\RG{116.904}{14.761}{10.9091}{rm}{105}\RG{119.871}{14.761}{10.9091}{rm}{100}\RT{136.785}{14.761}{10.9091}{为}\RG{152.000}{14.761}{10.9091}{mi}{86}\RT{164.907}{14.761}{10.9091}{的}\RT{176.253}{14.761}{10.9091}{恒}\RT{187.599}{14.761}{10.9091}{等}\RT{198.944}{14.761}{10.9091}{变}\RT{210.290}{14.761}{10.9091}{换}\RT{221.199}{14.761}{10.9091}{，}\RT{232.545}{14.761}{10.9091}{并}\RT{243.891}{14.761}{10.9091}{令}\RG{259.117}{14.761}{10.9091}{mi}{39}\RG{274.587}{14.761}{10.9091}{rm}{61}\RG{287.045}{14.761}{10.9091}{rm}{105}\RG{290.013}{14.761}{10.9091}{rm}{100}\RT{302.623}{14.761}{10.9091}{。}\RT{313.969}{14.761}{10.9091}{对}\RT{325.314}{14.761}{10.9091}{非}\RT{336.660}{14.761}{10.9091}{负}\RT{348.006}{14.761}{10.9091}{整}\RT{359.351}{14.761}{10.9091}{数}\RG{374.576}{14.761}{10.9091}{mi}{100}\RT{384.433}{14.761}{10.9091}{及}\RT{395.778}{14.761}{10.9091}{多}\RT{407.124}{14.761}{10.9091}{项}\RT{418.469}{14.761}{10.9091}{式}\RG{125.806}{13.125}{7.9701}{mi}{86}\RG{295.952}{13.125}{7.9701}{mi}{86}\RG{35.360}{7.270}{9.0910}{ex}{80}\RG{151.974}{7.270}{9.0910}{ex}{80}\RG{44.956}{4.392}{7.9701}{mi}{100}\RG{161.570}{4.392}{7.9701}{mi}{100}\RG{75.474}{3.957}{7.9701}{mi}{106}\RG{192.853}{3.957}{7.9701}{mi}{106}\RG{0.000}{-0.003}{10.9091}{mi}{102}\RG{6.430}{-0.003}{10.9091}{rm}{40}\RG{10.584}{-0.003}{10.9091}{mi}{21}\RG{16.815}{-0.003}{10.9091}{rm}{41}\RG{24.004}{-0.003}{10.9091}{rm}{61}\RG{60.647}{-0.003}{10.9091}{mi}{99}\RG{69.242}{-0.003}{10.9091}{mi}{21}\RT{79.473}{-0.003}{10.9091}{，}\RT{90.382}{-0.003}{10.9091}{定}\RT{101.291}{-0.003}{10.9091}{义}\RG{115.847}{-0.003}{10.9091}{mi}{102}\RG{122.278}{-0.003}{10.9091}{rm}{40}\RG{126.431}{-0.003}{10.9091}{mi}{39}\RG{133.429}{-0.003}{10.9091}{rm}{41}\RG{140.618}{-0.003}{10.9091}{rm}{61}\RG{177.260}{-0.003}{10.9091}{mi}{99}\RG{185.856}{-0.003}{10.9091}{mi}{39}\RT{196.853}{-0.003}{10.9091}{，}\RT{207.762}{-0.003}{10.9091}{并}\RT{218.671}{-0.003}{10.9091}{定}\RT{229.591}{-0.003}{10.9091}{义}\RT{240.501}{-0.003}{10.9091}{零}\RT{251.410}{-0.003}{10.9091}{多}\RT{262.330}{-0.003}{10.9091}{项}\RT{273.239}{-0.003}{10.9091}{式}\RT{284.160}{-0.003}{10.9091}{在}\RG{298.708}{-0.003}{10.9091}{mi}{39}\RT{309.345}{-0.003}{10.9091}{处}\RT{320.254}{-0.003}{10.9091}{的}\RT{331.163}{-0.003}{10.9091}{值}\RT{342.083}{-0.003}{10.9091}{为}\RT{352.993}{-0.003}{10.9091}{零}\RT{363.902}{-0.003}{10.9091}{变}\RT{374.822}{-0.003}{10.9091}{换}\RT{385.731}{-0.003}{10.9091}{。}\RT{396.641}{-0.003}{10.9091}{由}\RT{407.561}{-0.003}{10.9091}{线}\RT{418.470}{-0.003}{10.9091}{性}\RG{65.244}{-1.639}{7.9701}{mi}{106}\RG{181.857}{-1.639}{7.9701}{mi}{106}\RG{44.956}{-2.699}{7.9701}{mi}{106}\RG{48.500}{-2.699}{7.9701}{rm}{61}\RG{54.507}{-2.699}{7.9701}{rm}{48}\RG{161.570}{-2.699}{7.9701}{mi}{106}\RG{165.114}{-2.699}{7.9701}{rm}{61}\RG{171.121}{-2.699}{7.9701}{rm}{48}%
}

变换的复合运算，对任意\allowbreak{}\hskip.1em \RB{50.785}{8.680}{3.229}{%
\RG{0.000}{-0.002}{10.9091}{mi}{102}\RG{5.836}{-0.002}{10.9091}{mi}{59}\RG{10.626}{-0.002}{10.9091}{mi}{103}\RG{19.158}{-0.002}{10.9091}{sy}{50}\RG{29.461}{-0.002}{10.9091}{mi}{75}\RG{39.333}{-0.002}{10.9091}{rm}{91}\RG{42.363}{-0.002}{10.9091}{mi}{21}\RG{48.594}{-0.002}{10.9091}{rm}{93}%
}\hskip.1em\allowbreak{} ，有

\[
\RB{151.889}{8.682}{3.238}{%
\RG{0.000}{-0.000}{10.9091}{rm}{40}\RG{4.154}{-0.000}{10.9091}{mi}{102}\RG{10.579}{-0.000}{10.9091}{mi}{103}\RG{16.090}{-0.000}{10.9091}{rm}{41}\RG{20.247}{-0.000}{10.9091}{rm}{40}\RG{24.398}{-0.000}{10.9091}{mi}{39}\RG{31.396}{-0.000}{10.9091}{rm}{41}\RG{38.574}{-0.000}{10.9091}{rm}{61}\RG{49.919}{-0.000}{10.9091}{mi}{102}\RG{56.349}{-0.000}{10.9091}{rm}{40}\RG{60.502}{-0.000}{10.9091}{mi}{39}\RG{67.500}{-0.000}{10.9091}{rm}{41}\RG{71.654}{-0.000}{10.9091}{mi}{103}\RG{77.160}{-0.000}{10.9091}{rm}{40}\RG{81.314}{-0.000}{10.9091}{mi}{39}\RG{88.312}{-0.000}{10.9091}{rm}{41}\RG{95.490}{-0.000}{10.9091}{rm}{61}\RG{106.835}{-0.000}{10.9091}{mi}{103}\RG{112.341}{-0.000}{10.9091}{rm}{40}\RG{116.495}{-0.000}{10.9091}{mi}{39}\RG{123.492}{-0.000}{10.9091}{rm}{41}\RG{127.647}{-0.000}{10.9091}{mi}{102}\RG{134.076}{-0.000}{10.9091}{rm}{40}\RG{138.230}{-0.000}{10.9091}{mi}{39}\RG{145.228}{-0.000}{10.9091}{rm}{41}\RG{149.382}{-0.000}{10.9091}{mi}{58}%
}
\]
为使（3.4）包含零子空间，约定\allowbreak{}\hskip.1em \RB{64.180}{8.683}{3.237}{%
\RG{0.000}{0.001}{10.9091}{mi}{67}\RG{8.426}{0.001}{10.9091}{rm}{40}\RG{12.580}{0.001}{10.9091}{mi}{39}\RG{19.573}{0.001}{10.9091}{mi}{59}\RG{24.362}{0.001}{10.9091}{rm}{48}\RG{29.708}{0.001}{10.9091}{rm}{41}\RG{36.886}{0.001}{10.9091}{rm}{61}\RG{48.226}{0.001}{10.9091}{sy}{102}\RG{53.681}{0.001}{10.9091}{rm}{48}\RG{59.022}{0.001}{10.9091}{sy}{103}%
}\hskip.1em\allowbreak{} 。

% Original page 14.
（1）证明：固定\allowbreak{}\hskip.1em \RB{49.764}{8.314}{2.854}{%
\RG{0.000}{0.003}{10.9091}{rm}{48}\RG{8.984}{0.003}{10.9091}{sy}{54}\RG{8.984}{0.003}{10.9091}{rm}{61}\RG{20.935}{0.003}{10.9091}{mi}{118}\RG{30.134}{0.003}{10.9091}{sy}{50}\RG{41.049}{0.003}{10.9091}{mi}{86}%
}\hskip.1em\allowbreak{} 。由于\allowbreak{}\hskip.1em \RB{8.715}{7.954}{0.726}{%
\RG{0.000}{0.003}{10.9091}{mi}{86}%
}\hskip.1em\allowbreak{} 的维数为\allowbreak{}\hskip.1em \RB{6.555}{5.314}{0.606}{%
\RG{0.000}{0.003}{10.9091}{mi}{110}%
}\hskip.1em\allowbreak{} ，\allowbreak{}\hskip.1em \RB{24.993}{7.757}{1.337}{%
\RG{0.000}{0.003}{10.9091}{mi}{110}\RG{9.046}{0.003}{10.9091}{rm}{43}\RG{20.020}{0.003}{10.9091}{rm}{49}%
}\hskip.1em\allowbreak{} 个向量\allowbreak{}\hskip.1em \RB{80.556}{8.685}{3.235}{%
\RG{0.000}{0.003}{10.9091}{mi}{118}\RG{5.542}{0.003}{10.9091}{mi}{59}\RG{10.331}{0.003}{10.9091}{mi}{39}\RG{17.339}{0.003}{10.9091}{rm}{40}\RG{21.492}{0.003}{10.9091}{mi}{118}\RG{27.049}{0.003}{10.9091}{rm}{41}\RG{31.202}{0.003}{10.9091}{mi}{59}\RG{35.981}{0.003}{10.9091}{mi}{58}\RG{40.770}{0.003}{10.9091}{mi}{58}\RG{45.548}{0.003}{10.9091}{mi}{58}\RG{50.337}{0.003}{10.9091}{mi}{59}\RG{55.126}{0.003}{10.9091}{mi}{39}\RG{62.127}{3.961}{7.9701}{mi}{110}\RG{67.270}{0.003}{10.9091}{rm}{40}\RG{71.424}{0.003}{10.9091}{mi}{118}\RG{76.980}{0.003}{10.9091}{rm}{41}%
}\hskip.1em\allowbreak{} 线性相关。因此存在非零多项式\allowbreak{}\hskip.1em \RB{42.236}{8.685}{3.225}{%
\RG{0.000}{0.003}{10.9091}{mi}{102}\RG{10.035}{0.003}{10.9091}{sy}{50}\RG{20.912}{0.003}{10.9091}{mi}{75}\RG{30.783}{0.003}{10.9091}{rm}{91}\RG{33.813}{0.003}{10.9091}{mi}{21}\RG{40.045}{0.003}{10.9091}{rm}{93}%
}\hskip.1em\allowbreak{} ，使得\allowbreak{}\hskip.1em \RB{56.502}{8.685}{3.235}{%
\RG{0.000}{0.003}{10.9091}{mi}{102}\RG{6.430}{0.003}{10.9091}{rm}{40}\RG{10.583}{0.003}{10.9091}{mi}{39}\RG{17.581}{0.003}{10.9091}{rm}{41}\RG{21.738}{0.003}{10.9091}{rm}{40}\RG{25.889}{0.003}{10.9091}{mi}{118}\RG{31.445}{0.003}{10.9091}{rm}{41}\RG{39.202}{0.003}{10.9091}{rm}{61}\RG{51.115}{0.003}{10.9091}{rm}{48}%
}\hskip.1em\allowbreak{} 。在所有非零零化多项式中选取次数最小者，再除以其首项系数，得到首一多项式\allowbreak{}\hskip.1em \RB{13.458}{5.315}{2.205}{%
\RG{0.000}{0.004}{10.9091}{mi}{109}\RG{9.343}{-1.632}{7.9701}{mi}{118}%
}\hskip.1em\allowbreak{} 。非零常数多项式不能零化\allowbreak{}\hskip.1em \RB{5.486}{5.315}{0.605}{%
\RG{0.000}{0.004}{10.9091}{mi}{118}%
}\hskip.1em\allowbreak{} ，故

\[
\RB{63.103}{9.232}{2.732}{%
\RG{0.000}{0.005}{10.9091}{rm}{100}\RG{5.935}{0.005}{10.9091}{rm}{101}\RG{10.680}{0.005}{10.9091}{rm}{103}\RG{17.986}{0.005}{10.9091}{mi}{109}\RG{34.837}{0.005}{10.9091}{sy}{21}\RG{46.353}{0.005}{10.9091}{rm}{49}\RT{51.694}{0.005}{10.9091}{。}\RG{27.329}{-1.632}{7.9701}{mi}{118}%
}
\]
任取零化多项式\allowbreak{}\hskip.1em \RB{5.595}{5.309}{2.727}{%
\RG{0.000}{-0.002}{10.9091}{mi}{103}%
}\hskip.1em\allowbreak{} 。由多项式带余除法，存在\allowbreak{}\hskip.1em \RB{51.192}{8.680}{3.229}{%
\RG{0.000}{-0.002}{10.9091}{mi}{113}\RG{5.116}{-0.002}{10.9091}{mi}{59}\RG{9.906}{-0.002}{10.9091}{mi}{114}\RG{18.810}{-0.002}{10.9091}{sy}{50}\RG{29.868}{-0.002}{10.9091}{mi}{75}\RG{39.740}{-0.002}{10.9091}{rm}{91}\RG{42.770}{-0.002}{10.9091}{mi}{21}\RG{49.001}{-0.002}{10.9091}{rm}{93}%
}\hskip.1em\allowbreak{} ，使得\allowbreak{}\hskip.1em \RB{59.246}{6.782}{2.727}{%
\RG{0.000}{-0.002}{10.9091}{mi}{103}\RG{9.293}{-0.002}{10.9091}{rm}{61}\RG{21.387}{-0.002}{10.9091}{mi}{113}\RG{26.504}{-0.002}{10.9091}{mi}{109}\RG{35.858}{-1.638}{7.9701}{mi}{118}\RG{43.063}{-0.002}{10.9091}{rm}{43}\RG{54.098}{-0.002}{10.9091}{mi}{114}%
}\hskip.1em\allowbreak{} ，其中\allowbreak{}\hskip.1em \RB{24.864}{7.754}{0.730}{%
\RG{0.000}{-0.001}{10.9091}{mi}{114}\RG{8.141}{-0.001}{10.9091}{rm}{61}\RG{19.476}{-0.001}{10.9091}{rm}{48}%
}\hskip.1em\allowbreak{} 或\allowbreak{}\hskip.1em \RB{68.907}{8.070}{2.737}{%
\RG{0.000}{-0.001}{10.9091}{rm}{100}\RG{5.935}{-0.001}{10.9091}{rm}{101}\RG{10.680}{-0.001}{10.9091}{rm}{103}\RG{17.985}{-0.001}{10.9091}{mi}{114}\RG{26.112}{-0.001}{10.9091}{mi}{60}\RG{37.464}{-0.001}{10.9091}{rm}{100}\RG{43.399}{-0.001}{10.9091}{rm}{101}\RG{48.144}{-0.001}{10.9091}{rm}{103}\RG{55.449}{-0.001}{10.9091}{mi}{109}\RG{64.792}{-1.637}{7.9701}{mi}{118}%
}\hskip.1em\allowbreak{} 。于是

\[
\RB{159.504}{8.686}{18.216}{%
\RG{0.000}{0.004}{10.9091}{mi}{114}\RG{5.111}{0.004}{10.9091}{rm}{40}\RG{9.265}{0.004}{10.9091}{mi}{39}\RG{16.263}{0.004}{10.9091}{rm}{41}\RG{20.419}{0.004}{10.9091}{rm}{40}\RG{24.570}{0.004}{10.9091}{mi}{118}\RG{30.126}{0.004}{10.9091}{rm}{41}\RG{37.304}{0.004}{10.9091}{rm}{61}\RG{48.649}{0.004}{10.9091}{mi}{103}\RG{54.155}{0.004}{10.9091}{rm}{40}\RG{58.310}{0.004}{10.9091}{mi}{39}\RG{65.307}{0.004}{10.9091}{rm}{41}\RG{69.463}{0.004}{10.9091}{rm}{40}\RG{73.614}{0.004}{10.9091}{mi}{118}\RG{79.171}{0.004}{10.9091}{rm}{41}\RG{85.749}{0.004}{10.9091}{sy}{0}\RG{96.658}{0.004}{10.9091}{mi}{113}\RG{101.785}{0.004}{10.9091}{rm}{40}\RG{105.939}{0.004}{10.9091}{mi}{39}\RG{112.937}{0.004}{10.9091}{rm}{41}\RG{117.091}{0.004}{10.9091}{mi}{109}\RG{130.913}{0.004}{10.9091}{rm}{40}\RG{135.066}{0.004}{10.9091}{mi}{39}\RG{142.064}{0.004}{10.9091}{rm}{41}\RG{146.220}{0.004}{10.9091}{rm}{40}\RG{150.372}{0.004}{10.9091}{mi}{118}\RG{155.927}{0.004}{10.9091}{rm}{41}\RG{126.434}{-1.632}{7.9701}{mi}{118}\RG{37.310}{-17.487}{10.9091}{rm}{61}\RG{48.645}{-17.487}{10.9091}{rm}{48}\RG{53.990}{-17.487}{10.9091}{mi}{58}%
}
\]
若\allowbreak{}\hskip.1em \RB{24.864}{8.312}{2.855}{%
\RG{0.000}{0.001}{10.9091}{mi}{114}\RG{8.141}{0.001}{10.9091}{sy}{54}\RG{8.141}{0.001}{10.9091}{rm}{61}\RG{19.476}{0.001}{10.9091}{rm}{48}%
}\hskip.1em\allowbreak{} ，则与\allowbreak{}\hskip.1em \RB{13.458}{5.312}{2.208}{%
\RG{0.000}{0.001}{10.9091}{mi}{109}\RG{9.343}{-1.635}{7.9701}{mi}{118}%
}\hskip.1em\allowbreak{} 的次数最小性矛盾。因此\allowbreak{}\hskip.1em \RB{24.863}{7.756}{0.728}{%
\RG{0.000}{0.001}{10.9091}{mi}{114}\RG{8.141}{0.001}{10.9091}{rm}{61}\RG{19.475}{0.001}{10.9091}{rm}{48}%
}\hskip.1em\allowbreak{} ，即\allowbreak{}\hskip.1em \RB{28.508}{8.683}{3.226}{%
\RG{0.000}{0.001}{10.9091}{mi}{109}\RG{9.344}{-1.635}{7.9701}{mi}{118}\RG{16.852}{0.001}{10.9091}{sy}{106}\RG{22.913}{0.001}{10.9091}{mi}{103}%
}\hskip.1em\allowbreak{} 。

若\allowbreak{}\hskip.1em \RB{13.458}{8.105}{2.205}{%
\RG{0.000}{0.004}{10.9091}{mi}{109}\RG{0.126}{0.787}{9.0910}{ex}{102}\RG{9.343}{-1.632}{7.9701}{mi}{118}%
}\hskip.1em\allowbreak{} 也是满足定义的首一多项式，则上述整除性质给出\allowbreak{}\hskip.1em \RB{37.342}{8.686}{3.224}{%
\RG{0.000}{0.004}{10.9091}{mi}{109}\RG{9.343}{-1.632}{7.9701}{mi}{118}\RG{17.338}{0.004}{10.9091}{sy}{106}\RG{23.884}{0.004}{10.9091}{mi}{109}\RG{24.010}{0.787}{9.0910}{ex}{102}\RG{33.227}{-1.632}{7.9701}{mi}{118}%
}\hskip.1em\allowbreak{} 及\allowbreak{}\hskip.1em \RB{37.342}{8.686}{3.224}{%
\RG{0.000}{0.004}{10.9091}{mi}{109}\RG{0.126}{0.787}{9.0910}{ex}{102}\RG{9.343}{-1.632}{7.9701}{mi}{118}\RG{17.337}{0.004}{10.9091}{sy}{106}\RG{23.883}{0.004}{10.9091}{mi}{109}\RG{33.227}{-1.632}{7.9701}{mi}{118}%
}\hskip.1em\allowbreak{} 。二者首一，故\allowbreak{}\hskip.1em \RB{41.649}{8.106}{2.204}{%
\RG{0.000}{0.005}{10.9091}{mi}{109}\RG{0.126}{0.788}{9.0910}{ex}{102}\RG{9.343}{-1.632}{7.9701}{mi}{118}\RG{16.852}{0.005}{10.9091}{rm}{61}\RG{28.190}{0.005}{10.9091}{mi}{109}\RG{37.534}{-1.632}{7.9701}{mi}{118}%
}\hskip.1em\allowbreak{} 。这证明了存在性与唯一性。

（2）证明：记\allowbreak{}\hskip.1em \RB{51.362}{8.076}{2.732}{%
\RG{0.000}{0.005}{10.9091}{mi}{100}\RG{8.580}{0.005}{10.9091}{rm}{61}\RG{19.915}{0.005}{10.9091}{rm}{100}\RG{25.850}{0.005}{10.9091}{rm}{101}\RG{30.595}{0.005}{10.9091}{rm}{103}\RG{37.904}{0.005}{10.9091}{mi}{109}\RG{47.247}{-1.633}{7.9701}{mi}{118}%
}\hskip.1em\allowbreak{} ，并写成

\[
\RB{184.752}{10.052}{3.238}{%
\RG{48.961}{4.504}{7.9701}{mi}{100}\RG{91.681}{4.504}{7.9701}{mi}{100}\RG{95.657}{4.504}{7.9701}{sy}{0}\RG{101.668}{4.504}{7.9701}{rm}{49}\RG{0.000}{-0.000}{10.9091}{mi}{109}\RG{13.822}{-0.000}{10.9091}{rm}{40}\RG{17.976}{-0.000}{10.9091}{mi}{21}\RG{24.207}{-0.000}{10.9091}{rm}{41}\RG{31.385}{-0.000}{10.9091}{rm}{61}\RG{42.729}{-0.000}{10.9091}{mi}{21}\RG{55.816}{-0.000}{10.9091}{rm}{43}\RG{66.548}{-0.000}{10.9091}{mi}{99}\RG{85.450}{-0.000}{10.9091}{mi}{21}\RG{108.410}{-0.000}{10.9091}{rm}{43}\RG{119.142}{-0.000}{10.9091}{sy}{1}\RG{123.986}{-0.000}{10.9091}{sy}{1}\RG{128.841}{-0.000}{10.9091}{sy}{1}\RG{134.295}{-0.000}{10.9091}{rm}{43}\RG{145.027}{-0.000}{10.9091}{mi}{99}\RG{153.942}{-0.000}{10.9091}{mi}{21}\RG{162.598}{-0.000}{10.9091}{rm}{43}\RG{173.330}{-0.000}{10.9091}{mi}{99}\RG{182.245}{-0.000}{10.9091}{mi}{58}\RG{9.343}{-1.637}{7.9701}{mi}{118}\RG{71.145}{-1.637}{7.9701}{mi}{100}\RG{75.122}{-1.637}{7.9701}{sy}{0}\RG{81.132}{-1.637}{7.9701}{rm}{49}\RG{149.624}{-1.637}{7.9701}{rm}{49}\RG{177.927}{-1.637}{7.9701}{rm}{48}%
}
\]
由\allowbreak{}\hskip.1em \RB{62.748}{8.687}{3.234}{%
\RG{0.000}{0.005}{10.9091}{mi}{109}\RG{9.343}{-1.632}{7.9701}{mi}{118}\RG{13.821}{0.005}{10.9091}{rm}{40}\RG{17.976}{0.005}{10.9091}{mi}{39}\RG{24.973}{0.005}{10.9091}{rm}{41}\RG{29.129}{0.005}{10.9091}{rm}{40}\RG{33.281}{0.005}{10.9091}{mi}{118}\RG{38.837}{0.005}{10.9091}{rm}{41}\RG{46.015}{0.005}{10.9091}{rm}{61}\RG{57.360}{0.005}{10.9091}{rm}{48}%
}\hskip.1em\allowbreak{} ，有

\[
\RB{101.595}{17.363}{12.855}{%
\RG{49.964}{11.815}{7.9701}{mi}{100}\RG{53.940}{11.815}{7.9701}{sy}{0}\RG{59.950}{11.815}{7.9701}{rm}{49}\RG{50.323}{9.089}{9.0910}{ex}{88}\RG{6.998}{4.501}{7.9701}{mi}{100}\RG{81.226}{4.501}{7.9701}{mi}{106}\RG{0.000}{-0.003}{10.9091}{mi}{39}\RG{11.428}{-0.003}{10.9091}{rm}{40}\RG{15.582}{-0.003}{10.9091}{mi}{118}\RG{21.138}{-0.003}{10.9091}{rm}{41}\RG{28.317}{-0.003}{10.9091}{rm}{61}\RG{39.661}{-0.003}{10.9091}{sy}{0}\RG{65.633}{-0.003}{10.9091}{mi}{99}\RG{74.228}{-0.003}{10.9091}{mi}{39}\RG{85.224}{-0.003}{10.9091}{rm}{40}\RG{89.379}{-0.003}{10.9091}{mi}{118}\RG{94.934}{-0.003}{10.9091}{rm}{41}\RG{99.088}{-0.003}{10.9091}{mi}{58}\RG{70.230}{-1.639}{7.9701}{mi}{106}\RG{50.180}{-10.872}{7.9701}{mi}{106}\RG{53.724}{-10.872}{7.9701}{rm}{61}\RG{59.732}{-10.872}{7.9701}{rm}{48}%
}
\]
令\allowbreak{}\hskip.1em \RB{155.805}{9.503}{3.242}{%
\RG{0.000}{-0.004}{10.9091}{mi}{87}\RG{14.598}{-0.004}{10.9091}{rm}{61}\RG{25.933}{-0.004}{10.9091}{rm}{115}\RG{30.144}{-0.004}{10.9091}{rm}{112}\RG{36.079}{-0.004}{10.9091}{rm}{97}\RG{41.424}{-0.004}{10.9091}{rm}{110}\RG{47.413}{-2.580}{7.9701}{mi}{75}\RG{54.898}{-0.004}{10.9091}{sy}{102}\RG{60.353}{-0.004}{10.9091}{mi}{118}\RG{65.895}{-0.004}{10.9091}{mi}{59}\RG{70.684}{-0.004}{10.9091}{mi}{39}\RG{77.691}{-0.004}{10.9091}{rm}{40}\RG{81.845}{-0.004}{10.9091}{mi}{118}\RG{87.400}{-0.004}{10.9091}{rm}{41}\RG{91.554}{-0.004}{10.9091}{mi}{59}\RG{96.332}{-0.004}{10.9091}{mi}{58}\RG{101.122}{-0.004}{10.9091}{mi}{58}\RG{105.911}{-0.004}{10.9091}{mi}{58}\RG{110.689}{-0.004}{10.9091}{mi}{59}\RG{115.478}{-0.004}{10.9091}{mi}{39}\RG{122.479}{3.956}{7.9701}{mi}{100}\RG{126.455}{3.956}{7.9701}{sy}{0}\RG{132.465}{3.956}{7.9701}{rm}{49}\RG{136.784}{-0.004}{10.9091}{rm}{40}\RG{140.938}{-0.004}{10.9091}{mi}{118}\RG{146.493}{-0.004}{10.9091}{rm}{41}\RG{150.647}{-0.004}{10.9091}{sy}{103}%
}\hskip.1em\allowbreak{} 。上式表明\allowbreak{}\hskip.1em \RB{7.122}{5.307}{2.860}{%
\RG{0.000}{-0.004}{10.9091}{mi}{39}%
}\hskip.1em\allowbreak{} 将\allowbreak{}\hskip.1em \RB{11.693}{7.947}{0.733}{%
\RG{0.000}{-0.004}{10.9091}{mi}{87}%
}\hskip.1em\allowbreak{} 的每个生成向量映入\allowbreak{}\hskip.1em \RB{11.693}{7.947}{0.733}{%
\RG{0.000}{-0.004}{10.9091}{mi}{87}%
}\hskip.1em\allowbreak{} ，故\allowbreak{}\hskip.1em \RB{11.693}{7.947}{0.733}{%
\RG{0.000}{-0.004}{10.9091}{mi}{87}%
}\hskip.1em\allowbreak{} 是\allowbreak{}\hskip.1em \RB{7.122}{5.308}{2.859}{%
\RG{0.000}{-0.003}{10.9091}{mi}{39}%
}\hskip.1em\allowbreak{} –不变子空间。由于\allowbreak{}\hskip.1em \RB{30.891}{7.948}{0.939}{%
\RG{0.000}{-0.003}{10.9091}{mi}{118}\RG{8.741}{-0.003}{10.9091}{sy}{50}\RG{19.198}{-0.003}{10.9091}{mi}{87}%
}\hskip.1em\allowbreak{} ，对非负整数\allowbreak{}\hskip.1em \RB{4.766}{7.665}{2.728}{%
\RG{0.000}{-0.003}{10.9091}{mi}{106}%
}\hskip.1em\allowbreak{} 归纳可得\allowbreak{}\hskip.1em \RB{50.195}{9.270}{3.241}{%
\RG{0.000}{-0.003}{10.9091}{mi}{39}\RG{6.997}{3.956}{7.9701}{mi}{106}\RG{10.996}{-0.003}{10.9091}{rm}{40}\RG{15.150}{-0.003}{10.9091}{mi}{118}\RG{20.706}{-0.003}{10.9091}{rm}{41}\RG{28.044}{-0.003}{10.9091}{sy}{50}\RG{38.503}{-0.003}{10.9091}{mi}{87}%
}\hskip.1em\allowbreak{} ，因此\allowbreak{}\hskip.1em \RB{60.620}{8.679}{3.241}{%
\RG{0.000}{-0.003}{10.9091}{mi}{67}\RG{8.425}{-0.003}{10.9091}{rm}{40}\RG{12.580}{-0.003}{10.9091}{mi}{39}\RG{19.573}{-0.003}{10.9091}{mi}{59}\RG{24.351}{-0.003}{10.9091}{mi}{118}\RG{29.918}{-0.003}{10.9091}{rm}{41}\RG{37.257}{-0.003}{10.9091}{sy}{18}\RG{48.927}{-0.003}{10.9091}{mi}{87}%
}\hskip.1em\allowbreak{} 。由\allowbreak{}\hskip.1em \RB{11.693}{7.949}{0.731}{%
\RG{0.000}{-0.002}{10.9091}{mi}{87}%
}\hskip.1em\allowbreak{} 的定义又有\allowbreak{}\hskip.1em \RB{59.608}{8.680}{3.240}{%
\RG{0.000}{-0.002}{10.9091}{mi}{87}\RG{14.598}{-0.002}{10.9091}{sy}{18}\RG{26.113}{-0.002}{10.9091}{mi}{67}\RG{34.540}{-0.002}{10.9091}{rm}{40}\RG{38.693}{-0.002}{10.9091}{mi}{39}\RG{45.686}{-0.002}{10.9091}{mi}{59}\RG{50.464}{-0.002}{10.9091}{mi}{118}\RG{56.032}{-0.002}{10.9091}{rm}{41}%
}\hskip.1em\allowbreak{} ，故两者相等。

若存在不全为零的\allowbreak{}\hskip.1em \RB{75.370}{8.072}{2.605}{%
\RG{0.000}{0.001}{10.9091}{mi}{98}\RG{4.542}{-1.635}{7.9701}{rm}{48}\RG{8.860}{0.001}{10.9091}{mi}{59}\RG{13.638}{0.001}{10.9091}{mi}{58}\RG{18.427}{0.001}{10.9091}{mi}{58}\RG{23.216}{0.001}{10.9091}{mi}{58}\RG{27.994}{0.001}{10.9091}{mi}{59}\RG{32.784}{0.001}{10.9091}{mi}{98}\RG{37.328}{-1.635}{7.9701}{mi}{100}\RG{41.304}{-1.635}{7.9701}{sy}{0}\RG{47.315}{-1.635}{7.9701}{rm}{49}\RG{54.847}{0.001}{10.9091}{sy}{50}\RG{65.335}{0.001}{10.9091}{mi}{75}%
}\hskip.1em\allowbreak{} ，使得\allowbreak{}\hskip.1em \RB{79.247}{9.942}{4.680}{%
\RG{0.000}{7.274}{9.0910}{ex}{80}\RG{9.596}{4.395}{7.9701}{mi}{100}\RG{9.596}{-2.696}{7.9701}{mi}{106}\RG{13.140}{-2.696}{7.9701}{rm}{61}\RG{13.573}{4.395}{7.9701}{sy}{0}\RG{19.148}{-2.696}{7.9701}{rm}{48}\RG{19.583}{4.395}{7.9701}{rm}{49}\RG{25.719}{0.001}{10.9091}{mi}{98}\RG{30.261}{-1.635}{7.9701}{mi}{106}\RG{34.259}{0.001}{10.9091}{mi}{39}\RG{41.257}{3.959}{7.9701}{mi}{106}\RG{45.256}{0.001}{10.9091}{rm}{40}\RG{49.410}{0.001}{10.9091}{mi}{118}\RG{54.965}{0.001}{10.9091}{rm}{41}\RG{62.329}{0.001}{10.9091}{rm}{61}\RG{73.860}{0.001}{10.9091}{rm}{48}%
}\hskip.1em\allowbreak{} ，则非零多项式\allowbreak{}\hskip.1em \RB{43.995}{9.942}{4.680}{%
\RG{0.000}{7.274}{9.0910}{ex}{80}\RG{9.596}{4.395}{7.9701}{mi}{100}\RG{9.596}{-2.696}{7.9701}{mi}{106}\RG{13.140}{-2.696}{7.9701}{rm}{61}\RG{13.573}{4.395}{7.9701}{sy}{0}\RG{19.148}{-2.696}{7.9701}{rm}{48}\RG{19.583}{4.395}{7.9701}{rm}{49}\RG{25.719}{0.001}{10.9091}{mi}{98}\RG{30.261}{-1.635}{7.9701}{mi}{106}\RG{34.259}{0.001}{10.9091}{mi}{21}\RG{40.491}{3.959}{7.9701}{mi}{106}%
}\hskip.1em\allowbreak{}是次数小于\allowbreak{}\hskip.1em \RB{6.020}{8.073}{0.607}{%
\RG{0.000}{0.002}{10.9091}{mi}{100}%
}\hskip.1em\allowbreak{} 的零化多项式，与\allowbreak{}\hskip.1em \RB{13.459}{5.313}{2.208}{%
\RG{0.000}{0.002}{10.9091}{mi}{109}\RG{9.344}{-1.635}{7.9701}{mi}{118}%
}\hskip.1em\allowbreak{} 的定义矛盾。因此上述\allowbreak{}\hskip.1em \RB{6.020}{8.073}{0.607}{%
\RG{0.000}{0.002}{10.9091}{mi}{100}%
}\hskip.1em\allowbreak{} 个向量线性无关，从而\allowbreak{}\hskip.1em \RB{26.311}{8.073}{0.607}{%
\RG{0.000}{0.002}{10.9091}{mi}{107}\RG{8.952}{0.002}{10.9091}{rm}{61}\RG{20.291}{0.002}{10.9091}{mi}{100}%
}\hskip.1em\allowbreak{} ，且构成\allowbreak{}\hskip.1em \RB{33.495}{8.684}{3.236}{%
\RG{0.000}{0.002}{10.9091}{mi}{67}\RG{8.426}{0.002}{10.9091}{rm}{40}\RG{12.580}{0.002}{10.9091}{mi}{39}\RG{19.573}{0.002}{10.9091}{mi}{59}\RG{24.351}{0.002}{10.9091}{mi}{118}\RG{29.919}{0.002}{10.9091}{rm}{41}%
}\hskip.1em\allowbreak{} 的一组基。

题设关系给出的多项式\allowbreak{}\hskip.1em \RB{112.400}{9.502}{2.293}{%
\RG{0.000}{-0.004}{10.9091}{mi}{21}\RG{6.232}{3.955}{7.9701}{mi}{107}\RG{13.112}{-0.004}{10.9091}{rm}{43}\RG{23.629}{-0.004}{10.9091}{mi}{97}\RG{29.236}{-1.641}{7.9701}{mi}{107}\RG{33.453}{-1.641}{7.9701}{sy}{0}\RG{39.463}{-1.641}{7.9701}{rm}{49}\RG{43.781}{-0.004}{10.9091}{mi}{21}\RG{50.013}{3.955}{7.9701}{mi}{107}\RG{54.230}{3.955}{7.9701}{sy}{0}\RG{60.240}{3.955}{7.9701}{rm}{49}\RG{66.767}{-0.004}{10.9091}{rm}{43}\RG{77.283}{-0.004}{10.9091}{sy}{1}\RG{82.127}{-0.004}{10.9091}{sy}{1}\RG{86.982}{-0.004}{10.9091}{sy}{1}\RG{92.219}{-0.004}{10.9091}{rm}{43}\RG{102.736}{-0.004}{10.9091}{mi}{97}\RG{108.344}{-1.641}{7.9701}{rm}{48}%
}\hskip.1em\allowbreak{} 首一、次数为\allowbreak{}\hskip.1em \RB{5.966}{8.067}{0.613}{%
\RG{0.000}{-0.004}{10.9091}{mi}{107}%
}\hskip.1em\allowbreak{} ，并且零化\allowbreak{}\hskip.1em \RB{5.486}{5.307}{0.613}{%
\RG{0.000}{-0.004}{10.9091}{mi}{118}%
}\hskip.1em\allowbreak{} 。由（1）的整除性质及\allowbreak{}\hskip.1em \RB{52.173}{8.067}{2.740}{%
\RG{0.000}{-0.004}{10.9091}{rm}{100}\RG{5.935}{-0.004}{10.9091}{rm}{101}\RG{10.680}{-0.004}{10.9091}{rm}{103}\RG{17.985}{-0.004}{10.9091}{mi}{109}\RG{27.329}{-1.640}{7.9701}{mi}{118}\RG{34.853}{-0.004}{10.9091}{rm}{61}\RG{46.208}{-0.004}{10.9091}{mi}{107}%
}\hskip.1em\allowbreak{} ，它等于\allowbreak{}\hskip.1em \RB{13.458}{5.307}{2.213}{%
\RG{0.000}{-0.004}{10.9091}{mi}{109}\RG{9.343}{-1.640}{7.9701}{mi}{118}%
}\hskip.1em\allowbreak{} 。此外，\allowbreak{}\hskip.1em \RB{33.495}{8.678}{3.242}{%
\RG{0.000}{-0.004}{10.9091}{mi}{67}\RG{8.426}{-0.004}{10.9091}{rm}{40}\RG{12.580}{-0.004}{10.9091}{mi}{39}\RG{19.573}{-0.004}{10.9091}{mi}{59}\RG{24.351}{-0.004}{10.9091}{mi}{118}\RG{29.919}{-0.004}{10.9091}{rm}{41}%
}\hskip.1em\allowbreak{} 本身不变，而任何包含\allowbreak{}\hskip.1em \RB{5.486}{5.307}{0.613}{%
\RG{0.000}{-0.004}{10.9091}{mi}{118}%
}\hskip.1em\allowbreak{} 的不变子空间均包含所有\allowbreak{}\hskip.1em \RB{24.282}{9.271}{3.241}{%
\RG{0.000}{-0.003}{10.9091}{mi}{39}\RG{6.997}{3.957}{7.9701}{mi}{106}\RG{10.996}{-0.003}{10.9091}{rm}{40}\RG{15.150}{-0.003}{10.9091}{mi}{118}\RG{20.706}{-0.003}{10.9091}{rm}{41}%
}\hskip.1em\allowbreak{} ，故它是包含\allowbreak{}\hskip.1em \RB{5.486}{5.308}{0.612}{%
\RG{0.000}{-0.003}{10.9091}{mi}{118}%
}\hskip.1em\allowbreak{} 的最小不变子空间。

（3.1）证明：任取\allowbreak{}\hskip.1em \RB{54.757}{8.679}{3.241}{%
\RG{0.000}{-0.003}{10.9091}{mi}{119}\RG{10.959}{-0.003}{10.9091}{sy}{50}\RG{21.262}{-0.003}{10.9091}{mi}{67}\RG{29.689}{-0.003}{10.9091}{rm}{40}\RG{33.842}{-0.003}{10.9091}{mi}{39}\RG{40.835}{-0.003}{10.9091}{mi}{59}\RG{45.613}{-0.003}{10.9091}{mi}{118}\RG{51.181}{-0.003}{10.9091}{rm}{41}%
}\hskip.1em\allowbreak{} 。由（2）的基，存在唯一的\allowbreak{}\hskip.1em \RB{75.242}{8.068}{2.608}{%
\RG{0.000}{-0.003}{10.9091}{mi}{98}\RG{4.542}{-1.639}{7.9701}{rm}{48}\RG{8.860}{-0.003}{10.9091}{mi}{59}\RG{13.638}{-0.003}{10.9091}{mi}{58}\RG{18.427}{-0.003}{10.9091}{mi}{58}\RG{23.216}{-0.003}{10.9091}{mi}{58}\RG{27.994}{-0.003}{10.9091}{mi}{59}\RG{32.784}{-0.003}{10.9091}{mi}{98}\RG{37.328}{-1.639}{7.9701}{mi}{107}\RG{41.546}{-1.639}{7.9701}{sy}{0}\RG{47.556}{-1.639}{7.9701}{rm}{49}\RG{54.904}{-0.003}{10.9091}{sy}{50}\RG{65.207}{-0.003}{10.9091}{mi}{75}%
}\hskip.1em\allowbreak{} ，使得

\[
\RB{74.116}{17.369}{12.849}{%
\RG{22.298}{11.822}{7.9701}{mi}{107}\RG{26.516}{11.822}{7.9701}{sy}{0}\RG{32.527}{11.822}{7.9701}{rm}{49}\RG{22.778}{9.094}{9.0910}{ex}{88}\RG{53.746}{4.508}{7.9701}{mi}{106}\RG{0.000}{0.004}{10.9091}{mi}{119}\RG{10.960}{0.004}{10.9091}{rm}{61}\RG{38.208}{0.004}{10.9091}{mi}{98}\RG{46.749}{0.004}{10.9091}{mi}{39}\RG{57.745}{0.004}{10.9091}{rm}{40}\RG{61.899}{0.004}{10.9091}{mi}{118}\RG{67.455}{0.004}{10.9091}{rm}{41}\RG{71.609}{0.004}{10.9091}{mi}{58}\RG{42.749}{-1.633}{7.9701}{mi}{106}\RG{22.635}{-10.865}{7.9701}{mi}{106}\RG{26.179}{-10.865}{7.9701}{rm}{61}\RG{32.186}{-10.865}{7.9701}{rm}{48}%
}
\]
令\allowbreak{}\hskip.1em \RB{78.649}{9.946}{4.677}{%
\RG{0.000}{0.005}{10.9091}{mi}{103}\RG{5.506}{0.005}{10.9091}{rm}{40}\RG{9.660}{0.005}{10.9091}{mi}{21}\RG{15.891}{0.005}{10.9091}{rm}{41}\RG{23.070}{0.005}{10.9091}{rm}{61}\RG{34.414}{7.277}{9.0910}{ex}{80}\RG{44.010}{4.398}{7.9701}{mi}{107}\RG{44.010}{-2.693}{7.9701}{mi}{106}\RG{47.555}{-2.693}{7.9701}{rm}{61}\RG{48.227}{4.398}{7.9701}{sy}{0}\RG{53.562}{-2.693}{7.9701}{rm}{48}\RG{54.238}{4.398}{7.9701}{rm}{49}\RG{60.374}{0.004}{10.9091}{mi}{98}\RG{64.915}{-1.633}{7.9701}{mi}{106}\RG{68.915}{0.004}{10.9091}{mi}{21}\RG{75.145}{3.963}{7.9701}{mi}{106}%
}\hskip.1em\allowbreak{} ，则\allowbreak{}\hskip.1em \RB{56.395}{8.686}{3.234}{%
\RG{0.000}{0.004}{10.9091}{mi}{119}\RG{10.959}{0.004}{10.9091}{rm}{61}\RG{22.298}{0.004}{10.9091}{mi}{103}\RG{27.804}{0.004}{10.9091}{rm}{40}\RG{31.958}{0.004}{10.9091}{mi}{39}\RG{38.956}{0.004}{10.9091}{rm}{41}\RG{43.113}{0.004}{10.9091}{rm}{40}\RG{47.264}{0.004}{10.9091}{mi}{118}\RG{52.819}{0.004}{10.9091}{rm}{41}%
}\hskip.1em\allowbreak{} ，且\allowbreak{}\hskip.1em \RB{25.259}{7.759}{2.721}{%
\RG{0.000}{0.004}{10.9091}{mi}{103}\RG{8.537}{0.004}{10.9091}{rm}{61}\RG{19.871}{0.004}{10.9091}{rm}{48}%
}\hskip.1em\allowbreak{} 或\allowbreak{}\hskip.1em \RB{43.828}{8.076}{2.732}{%
\RG{0.000}{0.004}{10.9091}{rm}{100}\RG{5.935}{0.004}{10.9091}{rm}{101}\RG{10.680}{0.004}{10.9091}{rm}{103}\RG{17.986}{0.004}{10.9091}{mi}{103}\RG{26.517}{0.004}{10.9091}{mi}{60}\RG{37.862}{0.004}{10.9091}{mi}{107}%
}\hskip.1em\allowbreak{} 。

（3.2）证明：记\allowbreak{}\hskip.1em \RB{56.395}{8.686}{3.234}{%
\RG{0.000}{0.004}{10.9091}{mi}{119}\RG{10.959}{0.004}{10.9091}{rm}{61}\RG{22.298}{0.004}{10.9091}{mi}{103}\RG{27.804}{0.004}{10.9091}{rm}{40}\RG{31.958}{0.004}{10.9091}{mi}{39}\RG{38.956}{0.004}{10.9091}{rm}{41}\RG{43.113}{0.004}{10.9091}{rm}{40}\RG{47.264}{0.004}{10.9091}{mi}{118}\RG{52.819}{0.004}{10.9091}{rm}{41}%
}\hskip.1em\allowbreak{} 。由于\allowbreak{}\hskip.1em \RB{54.757}{8.686}{3.234}{%
\RG{0.000}{0.004}{10.9091}{mi}{119}\RG{10.960}{0.004}{10.9091}{sy}{50}\RG{21.263}{0.004}{10.9091}{mi}{67}\RG{29.689}{0.004}{10.9091}{rm}{40}\RG{33.842}{0.004}{10.9091}{mi}{39}\RG{40.835}{0.004}{10.9091}{mi}{59}\RG{45.613}{0.004}{10.9091}{mi}{118}\RG{51.181}{0.004}{10.9091}{rm}{41}%
}\hskip.1em\allowbreak{} 且该空间不变，有\allowbreak{}\hskip.1em \RB{84.486}{8.686}{3.234}{%
\RG{0.000}{0.004}{10.9091}{mi}{67}\RG{8.426}{0.004}{10.9091}{rm}{40}\RG{12.580}{0.004}{10.9091}{mi}{39}\RG{19.573}{0.004}{10.9091}{mi}{59}\RG{24.351}{0.004}{10.9091}{mi}{119}\RG{32.292}{0.004}{10.9091}{rm}{41}\RG{39.476}{0.004}{10.9091}{sy}{18}\RG{50.991}{0.004}{10.9091}{mi}{67}\RG{59.417}{0.004}{10.9091}{rm}{40}\RG{63.571}{0.004}{10.9091}{mi}{39}\RG{70.564}{0.004}{10.9091}{mi}{59}\RG{75.342}{0.004}{10.9091}{mi}{118}\RG{80.910}{0.004}{10.9091}{rm}{41}%
}\hskip.1em\allowbreak{} 。

若\allowbreak{}\hskip.1em \RB{51.764}{8.680}{3.240}{%
\RG{0.000}{-0.002}{10.9091}{rm}{40}\RG{4.154}{-0.002}{10.9091}{mi}{103}\RG{9.652}{-0.002}{10.9091}{mi}{59}\RG{14.441}{-0.002}{10.9091}{mi}{109}\RG{23.789}{-1.638}{7.9701}{mi}{118}\RG{28.267}{-0.002}{10.9091}{rm}{41}\RG{35.445}{-0.002}{10.9091}{rm}{61}\RG{46.791}{-0.002}{10.9091}{rm}{49}%
}\hskip.1em\allowbreak{} ，则由裴蜀等式，存在\allowbreak{}\hskip.1em \RB{49.591}{8.680}{3.229}{%
\RG{0.000}{-0.002}{10.9091}{mi}{97}\RG{5.607}{-0.002}{10.9091}{mi}{59}\RG{10.386}{-0.002}{10.9091}{mi}{98}\RG{17.965}{-0.002}{10.9091}{sy}{50}\RG{28.268}{-0.002}{10.9091}{mi}{75}\RG{38.139}{-0.002}{10.9091}{rm}{91}\RG{41.170}{-0.002}{10.9091}{mi}{21}\RG{47.400}{-0.002}{10.9091}{rm}{93}%
}\hskip.1em\allowbreak{} ，使得\allowbreak{}\hskip.1em \RB{61.972}{8.069}{2.727}{%
\RG{0.000}{-0.002}{10.9091}{mi}{97}\RG{5.607}{-0.002}{10.9091}{mi}{103}\RG{13.538}{-0.002}{10.9091}{rm}{43}\RG{24.271}{-0.002}{10.9091}{mi}{98}\RG{28.809}{-0.002}{10.9091}{mi}{109}\RG{38.156}{-1.638}{7.9701}{mi}{118}\RG{45.665}{-0.002}{10.9091}{rm}{61}\RG{57.000}{-0.002}{10.9091}{rm}{49}%
}\hskip.1em\allowbreak{} 。因此

\[
\RB{150.930}{8.685}{20.727}{%
\RG{0.000}{0.003}{10.9091}{mi}{118}\RG{8.586}{0.003}{10.9091}{rm}{61}\RG{19.925}{0.003}{10.9091}{mi}{97}\RG{25.532}{0.003}{10.9091}{rm}{40}\RG{29.686}{0.003}{10.9091}{mi}{39}\RG{36.684}{0.003}{10.9091}{rm}{41}\RG{40.838}{0.003}{10.9091}{mi}{103}\RG{46.344}{0.003}{10.9091}{rm}{40}\RG{50.498}{0.003}{10.9091}{mi}{39}\RG{57.496}{0.003}{10.9091}{rm}{41}\RG{61.652}{0.003}{10.9091}{rm}{40}\RG{65.803}{0.003}{10.9091}{mi}{118}\RG{71.359}{0.003}{10.9091}{rm}{41}\RG{77.938}{0.003}{10.9091}{rm}{43}\RG{88.670}{0.003}{10.9091}{mi}{98}\RG{93.211}{0.003}{10.9091}{rm}{40}\RG{97.365}{0.003}{10.9091}{mi}{39}\RG{104.363}{0.003}{10.9091}{rm}{41}\RG{108.517}{0.003}{10.9091}{mi}{109}\RG{122.339}{0.003}{10.9091}{rm}{40}\RG{126.493}{0.003}{10.9091}{mi}{39}\RG{133.491}{0.003}{10.9091}{rm}{41}\RG{137.647}{0.003}{10.9091}{rm}{40}\RG{141.798}{0.003}{10.9091}{mi}{118}\RG{147.354}{0.003}{10.9091}{rm}{41}\RG{117.860}{-1.634}{7.9701}{mi}{118}\RG{8.586}{-17.488}{10.9091}{rm}{61}\RG{19.925}{-17.488}{10.9091}{mi}{97}\RG{25.532}{-17.488}{10.9091}{rm}{40}\RG{29.686}{-17.488}{10.9091}{mi}{39}\RG{36.684}{-17.488}{10.9091}{rm}{41}\RG{40.840}{-17.488}{10.9091}{rm}{40}\RG{44.992}{-17.488}{10.9091}{mi}{119}\RG{52.921}{-17.488}{10.9091}{rm}{41}\RG{57.075}{-17.488}{10.9091}{mi}{58}%
}
\]
于是\allowbreak{}\hskip.1em \RB{54.757}{8.682}{3.238}{%
\RG{0.000}{-0.000}{10.9091}{mi}{118}\RG{8.586}{-0.000}{10.9091}{sy}{50}\RG{18.889}{-0.000}{10.9091}{mi}{67}\RG{27.314}{-0.000}{10.9091}{rm}{40}\RG{31.469}{-0.000}{10.9091}{mi}{39}\RG{38.462}{-0.000}{10.9091}{mi}{59}\RG{43.240}{-0.000}{10.9091}{mi}{119}\RG{51.181}{-0.000}{10.9091}{rm}{41}%
}\hskip.1em\allowbreak{} 。由循环子空间的最小性，\allowbreak{}\hskip.1em \RB{84.486}{8.682}{3.238}{%
\RG{0.000}{-0.000}{10.9091}{mi}{67}\RG{8.426}{-0.000}{10.9091}{rm}{40}\RG{12.580}{-0.000}{10.9091}{mi}{39}\RG{19.573}{-0.000}{10.9091}{mi}{59}\RG{24.351}{-0.000}{10.9091}{mi}{118}\RG{29.919}{-0.000}{10.9091}{rm}{41}\RG{37.103}{-0.000}{10.9091}{sy}{18}\RG{48.618}{-0.000}{10.9091}{mi}{67}\RG{57.044}{-0.000}{10.9091}{rm}{40}\RG{61.198}{-0.000}{10.9091}{mi}{39}\RG{68.190}{-0.000}{10.9091}{mi}{59}\RG{72.969}{-0.000}{10.9091}{mi}{119}\RG{80.910}{-0.000}{10.9091}{rm}{41}%
}\hskip.1em\allowbreak{} ，故两者相等。

反之，若两者相等，则存在\allowbreak{}\hskip.1em \RB{42.302}{8.685}{3.225}{%
\RG{0.000}{0.003}{10.9091}{mi}{104}\RG{9.927}{0.003}{10.9091}{sy}{50}\RG{20.979}{0.003}{10.9091}{mi}{75}\RG{30.850}{0.003}{10.9091}{rm}{91}\RG{33.881}{0.003}{10.9091}{mi}{21}\RG{40.111}{0.003}{10.9091}{rm}{93}%
}\hskip.1em\allowbreak{} ，使得\allowbreak{}\hskip.1em \RB{58.535}{8.685}{3.236}{%
\RG{0.000}{0.003}{10.9091}{mi}{118}\RG{9.334}{0.003}{10.9091}{rm}{61}\RG{21.421}{0.003}{10.9091}{mi}{104}\RG{27.571}{0.003}{10.9091}{rm}{40}\RG{31.724}{0.003}{10.9091}{mi}{39}\RG{38.722}{0.003}{10.9091}{rm}{41}\RG{42.878}{0.003}{10.9091}{rm}{40}\RG{47.030}{0.003}{10.9091}{mi}{119}\RG{54.959}{0.003}{10.9091}{rm}{41}%
}\hskip.1em\allowbreak{} 。因此\allowbreak{}\hskip.1em \RB{89.658}{8.685}{3.236}{%
\RG{0.000}{0.003}{10.9091}{rm}{40}\RG{4.154}{0.003}{10.9091}{mi}{104}\RG{10.307}{0.003}{10.9091}{mi}{103}\RG{18.533}{0.003}{10.9091}{sy}{0}\RG{29.742}{0.003}{10.9091}{rm}{49}\RG{35.087}{0.003}{10.9091}{rm}{41}\RG{39.244}{0.003}{10.9091}{rm}{40}\RG{43.390}{0.003}{10.9091}{mi}{39}\RG{50.388}{0.003}{10.9091}{rm}{41}\RG{54.544}{0.003}{10.9091}{rm}{40}\RG{58.696}{0.003}{10.9091}{mi}{118}\RG{64.252}{0.003}{10.9091}{rm}{41}\RG{72.183}{0.003}{10.9091}{rm}{61}\RG{84.270}{0.003}{10.9091}{rm}{48}%
}\hskip.1em\allowbreak{} ，由（1）得\allowbreak{}\hskip.1em \RB{55.902}{8.686}{3.224}{%
\RG{0.000}{0.004}{10.9091}{mi}{109}\RG{9.343}{-1.633}{7.9701}{mi}{118}\RG{17.933}{0.004}{10.9091}{sy}{106}\RG{25.075}{0.004}{10.9091}{mi}{104}\RG{31.228}{0.004}{10.9091}{mi}{103}\RG{39.588}{0.004}{10.9091}{sy}{0}\RG{50.929}{0.004}{10.9091}{rm}{49}%
}\hskip.1em\allowbreak{} 。任一同时整除\allowbreak{}\hskip.1em \RB{5.595}{5.315}{2.722}{%
\RG{0.000}{0.004}{10.9091}{mi}{103}%
}\hskip.1em\allowbreak{} 与\allowbreak{}\hskip.1em \RB{13.458}{5.315}{2.205}{%
\RG{0.000}{0.004}{10.9091}{mi}{109}\RG{9.343}{-1.633}{7.9701}{mi}{118}%
}\hskip.1em\allowbreak{} 的多项式均整除\allowbreak{}\hskip.1em \RB{86.864}{8.686}{3.235}{%
\RG{0.000}{0.004}{10.9091}{mi}{104}\RG{6.153}{0.004}{10.9091}{mi}{103}\RG{14.512}{0.004}{10.9091}{sy}{0}\RG{25.854}{0.004}{10.9091}{rm}{40}\RG{30.008}{0.004}{10.9091}{mi}{104}\RG{36.161}{0.004}{10.9091}{mi}{103}\RG{44.520}{0.004}{10.9091}{sy}{0}\RG{55.861}{0.004}{10.9091}{rm}{49}\RG{61.207}{0.004}{10.9091}{rm}{41}\RG{69.465}{0.004}{10.9091}{rm}{61}\RG{81.891}{0.004}{10.9091}{rm}{49}%
}\hskip.1em\allowbreak{} ，故

\[
\RB{63.541}{9.232}{3.234}{%
\RG{0.000}{0.004}{10.9091}{rm}{40}\RG{4.155}{0.004}{10.9091}{mi}{103}\RG{9.653}{0.004}{10.9091}{mi}{59}\RG{14.442}{0.004}{10.9091}{mi}{109}\RG{28.268}{0.004}{10.9091}{rm}{41}\RG{35.446}{0.004}{10.9091}{rm}{61}\RG{46.792}{0.004}{10.9091}{rm}{49}\RT{52.132}{0.004}{10.9091}{。}\RG{23.789}{-1.633}{7.9701}{mi}{118}%
}
\]
% Original page 15.
（3.3）证明：令\allowbreak{}\hskip.1em \RB{46.422}{8.685}{3.225}{%
\RG{0.000}{0.003}{10.9091}{mi}{113}\RG{8.869}{0.003}{10.9091}{rm}{61}\RG{20.918}{0.003}{10.9091}{mi}{109}\RG{30.263}{-1.634}{7.9701}{mi}{118}\RG{34.740}{0.003}{10.9091}{mi}{61}\RG{40.086}{0.003}{10.9091}{mi}{104}%
}\hskip.1em\allowbreak{} ，并取\allowbreak{}\hskip.1em \RB{57.439}{8.685}{3.235}{%
\RG{0.000}{0.003}{10.9091}{mi}{119}\RG{11.671}{0.003}{10.9091}{rm}{61}\RG{23.721}{0.003}{10.9091}{mi}{113}\RG{28.848}{0.003}{10.9091}{rm}{40}\RG{33.002}{0.003}{10.9091}{mi}{39}\RG{39.999}{0.003}{10.9091}{rm}{41}\RG{44.156}{0.003}{10.9091}{rm}{40}\RG{48.307}{0.003}{10.9091}{mi}{118}\RG{53.863}{0.003}{10.9091}{rm}{41}%
}\hskip.1em\allowbreak{} 。由于\allowbreak{}\hskip.1em \RB{6.336}{8.074}{0.606}{%
\RG{0.000}{0.003}{10.9091}{mi}{104}%
}\hskip.1em\allowbreak{} 非常数，有\allowbreak{}\hskip.1em \RB{70.347}{8.074}{2.734}{%
\RG{0.000}{0.003}{10.9091}{rm}{100}\RG{5.935}{0.003}{10.9091}{rm}{101}\RG{10.680}{0.003}{10.9091}{rm}{103}\RG{17.985}{0.003}{10.9091}{mi}{113}\RG{26.843}{0.003}{10.9091}{mi}{60}\RG{38.903}{0.003}{10.9091}{rm}{100}\RG{44.838}{0.003}{10.9091}{rm}{101}\RG{49.583}{0.003}{10.9091}{rm}{103}\RG{56.888}{0.003}{10.9091}{mi}{109}\RG{66.232}{-1.634}{7.9701}{mi}{118}%
}\hskip.1em\allowbreak{} ；又\allowbreak{}\hskip.1em \RB{24.880}{8.314}{2.854}{%
\RG{0.000}{0.003}{10.9091}{mi}{113}\RG{8.158}{0.003}{10.9091}{sy}{54}\RG{8.158}{0.003}{10.9091}{rm}{61}\RG{19.493}{0.003}{10.9091}{rm}{48}%
}\hskip.1em\allowbreak{} ，故由\allowbreak{}\hskip.1em \RB{13.459}{5.314}{2.206}{%
\RG{0.000}{0.003}{10.9091}{mi}{109}\RG{9.344}{-1.633}{7.9701}{mi}{118}%
}\hskip.1em\allowbreak{} 的次数最小性可得\allowbreak{}\hskip.1em \RB{27.681}{8.314}{2.854}{%
\RG{0.000}{0.003}{10.9091}{mi}{119}\RG{10.959}{0.003}{10.9091}{sy}{54}\RG{10.959}{0.003}{10.9091}{rm}{61}\RG{22.294}{0.003}{10.9091}{rm}{48}%
}\hskip.1em\allowbreak{} 。由构造，\allowbreak{}\hskip.1em \RB{54.757}{8.685}{3.235}{%
\RG{0.000}{0.003}{10.9091}{mi}{119}\RG{10.959}{0.003}{10.9091}{sy}{50}\RG{21.262}{0.003}{10.9091}{mi}{67}\RG{29.689}{0.003}{10.9091}{rm}{40}\RG{33.842}{0.003}{10.9091}{mi}{39}\RG{40.835}{0.003}{10.9091}{mi}{59}\RG{45.613}{0.003}{10.9091}{mi}{118}\RG{51.181}{0.003}{10.9091}{rm}{41}%
}\hskip.1em\allowbreak{} ，并且

\[
\RB{117.269}{8.683}{20.728}{%
\RG{0.000}{0.001}{10.9091}{mi}{104}\RG{6.149}{0.001}{10.9091}{rm}{40}\RG{10.303}{0.001}{10.9091}{mi}{39}\RG{17.300}{0.001}{10.9091}{rm}{41}\RG{21.457}{0.001}{10.9091}{rm}{40}\RG{25.609}{0.001}{10.9091}{mi}{119}\RG{33.538}{0.001}{10.9091}{rm}{41}\RG{40.716}{0.001}{10.9091}{rm}{61}\RG{52.061}{0.001}{10.9091}{mi}{104}\RG{58.209}{0.001}{10.9091}{rm}{40}\RG{62.364}{0.001}{10.9091}{mi}{39}\RG{69.361}{0.001}{10.9091}{rm}{41}\RG{73.515}{0.001}{10.9091}{mi}{113}\RG{78.643}{0.001}{10.9091}{rm}{40}\RG{82.797}{0.001}{10.9091}{mi}{39}\RG{89.794}{0.001}{10.9091}{rm}{41}\RG{93.950}{0.001}{10.9091}{rm}{40}\RG{98.102}{0.001}{10.9091}{mi}{118}\RG{103.658}{0.001}{10.9091}{rm}{41}\RG{40.722}{-17.490}{10.9091}{rm}{61}\RG{52.061}{-17.490}{10.9091}{mi}{109}\RG{65.883}{-17.490}{10.9091}{rm}{40}\RG{70.036}{-17.490}{10.9091}{mi}{39}\RG{77.033}{-17.490}{10.9091}{rm}{41}\RG{81.190}{-17.490}{10.9091}{rm}{40}\RG{85.342}{-17.490}{10.9091}{mi}{118}\RG{90.897}{-17.490}{10.9091}{rm}{41}\RG{98.076}{-17.490}{10.9091}{rm}{61}\RG{109.421}{-17.490}{10.9091}{rm}{48}\RG{114.761}{-17.490}{10.9091}{mi}{58}\RG{61.404}{-19.126}{7.9701}{mi}{118}%
}
\]
所以\allowbreak{}\hskip.1em \RB{30.966}{8.683}{3.226}{%
\RG{0.000}{0.001}{10.9091}{mi}{109}\RG{9.343}{-1.635}{7.9701}{mi}{119}\RG{18.569}{0.001}{10.9091}{sy}{106}\RG{24.630}{0.001}{10.9091}{mi}{104}%
}\hskip.1em\allowbreak{} 。另一方面，\allowbreak{}\hskip.1em \RB{66.839}{8.683}{3.237}{%
\RG{0.000}{0.001}{10.9091}{mi}{109}\RG{9.343}{-1.635}{7.9701}{mi}{119}\RG{15.539}{0.001}{10.9091}{rm}{40}\RG{19.693}{0.001}{10.9091}{mi}{39}\RG{26.690}{0.001}{10.9091}{rm}{41}\RG{30.847}{0.001}{10.9091}{rm}{40}\RG{34.999}{0.001}{10.9091}{mi}{119}\RG{42.927}{0.001}{10.9091}{rm}{41}\RG{50.106}{0.001}{10.9091}{rm}{61}\RG{61.451}{0.001}{10.9091}{rm}{48}%
}\hskip.1em\allowbreak{} 表明\allowbreak{}\hskip.1em \RB{43.762}{8.683}{3.226}{%
\RG{0.000}{0.001}{10.9091}{mi}{109}\RG{9.344}{-1.635}{7.9701}{mi}{118}\RG{16.852}{0.001}{10.9091}{sy}{106}\RG{22.913}{0.001}{10.9091}{mi}{109}\RG{32.256}{-1.635}{7.9701}{mi}{119}\RG{38.451}{0.001}{10.9091}{mi}{113}%
}\hskip.1em\allowbreak{} 。将\allowbreak{}\hskip.1em \RB{39.654}{8.072}{2.615}{%
\RG{0.000}{0.001}{10.9091}{mi}{109}\RG{9.343}{-1.635}{7.9701}{mi}{118}\RG{16.852}{0.001}{10.9091}{rm}{61}\RG{28.190}{0.001}{10.9091}{mi}{104}\RG{34.343}{0.001}{10.9091}{mi}{113}%
}\hskip.1em\allowbreak{} 代入，并在整环\allowbreak{}\hskip.1em \RB{21.324}{8.683}{3.226}{%
\RG{0.000}{0.001}{10.9091}{mi}{75}\RG{9.871}{0.001}{10.9091}{rm}{91}\RG{12.901}{0.001}{10.9091}{mi}{21}\RG{19.133}{0.001}{10.9091}{rm}{93}%
}\hskip.1em\allowbreak{}中约去非零因子\allowbreak{}\hskip.1em \RB{5.311}{5.313}{2.615}{%
\RG{0.000}{0.002}{10.9091}{mi}{113}%
}\hskip.1em\allowbreak{} ，得到\allowbreak{}\hskip.1em \RB{30.415}{8.684}{3.225}{%
\RG{0.000}{0.002}{10.9091}{mi}{104}\RG{9.180}{0.002}{10.9091}{sy}{106}\RG{15.240}{0.002}{10.9091}{mi}{109}\RG{24.584}{-1.635}{7.9701}{mi}{119}%
}\hskip.1em\allowbreak{} 。由于二者首一，故\allowbreak{}\hskip.1em \RB{36.244}{8.073}{2.208}{%
\RG{0.000}{0.002}{10.9091}{mi}{109}\RG{9.343}{-1.635}{7.9701}{mi}{119}\RG{18.569}{0.002}{10.9091}{rm}{61}\RG{29.908}{0.002}{10.9091}{mi}{104}%
}\hskip.1em\allowbreak{} 。

（3.4）证明：设\allowbreak{}\hskip.1em \RB{11.693}{7.951}{0.729}{%
\RG{0.000}{0.000}{10.9091}{mi}{87}%
}\hskip.1em\allowbreak{} 是\allowbreak{}\hskip.1em \RB{33.495}{8.681}{3.239}{%
\RG{0.000}{0.000}{10.9091}{mi}{67}\RG{8.426}{-0.001}{10.9091}{rm}{40}\RG{12.580}{0.000}{10.9091}{mi}{39}\RG{19.573}{0.000}{10.9091}{mi}{59}\RG{24.351}{0.000}{10.9091}{mi}{118}\RG{29.919}{-0.001}{10.9091}{rm}{41}%
}\hskip.1em\allowbreak{} 的任一\allowbreak{}\hskip.1em \RB{7.122}{5.310}{2.857}{%
\RG{0.000}{-0.001}{10.9091}{mi}{39}%
}\hskip.1em\allowbreak{} –不变子空间，定义

\[
\RB{152.707}{8.679}{3.241}{%
\RG{0.000}{-0.003}{10.9091}{mi}{73}\RG{16.559}{-0.003}{10.9091}{rm}{61}\RG{27.898}{-0.003}{10.9091}{sy}{102}\RG{33.353}{-0.003}{10.9091}{mi}{102}\RG{42.813}{-0.003}{10.9091}{sy}{50}\RG{53.116}{-0.003}{10.9091}{mi}{75}\RG{62.988}{-0.003}{10.9091}{rm}{91}\RG{66.018}{-0.003}{10.9091}{mi}{21}\RG{72.249}{-0.003}{10.9091}{rm}{93}\RG{78.238}{-0.003}{10.9091}{rm}{58}\RG{84.244}{-0.003}{10.9091}{mi}{102}\RG{90.675}{-0.003}{10.9091}{rm}{40}\RG{94.828}{-0.003}{10.9091}{mi}{39}\RG{101.825}{-0.003}{10.9091}{rm}{41}\RG{105.982}{-0.003}{10.9091}{rm}{40}\RG{110.134}{-0.003}{10.9091}{mi}{118}\RG{115.689}{-0.003}{10.9091}{rm}{41}\RG{122.873}{-0.003}{10.9091}{sy}{50}\RG{133.177}{-0.003}{10.9091}{mi}{87}\RG{144.746}{-0.003}{10.9091}{sy}{103}\RG{150.200}{-0.003}{10.9091}{mi}{58}\RG{4.711}{-1.639}{7.9701}{mi}{87}%
}
\]
由于\allowbreak{}\hskip.1em \RB{11.693}{7.946}{0.734}{%
\RG{0.000}{-0.005}{10.9091}{mi}{87}%
}\hskip.1em\allowbreak{} 是线性子空间，\allowbreak{}\hskip.1em \RB{13.284}{7.946}{2.293}{%
\RG{0.000}{-0.005}{10.9091}{mi}{73}\RG{4.711}{-1.641}{7.9701}{mi}{87}%
}\hskip.1em\allowbreak{} 对加减法封闭。由于\allowbreak{}\hskip.1em \RB{11.693}{7.946}{0.734}{%
\RG{0.000}{-0.005}{10.9091}{mi}{87}%
}\hskip.1em\allowbreak{} 对\allowbreak{}\hskip.1em \RB{7.122}{5.306}{2.861}{%
\RG{0.000}{-0.005}{10.9091}{mi}{39}%
}\hskip.1em\allowbreak{} 不变，它对每个\allowbreak{}\hskip.1em \RB{20.088}{8.677}{3.243}{%
\RG{0.000}{-0.005}{10.9091}{mi}{112}\RG{5.361}{-0.005}{10.9091}{rm}{40}\RG{9.515}{-0.005}{10.9091}{mi}{39}\RG{16.512}{-0.005}{10.9091}{rm}{41}%
}\hskip.1em\allowbreak{} （\allowbreak{}\hskip.1em \RB{40.018}{8.677}{3.232}{%
\RG{0.000}{-0.005}{10.9091}{mi}{112}\RG{8.391}{-0.005}{10.9091}{sy}{50}\RG{18.695}{-0.005}{10.9091}{mi}{75}\RG{28.566}{-0.005}{10.9091}{rm}{91}\RG{31.597}{-0.005}{10.9091}{mi}{21}\RG{37.827}{-0.005}{10.9091}{rm}{93}%
}\hskip.1em\allowbreak{} ）也不变，故\allowbreak{}\hskip.1em \RB{33.047}{8.176}{2.729}{%
\RG{0.000}{-0.004}{10.9091}{mi}{102}\RG{9.460}{-0.004}{10.9091}{sy}{50}\RG{19.763}{-0.004}{10.9091}{mi}{73}\RG{24.475}{-1.640}{7.9701}{mi}{87}%
}\hskip.1em\allowbreak{} 蕴含\allowbreak{}\hskip.1em \RB{38.408}{8.176}{2.729}{%
\RG{0.000}{-0.004}{10.9091}{mi}{112}\RG{5.356}{-0.004}{10.9091}{mi}{102}\RG{14.821}{-0.004}{10.9091}{sy}{50}\RG{25.124}{-0.004}{10.9091}{mi}{73}\RG{29.836}{-1.640}{7.9701}{mi}{87}%
}\hskip.1em\allowbreak{} 。因此\allowbreak{}\hskip.1em \RB{13.284}{7.947}{2.293}{%
\RG{0.000}{-0.004}{10.9091}{mi}{73}\RG{4.711}{-1.640}{7.9701}{mi}{87}%
}\hskip.1em\allowbreak{} 是\allowbreak{}\hskip.1em \RB{21.324}{8.678}{3.231}{%
\RG{0.000}{-0.004}{10.9091}{mi}{75}\RG{9.871}{-0.004}{10.9091}{rm}{91}\RG{12.901}{-0.004}{10.9091}{mi}{21}\RG{19.133}{-0.004}{10.9091}{rm}{93}%
}\hskip.1em\allowbreak{} 的理想。又\allowbreak{}\hskip.1em \RB{40.440}{7.947}{2.293}{%
\RG{0.000}{-0.004}{10.9091}{mi}{109}\RG{9.343}{-1.640}{7.9701}{mi}{118}\RG{16.852}{-0.004}{10.9091}{sy}{50}\RG{27.155}{-0.004}{10.9091}{mi}{73}\RG{31.867}{-1.640}{7.9701}{mi}{87}%
}\hskip.1em\allowbreak{} ，故该理想非零。

在\allowbreak{}\hskip.1em \RB{13.284}{7.950}{2.290}{%
\RG{0.000}{-0.001}{10.9091}{mi}{73}\RG{4.711}{-1.637}{7.9701}{mi}{87}%
}\hskip.1em\allowbreak{} 的非零元素中选取次数最小者，并将其化为首一多项式\allowbreak{}\hskip.1em \RB{5.595}{5.310}{2.726}{%
\RG{0.000}{-0.001}{10.9091}{mi}{103}%
}\hskip.1em\allowbreak{} 。任取\allowbreak{}\hskip.1em \RB{33.394}{8.179}{2.726}{%
\RG{0.000}{-0.001}{10.9091}{mi}{102}\RG{9.633}{-0.001}{10.9091}{sy}{50}\RG{20.110}{-0.001}{10.9091}{mi}{73}\RG{24.821}{-1.637}{7.9701}{mi}{87}%
}\hskip.1em\allowbreak{} ，将其带余除以\allowbreak{}\hskip.1em \RB{5.595}{5.311}{2.725}{%
\RG{0.000}{0.000}{10.9091}{mi}{103}%
}\hskip.1em\allowbreak{} ，写成\allowbreak{}\hskip.1em \RB{49.340}{8.180}{2.725}{%
\RG{0.000}{0.000}{10.9091}{mi}{102}\RG{9.460}{0.000}{10.9091}{rm}{61}\RG{20.799}{0.000}{10.9091}{mi}{112}\RG{26.155}{0.000}{10.9091}{mi}{103}\RG{33.775}{0.000}{10.9091}{rm}{43}\RG{44.193}{0.000}{10.9091}{mi}{114}%
}\hskip.1em\allowbreak{} ，其中\allowbreak{}\hskip.1em \RB{24.864}{7.755}{0.729}{%
\RG{0.000}{0.000}{10.9091}{mi}{114}\RG{8.141}{0.000}{10.9091}{rm}{61}\RG{19.476}{0.000}{10.9091}{rm}{48}%
}\hskip.1em\allowbreak{} 或\allowbreak{}\hskip.1em \RB{61.043}{8.071}{2.736}{%
\RG{0.000}{0.000}{10.9091}{rm}{100}\RG{5.935}{0.000}{10.9091}{rm}{101}\RG{10.680}{0.000}{10.9091}{rm}{103}\RG{17.985}{0.000}{10.9091}{mi}{114}\RG{26.112}{0.000}{10.9091}{mi}{60}\RG{37.464}{0.000}{10.9091}{rm}{100}\RG{43.399}{0.000}{10.9091}{rm}{101}\RG{48.144}{0.000}{10.9091}{rm}{103}\RG{55.449}{0.000}{10.9091}{mi}{103}%
}\hskip.1em\allowbreak{} 。由于\allowbreak{}\hskip.1em \RB{76.098}{8.180}{2.725}{%
\RG{0.000}{0.000}{10.9091}{mi}{114}\RG{8.141}{0.000}{10.9091}{rm}{61}\RG{19.479}{0.000}{10.9091}{mi}{102}\RG{28.019}{0.000}{10.9091}{sy}{0}\RG{38.613}{0.000}{10.9091}{mi}{112}\RG{43.969}{0.000}{10.9091}{mi}{103}\RG{52.511}{0.000}{10.9091}{sy}{50}\RG{62.814}{0.000}{10.9091}{mi}{73}\RG{67.525}{-1.636}{7.9701}{mi}{87}%
}\hskip.1em\allowbreak{} ，\allowbreak{}\hskip.1em \RB{5.595}{5.311}{2.725}{%
\RG{0.000}{0.000}{10.9091}{mi}{103}%
}\hskip.1em\allowbreak{} 的次数最小性迫使\allowbreak{}\hskip.1em \RB{25.666}{7.756}{0.728}{%
\RG{0.000}{0.001}{10.9091}{mi}{114}\RG{8.540}{0.001}{10.9091}{rm}{61}\RG{20.279}{0.001}{10.9091}{rm}{48}%
}\hskip.1em\allowbreak{} 。故\allowbreak{}\hskip.1em \RB{41.934}{8.683}{3.237}{%
\RG{0.000}{0.001}{10.9091}{mi}{73}\RG{4.711}{-1.636}{7.9701}{mi}{87}\RG{16.959}{0.001}{10.9091}{rm}{61}\RG{28.697}{0.001}{10.9091}{rm}{40}\RG{32.850}{0.001}{10.9091}{mi}{103}\RG{38.357}{0.001}{10.9091}{rm}{41}%
}\hskip.1em\allowbreak{} ，其中\allowbreak{}\hskip.1em \RB{101.881}{8.683}{3.237}{%
\RG{0.000}{0.001}{10.9091}{rm}{40}\RG{4.154}{0.001}{10.9091}{mi}{103}\RG{9.660}{0.001}{10.9091}{rm}{41}\RG{17.242}{0.001}{10.9091}{rm}{61}\RG{28.981}{0.001}{10.9091}{sy}{102}\RG{34.436}{0.001}{10.9091}{mi}{112}\RG{39.793}{0.001}{10.9091}{mi}{103}\RG{48.733}{0.001}{10.9091}{rm}{58}\RG{55.129}{0.001}{10.9091}{mi}{112}\RG{63.921}{0.001}{10.9091}{sy}{50}\RG{74.623}{0.001}{10.9091}{mi}{75}\RG{84.494}{0.001}{10.9091}{rm}{91}\RG{87.525}{0.001}{10.9091}{mi}{21}\RG{93.755}{0.001}{10.9091}{rm}{93}\RG{96.723}{0.001}{10.9091}{sy}{103}%
}\hskip.1em\allowbreak{} 。特别地，由\allowbreak{}\hskip.1em \RB{41.237}{7.952}{2.289}{%
\RG{0.000}{0.001}{10.9091}{mi}{109}\RG{9.343}{-1.636}{7.9701}{mi}{118}\RG{17.251}{0.001}{10.9091}{sy}{50}\RG{27.953}{0.001}{10.9091}{mi}{73}\RG{32.664}{-1.636}{7.9701}{mi}{87}%
}\hskip.1em\allowbreak{} 得

\[
\RB{39.829}{9.228}{3.226}{%
\RG{0.000}{0.001}{10.9091}{mi}{103}\RG{8.537}{0.001}{10.9091}{sy}{106}\RG{14.597}{0.001}{10.9091}{mi}{109}\RT{28.419}{0.001}{10.9091}{。}\RG{23.941}{-1.635}{7.9701}{mi}{118}%
}
\]
任意\allowbreak{}\hskip.1em \RB{32.956}{7.946}{0.941}{%
\RG{0.000}{-0.005}{10.9091}{mi}{119}\RG{10.960}{-0.005}{10.9091}{sy}{50}\RG{21.263}{-0.005}{10.9091}{mi}{87}%
}\hskip.1em\allowbreak{} 均可写成\allowbreak{}\hskip.1em \RB{57.320}{8.677}{3.243}{%
\RG{0.000}{-0.005}{10.9091}{mi}{119}\RG{10.959}{-0.005}{10.9091}{rm}{61}\RG{22.298}{-0.005}{10.9091}{mi}{102}\RG{28.729}{-0.005}{10.9091}{rm}{40}\RG{32.882}{-0.005}{10.9091}{mi}{39}\RG{39.880}{-0.005}{10.9091}{rm}{41}\RG{44.036}{-0.005}{10.9091}{rm}{40}\RG{48.188}{-0.005}{10.9091}{mi}{118}\RG{53.743}{-0.005}{10.9091}{rm}{41}%
}\hskip.1em\allowbreak{} 。此时\allowbreak{}\hskip.1em \RB{33.047}{8.176}{2.730}{%
\RG{0.000}{-0.005}{10.9091}{mi}{102}\RG{9.460}{-0.005}{10.9091}{sy}{50}\RG{19.763}{-0.005}{10.9091}{mi}{73}\RG{24.475}{-1.642}{7.9701}{mi}{87}%
}\hskip.1em\allowbreak{} ，故\allowbreak{}\hskip.1em \RB{31.750}{8.176}{2.730}{%
\RG{0.000}{-0.005}{10.9091}{mi}{102}\RG{9.460}{-0.005}{10.9091}{rm}{61}\RG{20.799}{-0.005}{10.9091}{mi}{112}\RG{26.155}{-0.005}{10.9091}{mi}{103}%
}\hskip.1em\allowbreak{} ，从而\allowbreak{}\hskip.1em \RB{87.530}{8.677}{3.243}{%
\RG{0.000}{-0.005}{10.9091}{mi}{119}\RG{10.959}{-0.005}{10.9091}{rm}{61}\RG{22.298}{-0.005}{10.9091}{mi}{112}\RG{27.659}{-0.005}{10.9091}{rm}{40}\RG{31.813}{-0.005}{10.9091}{mi}{39}\RG{38.810}{-0.005}{10.9091}{rm}{41}\RG{42.965}{-0.005}{10.9091}{mi}{103}\RG{48.471}{-0.005}{10.9091}{rm}{40}\RG{52.625}{-0.005}{10.9091}{mi}{39}\RG{59.623}{-0.005}{10.9091}{rm}{41}\RG{63.779}{-0.005}{10.9091}{rm}{40}\RG{67.930}{-0.005}{10.9091}{mi}{118}\RG{73.486}{-0.005}{10.9091}{rm}{41}\RG{80.670}{-0.005}{10.9091}{sy}{50}%
}\hskip.1em\allowbreak{}\allowbreak{}\hskip.1em \RB{62.617}{8.677}{3.243}{%
\RG{0.000}{-0.005}{10.9091}{mi}{67}\RG{8.426}{-0.005}{10.9091}{rm}{40}\RG{12.580}{-0.005}{10.9091}{mi}{39}\RG{19.573}{-0.005}{10.9091}{mi}{59}\RG{24.351}{-0.005}{10.9091}{mi}{103}\RG{29.869}{-0.005}{10.9091}{rm}{40}\RG{34.023}{-0.005}{10.9091}{mi}{39}\RG{41.020}{-0.005}{10.9091}{rm}{41}\RG{45.177}{-0.005}{10.9091}{rm}{40}\RG{49.329}{-0.005}{10.9091}{mi}{118}\RG{54.884}{-0.005}{10.9091}{rm}{41}\RG{59.040}{-0.005}{10.9091}{rm}{41}%
}\hskip.1em\allowbreak{} 。反向包含关系由\allowbreak{}\hskip.1em \RB{59.702}{8.677}{3.243}{%
\RG{0.000}{-0.005}{10.9091}{mi}{103}\RG{5.506}{-0.005}{10.9091}{rm}{40}\RG{9.661}{-0.005}{10.9091}{mi}{39}\RG{16.658}{-0.005}{10.9091}{rm}{41}\RG{20.814}{-0.005}{10.9091}{rm}{40}\RG{24.966}{-0.005}{10.9091}{mi}{118}\RG{30.522}{-0.005}{10.9091}{rm}{41}\RG{37.706}{-0.005}{10.9091}{sy}{50}\RG{48.009}{-0.005}{10.9091}{mi}{87}%
}\hskip.1em\allowbreak{} 及\allowbreak{}\hskip.1em \RB{11.693}{7.946}{0.734}{%
\RG{0.000}{-0.005}{10.9091}{mi}{87}%
}\hskip.1em\allowbreak{} 的不变性成立。因此

\[
\RB{91.636}{8.685}{3.235}{%
\RG{0.000}{0.003}{10.9091}{mi}{87}\RG{14.598}{0.003}{10.9091}{rm}{61}\RG{25.937}{0.003}{10.9091}{mi}{67}\RG{34.362}{0.003}{10.9091}{rm}{40}\RG{38.516}{0.003}{10.9091}{mi}{39}\RG{45.509}{0.003}{10.9091}{mi}{59}\RG{50.287}{0.003}{10.9091}{mi}{103}\RG{55.806}{0.003}{10.9091}{rm}{40}\RG{59.960}{0.003}{10.9091}{mi}{39}\RG{66.957}{0.003}{10.9091}{rm}{41}\RG{71.114}{0.003}{10.9091}{rm}{40}\RG{75.265}{0.003}{10.9091}{mi}{118}\RG{80.821}{0.003}{10.9091}{rm}{41}\RG{84.978}{0.003}{10.9091}{rm}{41}\RG{89.129}{0.003}{10.9091}{mi}{58}%
}
\]
反之，对\allowbreak{}\hskip.1em \RB{13.459}{5.315}{2.206}{%
\RG{0.000}{0.004}{10.9091}{mi}{109}\RG{9.344}{-1.633}{7.9701}{mi}{118}%
}\hskip.1em\allowbreak{} 的任一首一因式\allowbreak{}\hskip.1em \RB{5.595}{5.315}{2.722}{%
\RG{0.000}{0.004}{10.9091}{mi}{103}%
}\hskip.1em\allowbreak{} ，上述循环子空间都是\allowbreak{}\hskip.1em \RB{33.494}{8.686}{3.235}{%
\RG{0.000}{0.004}{10.9091}{mi}{67}\RG{8.425}{0.004}{10.9091}{rm}{40}\RG{12.580}{0.004}{10.9091}{mi}{39}\RG{19.573}{0.004}{10.9091}{mi}{59}\RG{24.351}{0.004}{10.9091}{mi}{118}\RG{29.918}{0.004}{10.9091}{rm}{41}%
}\hskip.1em\allowbreak{} 的不变子空间。其对应理想恰为\allowbreak{}\hskip.1em \RB{13.236}{8.686}{3.235}{%
\RG{0.000}{0.004}{10.9091}{rm}{40}\RG{4.154}{0.004}{10.9091}{mi}{103}\RG{9.660}{0.004}{10.9091}{rm}{41}%
}\hskip.1em\allowbreak{} ：若\allowbreak{}\hskip.1em \RB{105.186}{8.686}{3.235}{%
\RG{0.000}{0.004}{10.9091}{mi}{102}\RG{6.430}{0.004}{10.9091}{rm}{40}\RG{10.583}{0.004}{10.9091}{mi}{39}\RG{17.581}{0.004}{10.9091}{rm}{41}\RG{21.738}{0.004}{10.9091}{rm}{40}\RG{25.889}{0.004}{10.9091}{mi}{118}\RG{31.445}{0.004}{10.9091}{rm}{41}\RG{38.853}{0.004}{10.9091}{rm}{61}\RG{50.422}{0.004}{10.9091}{mi}{112}\RG{55.783}{0.004}{10.9091}{rm}{40}\RG{59.937}{0.004}{10.9091}{mi}{39}\RG{66.934}{0.004}{10.9091}{rm}{41}\RG{71.088}{0.004}{10.9091}{mi}{103}\RG{76.594}{0.004}{10.9091}{rm}{40}\RG{80.748}{0.004}{10.9091}{mi}{39}\RG{87.746}{0.004}{10.9091}{rm}{41}\RG{91.903}{0.004}{10.9091}{rm}{40}\RG{96.054}{0.004}{10.9091}{mi}{118}\RG{101.609}{0.004}{10.9091}{rm}{41}%
}\hskip.1em\allowbreak{} ，则\allowbreak{}\hskip.1em \RB{54.263}{8.686}{3.224}{%
\RG{0.000}{0.004}{10.9091}{mi}{109}\RG{9.343}{-1.632}{7.9701}{mi}{118}\RG{17.079}{0.004}{10.9091}{sy}{106}\RG{23.367}{0.004}{10.9091}{mi}{102}\RG{32.311}{0.004}{10.9091}{sy}{0}\RG{43.311}{0.004}{10.9091}{mi}{112}\RG{48.668}{0.004}{10.9091}{mi}{103}%
}\hskip.1em\allowbreak{} ；因\allowbreak{}\hskip.1em \RB{28.510}{8.686}{3.224}{%
\RG{0.000}{0.004}{10.9091}{mi}{103}\RG{8.764}{0.004}{10.9091}{sy}{106}\RG{15.051}{0.004}{10.9091}{mi}{109}\RG{24.395}{-1.632}{7.9701}{mi}{118}%
}\hskip.1em\allowbreak{} ，可得\allowbreak{}\hskip.1em \RB{21.507}{8.686}{3.224}{%
\RG{0.000}{0.004}{10.9091}{mi}{103}\RG{8.763}{0.004}{10.9091}{sy}{106}\RG{15.051}{0.004}{10.9091}{mi}{102}%
}\hskip.1em\allowbreak{} ，而反向包含关系直接来自循环子空间的定义。两个主理想相等时，其首一生成元互相整除，因而相等。因此不同首一因式给出不同不变子空间。

由多项式的唯一因式分解，\allowbreak{}\hskip.1em \RB{13.458}{5.310}{2.210}{%
\RG{0.000}{-0.001}{10.9091}{mi}{109}\RG{9.343}{-1.637}{7.9701}{mi}{118}%
}\hskip.1em\allowbreak{} 只有有限个首一因式，故上述分类给出的不变子空间只有有限个。其中\allowbreak{}\hskip.1em \RB{24.845}{7.755}{2.725}{%
\RG{0.000}{0.000}{10.9091}{mi}{103}\RG{8.537}{0.000}{10.9091}{rm}{61}\RG{19.872}{0.000}{10.9091}{rm}{49}%
}\hskip.1em\allowbreak{} 给出\allowbreak{}\hskip.1em \RB{33.494}{8.682}{3.238}{%
\RG{0.000}{0.000}{10.9091}{mi}{67}\RG{8.425}{0.000}{10.9091}{rm}{40}\RG{12.580}{0.000}{10.9091}{mi}{39}\RG{19.573}{0.000}{10.9091}{mi}{59}\RG{24.351}{0.000}{10.9091}{mi}{118}\RG{29.918}{0.000}{10.9091}{rm}{41}%
}\hskip.1em\allowbreak{} ，\allowbreak{}\hskip.1em \RB{33.333}{5.311}{2.725}{%
\RG{0.000}{0.000}{10.9091}{mi}{103}\RG{8.537}{0.000}{10.9091}{rm}{61}\RG{19.875}{0.000}{10.9091}{mi}{109}\RG{29.219}{-1.636}{7.9701}{mi}{118}%
}\hskip.1em\allowbreak{} 给出\allowbreak{}\hskip.1em \RB{15.954}{8.682}{3.227}{%
\RG{0.000}{0.000}{10.9091}{sy}{102}\RG{5.455}{0.000}{10.9091}{rm}{48}\RG{10.796}{0.000}{10.9091}{sy}{103}%
}\hskip.1em\allowbreak{} 。

（4.1）证明：记\allowbreak{}\hskip.1em \RB{34.462}{5.309}{2.212}{%
\RG{0.000}{-0.002}{10.9091}{mi}{97}\RG{8.835}{-0.002}{10.9091}{rm}{61}\RG{20.372}{-0.002}{10.9091}{mi}{109}\RG{29.715}{-1.639}{7.9701}{mi}{117}%
}\hskip.1em\allowbreak{} 、\allowbreak{}\hskip.1em \RB{32.764}{8.069}{2.212}{%
\RG{0.000}{-0.002}{10.9091}{mi}{98}\RG{7.770}{-0.002}{10.9091}{rm}{61}\RG{19.306}{-0.002}{10.9091}{mi}{109}\RG{28.649}{-1.639}{7.9701}{mi}{118}%
}\hskip.1em\allowbreak{} ，并记\allowbreak{}\hskip.1em \RB{57.176}{8.680}{3.240}{%
\RG{0.000}{-0.002}{10.9091}{mi}{85}\RG{11.622}{-0.002}{10.9091}{rm}{61}\RG{23.158}{-0.002}{10.9091}{mi}{67}\RG{31.584}{-0.002}{10.9091}{rm}{40}\RG{35.738}{-0.002}{10.9091}{mi}{39}\RG{42.730}{-0.002}{10.9091}{mi}{59}\RG{47.509}{-0.002}{10.9091}{mi}{117}\RG{53.600}{-0.002}{10.9091}{rm}{41}%
}\hskip.1em\allowbreak{} 、\allowbreak{}\hskip.1em \RB{59.826}{8.680}{3.240}{%
\RG{0.000}{-0.002}{10.9091}{mi}{87}\RG{14.795}{-0.002}{10.9091}{rm}{61}\RG{26.331}{-0.002}{10.9091}{mi}{67}\RG{34.758}{-0.002}{10.9091}{rm}{40}\RG{38.911}{-0.002}{10.9091}{mi}{39}\RG{45.904}{-0.002}{10.9091}{mi}{59}\RG{50.682}{-0.002}{10.9091}{mi}{118}\RG{56.250}{-0.002}{10.9091}{rm}{41}%
}\hskip.1em\allowbreak{} 。对任意\allowbreak{}\hskip.1em \RB{28.361}{7.949}{0.938}{%
\RG{0.000}{-0.002}{10.9091}{mi}{120}\RG{9.298}{-0.002}{10.9091}{sy}{50}\RG{19.799}{-0.002}{10.9091}{mi}{85}%
}\hskip.1em\allowbreak{} ，存在\allowbreak{}\hskip.1em \RB{40.018}{8.680}{3.229}{%
\RG{0.000}{-0.002}{10.9091}{mi}{112}\RG{8.391}{-0.002}{10.9091}{sy}{50}\RG{18.695}{-0.002}{10.9091}{mi}{75}\RG{28.566}{-0.002}{10.9091}{rm}{91}\RG{31.597}{-0.002}{10.9091}{mi}{21}\RG{37.827}{-0.002}{10.9091}{rm}{93}%
}\hskip.1em\allowbreak{} ，使得\allowbreak{}\hskip.1em \RB{54.915}{8.680}{3.240}{%
\RG{0.000}{-0.002}{10.9091}{mi}{120}\RG{9.100}{-0.002}{10.9091}{rm}{61}\RG{20.438}{-0.002}{10.9091}{mi}{112}\RG{25.800}{-0.002}{10.9091}{rm}{40}\RG{29.954}{-0.002}{10.9091}{mi}{39}\RG{36.951}{-0.002}{10.9091}{rm}{41}\RG{41.108}{-0.002}{10.9091}{rm}{40}\RG{45.260}{-0.002}{10.9091}{mi}{117}\RG{51.339}{-0.002}{10.9091}{rm}{41}%
}\hskip.1em\allowbreak{} ，因而

\[
\RB{127.845}{8.678}{3.242}{%
\RG{0.000}{-0.004}{10.9091}{mi}{97}\RG{5.607}{-0.004}{10.9091}{rm}{40}\RG{9.762}{-0.004}{10.9091}{mi}{39}\RG{16.759}{-0.004}{10.9091}{rm}{41}\RG{20.916}{-0.004}{10.9091}{rm}{40}\RG{25.067}{-0.004}{10.9091}{mi}{120}\RG{31.138}{-0.004}{10.9091}{rm}{41}\RG{38.316}{-0.004}{10.9091}{rm}{61}\RG{49.660}{-0.004}{10.9091}{mi}{112}\RG{55.021}{-0.004}{10.9091}{rm}{40}\RG{59.176}{-0.004}{10.9091}{mi}{39}\RG{66.173}{-0.004}{10.9091}{rm}{41}\RG{70.326}{-0.004}{10.9091}{mi}{97}\RG{75.934}{-0.004}{10.9091}{rm}{40}\RG{80.089}{-0.004}{10.9091}{mi}{39}\RG{87.086}{-0.004}{10.9091}{rm}{41}\RG{91.242}{-0.004}{10.9091}{rm}{40}\RG{95.393}{-0.004}{10.9091}{mi}{117}\RG{101.473}{-0.004}{10.9091}{rm}{41}\RG{108.652}{-0.004}{10.9091}{rm}{61}\RG{119.997}{-0.004}{10.9091}{rm}{48}\RG{125.337}{-0.004}{10.9091}{mi}{58}%
}
\]
同理，\allowbreak{}\hskip.1em \RB{19.270}{8.686}{3.235}{%
\RG{0.000}{0.004}{10.9091}{mi}{98}\RG{4.542}{0.004}{10.9091}{rm}{40}\RG{8.696}{0.004}{10.9091}{mi}{39}\RG{15.693}{0.004}{10.9091}{rm}{41}%
}\hskip.1em\allowbreak{} 在\allowbreak{}\hskip.1em \RB{11.693}{7.955}{0.725}{%
\RG{0.000}{0.004}{10.9091}{mi}{87}%
}\hskip.1em\allowbreak{} 上为零。由于\allowbreak{}\hskip.1em \RB{46.306}{8.686}{3.235}{%
\RG{0.000}{0.004}{10.9091}{rm}{40}\RG{4.155}{0.004}{10.9091}{mi}{97}\RG{9.762}{0.004}{10.9091}{mi}{59}\RG{14.540}{0.004}{10.9091}{mi}{98}\RG{19.089}{0.004}{10.9091}{rm}{41}\RG{28.133}{0.004}{10.9091}{rm}{61}\RG{41.333}{0.004}{10.9091}{rm}{49}%
}\hskip.1em\allowbreak{} ，存在\allowbreak{}\hskip.1em \RB{56.666}{8.686}{3.224}{%
\RG{0.000}{0.004}{10.9091}{mi}{11}\RG{6.862}{0.004}{10.9091}{mi}{59}\RG{11.640}{0.004}{10.9091}{mi}{12}\RG{23.177}{0.004}{10.9091}{sy}{50}\RG{35.342}{0.004}{10.9091}{mi}{75}\RG{45.213}{0.004}{10.9091}{rm}{91}\RG{48.244}{0.004}{10.9091}{mi}{21}\RG{54.475}{0.004}{10.9091}{rm}{93}%
}\hskip.1em\allowbreak{} ，使得\allowbreak{}\hskip.1em \RB{61.360}{8.195}{2.613}{%
\RG{0.000}{0.004}{10.9091}{mi}{11}\RG{6.862}{0.004}{10.9091}{mi}{97}\RG{15.640}{0.004}{10.9091}{rm}{43}\RG{27.118}{0.004}{10.9091}{mi}{12}\RG{33.751}{0.004}{10.9091}{mi}{98}\RG{43.187}{0.004}{10.9091}{rm}{61}\RG{56.388}{0.004}{10.9091}{rm}{49}%
}\hskip.1em\allowbreak{} 。若\allowbreak{}\hskip.1em \RB{51.613}{7.946}{0.941}{%
\RG{0.000}{-0.005}{10.9091}{mi}{120}\RG{9.100}{-0.005}{10.9091}{sy}{50}\RG{19.404}{-0.005}{10.9091}{mi}{85}\RG{30.222}{-0.005}{10.9091}{sy}{92}\RG{39.920}{-0.005}{10.9091}{mi}{87}%
}\hskip.1em\allowbreak{} ，则

\[
\RB{146.645}{8.685}{18.217}{%
\RG{0.000}{0.003}{10.9091}{mi}{120}\RG{9.100}{0.003}{10.9091}{rm}{61}\RG{20.439}{0.003}{10.9091}{mi}{11}\RG{27.303}{0.003}{10.9091}{rm}{40}\RG{31.457}{0.003}{10.9091}{mi}{39}\RG{38.454}{0.003}{10.9091}{rm}{41}\RG{42.609}{0.003}{10.9091}{mi}{97}\RG{48.216}{0.003}{10.9091}{rm}{40}\RG{52.370}{0.003}{10.9091}{mi}{39}\RG{59.367}{0.003}{10.9091}{rm}{41}\RG{63.524}{0.003}{10.9091}{rm}{40}\RG{67.676}{0.003}{10.9091}{mi}{120}\RG{73.746}{0.003}{10.9091}{rm}{41}\RG{80.324}{0.003}{10.9091}{rm}{43}\RG{91.056}{0.003}{10.9091}{mi}{12}\RG{97.692}{0.003}{10.9091}{rm}{40}\RG{101.846}{0.003}{10.9091}{mi}{39}\RG{108.843}{0.003}{10.9091}{rm}{41}\RG{112.997}{0.003}{10.9091}{mi}{98}\RG{117.539}{0.003}{10.9091}{rm}{40}\RG{121.693}{0.003}{10.9091}{mi}{39}\RG{128.690}{0.003}{10.9091}{rm}{41}\RG{132.846}{0.003}{10.9091}{rm}{40}\RG{136.998}{0.003}{10.9091}{mi}{120}\RG{143.068}{0.003}{10.9091}{rm}{41}\RG{9.100}{-17.488}{10.9091}{rm}{61}\RG{20.435}{-17.488}{10.9091}{rm}{48}\RG{25.780}{-17.488}{10.9091}{mi}{58}%
}
\]
因此\allowbreak{}\hskip.1em \RB{63.758}{8.684}{3.225}{%
\RG{0.000}{0.002}{10.9091}{mi}{85}\RG{11.011}{0.002}{10.9091}{sy}{92}\RG{20.902}{0.002}{10.9091}{mi}{87}\RG{35.983}{0.002}{10.9091}{rm}{61}\RG{47.805}{0.002}{10.9091}{sy}{102}\RG{53.259}{0.002}{10.9091}{rm}{48}\RG{58.600}{0.002}{10.9091}{sy}{103}%
}\hskip.1em\allowbreak{} ，所述和为直和。若\allowbreak{}\hskip.1em \RB{45.904}{7.757}{1.338}{%
\RG{0.000}{0.002}{10.9091}{mi}{117}\RG{8.697}{0.002}{10.9091}{rm}{43}\RG{19.623}{0.002}{10.9091}{mi}{118}\RG{28.691}{0.002}{10.9091}{rm}{61}\RG{40.517}{0.002}{10.9091}{rm}{48}%
}\hskip.1em\allowbreak{} ，则\allowbreak{}\hskip.1em \RB{82.349}{7.953}{0.934}{%
\RG{0.000}{0.002}{10.9091}{mi}{117}\RG{9.593}{0.002}{10.9091}{rm}{61}\RG{21.414}{0.002}{10.9091}{sy}{0}\RG{29.899}{0.002}{10.9091}{mi}{118}\RG{38.968}{0.002}{10.9091}{sy}{50}\RG{49.754}{0.002}{10.9091}{mi}{85}\RG{60.765}{0.002}{10.9091}{sy}{92}\RG{70.656}{0.002}{10.9091}{mi}{87}%
}\hskip.1em\allowbreak{} ，与\allowbreak{}\hskip.1em \RB{26.805}{8.313}{2.855}{%
\RG{0.000}{0.002}{10.9091}{mi}{117}\RG{9.592}{0.002}{10.9091}{sy}{54}\RG{9.592}{0.002}{10.9091}{rm}{61}\RG{21.418}{0.002}{10.9091}{rm}{48}%
}\hskip.1em\allowbreak{} 矛盾，故\allowbreak{}\hskip.1em \RB{44.545}{8.314}{2.854}{%
\RG{0.000}{0.003}{10.9091}{mi}{117}\RG{8.505}{0.003}{10.9091}{rm}{43}\RG{19.237}{0.003}{10.9091}{mi}{118}\RG{27.823}{0.003}{10.9091}{sy}{54}\RG{27.823}{0.003}{10.9091}{rm}{61}\RG{39.157}{0.003}{10.9091}{rm}{48}%
}\hskip.1em\allowbreak{} 。

多项式\allowbreak{}\hskip.1em \RB{10.504}{8.068}{0.612}{%
\RG{0.000}{-0.003}{10.9091}{mi}{97}\RG{5.607}{-0.003}{10.9091}{mi}{98}%
}\hskip.1em\allowbreak{} 零化\allowbreak{}\hskip.1em \RB{24.721}{6.781}{1.343}{%
\RG{0.000}{-0.003}{10.9091}{mi}{117}\RG{8.504}{-0.003}{10.9091}{rm}{43}\RG{19.236}{-0.003}{10.9091}{mi}{118}%
}\hskip.1em\allowbreak{} ，因为

\[
\RB{203.582}{8.686}{18.216}{%
\RG{0.000}{0.004}{10.9091}{rm}{40}\RG{4.155}{0.004}{10.9091}{mi}{97}\RG{9.762}{0.004}{10.9091}{mi}{98}\RG{14.304}{0.004}{10.9091}{rm}{41}\RG{18.460}{0.004}{10.9091}{rm}{40}\RG{22.611}{0.004}{10.9091}{mi}{39}\RG{29.609}{0.004}{10.9091}{rm}{41}\RG{33.766}{0.004}{10.9091}{rm}{40}\RG{37.917}{0.004}{10.9091}{mi}{117}\RG{46.420}{0.004}{10.9091}{rm}{43}\RG{57.153}{0.004}{10.9091}{mi}{118}\RG{62.709}{0.004}{10.9091}{rm}{41}\RG{69.887}{0.004}{10.9091}{rm}{61}\RG{81.232}{0.004}{10.9091}{mi}{98}\RG{85.773}{0.004}{10.9091}{rm}{40}\RG{89.927}{0.004}{10.9091}{mi}{39}\RG{96.924}{0.004}{10.9091}{rm}{41}\RG{101.079}{0.004}{10.9091}{mi}{97}\RG{106.686}{0.004}{10.9091}{rm}{40}\RG{110.840}{0.004}{10.9091}{mi}{39}\RG{117.837}{0.004}{10.9091}{rm}{41}\RG{121.994}{0.004}{10.9091}{rm}{40}\RG{126.146}{0.004}{10.9091}{mi}{117}\RG{132.225}{0.004}{10.9091}{rm}{41}\RG{138.803}{0.004}{10.9091}{rm}{43}\RG{149.536}{0.004}{10.9091}{mi}{97}\RG{155.144}{0.004}{10.9091}{rm}{40}\RG{159.297}{0.004}{10.9091}{mi}{39}\RG{166.295}{0.004}{10.9091}{rm}{41}\RG{170.449}{0.004}{10.9091}{mi}{98}\RG{174.991}{0.004}{10.9091}{rm}{40}\RG{179.145}{0.004}{10.9091}{mi}{39}\RG{186.142}{0.004}{10.9091}{rm}{41}\RG{190.298}{0.004}{10.9091}{rm}{40}\RG{194.450}{0.004}{10.9091}{mi}{118}\RG{200.005}{0.004}{10.9091}{rm}{41}\RG{69.893}{-17.487}{10.9091}{rm}{61}\RG{81.228}{-17.487}{10.9091}{rm}{48}\RG{86.573}{-17.487}{10.9091}{mi}{58}%
}
\]
% Original page 16.
任取\allowbreak{}\hskip.1em \RB{23.836}{6.786}{1.337}{%
\RG{0.000}{0.003}{10.9091}{mi}{117}\RG{8.061}{0.003}{10.9091}{rm}{43}\RG{18.351}{0.003}{10.9091}{mi}{118}%
}\hskip.1em\allowbreak{} 的零化多项式\allowbreak{}\hskip.1em \RB{6.456}{8.183}{2.723}{%
\RG{0.000}{0.003}{10.9091}{mi}{102}%
}\hskip.1em\allowbreak{} 。等式\allowbreak{}\hskip.1em \RB{93.998}{8.685}{3.235}{%
\RG{0.000}{0.003}{10.9091}{mi}{102}\RG{6.430}{0.003}{10.9091}{rm}{40}\RG{10.584}{0.003}{10.9091}{mi}{39}\RG{17.581}{0.003}{10.9091}{rm}{41}\RG{21.738}{0.003}{10.9091}{rm}{40}\RG{25.890}{0.003}{10.9091}{mi}{117}\RG{31.969}{0.003}{10.9091}{rm}{41}\RG{39.147}{0.003}{10.9091}{rm}{61}\RG{50.492}{0.003}{10.9091}{sy}{0}\RG{58.977}{0.003}{10.9091}{mi}{102}\RG{65.407}{0.003}{10.9091}{rm}{40}\RG{69.561}{0.003}{10.9091}{mi}{39}\RG{76.558}{0.003}{10.9091}{rm}{41}\RG{80.714}{0.003}{10.9091}{rm}{40}\RG{84.866}{0.003}{10.9091}{mi}{118}\RG{90.422}{0.003}{10.9091}{rm}{41}%
}\hskip.1em\allowbreak{} 的两端均属于\allowbreak{}\hskip.1em \RB{31.323}{7.954}{0.737}{%
\RG{0.000}{0.003}{10.9091}{mi}{85}\RG{10.375}{0.003}{10.9091}{sy}{92}\RG{19.630}{0.003}{10.9091}{mi}{87}%
}\hskip.1em\allowbreak{} ，故两端为零。由（1）得\allowbreak{}\hskip.1em \RB{21.154}{8.685}{3.225}{%
\RG{0.000}{0.003}{10.9091}{mi}{97}\RG{8.637}{0.003}{10.9091}{sy}{106}\RG{14.698}{0.003}{10.9091}{mi}{102}%
}\hskip.1em\allowbreak{} 、\allowbreak{}\hskip.1em \RB{20.088}{8.685}{3.225}{%
\RG{0.000}{0.003}{10.9091}{mi}{98}\RG{7.572}{0.003}{10.9091}{sy}{106}\RG{13.632}{0.003}{10.9091}{mi}{102}%
}\hskip.1em\allowbreak{} ，再由\allowbreak{}\hskip.1em \RB{42.586}{8.685}{3.235}{%
\RG{0.000}{0.003}{10.9091}{rm}{40}\RG{4.155}{0.003}{10.9091}{mi}{97}\RG{9.762}{0.003}{10.9091}{mi}{59}\RG{14.540}{0.003}{10.9091}{mi}{98}\RG{19.089}{0.003}{10.9091}{rm}{41}\RG{26.267}{0.003}{10.9091}{rm}{61}\RG{37.613}{0.003}{10.9091}{rm}{49}%
}\hskip.1em\allowbreak{} 得\allowbreak{}\hskip.1em \RB{25.697}{8.685}{3.225}{%
\RG{0.000}{0.003}{10.9091}{mi}{97}\RG{5.607}{0.003}{10.9091}{mi}{98}\RG{13.179}{0.003}{10.9091}{sy}{106}\RG{19.240}{0.003}{10.9091}{mi}{102}%
}\hskip.1em\allowbreak{} 。因此\allowbreak{}\hskip.1em \RB{10.504}{8.074}{0.606}{%
\RG{0.000}{0.003}{10.9091}{mi}{97}\RG{5.607}{0.003}{10.9091}{mi}{98}%
}\hskip.1em\allowbreak{} 是\allowbreak{}\hskip.1em \RB{24.722}{6.786}{1.337}{%
\RG{0.000}{0.003}{10.9091}{mi}{117}\RG{8.504}{0.003}{10.9091}{rm}{43}\RG{19.237}{0.003}{10.9091}{mi}{118}%
}\hskip.1em\allowbreak{} 的首一极小多项式。

此外，由（2）及直和的维数公式，

\[
\RB{149.320}{8.684}{38.218}{%
\RG{0.000}{0.002}{10.9091}{rm}{100}\RG{5.935}{0.002}{10.9091}{rm}{105}\RG{8.902}{0.002}{10.9091}{rm}{109}\RG{19.675}{0.002}{10.9091}{mi}{67}\RG{28.100}{0.002}{10.9091}{rm}{40}\RG{32.255}{0.002}{10.9091}{mi}{39}\RG{39.247}{0.002}{10.9091}{mi}{59}\RG{44.026}{0.002}{10.9091}{mi}{117}\RG{52.541}{0.002}{10.9091}{rm}{43}\RG{63.274}{0.002}{10.9091}{mi}{118}\RG{68.829}{0.002}{10.9091}{rm}{41}\RG{76.007}{0.002}{10.9091}{rm}{61}\RG{87.353}{0.002}{10.9091}{rm}{100}\RG{93.288}{0.002}{10.9091}{rm}{101}\RG{98.033}{0.002}{10.9091}{rm}{103}\RG{103.521}{0.002}{10.9091}{rm}{40}\RG{107.673}{0.002}{10.9091}{mi}{97}\RG{113.280}{0.002}{10.9091}{mi}{98}\RG{117.822}{0.002}{10.9091}{rm}{41}\RG{76.013}{-17.489}{10.9091}{rm}{61}\RG{87.348}{-17.489}{10.9091}{rm}{100}\RG{93.283}{-17.489}{10.9091}{rm}{101}\RG{98.028}{-17.489}{10.9091}{rm}{103}\RG{105.337}{-17.489}{10.9091}{mi}{97}\RG{113.369}{-17.489}{10.9091}{rm}{43}\RG{124.103}{-17.489}{10.9091}{rm}{100}\RG{130.038}{-17.489}{10.9091}{rm}{101}\RG{134.784}{-17.489}{10.9091}{rm}{103}\RG{142.086}{-17.489}{10.9091}{mi}{98}\RG{76.013}{-34.980}{10.9091}{rm}{61}\RG{87.348}{-34.980}{10.9091}{rm}{100}\RG{93.283}{-34.980}{10.9091}{rm}{105}\RG{96.250}{-34.980}{10.9091}{rm}{109}\RG{105.206}{-34.980}{10.9091}{rm}{40}\RG{109.362}{-34.980}{10.9091}{mi}{85}\RG{120.181}{-34.980}{10.9091}{sy}{8}\RG{131.090}{-34.980}{10.9091}{mi}{87}\RG{142.658}{-34.980}{10.9091}{rm}{41}\RG{146.812}{-34.980}{10.9091}{mi}{58}%
}
\]
结合\allowbreak{}\hskip.1em \RB{101.274}{8.683}{3.237}{%
\RG{0.000}{0.001}{10.9091}{mi}{67}\RG{8.425}{0.001}{10.9091}{rm}{40}\RG{12.580}{0.001}{10.9091}{mi}{39}\RG{19.573}{0.001}{10.9091}{mi}{59}\RG{24.351}{0.001}{10.9091}{mi}{117}\RG{32.866}{0.001}{10.9091}{rm}{43}\RG{43.599}{0.001}{10.9091}{mi}{118}\RG{49.154}{0.001}{10.9091}{rm}{41}\RG{56.339}{0.001}{10.9091}{sy}{18}\RG{67.854}{0.001}{10.9091}{mi}{85}\RG{78.672}{0.001}{10.9091}{sy}{8}\RG{89.581}{0.001}{10.9091}{mi}{87}%
}\hskip.1em\allowbreak{} ，得到\allowbreak{}\hskip.1em \RB{101.098}{8.683}{3.237}{%
\RG{0.000}{0.001}{10.9091}{mi}{67}\RG{8.426}{0.001}{10.9091}{rm}{40}\RG{12.580}{0.001}{10.9091}{mi}{39}\RG{19.573}{0.001}{10.9091}{mi}{59}\RG{24.351}{0.001}{10.9091}{mi}{117}\RG{32.867}{0.001}{10.9091}{rm}{43}\RG{43.599}{0.001}{10.9091}{mi}{118}\RG{49.154}{0.001}{10.9091}{rm}{41}\RG{56.333}{0.001}{10.9091}{rm}{61}\RG{67.677}{0.001}{10.9091}{mi}{85}\RG{78.496}{0.001}{10.9091}{sy}{8}\RG{89.405}{0.001}{10.9091}{mi}{87}%
}\hskip.1em\allowbreak{} 。

（4.2）证明：令\allowbreak{}\hskip.1em \RB{34.446}{5.316}{2.204}{%
\RG{0.000}{0.005}{10.9091}{mi}{97}\RG{8.827}{0.005}{10.9091}{rm}{61}\RG{20.355}{0.005}{10.9091}{mi}{109}\RG{29.699}{-1.632}{7.9701}{mi}{117}%
}\hskip.1em\allowbreak{} 、\allowbreak{}\hskip.1em \RB{32.748}{8.076}{2.204}{%
\RG{0.000}{0.005}{10.9091}{mi}{98}\RG{7.762}{0.005}{10.9091}{rm}{61}\RG{19.290}{0.005}{10.9091}{mi}{109}\RG{28.634}{-1.632}{7.9701}{mi}{118}%
}\hskip.1em\allowbreak{} 。由多项式的唯一因式分解，可取互异首一不可约多项式\allowbreak{}\hskip.1em \RB{42.828}{5.307}{2.620}{%
\RG{0.000}{-0.004}{10.9091}{mi}{112}\RG{5.362}{-1.641}{7.9701}{rm}{49}\RG{9.680}{-0.004}{10.9091}{mi}{59}\RG{14.458}{-0.004}{10.9091}{mi}{58}\RG{19.247}{-0.004}{10.9091}{mi}{58}\RG{24.026}{-0.004}{10.9091}{mi}{58}\RG{28.815}{-0.004}{10.9091}{mi}{59}\RG{33.604}{-0.004}{10.9091}{mi}{112}\RG{38.968}{-1.641}{7.9701}{mi}{114}%
}\hskip.1em\allowbreak{} 及非负整数\allowbreak{}\hskip.1em \RB{23.500}{8.187}{2.620}{%
\RG{0.000}{-0.004}{10.9091}{mi}{11}\RG{6.794}{-1.641}{7.9701}{mi}{105}\RG{9.879}{-0.004}{10.9091}{mi}{59}\RG{14.657}{-0.004}{10.9091}{mi}{12}\RG{20.695}{-1.641}{7.9701}{mi}{105}%
}\hskip.1em\allowbreak{} ，使得

\[
\RB{286.995}{15.525}{11.442}{%
\RG{24.377}{11.818}{7.9701}{mi}{114}\RG{97.879}{11.818}{7.9701}{mi}{114}\RG{20.420}{9.090}{9.0910}{ex}{89}\RG{93.921}{9.090}{9.0910}{ex}{89}\RG{113.161}{5.171}{7.9701}{mi}{12}\RG{39.660}{4.866}{7.9701}{mi}{11}\RG{117.497}{4.062}{5.9776}{mi}{105}\RG{44.623}{3.757}{5.9776}{mi}{105}\RG{0.000}{-0.000}{10.9091}{mi}{97}\RG{8.637}{-0.000}{10.9091}{rm}{61}\RG{34.299}{-0.000}{10.9091}{mi}{112}\RG{47.963}{-0.000}{10.9091}{mi}{59}\RG{74.559}{-0.000}{10.9091}{mi}{98}\RG{82.138}{-0.000}{10.9091}{rm}{61}\RG{107.800}{-0.000}{10.9091}{mi}{112}\RG{120.837}{-0.000}{10.9091}{mi}{59}\RG{147.440}{-0.000}{10.9091}{rm}{109}\RG{156.342}{-0.000}{10.9091}{rm}{97}\RG{161.688}{-0.000}{10.9091}{rm}{120}\RG{167.506}{-0.000}{10.9091}{sy}{102}\RG{172.961}{-0.000}{10.9091}{mi}{11}\RG{182.841}{-0.000}{10.9091}{mi}{59}\RG{187.619}{-0.000}{10.9091}{mi}{12}\RG{196.742}{-0.000}{10.9091}{sy}{103}\RG{205.219}{-0.000}{10.9091}{sy}{21}\RG{216.743}{-0.000}{10.9091}{rm}{49}\RG{232.987}{-0.000}{10.9091}{rm}{40}\RG{237.143}{-0.000}{10.9091}{rm}{49}\RG{245.518}{-0.000}{10.9091}{sy}{20}\RG{257.033}{-0.000}{10.9091}{mi}{105}\RG{263.708}{-0.000}{10.9091}{sy}{20}\RG{275.223}{-0.000}{10.9091}{mi}{114}\RG{280.333}{-0.000}{10.9091}{rm}{41}\RG{284.488}{-0.000}{10.9091}{mi}{58}\RG{179.756}{-1.637}{7.9701}{mi}{105}\RG{193.657}{-1.637}{7.9701}{mi}{105}\RG{39.660}{-2.505}{7.9701}{mi}{105}\RG{113.161}{-2.505}{7.9701}{mi}{105}\RG{19.976}{-10.869}{7.9701}{mi}{105}\RG{22.607}{-10.869}{7.9701}{rm}{61}\RG{28.614}{-10.869}{7.9701}{rm}{49}\RG{93.477}{-10.869}{7.9701}{mi}{105}\RG{96.108}{-10.869}{7.9701}{rm}{61}\RG{102.115}{-10.869}{7.9701}{rm}{49}%
}
\]
由于\allowbreak{}\hskip.1em \RB{15.282}{8.073}{2.604}{%
\RG{0.000}{0.002}{10.9091}{mi}{97}\RG{5.607}{0.002}{10.9091}{mi}{59}\RG{10.386}{0.002}{10.9091}{mi}{98}%
}\hskip.1em\allowbreak{} 均非常数，\allowbreak{}\hskip.1em \RB{26.316}{7.756}{1.993}{%
\RG{0.000}{0.002}{10.9091}{mi}{114}\RG{8.985}{0.002}{10.9091}{sy}{21}\RG{21.344}{0.002}{10.9091}{rm}{49}%
}\hskip.1em\allowbreak{} 。对每个\allowbreak{}\hskip.1em \RB{3.587}{7.669}{0.607}{%
\RG{0.000}{0.002}{10.9091}{mi}{105}%
}\hskip.1em\allowbreak{} ，记\allowbreak{}\hskip.1em \RB{78.553}{8.684}{3.226}{%
\RG{0.000}{0.002}{10.9091}{mi}{101}\RG{4.951}{-1.634}{7.9701}{mi}{105}\RG{11.910}{0.002}{10.9091}{rm}{61}\RG{24.096}{0.002}{10.9091}{rm}{109}\RG{32.998}{0.002}{10.9091}{rm}{97}\RG{38.343}{0.002}{10.9091}{rm}{120}\RG{44.158}{0.002}{10.9091}{sy}{102}\RG{49.613}{0.002}{10.9091}{mi}{11}\RG{56.407}{-1.634}{7.9701}{mi}{105}\RG{59.493}{0.002}{10.9091}{mi}{59}\RG{64.271}{0.002}{10.9091}{mi}{12}\RG{70.308}{-1.634}{7.9701}{mi}{105}\RG{73.394}{0.002}{10.9091}{sy}{103}%
}\hskip.1em\allowbreak{} 。当\allowbreak{}\hskip.1em \RB{34.949}{8.193}{2.615}{%
\RG{0.000}{0.002}{10.9091}{mi}{11}\RG{6.795}{-1.634}{7.9701}{mi}{105}\RG{13.753}{0.002}{10.9091}{sy}{21}\RG{26.112}{0.002}{10.9091}{mi}{12}\RG{32.143}{-1.634}{7.9701}{mi}{105}%
}\hskip.1em\allowbreak{} 时取\allowbreak{}\hskip.1em \RB{30.363}{5.313}{2.207}{%
\RG{0.000}{0.002}{10.9091}{mi}{122}\RG{4.962}{-1.634}{7.9701}{mi}{105}\RG{11.921}{0.002}{10.9091}{rm}{61}\RG{24.103}{0.002}{10.9091}{mi}{117}%
}\hskip.1em\allowbreak{} ，当\allowbreak{}\hskip.1em \RB{33.084}{8.194}{2.614}{%
\RG{0.000}{0.003}{10.9091}{mi}{11}\RG{6.795}{-1.634}{7.9701}{mi}{105}\RG{12.910}{0.003}{10.9091}{mi}{60}\RG{24.245}{0.003}{10.9091}{mi}{12}\RG{30.279}{-1.634}{7.9701}{mi}{105}%
}\hskip.1em\allowbreak{} 时取\allowbreak{}\hskip.1em \RB{27.901}{5.314}{2.207}{%
\RG{0.000}{0.003}{10.9091}{mi}{122}\RG{4.962}{-1.634}{7.9701}{mi}{105}\RG{11.078}{0.003}{10.9091}{rm}{61}\RG{22.416}{0.003}{10.9091}{mi}{118}%
}\hskip.1em\allowbreak{} ，并定义

\[
\RB{134.334}{12.689}{10.563}{%
\RG{23.282}{7.378}{10.9091}{mi}{109}\RG{32.625}{5.742}{7.9701}{mi}{122}\RG{36.219}{4.633}{5.9776}{mi}{105}\RG{0.000}{-0.002}{10.9091}{mi}{113}\RG{10.853}{-0.002}{10.9091}{rm}{61}\RG{40.650}{-0.002}{10.9091}{mi}{59}\RG{67.247}{-0.002}{10.9091}{mi}{119}\RG{81.006}{-0.002}{10.9091}{rm}{61}\RG{92.345}{-0.002}{10.9091}{mi}{113}\RG{100.166}{-0.002}{10.9091}{rm}{40}\RG{104.319}{-0.002}{10.9091}{mi}{39}\RG{111.317}{-0.002}{10.9091}{rm}{41}\RG{115.474}{-0.002}{10.9091}{rm}{40}\RG{119.625}{-0.002}{10.9091}{mi}{122}\RG{127.672}{-0.002}{10.9091}{rm}{41}\RG{131.827}{-0.002}{10.9091}{mi}{58}\RG{4.737}{-1.639}{7.9701}{mi}{105}\RG{74.890}{-1.639}{7.9701}{mi}{105}\RG{97.080}{-1.639}{7.9701}{mi}{105}\RG{124.587}{-1.639}{7.9701}{mi}{105}\RG{30.639}{-2.619}{7.9701}{mi}{101}\RG{34.224}{-3.728}{5.9776}{mi}{105}\RG{25.278}{-7.485}{10.9091}{mi}{112}\RG{30.639}{-9.990}{7.9701}{mi}{105}\RR{23.282}{2.725}{16.277}{0.364}%
}
\]
选取\allowbreak{}\hskip.1em \RB{7.768}{5.307}{2.213}{%
\RG{0.000}{-0.004}{10.9091}{mi}{122}\RG{4.962}{-1.640}{7.9701}{mi}{105}%
}\hskip.1em\allowbreak{} 的方式保证\allowbreak{}\hskip.1em \RB{44.646}{8.678}{3.231}{%
\RG{0.000}{-0.004}{10.9091}{mi}{113}\RG{4.737}{-1.640}{7.9701}{mi}{105}\RG{11.936}{-0.004}{10.9091}{sy}{50}\RG{23.322}{-0.004}{10.9091}{mi}{75}\RG{33.193}{-0.004}{10.9091}{rm}{91}\RG{36.223}{-0.004}{10.9091}{mi}{21}\RG{42.455}{-0.004}{10.9091}{rm}{93}%
}\hskip.1em\allowbreak{} 。由（3.3）的构造，\allowbreak{}\hskip.1em \RB{32.648}{8.307}{2.860}{%
\RG{0.000}{-0.004}{10.9091}{mi}{119}\RG{7.635}{-1.640}{7.9701}{mi}{105}\RG{14.834}{-0.004}{10.9091}{sy}{54}\RG{14.835}{-0.004}{10.9091}{rm}{61}\RG{27.260}{-0.004}{10.9091}{rm}{48}%
}\hskip.1em\allowbreak{} ，且\allowbreak{}\hskip.1em \RB{46.304}{8.570}{3.303}{%
\RG{0.000}{-0.004}{10.9091}{mi}{109}\RG{9.343}{-1.640}{7.9701}{mi}{119}\RG{14.889}{-2.748}{5.9776}{mi}{105}\RG{22.341}{-0.004}{10.9091}{rm}{61}\RG{34.763}{-0.004}{10.9091}{mi}{112}\RG{40.124}{4.863}{7.9701}{mi}{101}\RG{40.124}{-2.509}{7.9701}{mi}{105}\RG{43.709}{3.754}{5.9776}{mi}{105}%
}\hskip.1em\allowbreak{} 。各\allowbreak{}\hskip.1em \RB{10.441}{5.307}{2.213}{%
\RG{0.000}{-0.004}{10.9091}{mi}{119}\RG{7.636}{-1.640}{7.9701}{mi}{105}%
}\hskip.1em\allowbreak{} 均属于

\[
\RB{93.234}{9.225}{3.241}{%
\RG{0.000}{-0.003}{10.9091}{mi}{67}\RG{8.426}{-0.003}{10.9091}{rm}{40}\RG{12.580}{-0.003}{10.9091}{mi}{39}\RG{19.573}{-0.003}{10.9091}{mi}{59}\RG{24.351}{-0.003}{10.9091}{mi}{117}\RG{30.442}{-0.003}{10.9091}{rm}{41}\RG{37.021}{-0.003}{10.9091}{rm}{43}\RG{47.753}{-0.003}{10.9091}{mi}{67}\RG{56.179}{-0.003}{10.9091}{rm}{40}\RG{60.333}{-0.003}{10.9091}{mi}{39}\RG{67.325}{-0.003}{10.9091}{mi}{59}\RG{72.104}{-0.003}{10.9091}{mi}{118}\RG{77.671}{-0.003}{10.9091}{rm}{41}\RT{81.825}{-0.003}{10.9091}{。}%
}
\]
定义\allowbreak{}\hskip.1em \RB{58.214}{10.487}{3.620}{%
\RG{0.000}{-0.000}{10.9091}{mi}{115}\RG{5.032}{-1.636}{7.9701}{mi}{106}\RG{12.061}{-0.000}{10.9091}{rm}{61}\RG{23.400}{7.273}{9.0910}{ex}{80}\RG{32.996}{5.172}{7.9701}{mi}{106}\RG{32.996}{-2.505}{7.9701}{mi}{105}\RG{35.626}{-2.505}{7.9701}{rm}{61}\RG{41.634}{-2.505}{7.9701}{rm}{49}\RG{47.773}{-0.000}{10.9091}{mi}{119}\RG{55.409}{-1.636}{7.9701}{mi}{105}%
}\hskip.1em\allowbreak{} ，其中\allowbreak{}\hskip.1em \RB{44.612}{7.755}{2.725}{%
\RG{0.000}{-0.000}{10.9091}{rm}{49}\RG{8.370}{-0.000}{10.9091}{sy}{20}\RG{19.886}{-0.000}{10.9091}{mi}{106}\RG{27.949}{-0.000}{10.9091}{sy}{20}\RG{39.464}{-0.000}{10.9091}{mi}{114}%
}\hskip.1em\allowbreak{} 。对\allowbreak{}\hskip.1em \RB{4.766}{7.667}{2.725}{%
\RG{0.000}{-0.000}{10.9091}{mi}{106}%
}\hskip.1em\allowbreak{} 归纳证明

\[
\RB{115.761}{17.743}{11.438}{%
\RG{91.125}{12.428}{7.9701}{mi}{106}\RG{87.089}{9.095}{9.0910}{ex}{89}\RG{106.330}{4.870}{7.9701}{mi}{101}\RG{109.914}{3.761}{5.9776}{mi}{105}\RG{0.000}{0.004}{10.9091}{mi}{115}\RG{12.061}{0.004}{10.9091}{sy}{54}\RG{12.061}{0.004}{10.9091}{rm}{61}\RG{23.395}{0.004}{10.9091}{rm}{48}\RG{28.741}{0.004}{10.9091}{mi}{59}\RG{55.337}{0.004}{10.9091}{mi}{109}\RG{75.307}{0.004}{10.9091}{rm}{61}\RG{100.969}{0.004}{10.9091}{mi}{112}\RG{113.254}{0.004}{10.9091}{mi}{58}\RG{5.032}{-1.632}{7.9701}{mi}{106}\RG{64.687}{-1.632}{7.9701}{mi}{115}\RG{106.330}{-2.501}{7.9701}{mi}{105}\RG{68.261}{-2.741}{5.9776}{mi}{106}\RG{86.646}{-10.865}{7.9701}{mi}{105}\RG{89.277}{-10.865}{7.9701}{rm}{61}\RG{95.284}{-10.865}{7.9701}{rm}{49}%
}
\]
当\allowbreak{}\hskip.1em \RB{24.957}{7.753}{2.727}{%
\RG{0.000}{-0.002}{10.9091}{mi}{106}\RG{8.355}{-0.002}{10.9091}{rm}{61}\RG{19.984}{-0.002}{10.9091}{rm}{49}%
}\hskip.1em\allowbreak{} 时已由构造得到。设结论对\allowbreak{}\hskip.1em \RB{23.573}{7.753}{2.727}{%
\RG{0.000}{-0.002}{10.9091}{mi}{106}\RG{7.574}{-0.002}{10.9091}{sy}{0}\RG{18.600}{-0.002}{10.9091}{rm}{49}%
}\hskip.1em\allowbreak{} 成立，其中\allowbreak{}\hskip.1em \RB{45.779}{7.753}{2.727}{%
\RG{0.000}{-0.002}{10.9091}{rm}{50}\RG{8.662}{-0.002}{10.9091}{sy}{20}\RG{20.469}{-0.002}{10.9091}{mi}{106}\RG{28.824}{-0.002}{10.9091}{sy}{20}\RG{40.631}{-0.002}{10.9091}{mi}{114}%
}\hskip.1em\allowbreak{} 。由于\allowbreak{}\hskip.1em \RB{8.166}{5.309}{2.618}{%
\RG{0.000}{-0.002}{10.9091}{mi}{112}\RG{5.361}{-1.638}{7.9701}{mi}{105}%
}\hskip.1em\allowbreak{} 互异且不可约，\allowbreak{}\hskip.1em \RB{35.817}{10.507}{3.078}{%
\RG{0.000}{7.272}{9.0910}{ex}{81}\RG{8.586}{5.170}{7.9701}{mi}{106}\RG{8.586}{-2.506}{7.9701}{mi}{105}\RG{11.217}{-2.506}{7.9701}{rm}{61}\RG{12.130}{5.170}{7.9701}{sy}{0}\RG{17.224}{-2.506}{7.9701}{rm}{49}\RG{18.141}{5.170}{7.9701}{rm}{49}\RG{24.277}{-0.001}{10.9091}{mi}{112}\RG{29.638}{4.866}{7.9701}{mi}{101}\RG{29.638}{-2.506}{7.9701}{mi}{105}\RG{33.223}{3.757}{5.9776}{mi}{105}%
}\hskip.1em\allowbreak{} 与\allowbreak{}\hskip.1em \RB{11.977}{9.633}{4.489}{%
\RG{0.000}{-0.001}{10.9091}{mi}{112}\RG{5.361}{5.926}{7.9701}{mi}{101}\RG{5.361}{-2.506}{7.9701}{mi}{106}\RG{8.946}{4.817}{5.9776}{mi}{106}%
}\hskip.1em\allowbreak{} 互素。对\allowbreak{}\hskip.1em \RB{18.337}{5.310}{3.621}{%
\RG{0.000}{-0.001}{10.9091}{mi}{115}\RG{5.032}{-1.637}{7.9701}{mi}{106}\RG{8.576}{-1.637}{7.9701}{sy}{0}\RG{14.586}{-1.637}{7.9701}{rm}{49}%
}\hskip.1em\allowbreak{} 和\allowbreak{}\hskip.1em \RB{11.140}{5.310}{3.621}{%
\RG{0.000}{-0.001}{10.9091}{mi}{119}\RG{7.636}{-1.637}{7.9701}{mi}{106}%
}\hskip.1em\allowbreak{} 应用（4.1），得到\allowbreak{}\hskip.1em \RB{28.783}{8.310}{3.621}{%
\RG{0.000}{-0.001}{10.9091}{mi}{115}\RG{5.032}{-1.637}{7.9701}{mi}{106}\RG{12.061}{-0.001}{10.9091}{sy}{54}\RG{12.061}{-0.001}{10.9091}{rm}{61}\RG{23.395}{-0.001}{10.9091}{rm}{48}%
}\hskip.1em\allowbreak{} 及所述极小多项式公式。

取\allowbreak{}\hskip.1em \RB{31.190}{5.313}{2.207}{%
\RG{0.000}{0.002}{10.9091}{mi}{119}\RG{10.959}{0.002}{10.9091}{rm}{61}\RG{22.298}{0.002}{10.9091}{mi}{115}\RG{27.330}{-1.634}{7.9701}{mi}{114}%
}\hskip.1em\allowbreak{} ，则\allowbreak{}\hskip.1em \RB{27.681}{8.313}{2.855}{%
\RG{0.000}{0.002}{10.9091}{mi}{119}\RG{10.959}{0.002}{10.9091}{sy}{54}\RG{10.959}{0.002}{10.9091}{rm}{61}\RG{22.294}{0.002}{10.9091}{rm}{48}%
}\hskip.1em\allowbreak{} ，\allowbreak{}\hskip.1em \RB{102.510}{8.684}{3.236}{%
\RG{0.000}{0.002}{10.9091}{mi}{119}\RG{10.960}{0.002}{10.9091}{sy}{50}\RG{21.263}{0.002}{10.9091}{mi}{67}\RG{29.689}{0.002}{10.9091}{rm}{40}\RG{33.842}{0.002}{10.9091}{mi}{39}\RG{40.835}{0.002}{10.9091}{mi}{59}\RG{45.613}{0.002}{10.9091}{mi}{117}\RG{51.705}{0.002}{10.9091}{rm}{41}\RG{58.283}{0.002}{10.9091}{rm}{43}\RG{69.016}{0.002}{10.9091}{mi}{67}\RG{77.441}{0.002}{10.9091}{rm}{40}\RG{81.596}{0.002}{10.9091}{mi}{39}\RG{88.589}{0.002}{10.9091}{mi}{59}\RG{93.367}{0.002}{10.9091}{mi}{118}\RG{98.934}{0.002}{10.9091}{rm}{41}%
}\hskip.1em\allowbreak{} ，且

\[
\RB{127.260}{15.527}{11.440}{%
\RG{34.310}{11.820}{7.9701}{mi}{114}\RG{30.352}{9.093}{9.0910}{ex}{89}\RG{49.592}{5.577}{7.9701}{rm}{109}\RG{56.029}{5.577}{7.9701}{rm}{97}\RG{59.891}{5.577}{7.9701}{rm}{120}\RG{64.054}{5.577}{7.9701}{sy}{102}\RG{67.917}{5.577}{7.9701}{mi}{11}\RG{75.766}{5.577}{7.9701}{mi}{59}\RG{77.911}{5.577}{7.9701}{mi}{12}\RG{85.134}{5.577}{7.9701}{sy}{103}\RG{72.881}{4.469}{5.9776}{mi}{105}\RG{82.248}{4.469}{5.9776}{mi}{105}\RG{0.000}{0.002}{10.9091}{mi}{109}\RG{18.570}{0.002}{10.9091}{rm}{61}\RG{44.231}{0.002}{10.9091}{mi}{112}\RG{92.482}{0.002}{10.9091}{rm}{61}\RG{103.817}{0.002}{10.9091}{rm}{91}\RG{106.851}{0.002}{10.9091}{mi}{97}\RG{112.458}{0.002}{10.9091}{mi}{59}\RG{117.236}{0.002}{10.9091}{mi}{98}\RG{121.785}{0.002}{10.9091}{rm}{93}\RG{124.752}{0.002}{10.9091}{mi}{58}\RG{9.344}{-1.634}{7.9701}{mi}{119}\RG{49.592}{-2.503}{7.9701}{mi}{105}\RG{29.909}{-10.867}{7.9701}{mi}{105}\RG{32.539}{-10.867}{7.9701}{rm}{61}\RG{38.547}{-10.867}{7.9701}{rm}{49}%
}
\]
（5）证明：令\allowbreak{}\hskip.1em \RB{78.933}{8.684}{3.225}{%
\RG{0.000}{0.002}{10.9091}{mi}{76}\RG{10.298}{0.002}{10.9091}{rm}{61}\RG{21.632}{0.002}{10.9091}{rm}{91}\RG{24.667}{0.002}{10.9091}{mi}{109}\RG{34.010}{-1.634}{7.9701}{rm}{49}\RG{38.328}{0.002}{10.9091}{mi}{59}\RG{43.107}{0.002}{10.9091}{mi}{58}\RG{47.896}{0.002}{10.9091}{mi}{58}\RG{52.685}{0.002}{10.9091}{mi}{58}\RG{57.463}{0.002}{10.9091}{mi}{59}\RG{62.252}{0.002}{10.9091}{mi}{109}\RG{71.599}{-1.634}{7.9701}{mi}{110}\RG{76.742}{0.002}{10.9091}{rm}{93}%
}\hskip.1em\allowbreak{} 。由于\allowbreak{}\hskip.1em \RB{44.406}{8.684}{3.236}{%
\RG{0.000}{0.002}{10.9091}{mi}{109}\RG{9.343}{0.002}{10.9091}{rm}{40}\RG{13.498}{0.002}{10.9091}{mi}{39}\RG{20.495}{0.002}{10.9091}{rm}{41}\RG{27.673}{0.002}{10.9091}{rm}{61}\RG{39.019}{0.002}{10.9091}{rm}{48}%
}\hskip.1em\allowbreak{} ，对每个\allowbreak{}\hskip.1em \RB{3.587}{7.669}{0.607}{%
\RG{0.000}{0.002}{10.9091}{mi}{105}%
}\hskip.1em\allowbreak{} 均有\allowbreak{}\hskip.1em \RB{60.964}{8.684}{3.236}{%
\RG{0.000}{0.002}{10.9091}{mi}{109}\RG{9.343}{0.002}{10.9091}{rm}{40}\RG{13.497}{0.002}{10.9091}{mi}{39}\RG{20.495}{0.002}{10.9091}{rm}{41}\RG{24.651}{0.002}{10.9091}{rm}{40}\RG{28.803}{0.002}{10.9091}{mi}{118}\RG{33.967}{-1.634}{7.9701}{mi}{105}\RG{37.053}{0.002}{10.9091}{rm}{41}\RG{44.231}{0.002}{10.9091}{rm}{61}\RG{55.577}{0.002}{10.9091}{rm}{48}%
}\hskip.1em\allowbreak{} 。由（1）得\allowbreak{}\hskip.1em \RB{31.042}{8.685}{3.225}{%
\RG{0.000}{0.003}{10.9091}{mi}{109}\RG{9.343}{-1.633}{7.9701}{mi}{105}\RG{15.459}{0.003}{10.9091}{sy}{106}\RG{21.520}{0.003}{10.9091}{mi}{109}%
}\hskip.1em\allowbreak{} ，故\allowbreak{}\hskip.1em \RB{25.881}{8.685}{3.225}{%
\RG{0.000}{0.003}{10.9091}{mi}{76}\RG{10.298}{0.003}{10.9091}{sy}{106}\RG{16.359}{0.003}{10.9091}{mi}{109}%
}\hskip.1em\allowbreak{} 。

另一方面，\allowbreak{}\hskip.1em \RB{28.904}{8.678}{3.231}{%
\RG{0.000}{-0.004}{10.9091}{mi}{109}\RG{9.343}{-1.640}{7.9701}{mi}{105}\RG{15.459}{-0.004}{10.9091}{sy}{106}\RG{21.520}{-0.004}{10.9091}{mi}{76}%
}\hskip.1em\allowbreak{} 给出\allowbreak{}\hskip.1em \RB{58.888}{8.678}{3.242}{%
\RG{0.000}{-0.004}{10.9091}{mi}{76}\RG{7.268}{-0.004}{10.9091}{rm}{40}\RG{11.422}{-0.004}{10.9091}{mi}{39}\RG{18.419}{-0.004}{10.9091}{rm}{41}\RG{22.576}{-0.004}{10.9091}{rm}{40}\RG{26.727}{-0.004}{10.9091}{mi}{118}\RG{31.891}{-1.640}{7.9701}{mi}{105}\RG{34.977}{-0.004}{10.9091}{rm}{41}\RG{42.155}{-0.004}{10.9091}{rm}{61}\RG{53.501}{-0.004}{10.9091}{rm}{48}%
}\hskip.1em\allowbreak{} 。由于\allowbreak{}\hskip.1em \RB{43.525}{5.307}{2.609}{%
\RG{0.000}{-0.004}{10.9091}{mi}{118}\RG{5.164}{-1.640}{7.9701}{rm}{49}\RG{9.483}{-0.004}{10.9091}{mi}{59}\RG{14.261}{-0.004}{10.9091}{mi}{58}\RG{19.050}{-0.004}{10.9091}{mi}{58}\RG{23.829}{-0.004}{10.9091}{mi}{58}\RG{28.618}{-0.004}{10.9091}{mi}{59}\RG{33.407}{-0.004}{10.9091}{mi}{118}\RG{38.574}{-1.640}{7.9701}{mi}{110}%
}\hskip.1em\allowbreak{} 构成\allowbreak{}\hskip.1em \RB{8.715}{7.947}{0.733}{%
\RG{0.000}{-0.004}{10.9091}{mi}{86}%
}\hskip.1em\allowbreak{} 的一组基，线性变换\allowbreak{}\hskip.1em \RB{21.995}{8.678}{3.242}{%
\RG{0.000}{-0.004}{10.9091}{mi}{76}\RG{7.268}{-0.004}{10.9091}{rm}{40}\RG{11.422}{-0.004}{10.9091}{mi}{39}\RG{18.419}{-0.004}{10.9091}{rm}{41}%
}\hskip.1em\allowbreak{}在整个\allowbreak{}\hskip.1em \RB{8.715}{7.948}{0.732}{%
\RG{0.000}{-0.003}{10.9091}{mi}{86}%
}\hskip.1em\allowbreak{} 上为零。由变换极小多项式的次数最小性，\allowbreak{}\hskip.1em \RB{67.243}{8.068}{2.739}{%
\RG{0.000}{-0.003}{10.9091}{rm}{100}\RG{5.935}{-0.003}{10.9091}{rm}{101}\RG{10.680}{-0.003}{10.9091}{rm}{103}\RG{17.985}{-0.003}{10.9091}{mi}{109}\RG{30.358}{-0.003}{10.9091}{sy}{20}\RG{41.874}{-0.003}{10.9091}{rm}{100}\RG{47.809}{-0.003}{10.9091}{rm}{101}\RG{52.554}{-0.003}{10.9091}{rm}{103}\RG{59.859}{-0.003}{10.9091}{mi}{76}%
}\hskip.1em\allowbreak{} 。结合\allowbreak{}\hskip.1em \RB{25.881}{8.679}{3.230}{%
\RG{0.000}{-0.003}{10.9091}{mi}{76}\RG{10.298}{-0.003}{10.9091}{sy}{106}\RG{16.359}{-0.003}{10.9091}{mi}{109}%
}\hskip.1em\allowbreak{} 可知二者次数相等；再由二者首一，得到\allowbreak{}\hskip.1em \RB{31.095}{7.949}{0.611}{%
\RG{0.000}{-0.002}{10.9091}{mi}{109}\RG{12.374}{-0.002}{10.9091}{rm}{61}\RG{23.712}{-0.002}{10.9091}{mi}{76}%
}\hskip.1em\allowbreak{} 。

（6）证明：取\allowbreak{}\hskip.1em \RB{8.715}{7.953}{0.727}{%
\RG{0.000}{0.002}{10.9091}{mi}{86}%
}\hskip.1em\allowbreak{} 的一组基\allowbreak{}\hskip.1em \RB{43.525}{5.313}{2.604}{%
\RG{0.000}{0.002}{10.9091}{mi}{118}\RG{5.164}{-1.634}{7.9701}{rm}{49}\RG{9.483}{0.002}{10.9091}{mi}{59}\RG{14.261}{0.002}{10.9091}{mi}{58}\RG{19.050}{0.002}{10.9091}{mi}{58}\RG{23.829}{0.002}{10.9091}{mi}{58}\RG{28.618}{0.002}{10.9091}{mi}{59}\RG{33.407}{0.002}{10.9091}{mi}{118}\RG{38.574}{-1.634}{7.9701}{mi}{110}%
}\hskip.1em\allowbreak{} ，并记\allowbreak{}\hskip.1em \RB{43.555}{5.313}{3.298}{%
\RG{0.000}{0.002}{10.9091}{mi}{109}\RG{9.343}{-1.634}{7.9701}{mi}{105}\RG{15.987}{0.002}{10.9091}{rm}{61}\RG{27.854}{0.002}{10.9091}{mi}{109}\RG{37.197}{-1.634}{7.9701}{mi}{118}\RG{40.960}{-2.743}{5.9776}{mi}{105}%
}\hskip.1em\allowbreak{} 。令\allowbreak{}\hskip.1em \RB{33.621}{5.313}{2.134}{%
\RG{0.000}{0.002}{10.9091}{mi}{122}\RG{4.962}{-1.634}{7.9701}{rm}{49}\RG{12.839}{0.002}{10.9091}{rm}{61}\RG{24.705}{0.002}{10.9091}{mi}{118}\RG{29.870}{-1.634}{7.9701}{rm}{49}%
}\hskip.1em\allowbreak{} 。对\allowbreak{}\hskip.1em \RB{56.277}{7.756}{2.724}{%
\RG{0.000}{0.002}{10.9091}{mi}{106}\RG{8.591}{0.002}{10.9091}{rm}{61}\RG{20.461}{0.002}{10.9091}{rm}{50}\RG{25.799}{0.002}{10.9091}{mi}{59}\RG{30.577}{0.002}{10.9091}{mi}{58}\RG{35.366}{0.002}{10.9091}{mi}{58}\RG{40.155}{0.002}{10.9091}{mi}{58}\RG{44.933}{0.002}{10.9091}{mi}{59}\RG{49.723}{0.002}{10.9091}{mi}{110}%
}\hskip.1em\allowbreak{} ，按（4.2）的构造取非零向量

\[
\RB{123.110}{8.685}{3.617}{%
\RG{0.000}{0.003}{10.9091}{mi}{122}\RG{11.991}{0.003}{10.9091}{sy}{50}\RG{22.295}{0.003}{10.9091}{mi}{67}\RG{30.721}{0.003}{10.9091}{rm}{40}\RG{34.874}{0.003}{10.9091}{mi}{39}\RG{41.867}{0.003}{10.9091}{mi}{59}\RG{46.646}{0.003}{10.9091}{mi}{122}\RG{65.492}{0.003}{10.9091}{rm}{41}\RG{72.070}{0.003}{10.9091}{rm}{43}\RG{82.803}{0.003}{10.9091}{mi}{67}\RG{91.228}{0.003}{10.9091}{rm}{40}\RG{95.382}{0.003}{10.9091}{mi}{39}\RG{102.375}{0.003}{10.9091}{mi}{59}\RG{107.153}{0.003}{10.9091}{mi}{118}\RG{116.328}{0.003}{10.9091}{rm}{41}\RG{120.483}{0.003}{10.9091}{mi}{59}\RG{4.962}{-1.633}{7.9701}{mi}{106}\RG{51.619}{-1.633}{7.9701}{mi}{106}\RG{55.163}{-1.633}{7.9701}{sy}{0}\RG{61.174}{-1.633}{7.9701}{rm}{49}\RG{112.329}{-1.633}{7.9701}{mi}{106}%
}
\]
使得\allowbreak{}\hskip.1em \RB{80.211}{8.683}{4.357}{%
\RG{0.000}{0.001}{10.9091}{mi}{109}\RG{9.343}{-1.635}{7.9701}{mi}{122}\RG{12.938}{-2.744}{5.9776}{mi}{106}\RG{19.984}{0.001}{10.9091}{rm}{61}\RG{31.319}{0.001}{10.9091}{rm}{91}\RG{34.352}{0.001}{10.9091}{mi}{109}\RG{43.696}{-1.635}{7.9701}{mi}{122}\RG{47.290}{-2.744}{5.9776}{mi}{106}\RG{50.397}{-2.744}{5.9776}{sy}{0}\RG{55.650}{-2.744}{5.9776}{rm}{49}\RG{59.892}{0.001}{10.9091}{mi}{59}\RG{64.670}{0.001}{10.9091}{mi}{109}\RG{74.021}{-1.635}{7.9701}{mi}{106}\RG{78.020}{0.001}{10.9091}{rm}{93}%
}\hskip.1em\allowbreak{} 。由有限归纳，

\[
\RB{166.727}{8.678}{4.361}{%
\RG{0.000}{-0.004}{10.9091}{mi}{109}\RG{19.984}{-0.004}{10.9091}{rm}{61}\RG{31.319}{-0.004}{10.9091}{rm}{91}\RG{34.353}{-0.004}{10.9091}{mi}{109}\RG{48.015}{-0.004}{10.9091}{mi}{59}\RG{52.793}{-0.004}{10.9091}{mi}{58}\RG{57.582}{-0.004}{10.9091}{mi}{58}\RG{62.371}{-0.004}{10.9091}{mi}{58}\RG{67.149}{-0.004}{10.9091}{mi}{59}\RG{71.939}{-0.004}{10.9091}{mi}{109}\RG{85.284}{-0.004}{10.9091}{rm}{93}\RG{110.059}{-0.004}{10.9091}{rm}{40}\RG{114.216}{-0.004}{10.9091}{rm}{49}\RG{122.595}{-0.004}{10.9091}{sy}{20}\RG{134.110}{-0.004}{10.9091}{mi}{106}\RG{142.174}{-0.004}{10.9091}{sy}{20}\RG{153.689}{-0.004}{10.9091}{mi}{110}\RG{160.065}{-0.004}{10.9091}{rm}{41}\RG{164.219}{-0.004}{10.9091}{mi}{58}\RG{9.343}{-1.640}{7.9701}{mi}{122}\RG{43.697}{-1.640}{7.9701}{rm}{49}\RG{81.285}{-1.640}{7.9701}{mi}{106}\RG{12.938}{-2.749}{5.9776}{mi}{106}%
}
\]
于是\allowbreak{}\hskip.1em \RB{29.858}{8.312}{2.856}{%
\RG{0.000}{0.001}{10.9091}{mi}{122}\RG{4.961}{-1.635}{7.9701}{mi}{110}\RG{13.136}{0.001}{10.9091}{sy}{54}\RG{13.136}{0.001}{10.9091}{rm}{61}\RG{24.470}{0.001}{10.9091}{rm}{48}%
}\hskip.1em\allowbreak{} ，且由（5）得\allowbreak{}\hskip.1em \RB{41.936}{5.312}{3.099}{%
\RG{0.000}{0.001}{10.9091}{mi}{109}\RG{9.343}{-1.635}{7.9701}{mi}{122}\RG{12.938}{-2.544}{5.9776}{mi}{110}\RG{21.075}{0.001}{10.9091}{rm}{61}\RG{32.414}{0.001}{10.9091}{mi}{109}%
}\hskip.1em\allowbreak{} 。当\allowbreak{}\hskip.1em \RB{25.714}{7.755}{0.608}{%
\RG{0.000}{0.001}{10.9091}{mi}{110}\RG{9.406}{0.001}{10.9091}{rm}{61}\RG{20.741}{0.001}{10.9091}{rm}{49}%
}\hskip.1em\allowbreak{} 时，同一结论直接由\allowbreak{}\hskip.1em \RB{32.564}{5.312}{2.135}{%
\RG{0.000}{0.001}{10.9091}{mi}{122}\RG{4.962}{-1.635}{7.9701}{rm}{49}\RG{12.311}{0.001}{10.9091}{rm}{61}\RG{23.650}{0.001}{10.9091}{mi}{118}\RG{28.813}{-1.635}{7.9701}{rm}{49}%
}\hskip.1em\allowbreak{} 及（5）成立。

% Original page 17.
（7）证明：对\allowbreak{}\hskip.1em \RB{44.631}{7.757}{1.992}{%
\RG{0.000}{0.003}{10.9091}{rm}{49}\RG{8.371}{0.003}{10.9091}{sy}{20}\RG{19.887}{0.003}{10.9091}{mi}{105}\RG{26.561}{0.003}{10.9091}{sy}{20}\RG{38.076}{0.003}{10.9091}{mi}{110}%
}\hskip.1em\allowbreak{} ，定义拉格朗日插值多项式

\[
\RB{100.100}{15.451}{20.542}{%
\RG{42.727}{9.091}{9.0910}{ex}{89}\RG{65.045}{7.380}{10.9091}{mi}{21}\RG{73.700}{7.380}{10.9091}{sy}{0}\RG{84.609}{7.380}{10.9091}{mi}{21}\RG{90.841}{5.744}{7.9701}{mi}{106}\RG{0.000}{-0.000}{10.9091}{mi}{96}\RG{7.567}{-0.000}{10.9091}{rm}{40}\RG{11.722}{-0.000}{10.9091}{mi}{21}\RG{17.953}{-0.000}{10.9091}{rm}{41}\RG{25.131}{-0.000}{10.9091}{rm}{61}\RG{97.473}{-0.000}{10.9091}{mi}{59}\RG{4.482}{-1.636}{7.9701}{mi}{105}\RG{63.502}{-7.483}{10.9091}{mi}{21}\RG{75.243}{-7.483}{10.9091}{sy}{0}\RG{86.152}{-7.483}{10.9091}{mi}{21}\RG{69.733}{-9.120}{7.9701}{mi}{105}\RG{92.383}{-9.120}{7.9701}{mi}{106}\RG{36.476}{-10.923}{7.9701}{rm}{49}\RG{40.339}{-10.923}{7.9701}{sy}{20}\RG{46.349}{-10.923}{7.9701}{mi}{106}\RG{49.894}{-10.923}{7.9701}{sy}{20}\RG{55.904}{-10.923}{7.9701}{mi}{110}\RG{42.442}{-18.478}{7.9701}{mi}{106}\RG{45.986}{-18.478}{7.9701}{sy}{54}\RG{45.986}{-18.478}{7.9701}{rm}{61}\RG{51.996}{-18.478}{7.9701}{mi}{105}\RR{63.502}{2.727}{32.880}{0.364}%
}
\]
其中空积约定为\allowbreak{}\hskip.1em \RB{4.973}{7.754}{0.501}{%
\RG{0.000}{-0.001}{10.9091}{rm}{49}%
}\hskip.1em\allowbreak{} 。特征值互异，故所有分母非零。由\allowbreak{}\hskip.1em \RB{57.735}{8.681}{3.621}{%
\RG{0.000}{-0.001}{10.9091}{mi}{39}\RG{6.998}{-0.001}{10.9091}{rm}{40}\RG{11.152}{-0.001}{10.9091}{mi}{118}\RG{16.316}{-1.637}{7.9701}{mi}{106}\RG{20.314}{-0.001}{10.9091}{rm}{41}\RG{27.493}{-0.001}{10.9091}{rm}{61}\RG{38.837}{-0.001}{10.9091}{mi}{21}\RG{45.069}{-1.637}{7.9701}{mi}{106}\RG{49.068}{-0.001}{10.9091}{mi}{118}\RG{54.232}{-1.637}{7.9701}{mi}{106}%
}\hskip.1em\allowbreak{} ，对非负整数幂归纳并利用线性性，可得\allowbreak{}\hskip.1em \RB{88.664}{8.682}{3.620}{%
\RG{0.000}{-0.000}{10.9091}{mi}{102}\RG{6.430}{-0.000}{10.9091}{rm}{40}\RG{10.584}{-0.000}{10.9091}{mi}{39}\RG{17.581}{-0.000}{10.9091}{rm}{41}\RG{21.738}{-0.000}{10.9091}{rm}{40}\RG{25.890}{-0.000}{10.9091}{mi}{118}\RG{31.054}{-1.636}{7.9701}{mi}{106}\RG{35.052}{-0.000}{10.9091}{rm}{41}\RG{42.962}{-0.000}{10.9091}{rm}{61}\RG{55.028}{-0.000}{10.9091}{mi}{102}\RG{61.458}{-0.000}{10.9091}{rm}{40}\RG{65.612}{-0.000}{10.9091}{mi}{21}\RG{71.843}{-1.636}{7.9701}{mi}{106}\RG{75.842}{-0.000}{10.9091}{rm}{41}\RG{79.995}{-0.000}{10.9091}{mi}{118}\RG{85.160}{-1.636}{7.9701}{mi}{106}%
}\hskip.1em\allowbreak{} ，其中\allowbreak{}\hskip.1em \RB{42.539}{8.682}{3.227}{%
\RG{0.000}{-0.000}{10.9091}{mi}{102}\RG{10.186}{-0.000}{10.9091}{sy}{50}\RG{21.215}{-0.000}{10.9091}{mi}{75}\RG{31.087}{-0.000}{10.9091}{rm}{91}\RG{34.117}{-0.000}{10.9091}{mi}{21}\RG{40.348}{-0.000}{10.9091}{rm}{93}%
}\hskip.1em\allowbreak{} 。因此\allowbreak{}\hskip.1em \RB{63.222}{8.682}{3.238}{%
\RG{0.000}{-0.000}{10.9091}{mi}{96}\RG{4.482}{-1.636}{7.9701}{mi}{105}\RG{7.567}{-0.000}{10.9091}{rm}{40}\RG{11.722}{-0.000}{10.9091}{mi}{39}\RG{18.719}{-0.000}{10.9091}{rm}{41}\RG{22.876}{-0.000}{10.9091}{rm}{40}\RG{27.027}{-0.000}{10.9091}{mi}{118}\RG{32.192}{-1.636}{7.9701}{mi}{105}\RG{35.277}{-0.000}{10.9091}{rm}{41}\RG{43.186}{-0.000}{10.9091}{rm}{61}\RG{55.252}{-0.000}{10.9091}{mi}{118}\RG{60.417}{-1.636}{7.9701}{mi}{105}%
}\hskip.1em\allowbreak{} ，而\allowbreak{}\hskip.1em \RB{24.442}{8.311}{2.856}{%
\RG{0.000}{-0.000}{10.9091}{mi}{106}\RG{8.789}{-0.000}{10.9091}{sy}{54}\RG{8.789}{-0.000}{10.9091}{rm}{61}\RG{20.855}{-0.000}{10.9091}{mi}{105}%
}\hskip.1em\allowbreak{} 时

\[
\RB{71.464}{9.228}{3.620}{%
\RG{0.000}{0.001}{10.9091}{mi}{96}\RG{7.567}{0.001}{10.9091}{rm}{40}\RG{11.722}{0.001}{10.9091}{mi}{39}\RG{18.719}{0.001}{10.9091}{rm}{41}\RG{22.876}{0.001}{10.9091}{rm}{40}\RG{27.028}{0.001}{10.9091}{mi}{118}\RG{36.190}{0.001}{10.9091}{rm}{41}\RG{43.369}{0.001}{10.9091}{rm}{61}\RG{54.714}{0.001}{10.9091}{rm}{48}\RT{60.054}{0.001}{10.9091}{。}\RG{4.482}{-1.636}{7.9701}{mi}{105}\RG{32.192}{-1.636}{7.9701}{mi}{106}%
}
\]
若\allowbreak{}\hskip.1em \RB{62.799}{8.105}{4.677}{%
\RG{0.000}{7.276}{9.0910}{ex}{80}\RG{9.596}{4.397}{7.9701}{mi}{110}\RG{9.596}{-2.694}{7.9701}{mi}{106}\RG{13.140}{-2.694}{7.9701}{rm}{61}\RG{19.148}{-2.694}{7.9701}{rm}{49}\RG{25.287}{0.004}{10.9091}{mi}{99}\RG{29.884}{-1.632}{7.9701}{mi}{106}\RG{33.883}{0.004}{10.9091}{mi}{118}\RG{39.047}{-1.632}{7.9701}{mi}{106}\RG{46.076}{0.004}{10.9091}{rm}{61}\RG{57.411}{0.004}{10.9091}{rm}{48}%
}\hskip.1em\allowbreak{} ，其中\allowbreak{}\hskip.1em \RB{31.964}{7.955}{3.616}{%
\RG{0.000}{0.004}{10.9091}{mi}{99}\RG{4.597}{-1.632}{7.9701}{mi}{106}\RG{11.626}{0.004}{10.9091}{sy}{50}\RG{21.930}{0.004}{10.9091}{mi}{75}%
}\hskip.1em\allowbreak{} ，对等式施以\allowbreak{}\hskip.1em \RB{22.296}{8.686}{3.235}{%
\RG{0.000}{0.004}{10.9091}{mi}{96}\RG{4.482}{-1.632}{7.9701}{mi}{105}\RG{7.567}{0.004}{10.9091}{rm}{40}\RG{11.721}{0.004}{10.9091}{mi}{39}\RG{18.719}{0.004}{10.9091}{rm}{41}%
}\hskip.1em\allowbreak{} 得\allowbreak{}\hskip.1em \RB{35.685}{7.758}{2.205}{%
\RG{0.000}{0.004}{10.9091}{mi}{99}\RG{4.597}{-1.632}{7.9701}{mi}{105}\RG{7.682}{0.004}{10.9091}{mi}{118}\RG{12.846}{-1.632}{7.9701}{mi}{105}\RG{18.963}{0.004}{10.9091}{rm}{61}\RG{30.298}{0.004}{10.9091}{rm}{48}%
}\hskip.1em\allowbreak{} 。由于\allowbreak{}\hskip.1em \RB{28.002}{8.315}{2.853}{%
\RG{0.000}{0.004}{10.9091}{mi}{118}\RG{5.164}{-1.632}{7.9701}{mi}{105}\RG{11.280}{0.004}{10.9091}{sy}{54}\RG{11.280}{0.004}{10.9091}{rm}{61}\RG{22.614}{0.004}{10.9091}{rm}{48}%
}\hskip.1em\allowbreak{} ，有\allowbreak{}\hskip.1em \RB{27.435}{7.758}{2.205}{%
\RG{0.000}{0.004}{10.9091}{mi}{99}\RG{4.597}{-1.632}{7.9701}{mi}{105}\RG{10.713}{0.004}{10.9091}{rm}{61}\RG{22.048}{0.004}{10.9091}{rm}{48}%
}\hskip.1em\allowbreak{} 。这对每个\allowbreak{}\hskip.1em \RB{3.587}{7.672}{0.605}{%
\RG{0.000}{0.005}{10.9091}{mi}{105}%
}\hskip.1em\allowbreak{} 成立，故\allowbreak{}\hskip.1em \RB{43.525}{5.316}{2.601}{%
\RG{0.000}{0.005}{10.9091}{mi}{118}\RG{5.164}{-1.632}{7.9701}{rm}{49}\RG{9.483}{0.005}{10.9091}{mi}{59}\RG{14.261}{0.005}{10.9091}{mi}{58}\RG{19.050}{0.005}{10.9091}{mi}{58}\RG{23.829}{0.005}{10.9091}{mi}{58}\RG{28.618}{0.005}{10.9091}{mi}{59}\RG{33.407}{0.005}{10.9091}{mi}{118}\RG{38.574}{-1.632}{7.9701}{mi}{110}%
}\hskip.1em\allowbreak{} 线性无关；它们的数目为\allowbreak{}\hskip.1em \RB{6.555}{5.316}{0.605}{%
\RG{0.000}{0.005}{10.9091}{mi}{110}%
}\hskip.1em\allowbreak{} ，因此构成\allowbreak{}\hskip.1em \RB{8.715}{7.956}{0.725}{%
\RG{0.000}{0.005}{10.9091}{mi}{86}%
}\hskip.1em\allowbreak{} 的一组基。

任取\allowbreak{}\hskip.1em \RB{29.977}{7.949}{0.938}{%
\RG{0.000}{-0.002}{10.9091}{mi}{119}\RG{10.959}{-0.002}{10.9091}{sy}{50}\RG{21.262}{-0.002}{10.9091}{mi}{86}%
}\hskip.1em\allowbreak{} ，唯一地写成\allowbreak{}\hskip.1em \RB{62.324}{8.099}{3.271}{%
\RG{0.000}{-0.002}{10.9091}{mi}{119}\RG{10.959}{-0.002}{10.9091}{rm}{61}\RG{22.298}{7.271}{9.0910}{ex}{80}\RG{31.894}{4.392}{7.9701}{mi}{110}\RG{31.894}{-2.698}{7.9701}{mi}{105}\RG{34.525}{-2.699}{7.9701}{rm}{61}\RG{40.532}{-2.699}{7.9701}{rm}{49}\RG{46.671}{-0.002}{10.9091}{mi}{99}\RG{51.269}{-1.638}{7.9701}{mi}{105}\RG{54.355}{-0.002}{10.9091}{mi}{118}\RG{59.519}{-1.638}{7.9701}{mi}{105}%
}\hskip.1em\allowbreak{} 。若所有\allowbreak{}\hskip.1em \RB{7.403}{5.309}{2.211}{%
\RG{0.000}{-0.002}{10.9091}{mi}{99}\RG{4.597}{-1.638}{7.9701}{mi}{105}%
}\hskip.1em\allowbreak{} 均非零，则

\[
\RB{203.641}{9.844}{3.266}{%
\RG{27.216}{4.508}{7.9701}{sy}{0}\RG{33.226}{4.508}{7.9701}{rm}{49}\RG{0.000}{0.004}{10.9091}{mi}{118}\RG{11.280}{0.004}{10.9091}{rm}{61}\RG{22.618}{0.004}{10.9091}{mi}{99}\RG{37.544}{0.004}{10.9091}{mi}{96}\RG{45.112}{0.004}{10.9091}{rm}{40}\RG{49.267}{0.004}{10.9091}{mi}{39}\RG{56.264}{0.004}{10.9091}{rm}{41}\RG{60.420}{0.004}{10.9091}{rm}{40}\RG{64.572}{0.004}{10.9091}{mi}{119}\RG{72.501}{0.004}{10.9091}{rm}{41}\RG{79.685}{0.004}{10.9091}{sy}{50}\RG{89.988}{0.004}{10.9091}{mi}{67}\RG{98.414}{0.004}{10.9091}{rm}{40}\RG{102.568}{0.004}{10.9091}{mi}{39}\RG{109.561}{0.004}{10.9091}{mi}{59}\RG{114.339}{0.004}{10.9091}{mi}{119}\RG{122.280}{0.004}{10.9091}{rm}{41}\RG{148.255}{0.004}{10.9091}{rm}{40}\RG{152.412}{0.004}{10.9091}{rm}{49}\RG{160.778}{0.004}{10.9091}{sy}{20}\RG{172.293}{0.004}{10.9091}{mi}{105}\RG{178.967}{0.004}{10.9091}{sy}{20}\RG{190.483}{0.004}{10.9091}{mi}{110}\RG{196.859}{0.004}{10.9091}{rm}{41}\RG{201.014}{0.004}{10.9091}{mi}{59}\RG{5.164}{-1.633}{7.9701}{mi}{105}\RG{42.027}{-1.633}{7.9701}{mi}{105}\RG{27.216}{-2.694}{7.9701}{mi}{105}%
}
\]
因而\allowbreak{}\hskip.1em \RB{59.530}{8.681}{3.239}{%
\RG{0.000}{-0.001}{10.9091}{mi}{67}\RG{8.426}{-0.001}{10.9091}{rm}{40}\RG{12.580}{-0.001}{10.9091}{mi}{39}\RG{19.573}{-0.001}{10.9091}{mi}{59}\RG{24.351}{-0.001}{10.9091}{mi}{119}\RG{32.292}{-0.001}{10.9091}{rm}{41}\RG{39.470}{-0.001}{10.9091}{rm}{61}\RG{50.815}{-0.001}{10.9091}{mi}{86}%
}\hskip.1em\allowbreak{} 。若存在\allowbreak{}\hskip.1em \RB{3.587}{7.666}{0.610}{%
\RG{0.000}{-0.001}{10.9091}{mi}{105}%
}\hskip.1em\allowbreak{} 使得\allowbreak{}\hskip.1em \RB{27.436}{7.754}{2.210}{%
\RG{0.000}{-0.001}{10.9091}{mi}{99}\RG{4.597}{-1.637}{7.9701}{mi}{105}\RG{10.714}{-0.001}{10.9091}{rm}{61}\RG{22.049}{-0.001}{10.9091}{rm}{48}%
}\hskip.1em\allowbreak{} ，则对任意非负整数\allowbreak{}\hskip.1em \RB{5.147}{5.310}{0.610}{%
\RG{0.000}{-0.001}{10.9091}{mi}{114}%
}\hskip.1em\allowbreak{} ，

\[
\RB{192.789}{9.595}{20.538}{%
\RG{47.253}{9.095}{9.0910}{ex}{88}\RG{6.998}{4.509}{7.9701}{mi}{114}\RG{82.522}{4.509}{7.9701}{mi}{114}\RG{0.000}{0.005}{10.9091}{mi}{39}\RG{11.154}{0.005}{10.9091}{rm}{40}\RG{15.308}{0.005}{10.9091}{mi}{119}\RG{23.237}{0.005}{10.9091}{rm}{41}\RG{30.415}{0.005}{10.9091}{rm}{61}\RG{67.696}{0.005}{10.9091}{mi}{99}\RG{76.292}{0.005}{10.9091}{mi}{21}\RG{86.679}{0.005}{10.9091}{mi}{118}\RG{98.872}{0.005}{10.9091}{sy}{50}\RG{109.175}{0.005}{10.9091}{rm}{115}\RG{113.386}{0.005}{10.9091}{rm}{112}\RG{119.320}{0.005}{10.9091}{rm}{97}\RG{124.666}{0.005}{10.9091}{rm}{110}\RG{138.136}{0.005}{10.9091}{sy}{102}\RG{143.590}{0.005}{10.9091}{mi}{118}\RG{155.784}{0.005}{10.9091}{rm}{58}\RG{161.781}{0.005}{10.9091}{mi}{106}\RG{169.845}{0.005}{10.9091}{sy}{54}\RG{169.845}{0.005}{10.9091}{rm}{61}\RG{181.183}{0.005}{10.9091}{mi}{105}\RG{184.828}{0.005}{10.9091}{sy}{103}\RG{190.282}{0.005}{10.9091}{mi}{58}\RG{72.293}{-1.632}{7.9701}{mi}{106}\RG{91.843}{-1.632}{7.9701}{mi}{106}\RG{148.755}{-1.632}{7.9701}{mi}{106}\RG{130.651}{-2.571}{7.9701}{mi}{75}\RG{82.522}{-2.693}{7.9701}{mi}{106}\RG{41.760}{-10.919}{7.9701}{rm}{49}\RG{45.624}{-10.919}{7.9701}{sy}{20}\RG{51.635}{-10.919}{7.9701}{mi}{106}\RG{55.179}{-10.919}{7.9701}{sy}{20}\RG{61.189}{-10.919}{7.9701}{mi}{110}\RG{47.726}{-18.475}{7.9701}{mi}{106}\RG{51.270}{-18.475}{7.9701}{sy}{54}\RG{51.270}{-18.475}{7.9701}{rm}{61}\RG{57.280}{-18.475}{7.9701}{mi}{105}%
}
\]
该子空间维数为\allowbreak{}\hskip.1em \RB{24.683}{7.756}{0.607}{%
\RG{0.000}{0.002}{10.9091}{mi}{110}\RG{8.800}{0.002}{10.9091}{sy}{0}\RG{19.710}{0.002}{10.9091}{rm}{49}%
}\hskip.1em\allowbreak{} ，故\allowbreak{}\hskip.1em \RB{59.530}{8.684}{3.236}{%
\RG{0.000}{0.002}{10.9091}{mi}{67}\RG{8.426}{0.002}{10.9091}{rm}{40}\RG{12.580}{0.002}{10.9091}{mi}{39}\RG{19.573}{0.002}{10.9091}{mi}{59}\RG{24.351}{0.002}{10.9091}{mi}{119}\RG{32.292}{0.002}{10.9091}{rm}{41}\RG{39.476}{0.002}{10.9091}{sy}{54}\RG{39.476}{0.002}{10.9091}{rm}{61}\RG{50.815}{0.002}{10.9091}{mi}{86}%
}\hskip.1em\allowbreak{} 。因此全部循环向量恰为

\allowbreak{}\hskip.1em \RB{104.378}{16.865}{11.441}{%
\RG{0.000}{16.002}{9.0910}{ex}{40}\RG{7.324}{9.092}{9.0910}{ex}{88}\RG{7.637}{-10.868}{7.9701}{mi}{105}\RG{10.268}{-10.868}{7.9701}{rm}{61}\RG{11.545}{11.820}{7.9701}{mi}{110}\RG{16.275}{-10.868}{7.9701}{rm}{49}\RG{22.274}{0.001}{10.9091}{mi}{99}\RG{26.871}{-1.635}{7.9701}{mi}{105}\RG{29.956}{0.001}{10.9091}{mi}{118}\RG{35.120}{-1.635}{7.9701}{mi}{105}\RG{41.236}{0.001}{10.9091}{rm}{58}\RG{47.233}{0.001}{10.9091}{mi}{99}\RG{51.831}{-1.635}{7.9701}{mi}{105}\RG{57.946}{0.001}{10.9091}{sy}{50}\RG{68.250}{0.001}{10.9091}{mi}{75}\RG{80.545}{0.001}{10.9091}{sy}{110}\RG{88.421}{0.001}{10.9091}{sy}{102}\RG{93.878}{0.001}{10.9091}{rm}{48}\RG{99.219}{0.001}{10.9091}{sy}{103}%
}\hskip.1em\allowbreak{} 对所有\allowbreak{}\hskip.1em \RB{15.293}{16.865}{11.398}{%
\RG{0.000}{0.001}{10.9091}{mi}{105}\RG{3.645}{16.002}{9.0910}{ex}{41}\RG{12.786}{0.001}{10.9091}{mi}{58}%
}\hskip.1em\allowbreak{}

特别地，\allowbreak{}\hskip.1em \RB{58.639}{6.782}{2.211}{%
\RG{0.000}{-0.002}{10.9091}{mi}{118}\RG{5.164}{-1.638}{7.9701}{rm}{49}\RG{11.906}{-0.002}{10.9091}{rm}{43}\RG{22.639}{-0.002}{10.9091}{sy}{1}\RG{27.483}{-0.002}{10.9091}{sy}{1}\RG{32.338}{-0.002}{10.9091}{sy}{1}\RG{37.791}{-0.002}{10.9091}{rm}{43}\RG{48.523}{-0.002}{10.9091}{mi}{118}\RG{53.688}{-1.638}{7.9701}{mi}{110}%
}\hskip.1em\allowbreak{} 是循环向量，故\allowbreak{}\hskip.1em \RB{8.715}{7.949}{0.731}{%
\RG{0.000}{-0.002}{10.9091}{mi}{86}%
}\hskip.1em\allowbreak{} 是循环空间。\allowbreak{}\hskip.1em \RB{8.366}{7.993}{0.502}{%
\RG{0.000}{-0.002}{10.9091}{am}{3}%
}\hskip.1em\allowbreak{}


\end{document}
